\documentclass[10pt,a4paper]{article}
\usepackage[utf8]{inputenc}
\usepackage[margin=0.75in,top=0.8in,bottom=0.8in]{geometry}
\usepackage{xcolor}
\usepackage{titlesec}
\usepackage{enumitem}
\usepackage{hyperref}
\usepackage{tcolorbox}
\usepackage{tabularx}
\usepackage{fancyhdr}
\usepackage{microtype}
\usepackage{amssymb}

\tcbuselibrary{skins,breakable}

% -----------------------------------------------------------------------------
% COLOR PALETTE
% -----------------------------------------------------------------------------
\definecolor{brandPrimary}{RGB}{20, 45, 85}      % Deep Navy
\definecolor{brandSecondary}{RGB}{41, 128, 185}  % Slate Blue
\definecolor{brandDark}{RGB}{33, 37, 41}         % Charcoal Body Text
\definecolor{brandLight}{RGB}{95, 99, 230}       % Very Light Purple-Blue for What you'll learn
\definecolor{cardBg}{RGB}{248, 249, 250}         % Light Off-White
\definecolor{cardBorder}{RGB}{222, 226, 230}     % Subtle Border Gray
\definecolor{accentBg}{RGB}{235, 243, 250}       % Soft Highlight Blue

% -----------------------------------------------------------------------------
% TYPOGRAPHY & HYPERLINKS
% -----------------------------------------------------------------------------
\renewcommand{\familydefault}{\sfdefault}
\hypersetup{
    colorlinks=true,
    linkcolor=brandSecondary,
    urlcolor=brandSecondary,
    pdftitle={Skillzza Virtual Internship Catalog}
}

% -----------------------------------------------------------------------------
% HEADERS & FOOTERS
% -----------------------------------------------------------------------------
\pagestyle{fancy}
\fancyhf{}
\renewcommand{\headrulewidth}{0.4pt}
\renewcommand{\headrule}{\hbox to\headwidth{\color{cardBorder}\leaders\hrule height \headrulewidth\hfill}}
\lhead{\footnotesize \textcolor{brandSecondary}{\textbf{Skillzza Job Simulations}}}
\rhead{\footnotesize \textcolor{gray}{Virtual Internship Catalog}}
\cfoot{\footnotesize \textcolor{gray}{\thepage}}

% -----------------------------------------------------------------------------
% SECTION STYLES
% -----------------------------------------------------------------------------
\titleformat{\section}
  {\normalfont\Large\bfseries\color{brandPrimary}}
  {\thesection}{1em}{}[\vspace{2pt}\color{brandSecondary}\hrule height 1.5pt\vspace{4pt}]

\titlespacing*{\section}{0pt}{14pt}{10pt}

% -----------------------------------------------------------------------------
% CUSTOM TCOLORBOXES
% -----------------------------------------------------------------------------
\newtcolorbox{metabox}{
    enhanced,
    colback=accentBg,
    colframe=brandSecondary,
    arc=4mm,
    boxrule=1pt,
    left=12pt, right=12pt, top=10pt, bottom=10pt,
    before skip=6pt, after skip=10pt
}

\newtcolorbox{taskbox}[2][]{
    enhanced,
    breakable,
    colback=cardBg,
    colframe=cardBorder,
    arc=3mm,
    boxrule=0.8pt,
    left=12pt, right=12pt, top=12pt, bottom=12pt,
    before skip=12pt, after skip=12pt,
    title={\textbf{\large \textcolor{brandPrimary}{#2}}},
    coltitle=brandPrimary,
    fonttitle=\bfseries\Large,
    attach boxed title to top left={xshift=10pt, yshift=-3mm},
    boxed title style={
        colback=white,
        colframe=white,
        boxrule=0pt,
        arc=2mm,
        top=2pt, bottom=2pt, left=8pt, right=8pt
    },
    #1
}

\newtcolorbox{completionbox}{
    enhanced,
    colback=white,
    colframe=brandPrimary,
    arc=3mm,
    boxrule=1pt,
    left=10pt, right=10pt, top=8pt, bottom=8pt,
    before skip=12pt, after skip=8pt
}

\setlist[itemize]{leftmargin=1.4em, itemsep=2pt, topsep=3pt, parsep=0pt}
\setlength{\parindent}{0pt}
\setlength{\parskip}{8pt}

\begin{document}

\begin{titlepage}
    \centering
    \vspace*{2cm}
    {\Huge \bfseries \textcolor{brandPrimary}{Skillzza Virtual Internships}}\\[0.8cm]
    {\Large \textcolor{brandSecondary}{Task-Based Enterprise Job Simulations \& AI Curricula}}\\[2cm]
    
    \begin{tcolorbox}[enhanced, colback=accentBg, colframe=brandSecondary, width=0.85\textwidth, arc=4mm, boxrule=1pt, halign=center]
        \vspace{6pt}
        \textbf{\large Comprehensive Course Catalog \& Deliverables Guide}\\[6pt]
        \small Featuring curated simulations with industry partners including \textbf{NeuralFrontier}, \textbf{Stratagem Consultants}, \textbf{SynthAI Labs}, and \textbf{NovaCare Health Systems}.
        \vspace{6pt}
    \end{tcolorbox}
    
    \vfill
    {\small \textcolor{gray}{Generated with Skillzza Verified Simulation Engine $\bullet$ Self-Paced}}\\[1cm]
\end{titlepage}

\tableofcontents
\newpage

\newpage

% =============================================================================

% SIMULATION 1

% =============================================================================

\section*{AI Agent Designer – Build Your First Agent}

\addcontentsline{toc}{section}{AI Agent Designer – Build Your First Agent}

\begin{metabox}
    \begin{tabularx}{\textwidth}{@{}ll X@{}}
        \textbf{Industry:} & Technology & \textbf{Partner Enterprise:} \textbf{Skillzza Enterprise} \\
        \textbf{Duration:} & 12-18 hrs (Self-paced) & \textbf{Mode:} Individual, task-based simulation \\
        \textbf{Primary Skills:} & \multicolumn{2}{@{}l}{Python, AI, Data Analysis, Problem Solving}
    \end{tabularx}
\end{metabox}

\subsection*{Overview \& Problem Statement}

\textbf{Scenario:}  
As an AI Agent Designer, you are tasked with creating an intelligent virtual assistant capable of helping users navigate complex tasks efficiently. From understanding user goals to integrating memory and evaluating performance, your mission is to design an AI agent that balances functionality, usability, and ethical considerations. Imagine working on a groundbreaking AI product for a global tech company that aims to revolutionize customer service, productivity, and personal assistance.

\textbf{Mission:}  
Your mission is to design, build, and test an AI agent equipped with conversational capabilities, memory retention, and a goal-oriented problem-solving approach. The agent should be able to understand user intent, recall relevant information, and deliver solutions effectively. You will also focus on responsible AI practices, ensuring that your agent meets ethical standards while offering high usability.

\textbf{Final Challenge:}  
You will present your fully-functional AI agent prototype to a simulated product team. Your agent should demonstrate its ability to understand user goals, retrieve stored memories, and execute tasks with high precision while adhering to responsible AI guidelines. Your solution will be evaluated on technical accuracy, innovation, and ethical compliance.

---

\subsection*{Tasks \& Deliverables}

\begin{taskbox}{Task 1: Understand}

    {\small \textcolor{gray}{2--3 hrs $\cdot$ Intermediate}}\\[12pt]

      
You are tasked with designing an AI agent for a productivity app. The agent will assist users in scheduling meetings, tracking tasks, and answering queries.


\vspace{8pt}
{\large \textbf{\textcolor{brandLight}{What you'll learn}}}\\[4pt]  
Research and define the core functionalities of the agent. Create a user persona and list potential use cases for the agent.

\vspace{8pt}
{\large \textbf{\textcolor{brandLight}{What you'll do}}}\\[4pt]  
\begin{itemize}
    \item A document outlining the agent's purpose, target audience, and core functionalities.  
    \item A list of 5-7 use cases based on the user persona.\par
\end{itemize}

    \vspace{16pt}

    \textcolor{cardBorder}{\hrule height 1pt}\vspace{8pt}

    \small \textcolor{gray}{When you are done with this task, mark it as complete to proceed to the next one.} \hfill \textbf{\textcolor{brandDark}{$\checkmark$ Completed}}

\end{taskbox}

\begin{taskbox}{Task 2: Data/Setup}

    {\small \textcolor{gray}{2--3 hrs $\cdot$ Intermediate}}\\[12pt]

      
To enable memory and goal-oriented behavior, you need to set up the foundational architecture of the AI agent. This includes preparing the environment and connecting to the OpenAI API.


\vspace{8pt}
{\large \textbf{\textcolor{brandLight}{What you'll learn}}}\\[4pt]  
Set up Python, install necessary libraries (OpenAI, LangChain), and connect to the GPT model. Design a basic agent skeleton with input/output flow.

\vspace{8pt}
{\large \textbf{\textcolor{brandLight}{What you'll do}}}\\[4pt]  
\begin{itemize}
    \item Python script with API setup and a basic agent scaffold.  
    \item Screenshot of the agent successfully generating sample responses.\par
\end{itemize}

    \vspace{16pt}

    \textcolor{cardBorder}{\hrule height 1pt}\vspace{8pt}

    \small \textcolor{gray}{When you are done with this task, mark it as complete to proceed to the next one.} \hfill \textbf{\textcolor{brandDark}{$\checkmark$ Completed}}

\end{taskbox}

\begin{taskbox}{Task 3: Build/Execute}

    {\small \textcolor{gray}{2--3 hrs $\cdot$ Intermediate}}\\[12pt]

      
You will now integrate memory functionality so the agent can retain and recall user information across sessions. Add goal-setting logic to allow the agent to guide users toward task completion.


\vspace{8pt}
{\large \textbf{\textcolor{brandLight}{What you'll learn}}}\\[4pt]  
Use LangChain to implement memory and goal-setting modules. Train the agent to store user preferences and recall them in future interactions. Test the agent's ability to track goals across conversations.

\vspace{8pt}
{\large \textbf{\textcolor{brandLight}{What you'll do}}}\\[4pt]  
\begin{itemize}
    \item Python script with memory integration and goal-setting logic.  
    \item A set of test cases demonstrating that the agent can recall memory and guide users to complete tasks.\par
\end{itemize}

    \vspace{16pt}

    \textcolor{cardBorder}{\hrule height 1pt}\vspace{8pt}

    \small \textcolor{gray}{When you are done with this task, mark it as complete to proceed to the next one.} \hfill \textbf{\textcolor{brandDark}{$\checkmark$ Completed}}

\end{taskbox}

\begin{taskbox}{Task 4: GenAI/Explanation}

    {\small \textcolor{gray}{2--3 hrs $\cdot$ Intermediate}}\\[12pt]

      
Users often need step-by-step instructions for complex processes. Add a generative layer to the agent for explaining multi-step tasks clearly.


\vspace{8pt}
{\large \textbf{\textcolor{brandLight}{What you'll learn}}}\\[4pt]  
Integrate OpenAI's GPT capabilities to enable the agent to generate step-by-step instructions for tasks such as scheduling a meeting or organizing a to-do list.

\vspace{8pt}
{\large \textbf{\textcolor{brandLight}{What you'll do}}}\\[4pt]  
\begin{itemize}
    \item Python script demonstrating generative capabilities for step-by-step instructions.  
    \item A test log showing the agent successfully explaining three distinct workflows.\par
\end{itemize}

    \vspace{16pt}

    \textcolor{cardBorder}{\hrule height 1pt}\vspace{8pt}

    \small \textcolor{gray}{When you are done with this task, mark it as complete to proceed to the next one.} \hfill \textbf{\textcolor{brandDark}{$\checkmark$ Completed}}

\end{taskbox}

\begin{taskbox}{Task 5: Audit/Responsible AI}

    {\small \textcolor{gray}{2--3 hrs $\cdot$ Intermediate}}\\[12pt]

      
You must ensure the agent adheres to responsible AI practices. This includes avoiding biased responses, ensuring user data privacy, and providing disclaimers for ambiguous answers.


\vspace{8pt}
{\large \textbf{\textcolor{brandLight}{What you'll learn}}}\\[4pt]  
Implement filters to detect and manage biased or inappropriate responses. Add user consent mechanisms for data storage and integrate disclaimers for uncertain outputs.

\vspace{8pt}
{\large \textbf{\textcolor{brandLight}{What you'll do}}}\\[4pt]  
\begin{itemize}
    \item Python script with responsible AI features implemented.  
    \item A document outlining responsible AI measures and testing results.\par
\end{itemize}

    \vspace{16pt}

    \textcolor{cardBorder}{\hrule height 1pt}\vspace{8pt}

    \small \textcolor{gray}{When you are done with this task, mark it as complete to proceed to the next one.} \hfill \textbf{\textcolor{brandDark}{$\checkmark$ Completed}}

\end{taskbox}

\begin{taskbox}{Task 6: Present Recommendation}

    {\small \textcolor{gray}{2--3 hrs $\cdot$ Intermediate}}\\[12pt]

      
Your agent is ready for deployment, but you need to present its capabilities to the product team. Create a demo showcasing the agent's functionalities and explain its design choices.


\vspace{8pt}
{\large \textbf{\textcolor{brandLight}{What you'll learn}}}\\[4pt]  
Develop a short presentation and record a demo of the agent performing key tasks, showcasing its memory, goal-setting, and instructional capabilities. Highlight responsible AI features.

\vspace{8pt}
{\large \textbf{\textcolor{brandLight}{What you'll do}}}\\[4pt]  
\begin{itemize}
    \item Presentation slides (5-10 slides) covering agent design, features, and responsible AI considerations.  
    \item A video demo (3-5 minutes) showing the agent in action.\par
\end{itemize}

    \vspace{16pt}

    \textcolor{cardBorder}{\hrule height 1pt}\vspace{8pt}

    \small \textcolor{gray}{When you are done with this task, mark it as complete to proceed to the next one.} \hfill \textbf{\textcolor{brandDark}{$\checkmark$ Completed}}

\end{taskbox}

\begin{completionbox}
    \textbf{Knowledge Assessment \& Verification:} Complete the online knowledge check, submit your code and presentation files for AI verification, and claim your \textbf{Skillzza Verified Certificate}.
\end{completionbox}

\newpage

% =============================================================================

% SIMULATION 2

% =============================================================================

\section*{AI Research Analyst – Research a New Market}

\addcontentsline{toc}{section}{AI Research Analyst – Research a New Market}

\begin{metabox}
    \begin{tabularx}{\textwidth}{@{}ll X@{}}
        \textbf{Industry:} & Technology & \textbf{Partner Enterprise:} \textbf{Skillzza Enterprise} \\
        \textbf{Duration:} & 12-18 hrs (Self-paced) & \textbf{Mode:} Individual, task-based simulation \\
        \textbf{Primary Skills:} & \multicolumn{2}{@{}l}{Python, AI, Data Analysis, Problem Solving}
    \end{tabularx}
\end{metabox}

\subsection*{Overview \& Problem Statement}

\textbf{Scenario:}  
You have just been hired as an AI Research Analyst by a leading AI consulting firm. Your first assignment is to research a new market ripe for AI adoption. The firm is exploring opportunities to expand its AI offerings into industries that have traditionally relied on manual processes but could benefit from automation and data-driven decision-making. As part of your role, you will leverage AI tools to collect, analyze, and validate information and provide actionable recommendations to guide the firm's market entry strategy.

\textbf{Mission:}  
Conduct a detailed market research analysis using AI-powered tools. Your goal is to identify emerging trends, key players, challenges, and opportunities in the selected market. You will validate your data sources using responsible AI principles and synthesize your findings into actionable insights for decision-makers.

\textbf{Final Challenge:}  
Prepare a professional research report and presentation summarizing your findings. Your report should highlight the market potential, risks, and an AI-driven opportunity roadmap that will help your firm make informed investment and expansion decisions.

---

\subsection*{Tasks \& Deliverables}

\begin{taskbox}{Task 1: Understand the Scope of Research}

    {\small \textcolor{gray}{2--3 hrs $\cdot$ Intermediate}}\\[12pt]

      
Your manager has provided a vague request: "Find out how AI can transform the logistics industry." You need to refine the research scope to make it actionable.  


\vspace{8pt}
{\large \textbf{\textcolor{brandLight}{What you'll learn}}}\\[4pt]  
\begin{itemize}
    \item Break down the research question into sub-questions (e.g., "What are current pain points in logistics that AI can address-).  
    \item Research current AI applications in logistics using Google Scholar and ChatGPT.  
    \item Identify specific areas (e.g., route optimization, inventory management) for deeper research.  
\end{itemize}

\vspace{8pt}
{\large \textbf{\textcolor{brandLight}{What you'll do}}}\\[4pt]  
Submit a refined research scope document (1 page) outlining sub-questions, focus areas, and the rationale for choosing logistics as the target industry.\par

    \vspace{16pt}

    \textcolor{cardBorder}{\hrule height 1pt}\vspace{8pt}

    \small \textcolor{gray}{When you are done with this task, mark it as complete to proceed to the next one.} \hfill \textbf{\textcolor{brandDark}{$\checkmark$ Completed}}

\end{taskbox}

\begin{taskbox}{Task 2: Data Collection and Setup}

    {\small \textcolor{gray}{2--3 hrs $\cdot$ Intermediate}}\\[12pt]

      
You need reliable data to perform your analysis. Collect information from primary (official reports, industry statistics) and secondary (blogs, AI forums) sources.  


\vspace{8pt}
{\large \textbf{\textcolor{brandLight}{What you'll learn}}}\\[4pt]  
\begin{itemize}
    \item Use AI-powered web scraping tools (or manually search) to collect recent data on logistics trends and challenges.  
    \item Organize the data into categories (e.g., pain points, solutions, emerging technologies).  
    \item Ensure sources are credible by cross-checking for bias and accuracy.  
\end{itemize}

\vspace{8pt}
{\large \textbf{\textcolor{brandLight}{What you'll do}}}\\[4pt]  
Submit a structured dataset (Google Sheets or Excel) with at least 20 data points categorized and validated for your analysis.\par

    \vspace{16pt}

    \textcolor{cardBorder}{\hrule height 1pt}\vspace{8pt}

    \small \textcolor{gray}{When you are done with this task, mark it as complete to proceed to the next one.} \hfill \textbf{\textcolor{brandDark}{$\checkmark$ Completed}}

\end{taskbox}

\begin{taskbox}{Task 3: Build Insights Using AI Tools}

    {\small \textcolor{gray}{2--3 hrs $\cdot$ Intermediate}}\\[12pt]

      
You now have raw data and need to uncover trends and actionable insights. Leverage AI tools to synthesize information.  


\vspace{8pt}
{\large \textbf{\textcolor{brandLight}{What you'll learn}}}\\[4pt]  
\begin{itemize}
    \item Use ChatGPT to summarize industry trends and challenges based on your data.  
    \item Apply NLP techniques to identify recurring terms and themes in textual data.  
    \item Highlight correlations between AI adoption and efficiency gains in logistics.  
\end{itemize}

\vspace{8pt}
{\large \textbf{\textcolor{brandLight}{What you'll do}}}\\[4pt]  
Submit a concise insights report (2 pages) summarizing trends, challenges, and opportunities derived from the analysis.\par

    \vspace{16pt}

    \textcolor{cardBorder}{\hrule height 1pt}\vspace{8pt}

    \small \textcolor{gray}{When you are done with this task, mark it as complete to proceed to the next one.} \hfill \textbf{\textcolor{brandDark}{$\checkmark$ Completed}}

\end{taskbox}

\begin{taskbox}{Task 4: GenAI-Powered Explanation}

    {\small \textcolor{gray}{2--3 hrs $\cdot$ Intermediate}}\\[12pt]

      
You need to explain complex findings to non-experts. Use generative AI tools to craft an engaging summary.  


\vspace{8pt}
{\large \textbf{\textcolor{brandLight}{What you'll learn}}}\\[4pt]  
\begin{itemize}
    \item In ChatGPT, generate a simplified explanation of your insights for stakeholders unfamiliar with AI.  
    \item Create a visual representation (infographic) that showcases key points using tools like Canva or Google Slides.  
\end{itemize}

\vspace{8pt}
{\large \textbf{\textcolor{brandLight}{What you'll do}}}\\[4pt]  
Submit a one-page stakeholder-friendly summary and a visual infographic explaining your research findings.\par

    \vspace{16pt}

    \textcolor{cardBorder}{\hrule height 1pt}\vspace{8pt}

    \small \textcolor{gray}{When you are done with this task, mark it as complete to proceed to the next one.} \hfill \textbf{\textcolor{brandDark}{$\checkmark$ Completed}}

\end{taskbox}

\begin{taskbox}{Task 5: Audit and Responsible AI Practices}

    {\small \textcolor{gray}{2--3 hrs $\cdot$ Intermediate}}\\[12pt]

      
Your firm emphasizes responsible AI use. Ensure your research adheres to ethical standards and avoids biased data.  


\vspace{8pt}
{\large \textbf{\textcolor{brandLight}{What you'll learn}}}\\[4pt]  
\begin{itemize}
    \item Review your data sources for potential bias or misinformation.  
    \item Write a short audit report detailing how you ensured ethical AI practices throughout your research.  
    \item Identify any gaps or risks in your findings.  
\end{itemize}

\vspace{8pt}
{\large \textbf{\textcolor{brandLight}{What you'll do}}}\\[4pt]  
Submit a one-page Responsible AI Audit Report summarizing your validation process and ethical considerations.\par

    \vspace{16pt}

    \textcolor{cardBorder}{\hrule height 1pt}\vspace{8pt}

    \small \textcolor{gray}{When you are done with this task, mark it as complete to proceed to the next one.} \hfill \textbf{\textcolor{brandDark}{$\checkmark$ Completed}}

\end{taskbox}

\begin{taskbox}{Task 6: Present Your Recommendation}

    {\small \textcolor{gray}{2--3 hrs $\cdot$ Intermediate}}\\[12pt]

      
The firm’s leadership team needs your recommendation on whether logistics is a viable market for AI expansion. Prepare and pitch your proposal.  


\vspace{8pt}
{\large \textbf{\textcolor{brandLight}{What you'll learn}}}\\[4pt]  
\begin{itemize}
    \item Combine all your findings into a professional presentation.  
    \item Include key insights, visualizations, risks, and actionable recommendations.  
    \item Record a 5-minute video pitch summarizing your proposal.  
\end{itemize}

\vspace{8pt}
{\large \textbf{\textcolor{brandLight}{What you'll do}}}\\[4pt]  
Submit a presentation deck (PDF or PPT) and a video recording of your pitch.\par

    \vspace{16pt}

    \textcolor{cardBorder}{\hrule height 1pt}\vspace{8pt}

    \small \textcolor{gray}{When you are done with this task, mark it as complete to proceed to the next one.} \hfill \textbf{\textcolor{brandDark}{$\checkmark$ Completed}}

\end{taskbox}

\begin{completionbox}
    \textbf{Knowledge Assessment \& Verification:} Complete the online knowledge check, submit your code and presentation files for AI verification, and claim your \textbf{Skillzza Verified Certificate}.
\end{completionbox}

\newpage

% =============================================================================

% SIMULATION 3

% =============================================================================

\section*{GenAI Prompt Engineer – Enterprise Prompt Lab}

\addcontentsline{toc}{section}{GenAI Prompt Engineer – Enterprise Prompt Lab}

\begin{metabox}
    \begin{tabularx}{\textwidth}{@{}ll X@{}}
        \textbf{Industry:} & Technology & \textbf{Partner Enterprise:} \textbf{Skillzza Enterprise} \\
        \textbf{Duration:} & 12-18 hrs (Self-paced) & \textbf{Mode:} Individual, task-based simulation \\
        \textbf{Primary Skills:} & \multicolumn{2}{@{}l}{Python, AI, Data Analysis, Problem Solving}
    \end{tabularx}
\end{metabox}

\subsection*{Overview \& Problem Statement}

\#\#\# \textbf{Scenario}
In today’s enterprise landscape, organizations are leveraging Generative AI (GenAI) to automate tasks, enhance productivity, and generate structured outputs. As a \textbf{GenAI Prompt Engineer}, you will play a pivotal role in designing, testing, and optimizing AI prompts to ensure accurate, reliable, and structured responses tailored to business needs. The challenge lies in crafting prompts that minimize hallucinations (AI-generated inaccuracies), maximize productivity, and align with enterprise requirements.

\#\#\# \textbf{Mission}
Your mission is to design, evaluate, and refine prompts for a GenAI system that processes enterprise data to generate \textbf{structured outputs} like summaries, data classifications, or decision-support recommendations. You will optimize prompts to reduce hallucinations and ensure outputs are accurate, actionable, and contextually relevant.

\#\#\# \textbf{Final Challenge}
By the end of this simulation, you will deliver a polished \textbf{Prompt Engineering Playbook} for an enterprise GenAI model, including structured examples, evaluation metrics, and strategies for reducing hallucinations. Your playbook will be reviewed by industry experts to assess your readiness for real-world enterprise projects.

---

\subsection*{Tasks \& Deliverables}

\begin{taskbox}{Task 1: Understand the Role of Prompts}

    {\small \textcolor{gray}{2--3 hrs $\cdot$ Intermediate}}\\[12pt]

    \#\#\#\# \textbf{Scenario}
You have been hired by an enterprise that manages a large volume of customer service tickets. They want to use GenAI to classify tickets by issue type (e.g., billing, technical, feedback) and generate summaries for faster resolution.

\#\#\#\# \textbf{Student Assignment}
Research prompt engineering fundamentals. Learn how prompt structures (context, constraints, and examples) influence GenAI outputs. Study enterprise-centric use cases, focusing on structured data generation.

\#\#\#\# \textbf{Deliverable}
Submit a \textbf{1-page report} summarizing:
\begin{itemize}
    \item Key principles of prompt engineering.
    \item How structured prompts can benefit enterprise workflows.
    \item Examples of prompt patterns for structured outputs.\par
\end{itemize}

    \vspace{16pt}

    \textcolor{cardBorder}{\hrule height 1pt}\vspace{8pt}

    \small \textcolor{gray}{When you are done with this task, mark it as complete to proceed to the next one.} \hfill \textbf{\textcolor{brandDark}{$\checkmark$ Completed}}

\end{taskbox}

\begin{taskbox}{Task 2: Dataset Setup and Prompt Design}

    {\small \textcolor{gray}{2--3 hrs $\cdot$ Intermediate}}\\[12pt]

    \#\#\#\# \textbf{Scenario}
You’ve been provided with a sample dataset containing customer service tickets with columns like `Ticket ID`, `Issue Description`, and `Customer Priority`. Your goal is to design initial prompts to classify these tickets by issue type and summarize them.

\#\#\#\# \textbf{Student Assignment}
\begin{itemize}
    \item Analyze the provided dataset structure.
    \item Design 3 initial prompts:
  1. Classify tickets into predefined categories.
  2. Generate concise summaries of ticket issues.
  3. Identify high-priority customer concerns.
\end{itemize}
  
\#\#\#\# \textbf{Deliverable}
Submit a \textbf{document with your 3 prompts}, their intended outputs, and reasoning for their design.\par

    \vspace{16pt}

    \textcolor{cardBorder}{\hrule height 1pt}\vspace{8pt}

    \small \textcolor{gray}{When you are done with this task, mark it as complete to proceed to the next one.} \hfill \textbf{\textcolor{brandDark}{$\checkmark$ Completed}}

\end{taskbox}

\begin{taskbox}{Task 3: Build and Execute Prompts}

    {\small \textcolor{gray}{2--3 hrs $\cdot$ Intermediate}}\\[12pt]

    \#\#\#\# \textbf{Scenario}
Run your initial prompts on a GenAI platform (e.g., OpenAI GPT-4) and analyze the outputs. Refine your prompts iteratively to improve classification accuracy and summary quality.

\#\#\#\# \textbf{Student Assignment}
\begin{itemize}
    \item Test your prompts using the sample customer ticket dataset.
    \item Adjust context, constraints, or examples to improve outputs.
    \item Document the iterative refinements made to each prompt.
\end{itemize}

\#\#\#\# \textbf{Deliverable}
Submit:
1. \textbf{Initial outputs} from GenAI.
2. \textbf{Refined prompts} and their improved outputs.
3. A \textbf{brief explanation} of changes made and their impact.\par

    \vspace{16pt}

    \textcolor{cardBorder}{\hrule height 1pt}\vspace{8pt}

    \small \textcolor{gray}{When you are done with this task, mark it as complete to proceed to the next one.} \hfill \textbf{\textcolor{brandDark}{$\checkmark$ Completed}}

\end{taskbox}

\begin{taskbox}{Task 4: Analyze Hallucinations and Explain GenAI Behavior}

    {\small \textcolor{gray}{2--3 hrs $\cdot$ Intermediate}}\\[12pt]

    \#\#\#\# \textbf{Scenario}
Some GenAI outputs contain errors (e.g., misclassifications, irrelevant summaries). Your task is to identify and explain these hallucinations, then adjust your prompts to minimize such errors.

\#\#\#\# \textbf{Student Assignment}
\begin{itemize}
    \item Analyze common hallucinations in the outputs (e.g., incorrect ticket classifications).
    \item Adjust prompts to reduce errors (e.g., add constraints or improve examples).
    \item Explain the behavior behind the hallucinations (e.g., ambiguous wording, lack of context).
\end{itemize}

\#\#\#\# \textbf{Deliverable}
Submit a \textbf{hallucination analysis report} with:
\begin{itemize}
    \item Examples of errors.
    \item Adjusted prompts to reduce hallucinations.
    \item Explanations of GenAI behavior causing errors.\par
\end{itemize}

    \vspace{16pt}

    \textcolor{cardBorder}{\hrule height 1pt}\vspace{8pt}

    \small \textcolor{gray}{When you are done with this task, mark it as complete to proceed to the next one.} \hfill \textbf{\textcolor{brandDark}{$\checkmark$ Completed}}

\end{taskbox}

\begin{taskbox}{Task 5: Audit for Responsible AI}

    {\small \textcolor{gray}{2--3 hrs $\cdot$ Intermediate}}\\[12pt]

    \#\#\#\# \textbf{Scenario}
Ensure your prompts and outputs adhere to Responsible AI principles, including fairness, transparency, and accuracy. Avoid biases in ticket classifications and summaries.

\#\#\#\# \textbf{Student Assignment}
\begin{itemize}
    \item Perform a bias audit on GenAI outputs (e.g., verify fair classification across all categories).
    \item Document steps taken to ensure transparency and accuracy in prompts.
    \item Suggest improvements to align with Responsible AI principles.
\end{itemize}

\#\#\#\# \textbf{Deliverable}
Submit a \textbf{Responsible AI checklist} covering:
\begin{itemize}
    \item Bias detection results.
    \item Transparency steps taken.
    \item Suggestions for ethical improvements.\par
\end{itemize}

    \vspace{16pt}

    \textcolor{cardBorder}{\hrule height 1pt}\vspace{8pt}

    \small \textcolor{gray}{When you are done with this task, mark it as complete to proceed to the next one.} \hfill \textbf{\textcolor{brandDark}{$\checkmark$ Completed}}

\end{taskbox}

\begin{taskbox}{Task 6: Present Your Final Prompt Engineering Playbook}

    {\small \textcolor{gray}{2--3 hrs $\cdot$ Intermediate}}\\[12pt]

    \#\#\#\# \textbf{Scenario}
Compile your refined prompts, evaluation metrics, and Responsible AI findings into a \textbf{Prompt Engineering Playbook}. This will serve as a reusable resource for the enterprise.

\#\#\#\# \textbf{Student Assignment}
\begin{itemize}
    \item Summarize your prompt design process and key insights.
    \item Include examples of refined prompts and outputs.
    \item Provide actionable strategies for minimizing hallucinations and ensuring Responsible AI compliance.
\end{itemize}

\#\#\#\# \textbf{Deliverable}
Submit your \textbf{Prompt Engineering Playbook} (5-10 pages) with:
\begin{itemize}
    \item Overview of prompt design principles.
    \item Tested and refined prompts.
    \item Evaluation metrics and Responsible AI strategies.\par
\end{itemize}

    \vspace{16pt}

    \textcolor{cardBorder}{\hrule height 1pt}\vspace{8pt}

    \small \textcolor{gray}{When you are done with this task, mark it as complete to proceed to the next one.} \hfill \textbf{\textcolor{brandDark}{$\checkmark$ Completed}}

\end{taskbox}

\begin{completionbox}
    \textbf{Knowledge Assessment \& Verification:} Complete the online knowledge check, submit your code and presentation files for AI verification, and claim your \textbf{Skillzza Verified Certificate}.
\end{completionbox}

\newpage

% =============================================================================

% SIMULATION 4

% =============================================================================

\section*{Virtual Internship: Data Analyst – Turn Sales Data into Decisions}

\addcontentsline{toc}{section}{Virtual Internship: Data Analyst – Turn Sales Data into Decisions}

\begin{metabox}
    \begin{tabularx}{\textwidth}{@{}ll X@{}}
        \textbf{Industry:} & Technology & \textbf{Partner Enterprise:} \textbf{Skillzza Enterprise} \\
        \textbf{Duration:} & 12-18 hrs (Self-paced) & \textbf{Mode:} Individual, task-based simulation \\
        \textbf{Primary Skills:} & \multicolumn{2}{@{}l}{Python, AI, Data Analysis, Problem Solving}
    \end{tabularx}
\end{metabox}

\subsection*{Overview \& Problem Statement}

\#\#\# Scenario  
You’ve just been hired as a \textbf{Data Analyst} for a retail company that specializes in selling consumer electronics. The company is grappling with growing sales data across multiple regions and product categories but lacks insights into key performance metrics like revenue growth, customer segmentation, and product performance. Your role is to analyze raw sales data, clean it, derive actionable insights, and present recommendations to the management team for improving sales strategies.

\#\#\# Mission  
Your mission is to transform messy, unstructured sales data into clean datasets, calculate KPIs (key performance indicators) to measure business performance, and apply advanced segmentation techniques to identify high-value customers. By the end of this internship, you will create a report with clear recommendations that the management team can use to make data-driven decisions.

\#\#\# Final Challenge  
In the final task, you will compile your findings into a presentation for the executive team. You will use both visualizations and narrative storytelling to explain sales trends, customer behavior, and actionable recommendations.

---

\subsection*{Tasks \& Deliverables}

\begin{taskbox}{Task 1: Understand the Sales Dataset}

    {\small \textcolor{gray}{2--3 hrs $\cdot$ Intermediate}}\\[12pt]

    \#\#\#\# Scenario  
You’ve been handed a raw sales dataset from the company’s CRM system. The dataset contains fields such as Order ID, Product Category, Purchase Date, Customer ID, Region, Quantity Sold, Revenue, and Discount Applied. Unfortunately, the data has missing values and inconsistencies.  

\#\#\#\# Student Assignment  
\begin{itemize}
    \item Explore the raw dataset and identify key fields relevant for analysis.  
    \item Document challenges in the dataset, such as missing values, duplicate entries, and inconsistent formatting.  
\end{itemize}

\#\#\#\# Deliverable  
\begin{itemize}
    \item Submit your exploratory data analysis summary (a brief document or spreadsheet) highlighting key patterns and challenges in the dataset.\par
\end{itemize}

    \vspace{16pt}

    \textcolor{cardBorder}{\hrule height 1pt}\vspace{8pt}

    \small \textcolor{gray}{When you are done with this task, mark it as complete to proceed to the next one.} \hfill \textbf{\textcolor{brandDark}{$\checkmark$ Completed}}

\end{taskbox}

\begin{taskbox}{Task 2: Data Cleaning and Setup}

    {\small \textcolor{gray}{2--3 hrs $\cdot$ Intermediate}}\\[12pt]

    \#\#\#\# Scenario  
Before diving into deeper analysis, you need to clean the data to ensure accuracy. The management relies on precise insights, and errors in the dataset could mislead their decision-making.  

\#\#\#\# Student Assignment  
\begin{itemize}
    \item Remove duplicate entries and handle missing values.  
    \item Standardize column names and formats (e.g., date formats, currency).  
    \item Validate the cleaned dataset by spot-checking for errors.  
\end{itemize}

\#\#\#\# Deliverable  
\begin{itemize}
    \item Submit a cleaned dataset (CSV file) with a brief document explaining the cleaning steps and decisions.\par
\end{itemize}

    \vspace{16pt}

    \textcolor{cardBorder}{\hrule height 1pt}\vspace{8pt}

    \small \textcolor{gray}{When you are done with this task, mark it as complete to proceed to the next one.} \hfill \textbf{\textcolor{brandDark}{$\checkmark$ Completed}}

\end{taskbox}

\begin{taskbox}{Task 3: Calculate Key Performance Indicators (KPIs)}

    {\small \textcolor{gray}{2--3 hrs $\cdot$ Intermediate}}\\[12pt]

    \#\#\#\# Scenario  
Management has requested a summary of sales performance based on key metrics. They need insights into total revenue, profit margin, average order value (AOV), and customer retention rate to evaluate current strategies.  

\#\#\#\# Student Assignment  
\begin{itemize}
    \item Use Python or Excel to calculate:  
    \item Total Revenue  
    \item Average Order Value (Revenue / Number of Orders)  
    \item Monthly Revenue Growth (\%)  
    \item Customer Retention Rate (returning customers over total customers)  
    \item Visualize the trends using line charts or bar charts.  
\end{itemize}

\#\#\#\# Deliverable  
\begin{itemize}
    \item Submit a spreadsheet or Python script with calculated KPIs.  
    \item Include visualizations showing trends over time.\par
\end{itemize}

    \vspace{16pt}

    \textcolor{cardBorder}{\hrule height 1pt}\vspace{8pt}

    \small \textcolor{gray}{When you are done with this task, mark it as complete to proceed to the next one.} \hfill \textbf{\textcolor{brandDark}{$\checkmark$ Completed}}

\end{taskbox}

\begin{taskbox}{Task 4: Segment Customers Using Behavioral Data}

    {\small \textcolor{gray}{2--3 hrs $\cdot$ Intermediate}}\\[12pt]

    \#\#\#\# Scenario  
The marketing team wants to launch targeted campaigns but doesn’t know which customers to prioritize. You need to segment the customers based on their purchase behavior, such as frequency (number of purchases), recency (time since last purchase), and monetary value (total revenue contributed).  

\#\#\#\# Student Assignment  
\begin{itemize}
    \item Apply RFM (Recency, Frequency, Monetary) analysis using Python or Excel.  
    \item Create 3-5 customer segments (e.g., High-value customers, New customers, At-risk customers).  
    \item Visualize the segments using scatter plots or bar charts.  
\end{itemize}

\#\#\#\# Deliverable  
\begin{itemize}
    \item Submit the segmented dataset (CSV file) with clear labels for each segment.  
    \item Include a chart summarizing the segmentation insights.\par
\end{itemize}

    \vspace{16pt}

    \textcolor{cardBorder}{\hrule height 1pt}\vspace{8pt}

    \small \textcolor{gray}{When you are done with this task, mark it as complete to proceed to the next one.} \hfill \textbf{\textcolor{brandDark}{$\checkmark$ Completed}}

\end{taskbox}

\begin{taskbox}{Task 5: Evaluate Recommendations Using Responsible Analytics}

    {\small \textcolor{gray}{2--3 hrs $\cdot$ Intermediate}}\\[12pt]

    \#\#\#\# Scenario  
The management team wants to ensure that your recommendations are fair and ethical. They are concerned about biases in the segmentation or KPI analysis that may disproportionately affect certain regions or customer groups.  

\#\#\#\# Student Assignment  
\begin{itemize}
    \item Audit your analysis for potential biases (e.g., over-prioritizing one region or product category).  
    \item Suggest adjustments to the segmentation or KPI calculations to ensure fairness.  
\end{itemize}

\#\#\#\# Deliverable  
\begin{itemize}
    \item Submit an audit report highlighting potential biases and steps taken to address them.\par
\end{itemize}

    \vspace{16pt}

    \textcolor{cardBorder}{\hrule height 1pt}\vspace{8pt}

    \small \textcolor{gray}{When you are done with this task, mark it as complete to proceed to the next one.} \hfill \textbf{\textcolor{brandDark}{$\checkmark$ Completed}}

\end{taskbox}

\begin{taskbox}{Task 6: Present Recommendations to Management}

    {\small \textcolor{gray}{2--3 hrs $\cdot$ Intermediate}}\\[12pt]

    \#\#\#\# Scenario  
You are required to compile your findings and present actionable recommendations to the executive team. Your presentation should include visualizations, key insights, and suggested strategies for improving sales.  

\#\#\#\# Student Assignment  
\begin{itemize}
    \item Create a slide deck summarizing:  
    \item Key findings from the KPI analysis.  
    \item Customer segmentation results.  
    \item Recommendations for improving sales strategies (e.g., targeted campaigns, product focus, regional strategies).  
    \item Practice delivering a concise, data-driven presentation.  
\end{itemize}

\#\#\#\# Deliverable  
\begin{itemize}
    \item Submit your slide deck (PDF or PowerPoint).  
    \item Record a short video (3–5 minutes) presenting your findings.\par
\end{itemize}

    \vspace{16pt}

    \textcolor{cardBorder}{\hrule height 1pt}\vspace{8pt}

    \small \textcolor{gray}{When you are done with this task, mark it as complete to proceed to the next one.} \hfill \textbf{\textcolor{brandDark}{$\checkmark$ Completed}}

\end{taskbox}

\begin{completionbox}
    \textbf{Knowledge Assessment \& Verification:} Complete the online knowledge check, submit your code and presentation files for AI verification, and claim your \textbf{Skillzza Verified Certificate}.
\end{completionbox}

\newpage

% =============================================================================

% SIMULATION 5

% =============================================================================

\section*{Data Visualization Specialist – Tell a Story with Data}

\addcontentsline{toc}{section}{Data Visualization Specialist – Tell a Story with Data}

\begin{metabox}
    \begin{tabularx}{\textwidth}{@{}ll X@{}}
        \textbf{Industry:} & Technology & \textbf{Partner Enterprise:} \textbf{Skillzza Enterprise} \\
        \textbf{Duration:} & 12-18 hrs (Self-paced) & \textbf{Mode:} Individual, task-based simulation \\
        \textbf{Primary Skills:} & \multicolumn{2}{@{}l}{Python, AI, Data Analysis, Problem Solving}
    \end{tabularx}
\end{metabox}

\subsection*{Overview \& Problem Statement}

\#\#\# \textbf{Scenario}
In today's data-driven world, raw data often fails to communicate insights effectively to decision-makers. As a Data Visualization Specialist, your job is to transform complex datasets into compelling visuals that tell a clear and actionable story. Whether you're presenting to a business client, a management team, or a public audience, your ability to analyze data and craft impactful visual narratives can drive meaningful decisions.

\#\#\# \textbf{Mission}
Your mission in this virtual internship is to take on the role of a Data Visualization Specialist for a fictional retail company, \textbf{TrendMart}, which is struggling to understand customer purchasing patterns. You will analyze a dataset containing 10,000+ rows of transaction data, identify key metrics, design a visually engaging dashboard, and present your findings in a story format to stakeholders.

\#\#\# \textbf{Final Challenge}
By the end of this simulation, you will present a fully interactive dashboard that visually explains customer behavior trends, product sales performance, and seasonal purchase patterns. You will also deliver a concise narrative explaining how your visualizations support decision-making.

---

\subsection*{Tasks \& Deliverables}

\begin{taskbox}{Task 1: Understand}

    {\small \textcolor{gray}{2--3 hrs $\cdot$ Intermediate}}\\[12pt]

    \#\#\#\# Scenario:
TrendMart has provided you with a dataset containing transaction-level details, customer demographics, and product categories. The management team is unsure which metrics to focus on and needs your help to identify actionable insights.

\#\#\#\# Student Assignment:
\begin{itemize}
    \item Review the dataset fields, which include:
    \item `Transaction\_ID`, `Customer\_ID`, `Age`, `Gender`, `Region`, `Product\_Category`, `Purchase\_Amount`, `Purchase\_Date`.
    \item Identify potential KPIs such as total revenue, average order value (AOV), product popularity, and seasonal trends.
\end{itemize}

\#\#\#\# Deliverable:
Submit a one-page document outlining:
\begin{itemize}
    \item Key KPIs to focus on.
    \item The rationale behind your selection.
    \item Any assumptions made about the dataset.\par
\end{itemize}

    \vspace{16pt}

    \textcolor{cardBorder}{\hrule height 1pt}\vspace{8pt}

    \small \textcolor{gray}{When you are done with this task, mark it as complete to proceed to the next one.} \hfill \textbf{\textcolor{brandDark}{$\checkmark$ Completed}}

\end{taskbox}

\begin{taskbox}{Task 2: Data/Setup}

    {\small \textcolor{gray}{2--3 hrs $\cdot$ Intermediate}}\\[12pt]

    \#\#\#\# Scenario:
TrendMart’s dataset has inconsistencies, including missing values and duplicate entries. Before creating visuals, you need to clean and structure the data to ensure reliable analysis.

\#\#\#\# Student Assignment:
\begin{itemize}
    \item Perform data cleaning:
    \item Remove duplicates.
    \item Handle missing values (e.g., imputing median values for `Age`).
    \item Format `Purchase\_Date` into a standard date format.
    \item Create meaningful groupings (e.g., aggregate revenue by product categories and regions).
\end{itemize}

\#\#\#\# Deliverable:
Submit the cleaned dataset (in Excel or CSV format) and a summary of the cleaning steps performed.\par

    \vspace{16pt}

    \textcolor{cardBorder}{\hrule height 1pt}\vspace{8pt}

    \small \textcolor{gray}{When you are done with this task, mark it as complete to proceed to the next one.} \hfill \textbf{\textcolor{brandDark}{$\checkmark$ Completed}}

\end{taskbox}

\begin{taskbox}{Task 3: Build/Execute}

    {\small \textcolor{gray}{2--3 hrs $\cdot$ Intermediate}}\\[12pt]

    \#\#\#\# Scenario:
The cleaned dataset is ready, and you need to choose appropriate visualizations to represent the KPIs identified earlier. TrendMart’s management requests a dashboard that can display:
\begin{itemize}
    \item Regional revenue distribution.
    \item Top 5 product categories by sales.
    \item Seasonal revenue trends (monthly).
\end{itemize}

\#\#\#\# Student Assignment:
\begin{itemize}
    \item Design an interactive dashboard using Tableau Public or Power BI.
    \item Ensure the dashboard includes:
    \item A bar chart for top product categories.
    \item A line chart for monthly trends.
    \item A heatmap for regional revenue distribution.
\end{itemize}

\#\#\#\# Deliverable:
Submit the dashboard file and screenshots of key charts, along with a brief explanation of your design choices.\par

    \vspace{16pt}

    \textcolor{cardBorder}{\hrule height 1pt}\vspace{8pt}

    \small \textcolor{gray}{When you are done with this task, mark it as complete to proceed to the next one.} \hfill \textbf{\textcolor{brandDark}{$\checkmark$ Completed}}

\end{taskbox}

\begin{taskbox}{Task 4: GenAI/Explanation}

    {\small \textcolor{gray}{2--3 hrs $\cdot$ Intermediate}}\\[12pt]

    \#\#\#\# Scenario:
TrendMart’s leadership team wants an automated summary of customer purchasing patterns to save time during presentations. Use generative AI tools to generate a narrative based on your dashboard insights.

\#\#\#\# Student Assignment:
\begin{itemize}
    \item Input the cleaned dataset or dashboard insights into an AI tool (e.g., ChatGPT or Tableau AI).
    \item Generate a text-based summary, including:
    \item Key trends (e.g., “Revenue peaked during December due to holiday sales”).
    \item Anomalies (e.g., “Product Category X showed unusually high sales in Region Y”).
\end{itemize}

\#\#\#\# Deliverable:
Submit the AI-generated narrative and refine it into a concise 300-word summary suitable for a presentation.\par

    \vspace{16pt}

    \textcolor{cardBorder}{\hrule height 1pt}\vspace{8pt}

    \small \textcolor{gray}{When you are done with this task, mark it as complete to proceed to the next one.} \hfill \textbf{\textcolor{brandDark}{$\checkmark$ Completed}}

\end{taskbox}

\begin{taskbox}{Task 5: Audit/Responsible AI}

    {\small \textcolor{gray}{2--3 hrs $\cdot$ Intermediate}}\\[12pt]

    \#\#\#\# Scenario:
Data visualizations can unintentionally mislead viewers if not presented responsibly. Your task is to audit your dashboard for potential bias or misrepresentation.

\#\#\#\# Student Assignment:
\begin{itemize}
    \item Check for:
    \item Cherry-picking (e.g., excluding low-performing regions or products).
    \item Misleading scales or axes in charts.
    \item Overcomplicated visuals that may confuse viewers.
    \item Make corrections where necessary.
\end{itemize}

\#\#\#\# Deliverable:
Submit a 1-page audit report explaining:
\begin{itemize}
    \item Potential biases identified.
    \item Changes made to improve ethical visualization.\par
\end{itemize}

    \vspace{16pt}

    \textcolor{cardBorder}{\hrule height 1pt}\vspace{8pt}

    \small \textcolor{gray}{When you are done with this task, mark it as complete to proceed to the next one.} \hfill \textbf{\textcolor{brandDark}{$\checkmark$ Completed}}

\end{taskbox}

\begin{taskbox}{Task 6: Present Recommendation}

    {\small \textcolor{gray}{2--3 hrs $\cdot$ Intermediate}}\\[12pt]

    \#\#\#\# Scenario:
TrendMart’s management team is ready to see your findings. Present your dashboard and explain the insights through storytelling.

\#\#\#\# Student Assignment:
\begin{itemize}
    \item Prepare a 5-minute video or slide deck presentation that:
    \item Highlights the key KPIs and trends.
    \item Explains the visuals in the dashboard.
    \item Recommends actionable business decisions (e.g., “Focus marketing efforts on Region X during Q4”).
\end{itemize}

\#\#\#\# Deliverable:
Submit the presentation file (video link or slide deck) and a short summary of your recommendations.\par

    \vspace{16pt}

    \textcolor{cardBorder}{\hrule height 1pt}\vspace{8pt}

    \small \textcolor{gray}{When you are done with this task, mark it as complete to proceed to the next one.} \hfill \textbf{\textcolor{brandDark}{$\checkmark$ Completed}}

\end{taskbox}

\begin{completionbox}
    \textbf{Knowledge Assessment \& Verification:} Complete the online knowledge check, submit your code and presentation files for AI verification, and claim your \textbf{Skillzza Verified Certificate}.
\end{completionbox}

\newpage

% =============================================================================

% SIMULATION 6

% =============================================================================

\section*{Predictive Analytics Analyst – Forecast Business Demand}

\addcontentsline{toc}{section}{Predictive Analytics Analyst – Forecast Business Demand}

\begin{metabox}
    \begin{tabularx}{\textwidth}{@{}ll X@{}}
        \textbf{Industry:} & Technology & \textbf{Partner Enterprise:} \textbf{Skillzza Enterprise} \\
        \textbf{Duration:} & 12-18 hrs (Self-paced) & \textbf{Mode:} Individual, task-based simulation \\
        \textbf{Primary Skills:} & \multicolumn{2}{@{}l}{Python, AI, Data Analysis, Problem Solving}
    \end{tabularx}
\end{metabox}

\subsection*{Overview \& Problem Statement}

\textbf{Scenario:}  
You are hired as a Predictive Analytics Analyst by a retail company struggling with demand forecasting. They often face stockouts or overstock issues due to poor planning, which affects both revenue and customer satisfaction. Your role is to analyze historical sales data, identify demand patterns, and build a predictive model for accurate demand forecasting.

\textbf{Mission:}  
Your mission is to leverage predictive analytics and machine learning techniques to forecast business demand for the next quarter. By identifying trends and contributing actionable insights, you will help the company optimize inventory levels, improve customer experience, and reduce costs.

\textbf{Final Challenge:}  
Deliver a business-ready solution that includes a demand forecast model, its evaluation metrics, and a presentation of actionable recommendations for the company's inventory planning team.

---

\subsection*{Tasks \& Deliverables}

\begin{taskbox}{Task 1: Understand the Business Context}

    {\small \textcolor{gray}{2--3 hrs $\cdot$ Intermediate}}\\[12pt]

      
You are given historical sales data from the past three years for a retail company. This dataset includes daily sales, product categories, promotions, and external factors like holidays. Your first step is to understand the industry-specific factors that influence demand and define the forecasting problem clearly.  


\vspace{8pt}
{\large \textbf{\textcolor{brandLight}{What you'll learn}}}\\[4pt]  
\begin{itemize}
    \item Explore the dataset fields such as `date`, `sales`, `product\_category`, `promotion`, `holiday\_flag`.
    \item Identify key demand influencers like seasonality, trends, and external events.  
    \item Create a written summary of the forecasting problem.  
\end{itemize}

\vspace{8pt}
{\large \textbf{\textcolor{brandLight}{What you'll do}}}\\[4pt]  
Submit a one-page document summarizing the dataset structure, key factors affecting demand, and the problem statement.\par

    \vspace{16pt}

    \textcolor{cardBorder}{\hrule height 1pt}\vspace{8pt}

    \small \textcolor{gray}{When you are done with this task, mark it as complete to proceed to the next one.} \hfill \textbf{\textcolor{brandDark}{$\checkmark$ Completed}}

\end{taskbox}

\begin{taskbox}{Task 2: Data Preparation and Setup}

    {\small \textcolor{gray}{2--3 hrs $\cdot$ Intermediate}}\\[12pt]

      
Raw data often has missing values, outliers, or inconsistencies that can affect model performance. Your next step is to clean and preprocess the data for analysis.  


\vspace{8pt}
{\large \textbf{\textcolor{brandLight}{What you'll learn}}}\\[4pt]  
\begin{itemize}
    \item Handle missing values in `sales` using interpolation or other techniques.  
    \item Remove or impute outliers in the dataset.  
    \item Engineer new features like `weekday`, `month`, and `season` from the `date` field.  
    \item Split the data into training and test sets (80\%-20\%).  
\end{itemize}

\vspace{8pt}
{\large \textbf{\textcolor{brandLight}{What you'll do}}}\\[4pt]  
Submit a clean dataset (.csv) with engineered features and a short report explaining the preprocessing steps.\par

    \vspace{16pt}

    \textcolor{cardBorder}{\hrule height 1pt}\vspace{8pt}

    \small \textcolor{gray}{When you are done with this task, mark it as complete to proceed to the next one.} \hfill \textbf{\textcolor{brandDark}{$\checkmark$ Completed}}

\end{taskbox}

\begin{taskbox}{Task 3: Build and Execute Forecasting Models}

    {\small \textcolor{gray}{2--3 hrs $\cdot$ Intermediate}}\\[12pt]

      
You will now build predictive models for demand forecasting using machine learning algorithms. Compare traditional methods like ARIMA with advanced techniques like XGBoost or Random Forest.  


\vspace{8pt}
{\large \textbf{\textcolor{brandLight}{What you'll learn}}}\\[4pt]  
\begin{itemize}
    \item Train three different forecasting models: ARIMA, Random Forest, and XGBoost.  
    \item Evaluate each model using metrics like RMSE, MAE, and MAPE.  
    \item Document the hyperparameter settings for each model.  
\end{itemize}

\vspace{8pt}
{\large \textbf{\textcolor{brandLight}{What you'll do}}}\\[4pt]  
Submit the trained model files and a performance comparison table of the three models.\par

    \vspace{16pt}

    \textcolor{cardBorder}{\hrule height 1pt}\vspace{8pt}

    \small \textcolor{gray}{When you are done with this task, mark it as complete to proceed to the next one.} \hfill \textbf{\textcolor{brandDark}{$\checkmark$ Completed}}

\end{taskbox}

\begin{taskbox}{Task 4: Explain Model Decisions with GenAI}

    {\small \textcolor{gray}{2--3 hrs $\cdot$ Intermediate}}\\[12pt]

      
Management needs to understand how your model predicts demand. Use generative AI tools to summarize insights from feature importance and model outputs.  


\vspace{8pt}
{\large \textbf{\textcolor{brandLight}{What you'll learn}}}\\[4pt]  
\begin{itemize}
    \item Use tools like ChatGPT or LIME to generate explainable insights from the model outputs.  
    \item Identify key features driving demand predictions (e.g., holidays, promotions).  
    \item Create a summary of insights for non-technical stakeholders.  
\end{itemize}

\vspace{8pt}
{\large \textbf{\textcolor{brandLight}{What you'll do}}}\\[4pt]  
Submit a one-page document explaining model predictions in layman’s terms, including visualizations.\par

    \vspace{16pt}

    \textcolor{cardBorder}{\hrule height 1pt}\vspace{8pt}

    \small \textcolor{gray}{When you are done with this task, mark it as complete to proceed to the next one.} \hfill \textbf{\textcolor{brandDark}{$\checkmark$ Completed}}

\end{taskbox}

\begin{taskbox}{Task 5: Audit for Responsible AI}

    {\small \textcolor{gray}{2--3 hrs $\cdot$ Intermediate}}\\[12pt]

      
Your forecasting model must be unbiased and fair. Conduct a Responsible AI audit to ensure that predictions are not skewed by factors like product category or promotion bias.  


\vspace{8pt}
{\large \textbf{\textcolor{brandLight}{What you'll learn}}}\\[4pt]  
\begin{itemize}
    \item Analyze prediction bias across different product categories.  
    \item Check if external factors like promotions disproportionately affect predictions.  
    \item Suggest steps for mitigating bias (e.g., reweighting training samples).  
\end{itemize}

\vspace{8pt}
{\large \textbf{\textcolor{brandLight}{What you'll do}}}\\[4pt]  
Submit a Responsible AI audit report highlighting bias findings and mitigation strategies.\par

    \vspace{16pt}

    \textcolor{cardBorder}{\hrule height 1pt}\vspace{8pt}

    \small \textcolor{gray}{When you are done with this task, mark it as complete to proceed to the next one.} \hfill \textbf{\textcolor{brandDark}{$\checkmark$ Completed}}

\end{taskbox}

\begin{taskbox}{Task 6: Present Forecast and Business Recommendations}

    {\small \textcolor{gray}{2--3 hrs $\cdot$ Intermediate}}\\[12pt]

      
You are ready to present your findings and recommendations to the inventory planning team. Prepare actionable insights from your demand forecast to drive business decisions.  


\vspace{8pt}
{\large \textbf{\textcolor{brandLight}{What you'll learn}}}\\[4pt]  
\begin{itemize}
    \item Create a dashboard or presentation summarizing demand forecasts per product category and time period.  
    \item Recommend inventory adjustments based on forecasted demand.  
    \item Include an executive summary highlighting cost-saving opportunities.  
\end{itemize}

\vspace{8pt}
{\large \textbf{\textcolor{brandLight}{What you'll do}}}\\[4pt]  
Submit a presentation (PDF or PPT) and an optional Tableau dashboard file.\par

    \vspace{16pt}

    \textcolor{cardBorder}{\hrule height 1pt}\vspace{8pt}

    \small \textcolor{gray}{When you are done with this task, mark it as complete to proceed to the next one.} \hfill \textbf{\textcolor{brandDark}{$\checkmark$ Completed}}

\end{taskbox}

\begin{completionbox}
    \textbf{Knowledge Assessment \& Verification:} Complete the online knowledge check, submit your code and presentation files for AI verification, and claim your \textbf{Skillzza Verified Certificate}.
\end{completionbox}

\newpage

% =============================================================================

% SIMULATION 7

% =============================================================================

\section*{Data Scientist – Customer Churn Prediction}

\addcontentsline{toc}{section}{Data Scientist – Customer Churn Prediction}

\begin{metabox}
    \begin{tabularx}{\textwidth}{@{}ll X@{}}
        \textbf{Industry:} & Technology & \textbf{Partner Enterprise:} \textbf{Skillzza Enterprise} \\
        \textbf{Duration:} & 12-18 hrs (Self-paced) & \textbf{Mode:} Individual, task-based simulation \\
        \textbf{Primary Skills:} & \multicolumn{2}{@{}l}{Python, AI, Data Analysis, Problem Solving}
    \end{tabularx}
\end{metabox}

\subsection*{Overview \& Problem Statement}

\textbf{Scenario:}  
Customer churn is one of the most pressing challenges faced by subscription-based businesses, SaaS companies, and service providers. Losing customers not only impacts revenue but also increases acquisition costs for new customers to replace them. As a Data Scientist specializing in churn prediction, you will help businesses identify customers at risk of leaving and provide actionable insights to improve retention.

\textbf{Mission:}  
Your mission is to develop a predictive model that identifies customers likely to churn using a dataset containing customer demographics, usage patterns, and service feedback. You will explore churn drivers, validate predictions, and propose a retention strategy based on actionable insights.

\textbf{Final Challenge:}  
Deliver a detailed analytical report including your churn prediction model, key drivers analysis, and a retention strategy that a business can implement to reduce churn rates by at least 15\%.

---

\subsection*{Tasks \& Deliverables}

\begin{taskbox}{Task 1: Understand the Problem}

    {\small \textcolor{gray}{2--3 hrs $\cdot$ Intermediate}}\\[12pt]

      
You are hired by a subscription-based video streaming company experiencing increased customer churn. Your manager has provided a dataset containing customer profiles, subscription details, usage patterns, and churn labels. You need to start by understanding the problem and defining the business objectives.


\vspace{8pt}
{\large \textbf{\textcolor{brandLight}{What you'll learn}}}\\[4pt]  
\begin{itemize}
    \item Define the business problem: What does churn mean in this context, and why is it important?  
    \item Explore the dataset structure: Identify key features, null values, and overall data quality.  
    \item Document initial hypotheses about churn drivers based on domain knowledge (e.g., high service usage might reduce churn).  
\end{itemize}

\vspace{8pt}
{\large \textbf{\textcolor{brandLight}{What you'll do}}}\\[4pt]  
Submit a 1-2 page report explaining the business problem, dataset overview, and initial hypotheses about churn factors.\par

    \vspace{16pt}

    \textcolor{cardBorder}{\hrule height 1pt}\vspace{8pt}

    \small \textcolor{gray}{When you are done with this task, mark it as complete to proceed to the next one.} \hfill \textbf{\textcolor{brandDark}{$\checkmark$ Completed}}

\end{taskbox}

\begin{taskbox}{Task 2: Data Preparation and Exploratory Data Analysis}

    {\small \textcolor{gray}{2--3 hrs $\cdot$ Intermediate}}\\[12pt]

      
Before building the model, you need to clean the dataset and perform exploratory data analysis to uncover trends and outliers.


\vspace{8pt}
{\large \textbf{\textcolor{brandLight}{What you'll learn}}}\\[4pt]  
1. Handle missing values (imputation or removal).  
2. Perform feature engineering to create new variables such as "average monthly usage."  
3. Conduct exploratory data analysis (EDA) using statistical summaries and visualizations (e.g., histograms, box plots, correlation heatmaps).  

\vspace{8pt}
{\large \textbf{\textcolor{brandLight}{What you'll do}}}\\[4pt]  
Submit a Jupyter Notebook containing cleaned data, visualizations, and a summary of insights from EDA.\par

    \vspace{16pt}

    \textcolor{cardBorder}{\hrule height 1pt}\vspace{8pt}

    \small \textcolor{gray}{When you are done with this task, mark it as complete to proceed to the next one.} \hfill \textbf{\textcolor{brandDark}{$\checkmark$ Completed}}

\end{taskbox}

\begin{taskbox}{Task 3: Build Churn Prediction Model}

    {\small \textcolor{gray}{2--3 hrs $\cdot$ Intermediate}}\\[12pt]

      
Now that the data is clean, you need to build a machine learning classification model to predict churn. You aim to optimize predictive accuracy while balancing bias and variance.


\vspace{8pt}
{\large \textbf{\textcolor{brandLight}{What you'll learn}}}\\[4pt]  
1. Split the dataset into training and testing subsets.  
2. Train models like logistic regression, random forest, and XGBoost.  
3. Evaluate model performance using metrics such as accuracy, precision, recall, and F1-score.  
4. Select the best-performing model and justify your choice.  

\vspace{8pt}
{\large \textbf{\textcolor{brandLight}{What you'll do}}}\\[4pt]  
Submit a Jupyter Notebook detailing the models, evaluation metrics, and your selection rationale. Include visualizations of model performance (e.g., ROC curves).\par

    \vspace{16pt}

    \textcolor{cardBorder}{\hrule height 1pt}\vspace{8pt}

    \small \textcolor{gray}{When you are done with this task, mark it as complete to proceed to the next one.} \hfill \textbf{\textcolor{brandDark}{$\checkmark$ Completed}}

\end{taskbox}

\begin{taskbox}{Task 4: Explain Predictions Using GenAI Tools}

    {\small \textcolor{gray}{2--3 hrs $\cdot$ Intermediate}}\\[12pt]

      
Your manager wants to understand why certain customers are predicted to churn. Use SHAP to explain the predictions and identify key drivers of churn.  


\vspace{8pt}
{\large \textbf{\textcolor{brandLight}{What you'll learn}}}\\[4pt]  
1. Apply SHAP to analyze feature importance for individual predictions.  
2. Create visualizations to highlight the most significant features influencing churn.  
3. Write a brief explanation of the results and insights.  

\vspace{8pt}
{\large \textbf{\textcolor{brandLight}{What you'll do}}}\\[4pt]  
Submit SHAP visualizations and a 1-page report summarizing the key churn drivers revealed by your analysis.\par

    \vspace{16pt}

    \textcolor{cardBorder}{\hrule height 1pt}\vspace{8pt}

    \small \textcolor{gray}{When you are done with this task, mark it as complete to proceed to the next one.} \hfill \textbf{\textcolor{brandDark}{$\checkmark$ Completed}}

\end{taskbox}

\begin{taskbox}{Task 5: Audit Model for Responsible AI}

    {\small \textcolor{gray}{2--3 hrs $\cdot$ Intermediate}}\\[12pt]

      
Your client wants assurance that your model is fair and unbiased. Audit the model for potential bias based on customer demographics (e.g., age, gender, location).


\vspace{8pt}
{\large \textbf{\textcolor{brandLight}{What you'll learn}}}\\[4pt]  
1. Analyze discrepancies in prediction rates across demographic groups.  
2. Use fairness metrics like demographic parity and equalized odds to evaluate bias.  
3. Adjust the model or provide recommendations to mitigate bias if present.  

\vspace{8pt}
{\large \textbf{\textcolor{brandLight}{What you'll do}}}\\[4pt]  
Submit a Jupyter Notebook with fairness analysis results and a 1-page report detailing your findings and recommendations.\par

    \vspace{16pt}

    \textcolor{cardBorder}{\hrule height 1pt}\vspace{8pt}

    \small \textcolor{gray}{When you are done with this task, mark it as complete to proceed to the next one.} \hfill \textbf{\textcolor{brandDark}{$\checkmark$ Completed}}

\end{taskbox}

\begin{taskbox}{Task 6: Present Retention Strategy}

    {\small \textcolor{gray}{2--3 hrs $\cdot$ Intermediate}}\\[12pt]

      
Your manager is impressed with your model and insights but wants a clear retention strategy based on your findings. Propose actionable steps the business can take to reduce churn.


\vspace{8pt}
{\large \textbf{\textcolor{brandLight}{What you'll learn}}}\\[4pt]  
1. Summarize key churn drivers and customer segments most at risk.  
2. Recommend retention strategies (e.g., personalized offers, improved customer support, loyalty programs).  
3. Create a presentation slide deck to communicate findings and strategies effectively.  

\vspace{8pt}
{\large \textbf{\textcolor{brandLight}{What you'll do}}}\\[4pt]  
Submit a 5-slide presentation summarizing your analytics, insights, and retention strategy.\par

    \vspace{16pt}

    \textcolor{cardBorder}{\hrule height 1pt}\vspace{8pt}

    \small \textcolor{gray}{When you are done with this task, mark it as complete to proceed to the next one.} \hfill \textbf{\textcolor{brandDark}{$\checkmark$ Completed}}

\end{taskbox}

\begin{completionbox}
    \textbf{Knowledge Assessment \& Verification:} Complete the online knowledge check, submit your code and presentation files for AI verification, and claim your \textbf{Skillzza Verified Certificate}.
\end{completionbox}

\newpage

% =============================================================================

% SIMULATION 8

% =============================================================================

\section*{Fraud Analytics Analyst – Detect Suspicious Transactions}

\addcontentsline{toc}{section}{Fraud Analytics Analyst – Detect Suspicious Transactions}

\begin{metabox}
    \begin{tabularx}{\textwidth}{@{}ll X@{}}
        \textbf{Industry:} & Technology & \textbf{Partner Enterprise:} \textbf{Skillzza Enterprise} \\
        \textbf{Duration:} & 12-18 hrs (Self-paced) & \textbf{Mode:} Individual, task-based simulation \\
        \textbf{Primary Skills:} & \multicolumn{2}{@{}l}{Python, AI, Data Analysis, Problem Solving}
    \end{tabularx}
\end{metabox}

\subsection*{Overview \& Problem Statement}

\#\#\# \textbf{Scenario}  
Financial institutions face growing threats from fraudulent activities, including identity theft, money laundering, and transaction tampering. As a Fraud Analytics Analyst, you will step into the shoes of a specialist tasked with detecting suspicious activity within transaction datasets and designing fraud detection mechanisms using AI-powered tools and analytics frameworks.

\#\#\# \textbf{Mission}  
Your mission is to analyze transaction data, uncover patterns of fraudulent behavior, and build a fraud detection model capable of identifying anomalies. You will also define fraud rules, audit the AI system for fairness, and present actionable recommendations to reduce fraud risks.

\#\#\# \textbf{Final Challenge}  
At the end of this track, you will submit a detailed fraud investigation report, complete with AI-generated insights, fraud rule recommendations, and ethical considerations for the deployed system.

---

\subsection*{Tasks \& Deliverables}

\begin{taskbox}{Task 1: Understand Fraud Analytics in BFSI}

    {\small \textcolor{gray}{2--3 hrs $\cdot$ Intermediate}}\\[12pt]

    \#\#\#\# \textbf{Scenario}  
A bank has flagged a dataset containing transactional data for potential fraud. You are required to understand fraud patterns, regulatory requirements, and common techniques used in fraud detection.  

\#\#\#\# \textbf{Student Assignment}  
\begin{itemize}
    \item Review the dataset fields and understand their significance (e.g., transaction amount, account balance, merchant category).  
    \item Research common fraud techniques like phishing, synthetic identity fraud, and transaction tampering.  
    \item Explore fraud metrics like True Positive Rate and False Positive Rate.  
\end{itemize}

\#\#\#\# \textbf{Deliverable}  
A summary report (500 words) detailing:  
\begin{itemize}
    \item Types of fraud in BFSI  
    \item Key dataset fields relevant to fraud detection  
    \item Fraud detection challenges in AI systems\par
\end{itemize}

    \vspace{16pt}

    \textcolor{cardBorder}{\hrule height 1pt}\vspace{8pt}

    \small \textcolor{gray}{When you are done with this task, mark it as complete to proceed to the next one.} \hfill \textbf{\textcolor{brandDark}{$\checkmark$ Completed}}

\end{taskbox}

\begin{taskbox}{Task 2: Dataset Preprocessing and Setup}

    {\small \textcolor{gray}{2--3 hrs $\cdot$ Intermediate}}\\[12pt]

    \#\#\#\# \textbf{Scenario}  
You’ve been provided with a raw transaction dataset containing fields like `Transaction\_ID`, `Account\_ID`, `Transaction\_Amount`, `Merchant\_Category`, `Transaction\_Timestamp`, and `Fraud\_Flag`. Your task is to preprocess the dataset for analysis.  

\#\#\#\# \textbf{Student Assignment}  
\begin{itemize}
    \item Handle missing values and outliers in the dataset.  
    \item Convert timestamps into features such as `Hour\_of\_Day` and `Day\_of\_Week`.  
    \item Normalize numerical fields like `Transaction\_Amount`.  
    \item Visualize transaction patterns using heatmaps and histograms.  
\end{itemize}

\#\#\#\# \textbf{Deliverable}  
A cleaned and preprocessed dataset (CSV format) along with a Jupyter Notebook showing EDA results.\par

    \vspace{16pt}

    \textcolor{cardBorder}{\hrule height 1pt}\vspace{8pt}

    \small \textcolor{gray}{When you are done with this task, mark it as complete to proceed to the next one.} \hfill \textbf{\textcolor{brandDark}{$\checkmark$ Completed}}

\end{taskbox}

\begin{taskbox}{Task 3: Build an Anomaly Detection Model}

    {\small \textcolor{gray}{2--3 hrs $\cdot$ Intermediate}}\\[12pt]

    \#\#\#\# \textbf{Scenario}  
Using the preprocessed dataset, you need to build an AI model to identify anomalous transactions that could indicate fraud.  

\#\#\#\# \textbf{Student Assignment}  
\begin{itemize}
    \item Use clustering algorithms (e.g., k-means or DBSCAN) for unsupervised anomaly detection.  
    \item Train a supervised learning model (e.g., Random Forest or XGBoost) using `Fraud\_Flag` as a target variable.  
    \item Evaluate model performance using metrics like precision, recall, and F1-score.  
\end{itemize}

\#\#\#\# \textbf{Deliverable}  
A trained anomaly detection model and a summary report containing:  
\begin{itemize}
    \item Model architecture and parameters  
    \item Performance metrics and confusion matrix\par
\end{itemize}

    \vspace{16pt}

    \textcolor{cardBorder}{\hrule height 1pt}\vspace{8pt}

    \small \textcolor{gray}{When you are done with this task, mark it as complete to proceed to the next one.} \hfill \textbf{\textcolor{brandDark}{$\checkmark$ Completed}}

\end{taskbox}

\begin{taskbox}{Task 4: Use GenAI to Explain Fraud Predictions}

    {\small \textcolor{gray}{2--3 hrs $\cdot$ Intermediate}}\\[12pt]

    \#\#\#\# \textbf{Scenario}  
Your fraud detection model flagged 200 transactions as suspicious. The bank’s compliance team requires explanations for why these transactions were flagged.  

\#\#\#\# \textbf{Student Assignment}  
\begin{itemize}
    \item Use a GenAI tool (e.g., ChatGPT) to generate explanations for flagged transactions based on their features.  
    \item Identify patterns in flagged transactions (e.g., high transaction amounts at odd hours).  
    \item Summarize findings into an actionable report for non-technical stakeholders.  
\end{itemize}

\#\#\#\# \textbf{Deliverable}  
A transaction explanation report (PDF format) containing:  
\begin{itemize}
    \item Top 5 fraud patterns observed  
    \item GenAI-generated explanations for flagged transactions\par
\end{itemize}

    \vspace{16pt}

    \textcolor{cardBorder}{\hrule height 1pt}\vspace{8pt}

    \small \textcolor{gray}{When you are done with this task, mark it as complete to proceed to the next one.} \hfill \textbf{\textcolor{brandDark}{$\checkmark$ Completed}}

\end{taskbox}

\begin{taskbox}{Task 5: Audit the AI System for Responsible Fraud Detection}

    {\small \textcolor{gray}{2--3 hrs $\cdot$ Intermediate}}\\[12pt]

    \#\#\#\# \textbf{Scenario}  
The bank is concerned about potential biases in the fraud detection model, especially regarding merchant categories and transaction amounts.  

\#\#\#\# \textbf{Student Assignment}  
\begin{itemize}
    \item Analyze the model for biased predictions (e.g., disproportionately flagging certain merchant categories).  
    \item Use statistical tests to check for fairness across demographic fields (if available, e.g., `Account\_Holder\_Age`).  
    \item Recommend strategies to mitigate bias and improve fairness.  
\end{itemize}

\#\#\#\# \textbf{Deliverable}  
An AI audit report (PDF format) containing:  
\begin{itemize}
    \item Bias analysis results  
    \item Recommendations for ethical AI deployment\par
\end{itemize}

    \vspace{16pt}

    \textcolor{cardBorder}{\hrule height 1pt}\vspace{8pt}

    \small \textcolor{gray}{When you are done with this task, mark it as complete to proceed to the next one.} \hfill \textbf{\textcolor{brandDark}{$\checkmark$ Completed}}

\end{taskbox}

\begin{taskbox}{Task 6: Present Fraud Risk Mitigation Recommendations}

    {\small \textcolor{gray}{2--3 hrs $\cdot$ Intermediate}}\\[12pt]

    \#\#\#\# \textbf{Scenario}  
Your final task is to present actionable recommendations to reduce fraud risks based on your analysis and model findings.  

\#\#\#\# \textbf{Student Assignment}  
\begin{itemize}
    \item Summarize findings from previous tasks into a cohesive presentation.  
    \item Highlight fraud prevention strategies (e.g., real-time monitoring, rule-based interventions).  
    \item Provide recommendations for improving the fraud detection model.  
\end{itemize}

\#\#\#\# \textbf{Deliverable}  
A presentation slide deck (10-15 slides) covering:  
\begin{itemize}
    \item Fraud trends and findings from the dataset  
    \item Model insights and fraud rules  
    \item Risk mitigation strategies\par
\end{itemize}

    \vspace{16pt}

    \textcolor{cardBorder}{\hrule height 1pt}\vspace{8pt}

    \small \textcolor{gray}{When you are done with this task, mark it as complete to proceed to the next one.} \hfill \textbf{\textcolor{brandDark}{$\checkmark$ Completed}}

\end{taskbox}

\begin{completionbox}
    \textbf{Knowledge Assessment \& Verification:} Complete the online knowledge check, submit your code and presentation files for AI verification, and claim your \textbf{Skillzza Verified Certificate}.
\end{completionbox}

\newpage

% =============================================================================

% SIMULATION 9

% =============================================================================

\section*{Data Quality Analyst – Clean a Messy Enterprise Dataset}

\addcontentsline{toc}{section}{Data Quality Analyst – Clean a Messy Enterprise Dataset}

\begin{metabox}
    \begin{tabularx}{\textwidth}{@{}ll X@{}}
        \textbf{Industry:} & Technology & \textbf{Partner Enterprise:} \textbf{Skillzza Enterprise} \\
        \textbf{Duration:} & 12-18 hrs (Self-paced) & \textbf{Mode:} Individual, task-based simulation \\
        \textbf{Primary Skills:} & \multicolumn{2}{@{}l}{Python, AI, Data Analysis, Problem Solving}
    \end{tabularx}
\end{metabox}

\subsection*{Overview \& Problem Statement}

\#\#\# \textbf{Scenario}  
In the fast-paced world of enterprise data management, businesses rely on clean, accurate, and complete datasets to drive strategy. Imagine you're part of the Data Governance team at a large retail company. You've been handed a massive dataset containing customer purchase records, but it's riddled with inconsistencies: duplicates, missing values, and formatting issues. Your mission is to clean the dataset and provide actionable insights to improve its quality score.

\#\#\# \textbf{Mission}  
Your goal is to analyze, clean, and enhance the quality of a messy enterprise dataset, ensuring it meets industry standards for data integrity. You'll perform profiling and validation, clean anomalies, deduplicate records, handle missing data, and calculate a final quality score. By the end, you’ll present recommendations to stakeholders on how clean data can improve business operations.

\#\#\# \textbf{Final Challenge}  
Deliver a cleaned dataset with a detailed report explaining the data issues resolved, your methodology, and the resulting quality score. Additionally, provide strategic insights on how maintaining clean data can benefit enterprise-level decision-making.

---

\subsection*{Tasks \& Deliverables}

\begin{taskbox}{Task 1: Understand}

    {\small \textcolor{gray}{2--3 hrs $\cdot$ Intermediate}}\\[12pt]

    \#\#\#\#   
You've just received a messy dataset containing customer purchase records for the last quarter. Before diving into cleaning, you need to understand the structure and identify key data issues.

\#\#\#\# \\[10pt]
\vspace{8pt}
{\large \textbf{\textcolor{brandLight}{What you'll learn}}}\\[4pt]  
Perform initial data profiling to identify duplicates, missing values, and inconsistencies. Create a summary report highlighting the major quality issues in the dataset.

\#\#\#\# \vspace{8pt}
{\large \textbf{\textcolor{brandLight}{What you'll do}}}\\[4pt]  
A document summarizing:  
\begin{itemize}
    \item Key columns and their data types.  
    \item Percentage of duplicates and missing values per column.  
    \item Any obvious formatting issues (e.g., inconsistent date formats).\par
\end{itemize}

    \vspace{16pt}

    \textcolor{cardBorder}{\hrule height 1pt}\vspace{8pt}

    \small \textcolor{gray}{When you are done with this task, mark it as complete to proceed to the next one.} \hfill \textbf{\textcolor{brandDark}{$\checkmark$ Completed}}

\end{taskbox}

\begin{taskbox}{Task 2: Data/Setup}

    {\small \textcolor{gray}{2--3 hrs $\cdot$ Intermediate}}\\[12pt]

    \#\#\#\#   
Now that you’ve identified the issues, it’s time to prepare the workspace and tools for cleaning. You'll load the messy dataset into Excel and Python and set up the cleaning pipeline.

\#\#\#\# \\[10pt]
\vspace{8pt}
{\large \textbf{\textcolor{brandLight}{What you'll learn}}}\\[4pt]  
\begin{itemize}
    \item Import the dataset into Excel and Python.  
    \item Set up the environment for cleaning tasks (e.g., install Pandas if needed).  
    \item Save the dataset in a format ready for cleaning (e.g., CSV).  
\end{itemize}

\#\#\#\# \vspace{8pt}
{\large \textbf{\textcolor{brandLight}{What you'll do}}}\\[4pt]  
Submit:  
1. Screenshot of the dataset loaded in Excel and Python.  
2. A brief report describing your setup process.\par

    \vspace{16pt}

    \textcolor{cardBorder}{\hrule height 1pt}\vspace{8pt}

    \small \textcolor{gray}{When you are done with this task, mark it as complete to proceed to the next one.} \hfill \textbf{\textcolor{brandDark}{$\checkmark$ Completed}}

\end{taskbox}

\begin{taskbox}{Task 3: Build/Execute}

    {\small \textcolor{gray}{2--3 hrs $\cdot$ Intermediate}}\\[12pt]

    \#\#\#\#   
It’s time to clean the dataset! You’ll handle duplicates, missing values, and apply validation rules to ensure data accuracy.  

\#\#\#\# \\[10pt]
\vspace{8pt}
{\large \textbf{\textcolor{brandLight}{What you'll learn}}}\\[4pt]  
\begin{itemize}
    \item Deduplicate the dataset using both Excel (Remove Duplicates) and Python (drop\_duplicates in Pandas).  
    \item Handle missing values by imputing or removing incomplete rows.  
    \item Apply validation rules to standardize inconsistent formats (e.g., date formats, phone numbers).  
\end{itemize}

\#\#\#\# \vspace{8pt}
{\large \textbf{\textcolor{brandLight}{What you'll do}}}\\[4pt]  
Submit the cleaned dataset with:  
1. Before and after statistics for duplicates and missing values.  
2. A list of validation rules applied.\par

    \vspace{16pt}

    \textcolor{cardBorder}{\hrule height 1pt}\vspace{8pt}

    \small \textcolor{gray}{When you are done with this task, mark it as complete to proceed to the next one.} \hfill \textbf{\textcolor{brandDark}{$\checkmark$ Completed}}

\end{taskbox}

\begin{taskbox}{Task 4: GenAI/Explanation}

    {\small \textcolor{gray}{2--3 hrs $\cdot$ Intermediate}}\\[12pt]

    \#\#\#\#   
Use Generative AI tools to explain data quality metrics to non-technical stakeholders. Simulate the use of AI-powered profiling tools to generate insights.

\#\#\#\# \\[10pt]
\vspace{8pt}
{\large \textbf{\textcolor{brandLight}{What you'll learn}}}\\[4pt]  
\begin{itemize}
    \item Use ChatGPT or similar tools to generate a plain-English explanation of the cleaning process and its impact on data quality.  
    \item Create a summary of key metrics (e.g., \% of duplicate records removed, \% of missing values handled) for stakeholders.  
\end{itemize}

\#\#\#\# \vspace{8pt}
{\large \textbf{\textcolor{brandLight}{What you'll do}}}\\[4pt]  
Submit:  
1. AI-generated explanation of data cleaning results.  
2. A stakeholder-friendly summary of metrics.\par

    \vspace{16pt}

    \textcolor{cardBorder}{\hrule height 1pt}\vspace{8pt}

    \small \textcolor{gray}{When you are done with this task, mark it as complete to proceed to the next one.} \hfill \textbf{\textcolor{brandDark}{$\checkmark$ Completed}}

\end{taskbox}

\begin{taskbox}{Task 5: Audit/Responsible AI}

    {\small \textcolor{gray}{2--3 hrs $\cdot$ Intermediate}}\\[12pt]

    \#\#\#\#   
Ensuring data cleaning complies with ethical guidelines is crucial. Evaluate the cleaned dataset for fairness and audit your cleaning decisions.  

\#\#\#\# \\[10pt]
\vspace{8pt}
{\large \textbf{\textcolor{brandLight}{What you'll learn}}}\\[4pt]  
\begin{itemize}
    \item Audit the dataset for biases introduced during cleaning (e.g., removing rows based on incomplete demographic data).  
    \item Validate that no critical information was lost during deduplication or missing value handling.  
\end{itemize}

\#\#\#\# \vspace{8pt}
{\large \textbf{\textcolor{brandLight}{What you'll do}}}\\[4pt]  
Submit an audit report covering:  
1. Bias detected (if any) and mitigation strategies.  
2. Validation of dataset integrity.\par

    \vspace{16pt}

    \textcolor{cardBorder}{\hrule height 1pt}\vspace{8pt}

    \small \textcolor{gray}{When you are done with this task, mark it as complete to proceed to the next one.} \hfill \textbf{\textcolor{brandDark}{$\checkmark$ Completed}}

\end{taskbox}

\begin{taskbox}{Task 6: Present Recommendation}

    {\small \textcolor{gray}{2--3 hrs $\cdot$ Intermediate}}\\[12pt]

    \#\#\#\#   
Your cleaned dataset is ready, but stakeholders need actionable insights. Prepare a final report and presentation summarizing your work and recommendations for maintaining data quality in the future.

\#\#\#\# \\[10pt]
\vspace{8pt}
{\large \textbf{\textcolor{brandLight}{What you'll learn}}}\\[4pt]  
\begin{itemize}
    \item Create a presentation summarizing the cleaning process, data quality improvement, and insights gained.  
    \item Propose strategies for ensuring continuous data quality in enterprise operations.
\end{itemize}

\#\#\#\# \vspace{8pt}
{\large \textbf{\textcolor{brandLight}{What you'll do}}}\\[4pt]  
Submit:  
1. Final cleaned dataset (.csv or .xlsx).  
2. Presentation slides (PDF or PowerPoint).\par

    \vspace{16pt}

    \textcolor{cardBorder}{\hrule height 1pt}\vspace{8pt}

    \small \textcolor{gray}{When you are done with this task, mark it as complete to proceed to the next one.} \hfill \textbf{\textcolor{brandDark}{$\checkmark$ Completed}}

\end{taskbox}

\begin{completionbox}
    \textbf{Knowledge Assessment \& Verification:} Complete the online knowledge check, submit your code and presentation files for AI verification, and claim your \textbf{Skillzza Verified Certificate}.
\end{completionbox}

\newpage

% =============================================================================

% SIMULATION 10

% =============================================================================

\section*{Virtual Internship: AI Healthcare Analyst – Hospital Operations Optimisation}

\addcontentsline{toc}{section}{Virtual Internship: AI Healthcare Analyst – Hospital Operations Optimisation}

\begin{metabox}
    \begin{tabularx}{\textwidth}{@{}ll X@{}}
        \textbf{Industry:} & Technology & \textbf{Partner Enterprise:} \textbf{Skillzza Enterprise} \\
        \textbf{Duration:} & 12-18 hrs (Self-paced) & \textbf{Mode:} Individual, task-based simulation \\
        \textbf{Primary Skills:} & \multicolumn{2}{@{}l}{Python, AI, Data Analysis, Problem Solving}
    \end{tabularx}
\end{metabox}

\subsection*{Overview \& Problem Statement}

\textbf{Scenario:}  
Hospitals are complex systems with numerous moving parts — from patient care to staff scheduling to inventory management. As an AI Healthcare Analyst, your mission is to analyze operational bottlenecks and propose AI-driven solutions to optimize hospital performance. In this simulation, you will work with real-world hospital data to identify inefficiencies, leverage predictive analytics, and recommend interventions to improve patient care and operational efficiency.

\textbf{Mission:}  
Your mission is to use AI tools to analyze hospital operations, uncover bottlenecks, and evaluate opportunities for automation or optimization. You will identify areas where AI solutions can streamline processes such as patient flow, staff allocation, and equipment utilization. Your recommendations will aim to enhance hospital efficiency while adhering to ethical AI principles.  

\textbf{Final Challenge:}  
The culmination of this internship will require you to present a comprehensive "AI Intervention Plan" to the hospital's leadership team. Your report will include detailed analysis, AI-driven recommendations, and an implementation roadmap to address the identified bottlenecks.  

---

\begin{completionbox}
    \textbf{Knowledge Assessment \& Verification:} Complete the online knowledge check, submit your code and presentation files for AI verification, and claim your \textbf{Skillzza Verified Certificate}.
\end{completionbox}

\newpage

% =============================================================================

% SIMULATION 11

% =============================================================================

\section*{AI Healthcare – Medical Appointment No-Show Predictor}

\addcontentsline{toc}{section}{AI Healthcare – Medical Appointment No-Show Predictor}

\begin{metabox}
    \begin{tabularx}{\textwidth}{@{}ll X@{}}
        \textbf{Industry:} & Technology & \textbf{Partner Enterprise:} \textbf{Skillzza Enterprise} \\
        \textbf{Duration:} & 12-18 hrs (Self-paced) & \textbf{Mode:} Individual, task-based simulation \\
        \textbf{Primary Skills:} & \multicolumn{2}{@{}l}{Python, AI, Data Analysis, Problem Solving}
    \end{tabularx}
\end{metabox}

\subsection*{Overview \& Problem Statement}

\#\#\# \textbf{Scenario}  
In healthcare systems, missed medical appointments (commonly referred to as "no-shows") contribute to inefficiencies, wasted resources, and delayed care for patients. Healthcare organizations are increasingly leveraging AI to predict which patients are likely to skip appointments and proactively intervene to reduce no-show rates.  

As an \textbf{AI Healthcare Analyst}, you will be tasked with developing and deploying a machine learning model to predict medical appointment no-shows using historical appointment data. Using your insights, you'll segment patients by risk level and propose intervention strategies to improve attendance rates.  

\#\#\# \textbf{Mission}  
Your mission is to analyze medical appointment data, build an AI-powered no-show predictor model, and deliver actionable recommendations to reduce no-show rates through segmentation and intervention strategies.  

\#\#\# \textbf{Final Challenge}  
At the end of this simulation, you will present a comprehensive report to stakeholders that includes:  
1. Predictive model performance metrics.  
2. Patient segmentation based on no-show risk levels.  
3. Evidence-based recommendations for reducing no-shows.  

---

\subsection*{Tasks \& Deliverables}

\begin{taskbox}{Task 1: Understand – Analyze the Problem and Dataset}

    {\small \textcolor{gray}{2--3 hrs $\cdot$ Intermediate}}\\[12pt]

    \#\#\#\# \textbf{Scenario}  
You’ve just joined a healthcare analytics team tasked with reducing no-shows for medical appointments. Your first step is to understand the scope of the problem and analyze the provided dataset.  

\#\#\#\# \textbf{Student Assignment}  
\begin{itemize}
    \item Analyze the impact of no-shows on healthcare efficiency (research online).  
    \item Load the provided appointment dataset and perform an initial review of its fields, data types, and missing values.  
    \item Summarize the key trends in no-show rates (e.g., by gender, age, and health condition).  
\end{itemize}

\#\#\#\# \textbf{Deliverable}  
Submit a report containing:  
\begin{itemize}
    \item Key insights about no-show trends (graphs and tables).  
    \item Summary of missing values and potential data quality issues.\par
\end{itemize}

    \vspace{16pt}

    \textcolor{cardBorder}{\hrule height 1pt}\vspace{8pt}

    \small \textcolor{gray}{When you are done with this task, mark it as complete to proceed to the next one.} \hfill \textbf{\textcolor{brandDark}{$\checkmark$ Completed}}

\end{taskbox}

\begin{taskbox}{Task 2: Data/Setup – Preprocess and Prepare the Dataset}

    {\small \textcolor{gray}{2--3 hrs $\cdot$ Intermediate}}\\[12pt]

    \#\#\#\# \textbf{Scenario}  
The dataset needs to be cleaned and prepared for predictive modeling. You’ll handle missing values, encode categorical variables, and address class imbalance.  

\#\#\#\# \textbf{Student Assignment}  
\begin{itemize}
    \item Handle missing values using appropriate imputation techniques.  
    \item Encode categorical variables like `Gender` and `HealthCondition`.  
    \item Address class imbalance in the `NoShow` field using techniques such as oversampling or undersampling.  
\end{itemize}

\#\#\#\# \textbf{Deliverable}  
Submit a clean and preprocessed dataset ready for modeling. Include a summary of preprocessing steps taken.\par

    \vspace{16pt}

    \textcolor{cardBorder}{\hrule height 1pt}\vspace{8pt}

    \small \textcolor{gray}{When you are done with this task, mark it as complete to proceed to the next one.} \hfill \textbf{\textcolor{brandDark}{$\checkmark$ Completed}}

\end{taskbox}

\begin{taskbox}{Task 3: Build/Execute – Train and Evaluate the No-Show Predictor}

    {\small \textcolor{gray}{2--3 hrs $\cdot$ Intermediate}}\\[12pt]

    \#\#\#\# \textbf{Scenario}  
Using the preprocessed dataset, you’ll train a classification model to predict whether a patient is likely to no-show an appointment.  

\#\#\#\# \textbf{Student Assignment}  
\begin{itemize}
    \item Split the dataset into training and testing sets (e.g., 80/20 split).  
    \item Train at least two models (e.g., Logistic Regression, Random Forest).  
    \item Evaluate model performance using metrics such as accuracy, precision, recall, and F1-score.  
\end{itemize}

\#\#\#\# \textbf{Deliverable}  
Submit a summary comparing the performance of different models. Include code snippets and visualizations of key metrics (e.g., confusion matrix).\par

    \vspace{16pt}

    \textcolor{cardBorder}{\hrule height 1pt}\vspace{8pt}

    \small \textcolor{gray}{When you are done with this task, mark it as complete to proceed to the next one.} \hfill \textbf{\textcolor{brandDark}{$\checkmark$ Completed}}

\end{taskbox}

\begin{taskbox}{Task 4: GenAI/Explanation – Interpret Predictions}

    {\small \textcolor{gray}{2--3 hrs $\cdot$ Intermediate}}\\[12pt]

    \#\#\#\# \textbf{Scenario}  
Stakeholders need to understand the factors influencing no-show predictions to trust the model's recommendations. You’ll use explainable AI techniques to interpret predictions.  

\#\#\#\# \textbf{Student Assignment}  
\begin{itemize}
    \item Choose one model and explain its predictions using SHAP or LIME.  
    \item Identify which features contribute most to the likelihood of a no-show (e.g., age, SMS reminder).  
\end{itemize}

\#\#\#\# \textbf{Deliverable}  
Submit a document with feature importance graphs and written explanations of model predictions.\par

    \vspace{16pt}

    \textcolor{cardBorder}{\hrule height 1pt}\vspace{8pt}

    \small \textcolor{gray}{When you are done with this task, mark it as complete to proceed to the next one.} \hfill \textbf{\textcolor{brandDark}{$\checkmark$ Completed}}

\end{taskbox}

\begin{taskbox}{Task 5: Audit/Responsible AI – Ensure Ethical Deployment}

    {\small \textcolor{gray}{2--3 hrs $\cdot$ Intermediate}}\\[12pt]

    \#\#\#\# \textbf{Scenario}  
AI models in healthcare must adhere to ethical principles. You are tasked with auditing the no-show predictor for biases.  

\#\#\#\# \textbf{Student Assignment}  
\begin{itemize}
    \item Analyze model predictions for potential biases (e.g., gender or age bias).  
    \item Propose strategies to mitigate any identified biases (e.g., reweighting samples or adding fairness constraints).  
\end{itemize}

\#\#\#\# \textbf{Deliverable}  
Submit an audit report on model fairness and proposed mitigation strategies.\par

    \vspace{16pt}

    \textcolor{cardBorder}{\hrule height 1pt}\vspace{8pt}

    \small \textcolor{gray}{When you are done with this task, mark it as complete to proceed to the next one.} \hfill \textbf{\textcolor{brandDark}{$\checkmark$ Completed}}

\end{taskbox}

\begin{taskbox}{Task 6: Present Recommendation – Segment Patients and Propose Interventions}

    {\small \textcolor{gray}{2--3 hrs $\cdot$ Intermediate}}\\[12pt]

    \#\#\#\# \textbf{Scenario}  
Based on your predictions, segment patients by no-show risk levels (e.g., high, medium, low). Develop actionable strategies to reduce no-shows for each segment.  

\#\#\#\# \textbf{Student Assignment}  
\begin{itemize}
    \item Use segmentation techniques to group patients by risk levels.  
    \item Propose intervention strategies for each segment (e.g., SMS reminders, telehealth options for high-risk patients).  
\end{itemize}

\#\#\#\# \textbf{Deliverable}  
Submit a stakeholder presentation with:  
\begin{itemize}
    \item Patient segmentation visualizations.  
    \item Intervention strategies tailored to each risk group.\par
\end{itemize}

    \vspace{16pt}

    \textcolor{cardBorder}{\hrule height 1pt}\vspace{8pt}

    \small \textcolor{gray}{When you are done with this task, mark it as complete to proceed to the next one.} \hfill \textbf{\textcolor{brandDark}{$\checkmark$ Completed}}

\end{taskbox}

\begin{completionbox}
    \textbf{Knowledge Assessment \& Verification:} Complete the online knowledge check, submit your code and presentation files for AI verification, and claim your \textbf{Skillzza Verified Certificate}.
\end{completionbox}

\newpage

% =============================================================================

% SIMULATION 12

% =============================================================================

\section*{AI Pharma Analyst – Drug Discovery Intelligence}

\addcontentsline{toc}{section}{AI Pharma Analyst – Drug Discovery Intelligence}

\begin{metabox}
    \begin{tabularx}{\textwidth}{@{}ll X@{}}
        \textbf{Industry:} & Technology & \textbf{Partner Enterprise:} \textbf{Skillzza Enterprise} \\
        \textbf{Duration:} & 12-18 hrs (Self-paced) & \textbf{Mode:} Individual, task-based simulation \\
        \textbf{Primary Skills:} & \multicolumn{2}{@{}l}{Python, AI, Data Analysis, Problem Solving}
    \end{tabularx}
\end{metabox}

\subsection*{Overview \& Problem Statement}

\#\#\# \textbf{Scenario}
In the rapidly evolving pharmaceutical industry, artificial intelligence is revolutionizing drug discovery and development. As an \textbf{AI Pharma Analyst}, you will work on analyzing complex biological datasets, identifying drug targets, and leveraging AI tools to prioritize drug candidates effectively. Your mission is to simulate the role of an advanced AI researcher working in a pharmaceutical company to facilitate faster and more accurate drug discovery.

\#\#\# \textbf{Mission}
Your virtual internship revolves around helping a biotech company identify promising drug targets for treating a rare neurodegenerative disease. Using AI-powered tools, you will extract insights from large-scale datasets, analyze genomic research, and prioritize drug candidates based on efficacy and safety predictions.

\#\#\# \textbf{Final Challenge}
At the end of the internship, you will deliver a comprehensive report detailing your findings, methods, and AI-driven prioritization of drug candidates, along with actionable recommendations for the next stages of research.

---

\subsection*{Tasks \& Deliverables}

\begin{taskbox}{Task 1: Understand the Drug Discovery Pipeline}

    {\small \textcolor{gray}{2--3 hrs $\cdot$ Intermediate}}\\[12pt]

    \#\#\#\# \textbf{Scenario}
You are tasked with understanding how AI fits into the modern drug discovery pipeline, including target identification, validation, and candidate prioritization.

\#\#\#\# \textbf{Student Assignment}
\begin{itemize}
    \item Research the stages of drug discovery.
    \item Study how AI techniques like deep learning and NLP are used in pharma research.
    \item Review the provided case study explaining AI-assisted drug discovery.
\end{itemize}

\#\#\#\# \textbf{Deliverable}
A one-page summary explaining the role of AI in drug discovery with references to the case study provided.\par

    \vspace{16pt}

    \textcolor{cardBorder}{\hrule height 1pt}\vspace{8pt}

    \small \textcolor{gray}{When you are done with this task, mark it as complete to proceed to the next one.} \hfill \textbf{\textcolor{brandDark}{$\checkmark$ Completed}}

\end{taskbox}

\begin{taskbox}{Task 2: Explore the Dataset and Set Up Your Environment}

    {\small \textcolor{gray}{2--3 hrs $\cdot$ Intermediate}}\\[12pt]

    \#\#\#\# \textbf{Scenario}
You have received a dataset containing genomic sequences, proteomic profiles, and experimental drug data. Your first step is to understand the structure of the data and set up your analysis environment.

\#\#\#\# \textbf{Student Assignment}
\begin{itemize}
    \item Load and explore the dataset provided (CSV format).
    \item Identify key fields such as gene names, protein structures, efficacy scores, and toxicity data.
    \item Set up your Python environment with necessary libraries.
\end{itemize}

\#\#\#\# \textbf{Deliverable}
A Python notebook with exploratory data analysis (EDA), including summary statistics, missing value handling, and visualizations.\par

    \vspace{16pt}

    \textcolor{cardBorder}{\hrule height 1pt}\vspace{8pt}

    \small \textcolor{gray}{When you are done with this task, mark it as complete to proceed to the next one.} \hfill \textbf{\textcolor{brandDark}{$\checkmark$ Completed}}

\end{taskbox}

\begin{taskbox}{Task 3: Build AI Models for Target Identification}

    {\small \textcolor{gray}{2--3 hrs $\cdot$ Intermediate}}\\[12pt]

    \#\#\#\# \textbf{Scenario}
Using the dataset, your next task is to create predictive models to identify drug targets that are most likely to be effective against the neurodegenerative disease.

\#\#\#\# \textbf{Student Assignment}
\begin{itemize}
    \item Preprocess the data and prepare it for model training.
    \item Train AI models (e.g., random forest, deep learning) to predict potential drug targets.
    \item Evaluate model performance using metrics like accuracy, precision, and recall.
\end{itemize}

\#\#\#\# \textbf{Deliverable}
Python scripts and trained model files, along with an evaluation report containing key model metrics.\par

    \vspace{16pt}

    \textcolor{cardBorder}{\hrule height 1pt}\vspace{8pt}

    \small \textcolor{gray}{When you are done with this task, mark it as complete to proceed to the next one.} \hfill \textbf{\textcolor{brandDark}{$\checkmark$ Completed}}

\end{taskbox}

\begin{taskbox}{Task 4: Use GenAI for Literature Mining and Hypothesis Generation}

    {\small \textcolor{gray}{2--3 hrs $\cdot$ Intermediate}}\\[12pt]

    \#\#\#\# \textbf{Scenario}
To complement your AI model findings, use generative AI tools to mine scientific literature and generate hypotheses about the identified drug targets.

\#\#\#\# \textbf{Student Assignment}
\begin{itemize}
    \item Use PubMed API or a GPT-based tool to extract insights related to the disease and potential targets.
    \item Summarize relevant research papers and identify gaps in existing studies.
    \item Generate hypotheses for further experimental validation.
\end{itemize}

\#\#\#\# \textbf{Deliverable}
A document summarizing literature findings and hypotheses generated using GenAI tools.\par

    \vspace{16pt}

    \textcolor{cardBorder}{\hrule height 1pt}\vspace{8pt}

    \small \textcolor{gray}{When you are done with this task, mark it as complete to proceed to the next one.} \hfill \textbf{\textcolor{brandDark}{$\checkmark$ Completed}}

\end{taskbox}

\begin{taskbox}{Task 5: Audit Your AI Model for Responsible Use}

    {\small \textcolor{gray}{2--3 hrs $\cdot$ Intermediate}}\\[12pt]

    \#\#\#\# \textbf{Scenario}
To ensure ethical use of AI in drug discovery, you must audit your model for biases and explainability.

\#\#\#\# \textbf{Student Assignment}
\begin{itemize}
    \item Analyze your model using SHAP or LIME to understand feature importance.
    \item Check for biases in predictions (e.g., against specific genomic profiles).
    \item Propose strategies to improve the fairness and reliability of the model.
\end{itemize}

\#\#\#\# \textbf{Deliverable}
A report detailing your model audit findings and proposed improvements.\par

    \vspace{16pt}

    \textcolor{cardBorder}{\hrule height 1pt}\vspace{8pt}

    \small \textcolor{gray}{When you are done with this task, mark it as complete to proceed to the next one.} \hfill \textbf{\textcolor{brandDark}{$\checkmark$ Completed}}

\end{taskbox}

\begin{taskbox}{Task 6: Present Your Recommendations}

    {\small \textcolor{gray}{2--3 hrs $\cdot$ Intermediate}}\\[12pt]

    \#\#\#\# \textbf{Scenario}
As the final step, you will present your findings to the biotech company’s research team, including your AI-driven prioritization of drug candidates.

\#\#\#\# \textbf{Student Assignment}
\begin{itemize}
    \item Synthesize your findings into a concise slide deck.
    \item Provide actionable recommendations based on your analysis and AI model results.
    \item Highlight the most promising drug candidates for further development.
\end{itemize}

\#\#\#\# \textbf{Deliverable}
A PDF slide deck (6–8 slides) summarizing key findings, methods, and recommendations.\par

    \vspace{16pt}

    \textcolor{cardBorder}{\hrule height 1pt}\vspace{8pt}

    \small \textcolor{gray}{When you are done with this task, mark it as complete to proceed to the next one.} \hfill \textbf{\textcolor{brandDark}{$\checkmark$ Completed}}

\end{taskbox}

\begin{completionbox}
    \textbf{Knowledge Assessment \& Verification:} Complete the online knowledge check, submit your code and presentation files for AI verification, and claim your \textbf{Skillzza Verified Certificate}.
\end{completionbox}

\newpage

% =============================================================================

% SIMULATION 13

% =============================================================================

\section*{Manufacturing Data Analyst – Quality Defect Detection}

\addcontentsline{toc}{section}{Manufacturing Data Analyst – Quality Defect Detection}

\begin{metabox}
    \begin{tabularx}{\textwidth}{@{}ll X@{}}
        \textbf{Industry:} & Technology & \textbf{Partner Enterprise:} \textbf{Skillzza Enterprise} \\
        \textbf{Duration:} & 12-18 hrs (Self-paced) & \textbf{Mode:} Individual, task-based simulation \\
        \textbf{Primary Skills:} & \multicolumn{2}{@{}l}{Python, AI, Data Analysis, Problem Solving}
    \end{tabularx}
\end{metabox}

\subsection*{Overview \& Problem Statement}

\#\#\# \textbf{Scenario}  
You have been hired as a Manufacturing Data Analyst at a global consumer goods company that produces high-volume products such as electronics, appliances, and tools. Your job is to analyze production data to identify recurring quality defects, quantify their impact, and recommend corrective actions to improve manufacturing processes.

\#\#\# \textbf{Mission}  
Your mission is to work with raw production data, uncover patterns in quality defects, and use data-driven techniques to prioritize issues using Pareto principles. By leveraging analytical tools and techniques, you will help the organization reduce defects, enhance product reliability, and save costs.

\#\#\# \textbf{Final Challenge}  
For the final challenge, you will present a comprehensive report summarizing your findings. This report will include defect trends, Pareto analysis, and actionable recommendations for process improvement based on data insights.

---

\begin{completionbox}
    \textbf{Knowledge Assessment \& Verification:} Complete the online knowledge check, submit your code and presentation files for AI verification, and claim your \textbf{Skillzza Verified Certificate}.
\end{completionbox}

\newpage

% =============================================================================

% SIMULATION 14

% =============================================================================

\section*{AI Supply Chain Analyst – Demand \& Inventory Prediction}

\addcontentsline{toc}{section}{AI Supply Chain Analyst – Demand \& Inventory Prediction}

\begin{metabox}
    \begin{tabularx}{\textwidth}{@{}ll X@{}}
        \textbf{Industry:} & Technology & \textbf{Partner Enterprise:} \textbf{Skillzza Enterprise} \\
        \textbf{Duration:} & 12-18 hrs (Self-paced) & \textbf{Mode:} Individual, task-based simulation \\
        \textbf{Primary Skills:} & \multicolumn{2}{@{}l}{Python, AI, Data Analysis, Problem Solving}
    \end{tabularx}
\end{metabox}

\subsection*{Overview \& Problem Statement}

\#\#\# \textbf{Scenario}  
As global businesses navigate increasingly complex supply chain challenges, the role of AI in predicting demand and optimizing inventory has become crucial. In this virtual internship, you will step into the role of an AI Supply Chain Analyst at "OptiFlow Logistics," a leading supply chain solutions company specializing in AI-driven inventory management. Your mission is to analyze historical sales data, forecast future demand, and optimize inventory levels using advanced AI techniques.  

\#\#\# \textbf{Mission}  
Your goal is to create an AI-powered system that accurately predicts product demand while minimizing inventory costs, reducing stockouts, and improving supply chain efficiency. The final challenge is to present a business case to leadership, showcasing how your solution can drive operational excellence and tangible business outcomes.  

\#\#\# \textbf{Final Challenge}  
\begin{itemize}
    \item Develop a predictive demand forecasting model using real-world supply chain data.  
    \item Optimize inventory levels based on predictions and business constraints.  
    \item Present an executive-level report highlighting the financial and operational impact of your solution.  
\end{itemize}

---

\subsection*{Tasks \& Deliverables}

\begin{taskbox}{Task 1: Understand the Supply Chain Context}

    {\small \textcolor{gray}{2--3 hrs $\cdot$ Intermediate}}\\[12pt]

    \#\#\#\# \textbf{Scenario}  
OptiFlow Logistics is experiencing issues with overstocking certain product categories while facing stockouts in others. The leadership team needs a detailed analysis of historical sales data to understand patterns and drivers of demand fluctuations.

\#\#\#\# \textbf{Student Assignment}  
\begin{itemize}
    \item Study the provided dataset containing 3 years of demand history for 10 product categories.  
    \item Identify seasonality, trends, and anomalies using visualizations and statistical techniques.  
    \item Document key insights in a report.  
\end{itemize}

\#\#\#\# \textbf{Deliverable}  
\begin{itemize}
    \item A detailed EDA report with visualizations highlighting patterns and anomalies.\par
\end{itemize}

    \vspace{16pt}

    \textcolor{cardBorder}{\hrule height 1pt}\vspace{8pt}

    \small \textcolor{gray}{When you are done with this task, mark it as complete to proceed to the next one.} \hfill \textbf{\textcolor{brandDark}{$\checkmark$ Completed}}

\end{taskbox}

\begin{taskbox}{Task 2: Data Preparation and Setup}

    {\small \textcolor{gray}{2--3 hrs $\cdot$ Intermediate}}\\[12pt]

    \#\#\#\# \textbf{Scenario}  
The raw dataset contains missing values, outliers, and inconsistent formats. Before building a forecasting model, you need to clean and preprocess the data to ensure accuracy.

\#\#\#\# \textbf{Student Assignment}  
\begin{itemize}
    \item Handle missing values using imputation techniques (e.g., mean/median or forward fill).  
    \item Remove outliers using statistical thresholds or visualization methods.  
    \item Normalize numerical data and encode categorical fields.  
\end{itemize}

\#\#\#\# \textbf{Deliverable}  
\begin{itemize}
    \item A clean, structured dataset ready for modeling.\par
\end{itemize}

    \vspace{16pt}

    \textcolor{cardBorder}{\hrule height 1pt}\vspace{8pt}

    \small \textcolor{gray}{When you are done with this task, mark it as complete to proceed to the next one.} \hfill \textbf{\textcolor{brandDark}{$\checkmark$ Completed}}

\end{taskbox}

\begin{taskbox}{Task 3: Build and Execute the Demand Forecasting Model}

    {\small \textcolor{gray}{2--3 hrs $\cdot$ Intermediate}}\\[12pt]

    \#\#\#\# \textbf{Scenario}  
Using the cleaned dataset, you will train an AI model to predict monthly demand for each product category over the next 12 months. Your model needs to account for seasonality and trends.

\#\#\#\# \textbf{Student Assignment}  
\begin{itemize}
    \item Split the dataset into training and test sets.  
    \item Build and train a time series forecasting model using ARIMA, Prophet, or LSTM.  
    \item Evaluate model performance using metrics like MAPE (Mean Absolute Percentage Error) and RMSE (Root Mean Square Error).  
\end{itemize}

\#\#\#\# \textbf{Deliverable}  
\begin{itemize}
    \item A trained and tested demand forecasting model with performance metrics.\par
\end{itemize}

    \vspace{16pt}

    \textcolor{cardBorder}{\hrule height 1pt}\vspace{8pt}

    \small \textcolor{gray}{When you are done with this task, mark it as complete to proceed to the next one.} \hfill \textbf{\textcolor{brandDark}{$\checkmark$ Completed}}

\end{taskbox}

\begin{taskbox}{Task 4: Explain Insights Using Generative AI}

    {\small \textcolor{gray}{2--3 hrs $\cdot$ Intermediate}}\\[12pt]

    \#\#\#\# \textbf{Scenario}  
The leadership team at OptiFlow Logistics wants to understand the model’s predictions and implications for their business. Use Generative AI tools to summarize findings in a clear and concise way.

\#\#\#\# \textbf{Student Assignment}  
\begin{itemize}
    \item Use an AI tool like ChatGPT or Bard to generate a business summary of your model’s predictions.  
    \item Include key insights on product categories with high or low demand and explain trends.  
\end{itemize}

\#\#\#\# \textbf{Deliverable}  
\begin{itemize}
    \item A leadership-ready summary report generated using Generative AI.\par
\end{itemize}

    \vspace{16pt}

    \textcolor{cardBorder}{\hrule height 1pt}\vspace{8pt}

    \small \textcolor{gray}{When you are done with this task, mark it as complete to proceed to the next one.} \hfill \textbf{\textcolor{brandDark}{$\checkmark$ Completed}}

\end{taskbox}

\begin{taskbox}{Task 5: Audit and Ensure Responsible AI Practices}

    {\small \textcolor{gray}{2--3 hrs $\cdot$ Intermediate}}\\[12pt]

    \#\#\#\# \textbf{Scenario}  
Decision-making based on AI models can have ethical implications. You need to ensure your forecasting model is transparent, unbiased, and interpretable.

\#\#\#\# \textbf{Student Assignment}  
\begin{itemize}
    \item Conduct a fairness audit of your model’s predictions across different product categories.  
    \item Use SHAP (SHapley Additive exPlanations) or LIME to explain model decisions.  
    \item Document areas for improvement in the model.  
\end{itemize}

\#\#\#\# \textbf{Deliverable}  
\begin{itemize}
    \item A Responsible AI audit report with fairness metrics and interpretability analysis.\par
\end{itemize}

    \vspace{16pt}

    \textcolor{cardBorder}{\hrule height 1pt}\vspace{8pt}

    \small \textcolor{gray}{When you are done with this task, mark it as complete to proceed to the next one.} \hfill \textbf{\textcolor{brandDark}{$\checkmark$ Completed}}

\end{taskbox}

\begin{taskbox}{Task 6: Present Your Business Case and Recommendations}

    {\small \textcolor{gray}{2--3 hrs $\cdot$ Intermediate}}\\[12pt]

    \#\#\#\# \textbf{Scenario}  
You are presenting your solution to the leadership team. Showcase how your demand forecasting model and inventory optimization plan can drive cost savings and operational efficiency.

\#\#\#\# \textbf{Student Assignment}  
\begin{itemize}
    \item Create a PowerPoint or dashboard summarizing your findings.  
    \item Highlight the impact of your solution on reducing stockouts, minimizing overstock, and improving forecasting accuracy.  
    \item Include financial projections supported by data.  
\end{itemize}

\#\#\#\# \textbf{Deliverable}  
\begin{itemize}
    \item A visually compelling business case presentation or dashboard.\par
\end{itemize}

    \vspace{16pt}

    \textcolor{cardBorder}{\hrule height 1pt}\vspace{8pt}

    \small \textcolor{gray}{When you are done with this task, mark it as complete to proceed to the next one.} \hfill \textbf{\textcolor{brandDark}{$\checkmark$ Completed}}

\end{taskbox}

\begin{completionbox}
    \textbf{Knowledge Assessment \& Verification:} Complete the online knowledge check, submit your code and presentation files for AI verification, and claim your \textbf{Skillzza Verified Certificate}.
\end{completionbox}

\newpage

% =============================================================================

% SIMULATION 15

% =============================================================================

\section*{Computer Vision Engineer – Factory Quality Inspector}

\addcontentsline{toc}{section}{Computer Vision Engineer – Factory Quality Inspector}

\begin{metabox}
    \begin{tabularx}{\textwidth}{@{}ll X@{}}
        \textbf{Industry:} & Technology & \textbf{Partner Enterprise:} \textbf{Skillzza Enterprise} \\
        \textbf{Duration:} & 12-18 hrs (Self-paced) & \textbf{Mode:} Individual, task-based simulation \\
        \textbf{Primary Skills:} & \multicolumn{2}{@{}l}{Python, AI, Data Analysis, Problem Solving}
    \end{tabularx}
\end{metabox}

\subsection*{Overview \& Problem Statement}

\#\#\# \textbf{Scenario}  
Imagine you are working as a Computer Vision Engineer for a high-tech manufacturing company specializing in the production of industrial parts. The company aims to improve its quality inspection process by automating defect detection using computer vision. Traditional manual inspection is both time-consuming and prone to errors, leading to inconsistent product quality that impacts customer satisfaction. Your role is critical in designing, deploying, and optimizing an AI-powered defect classification system.  

\#\#\# \textbf{Mission}  
Your mission is to create a defect classification model that can accurately identify faulty parts from factory images, integrate this model into a workflow, and design an inspection dashboard for real-time monitoring. You will also ensure the AI system adheres to responsible AI practices, including accuracy, fairness, and transparency.  

\#\#\# \textbf{Final Challenge}  
By the end of this simulation, you will present your defect classification pipeline and an interactive dashboard to factory stakeholders. Your deliverable should demonstrate how the AI system improves accuracy, reduces costs, and integrates seamlessly into the factory workflow.  

---

\subsection*{Tasks \& Deliverables}

\begin{taskbox}{Task 1: Understand the Problem Domain}

    {\small \textcolor{gray}{2--3 hrs $\cdot$ Intermediate}}\\[12pt]

    \#\#\#\# \textbf{Scenario}  
Your factory manager has provided you with a dataset of images containing industrial parts marked as "defective" or "non-defective." The defects include scratches, dents, and misalignments. Your first task is to analyze the dataset and understand the patterns that differentiate defective from non-defective parts.  

\#\#\#\# \textbf{Student Assignment}  
1. Explore the dataset (CSV file with image paths and labels).  
2. Visualize the sample images of defective and non-defective parts.  
3. Identify unique challenges in detecting defects (e.g., small scratches, varied lighting conditions).  

\#\#\#\# \textbf{Deliverable}  
A short report (1-2 pages) summarizing your observations about the dataset, including key challenges and visual samples of defects.\par

    \vspace{16pt}

    \textcolor{cardBorder}{\hrule height 1pt}\vspace{8pt}

    \small \textcolor{gray}{When you are done with this task, mark it as complete to proceed to the next one.} \hfill \textbf{\textcolor{brandDark}{$\checkmark$ Completed}}

\end{taskbox}

\begin{taskbox}{Task 2: Preprocess the Dataset}

    {\small \textcolor{gray}{2--3 hrs $\cdot$ Intermediate}}\\[12pt]

    \#\#\#\# \textbf{Scenario}  
Factory images often have inconsistent lighting, variable angles, and noise. To train an accurate model, you need to preprocess the dataset by resizing, normalizing, and augmenting the images.  

\#\#\#\# \textbf{Student Assignment}  
1. Resize all images to a standard input size (e.g., 224x224 pixels).  
2. Normalize pixel values to a range of 0-1.  
3. Apply data augmentation techniques: rotation, flipping, and brightness adjustment.  

\#\#\#\# \textbf{Deliverable}  
A Python script that preprocesses the dataset and outputs a set of augmented images ready for model training.\par

    \vspace{16pt}

    \textcolor{cardBorder}{\hrule height 1pt}\vspace{8pt}

    \small \textcolor{gray}{When you are done with this task, mark it as complete to proceed to the next one.} \hfill \textbf{\textcolor{brandDark}{$\checkmark$ Completed}}

\end{taskbox}

\begin{taskbox}{Task 3: Build and Train the Defect Classification Model}

    {\small \textcolor{gray}{2--3 hrs $\cdot$ Intermediate}}\\[12pt]

    \#\#\#\# \textbf{Scenario}  
Using transfer learning, you will train a convolutional neural network (CNN) to classify images as defective or non-defective. You aim for a minimum accuracy of 90\% on the test set.  

\#\#\#\# \textbf{Student Assignment}  
1. Select a pre-trained model (e.g., ResNet50, MobileNet).  
2. Fine-tune the model with your preprocessed dataset.  
3. Evaluate the model using metrics like precision, recall, and F1 score.  

\#\#\#\# \textbf{Deliverable}  
A Jupyter Notebook with the trained model, performance metrics, and sample predictions.\par

    \vspace{16pt}

    \textcolor{cardBorder}{\hrule height 1pt}\vspace{8pt}

    \small \textcolor{gray}{When you are done with this task, mark it as complete to proceed to the next one.} \hfill \textbf{\textcolor{brandDark}{$\checkmark$ Completed}}

\end{taskbox}

\begin{taskbox}{Task 4: Explain Model Predictions with GenAI Tools}

    {\small \textcolor{gray}{2--3 hrs $\cdot$ Intermediate}}\\[12pt]

    \#\#\#\# \textbf{Scenario}  
To maintain transparency, you need to explain why your model is classifying certain parts as defective or non-defective. This will help factory staff understand and trust the AI system.  

\#\#\#\# \textbf{Student Assignment}  
1. Use SHAP or LIME to analyze the predictions made by your model.  
2. Visualize the regions of the image that contributed to the defect classification.  
3. Write a brief explanation of the results for stakeholders.  

\#\#\#\# \textbf{Deliverable}  
A set of visual explanations (heatmaps) showing defect regions and a summary document explaining the model’s decision-making process.\par

    \vspace{16pt}

    \textcolor{cardBorder}{\hrule height 1pt}\vspace{8pt}

    \small \textcolor{gray}{When you are done with this task, mark it as complete to proceed to the next one.} \hfill \textbf{\textcolor{brandDark}{$\checkmark$ Completed}}

\end{taskbox}

\begin{taskbox}{Task 5: Audit the Model for Responsible AI Practices}

    {\small \textcolor{gray}{2--3 hrs $\cdot$ Intermediate}}\\[12pt]

    \#\#\#\# \textbf{Scenario}  
The factory’s leadership wants to ensure the AI system is fair and unbiased. You need to audit the model for potential biases (e.g., bias against certain defect types) and retrain if necessary.  

\#\#\#\# \textbf{Student Assignment}  
1. Analyze the model’s performance across all defect categories.  
2. Identify any bias in predictions (e.g., higher recall for scratches but lower for dents).  
3. Retrain the model or adjust hyperparameters to address bias.  

\#\#\#\# \textbf{Deliverable}  
An audit report detailing the bias analysis and steps taken to mitigate it. Include updated performance metrics.\par

    \vspace{16pt}

    \textcolor{cardBorder}{\hrule height 1pt}\vspace{8pt}

    \small \textcolor{gray}{When you are done with this task, mark it as complete to proceed to the next one.} \hfill \textbf{\textcolor{brandDark}{$\checkmark$ Completed}}

\end{taskbox}

\begin{taskbox}{Task 6: Present Your Recommendation}

    {\small \textcolor{gray}{2--3 hrs $\cdot$ Intermediate}}\\[12pt]

    \#\#\#\# \textbf{Scenario}  
You are tasked with presenting the defect classification system and an interactive dashboard to factory stakeholders. The dashboard should show real-time predictions and model performance metrics.  

\#\#\#\# \textbf{Student Assignment}  
1. Create a Streamlit dashboard showing sample predictions, confidence scores, and overall model metrics.  
2. Prepare a brief presentation (5-10 slides) summarizing your work, including dataset challenges, model performance, and dashboard functionality.  

\#\#\#\# \textbf{Deliverable}  
A Streamlit dashboard URL and a presentation file (PDF or PPT).\par

    \vspace{16pt}

    \textcolor{cardBorder}{\hrule height 1pt}\vspace{8pt}

    \small \textcolor{gray}{When you are done with this task, mark it as complete to proceed to the next one.} \hfill \textbf{\textcolor{brandDark}{$\checkmark$ Completed}}

\end{taskbox}

\begin{completionbox}
    \textbf{Knowledge Assessment \& Verification:} Complete the online knowledge check, submit your code and presentation files for AI verification, and claim your \textbf{Skillzza Verified Certificate}.
\end{completionbox}

\newpage

% =============================================================================

% SIMULATION 16

% =============================================================================

\section*{Precision Agriculture Analyst – Smart Irrigation}

\addcontentsline{toc}{section}{Precision Agriculture Analyst – Smart Irrigation}

\begin{metabox}
    \begin{tabularx}{\textwidth}{@{}ll X@{}}
        \textbf{Industry:} & Technology & \textbf{Partner Enterprise:} \textbf{Skillzza Enterprise} \\
        \textbf{Duration:} & 12-18 hrs (Self-paced) & \textbf{Mode:} Individual, task-based simulation \\
        \textbf{Primary Skills:} & \multicolumn{2}{@{}l}{Python, AI, Data Analysis, Problem Solving}
    \end{tabularx}
\end{metabox}

\subsection*{Overview \& Problem Statement}

\#\#\# \textbf{Scenario}  
The world is facing increasing challenges with water scarcity and food security, driving the need for smarter agricultural practices. As a Precision Agriculture Analyst specializing in Smart Irrigation, you will work to optimize water usage for crop fields by leveraging weather, soil, and crop data. Your goal is to design predictive models to estimate water demand and create efficient irrigation plans that minimize waste while maximizing crop yield.

\#\#\# \textbf{Mission}  
Your mission in this simulation is to analyze data from sensors and weather stations to develop an AI-driven irrigation system. You will create a predictive model for water demand based on environmental factors and crop requirements, ensure the system is scalable and responsible, and present your findings to stakeholders in the agricultural sector.

\#\#\# \textbf{Final Challenge}  
Build a predictive model that forecasts water demand for a given crop field based on historical weather, soil moisture, and crop data. Use this model to create a smart irrigation plan, ensuring sustainable water use while maximizing agricultural yield. Present your solution along with key analytics and insights to hypothetical investors and stakeholders.

---

\subsection*{Tasks \& Deliverables}

\begin{taskbox}{Task 1: Understand the Problem}

    {\small \textcolor{gray}{2--3 hrs $\cdot$ Intermediate}}\\[12pt]

    \#\#\#\# \textbf{Scenario}  
You are hired by an AgriTech company to optimize water usage for a wheat field in a semi-arid region. The region frequently experiences unpredictable weather patterns, and soil moisture levels vary significantly across the field.

\#\#\#\# \textbf{Student Assignment}  
\begin{itemize}
    \item Research the principles of precision agriculture and smart irrigation systems.
    \item Study the provided dataset containing historical weather, soil moisture, and crop yield data.
    \item Identify key factors that influence water demand for wheat crops.
\end{itemize}

\#\#\#\# \textbf{Deliverable}  
Submit a 500-word report summarizing:
\begin{itemize}
    \item Key factors affecting water usage in agriculture.
    \item The role of predictive modeling in smart irrigation systems.
    \item Initial insights from the dataset.\par
\end{itemize}

    \vspace{16pt}

    \textcolor{cardBorder}{\hrule height 1pt}\vspace{8pt}

    \small \textcolor{gray}{When you are done with this task, mark it as complete to proceed to the next one.} \hfill \textbf{\textcolor{brandDark}{$\checkmark$ Completed}}

\end{taskbox}

\begin{taskbox}{Task 2: Preprocess Data and Set Up Environment}

    {\small \textcolor{gray}{2--3 hrs $\cdot$ Intermediate}}\\[12pt]

    \#\#\#\# \textbf{Scenario}  
You receive raw sensor data from soil moisture monitors and weather stations. However, the data contains missing values, outliers, and inconsistencies that need to be cleaned before analysis.

\#\#\#\# \textbf{Student Assignment}  
\begin{itemize}
    \item Import the provided dataset into a Jupyter Notebook.
    \item Perform data cleaning: handle missing values, remove outliers, and standardize units (e.g., temperature in °C, rainfall in mm).
    \item Conduct exploratory data analysis (EDA) and visualize key trends in the data.
\end{itemize}

\#\#\#\# \textbf{Deliverable}  
Submit an annotated Jupyter Notebook containing:
\begin{itemize}
    \item Cleaned dataset.
    \item Summary statistics for key variables.
    \item Visualizations (e.g., histograms, scatter plots, and time-series graphs).\par
\end{itemize}

    \vspace{16pt}

    \textcolor{cardBorder}{\hrule height 1pt}\vspace{8pt}

    \small \textcolor{gray}{When you are done with this task, mark it as complete to proceed to the next one.} \hfill \textbf{\textcolor{brandDark}{$\checkmark$ Completed}}

\end{taskbox}

\begin{taskbox}{Task 3: Build Predictive Water Demand Model}

    {\small \textcolor{gray}{2--3 hrs $\cdot$ Intermediate}}\\[12pt]

    \#\#\#\# \textbf{Scenario}  
Using the cleaned data, you will build a predictive model to estimate daily water demand for the wheat field based on weather and soil conditions.

\#\#\#\# \textbf{Student Assignment}  
\begin{itemize}
    \item Split the dataset into training and testing sets.
    \item Train a regression model (e.g., Random Forest or Gradient Boosting) to predict daily water demand.
    \item Evaluate model performance using metrics like RMSE and R².
\end{itemize}

\#\#\#\# \textbf{Deliverable}  
Submit a Python script and report containing:
\begin{itemize}
    \item Model training code.
    \item Evaluation metrics (e.g., RMSE, R²).
    \item Insights into the model’s feature importance and accuracy.\par
\end{itemize}

    \vspace{16pt}

    \textcolor{cardBorder}{\hrule height 1pt}\vspace{8pt}

    \small \textcolor{gray}{When you are done with this task, mark it as complete to proceed to the next one.} \hfill \textbf{\textcolor{brandDark}{$\checkmark$ Completed}}

\end{taskbox}

\begin{taskbox}{Task 4: Use Generative AI for Explanation}

    {\small \textcolor{gray}{2--3 hrs $\cdot$ Intermediate}}\\[12pt]

    \#\#\#\# \textbf{Scenario}  
Your team needs to explain the predictive model to non-technical stakeholders. Use Generative AI tools to create accessible explanations and visualizations.

\#\#\#\# \textbf{Student Assignment}  
\begin{itemize}
    \item Use a Generative AI tool (e.g., ChatGPT or Jasper) to summarize the methodology and results of your predictive model.
    \item Create visualizations to explain how weather and soil factors influence water demand.
    \item Generate a brief video script explaining the smart irrigation solution.
\end{itemize}

\#\#\#\# \textbf{Deliverable}  
Submit:
\begin{itemize}
    \item A 300-word text explanation generated by AI.
    \item Two visualizations explaining the model.
    \item A 2-minute video script for stakeholder presentation.\par
\end{itemize}

    \vspace{16pt}

    \textcolor{cardBorder}{\hrule height 1pt}\vspace{8pt}

    \small \textcolor{gray}{When you are done with this task, mark it as complete to proceed to the next one.} \hfill \textbf{\textcolor{brandDark}{$\checkmark$ Completed}}

\end{taskbox}

\begin{taskbox}{Task 5: Audit for Responsible AI Practices}

    {\small \textcolor{gray}{2--3 hrs $\cdot$ Intermediate}}\\[12pt]

    \#\#\#\# \textbf{Scenario}  
Your irrigation system might be deployed across diverse agricultural regions. Conduct an audit to ensure the predictive model is fair, unbiased, and adaptable to different crops and climates.

\#\#\#\# \textbf{Student Assignment}  
\begin{itemize}
    \item Test the model on a new dataset containing data for a cornfield.
    \item Check for bias or underperformance in certain environmental conditions.
    \item Propose adjustments or improvements to ensure the model’s fairness.
\end{itemize}

\#\#\#\# \textbf{Deliverable}  
Submit a report containing:
\begin{itemize}
    \item Bias audit results (e.g., performance comparison across datasets).
    \item Proposed improvements for model adaptability and fairness.
    \item Justification for Responsible AI practices in agriculture.\par
\end{itemize}

    \vspace{16pt}

    \textcolor{cardBorder}{\hrule height 1pt}\vspace{8pt}

    \small \textcolor{gray}{When you are done with this task, mark it as complete to proceed to the next one.} \hfill \textbf{\textcolor{brandDark}{$\checkmark$ Completed}}

\end{taskbox}

\begin{taskbox}{Task 6: Present Recommendation}

    {\small \textcolor{gray}{2--3 hrs $\cdot$ Intermediate}}\\[12pt]

    \#\#\#\# \textbf{Scenario}  
You are presenting your irrigation solution to a panel of investors and agricultural experts. They need a clear understanding of the technology, its benefits, and its impact on sustainability.

\#\#\#\# \textbf{Student Assignment}  
\begin{itemize}
    \item Create a presentation summarizing your work, including challenges, methodology, results, and recommendations.
    \item Include visuals (e.g., charts, graphs, and images) to support your arguments.
    \item Highlight the economic and environmental benefits of the smart irrigation system.
\end{itemize}

\#\#\#\# \textbf{Deliverable}  
Submit:
\begin{itemize}
    \item A PowerPoint/Google Slides presentation (10 slides max).
    \item A video recording of your presentation (3-5 minutes).
    \item A written summary of key points (400 words).\par
\end{itemize}

    \vspace{16pt}

    \textcolor{cardBorder}{\hrule height 1pt}\vspace{8pt}

    \small \textcolor{gray}{When you are done with this task, mark it as complete to proceed to the next one.} \hfill \textbf{\textcolor{brandDark}{$\checkmark$ Completed}}

\end{taskbox}

\begin{completionbox}
    \textbf{Knowledge Assessment \& Verification:} Complete the online knowledge check, submit your code and presentation files for AI verification, and claim your \textbf{Skillzza Verified Certificate}.
\end{completionbox}

\newpage

% =============================================================================

% SIMULATION 17

% =============================================================================

\section*{Credit Analyst – AI Credit Risk Assessment}

\addcontentsline{toc}{section}{Credit Analyst – AI Credit Risk Assessment}

\begin{metabox}
    \begin{tabularx}{\textwidth}{@{}ll X@{}}
        \textbf{Industry:} & Technology & \textbf{Partner Enterprise:} \textbf{Skillzza Enterprise} \\
        \textbf{Duration:} & 12-18 hrs (Self-paced) & \textbf{Mode:} Individual, task-based simulation \\
        \textbf{Primary Skills:} & \multicolumn{2}{@{}l}{Python, AI, Data Analysis, Problem Solving}
    \end{tabularx}
\end{metabox}

\subsection*{Overview \& Problem Statement}

\#\#\# Scenario  
As a Credit Analyst in a leading financial institution, your role is critical to ensuring responsible lending decisions. You are tasked with analyzing financial profiles, assessing credit risks using AI models, and making informed recommendations for loan approvals or denials. In this simulation, you will step into the shoes of a modern-day credit analyst, leveraging advanced AI tools to streamline and refine credit risk assessment processes.

\#\#\# Mission  
Your mission is to utilize AI-powered techniques to develop a credit risk scoring engine, analyze borrower profiles, and provide transparent explanations for loan decisions. You will also ensure compliance with responsible AI practices and present your recommendations to the financial advisory board for final review.

\#\#\# Final Challenge  
The final challenge involves creating an end-to-end credit risk assessment report for a real-world dataset, including:
1. Financial profile analysis.
2. AI-driven risk scoring with key metrics.
3. Transparent explanation of risk scoring.
4. Recommendations for credit approval or denial.
5. A presentation to simulate board-level decision-making.

---

\subsection*{Tasks \& Deliverables}

\begin{taskbox}{Task 1: Understand}

    {\small \textcolor{gray}{2--3 hrs $\cdot$ Intermediate}}\\[12pt]

    \#\#\#\# Scenario  
You’ve just joined the credit risk team of a financial institution. To start, you need to familiarize yourself with the dataset provided and understand the key variables that influence credit risk scoring.  

\#\#\#\# Student Assignment  
\begin{itemize}
    \item Explore the dataset fields, including `Income`, `Debt`, `Credit\_Score`, `Loan\_Amount`, `Employment\_Status`, and `Loan\_Status`.  
    \item Identify missing values and outliers.  
    \item Research the importance of each variable in credit scoring models.  
\end{itemize}

\#\#\#\# Deliverable  
A report summarizing:  
1. Description of dataset fields.  
2. Identified data issues like missing values or extreme outliers.  
3. Key insights about important variables for risk scoring.\par

    \vspace{16pt}

    \textcolor{cardBorder}{\hrule height 1pt}\vspace{8pt}

    \small \textcolor{gray}{When you are done with this task, mark it as complete to proceed to the next one.} \hfill \textbf{\textcolor{brandDark}{$\checkmark$ Completed}}

\end{taskbox}

\begin{taskbox}{Task 2: Data/Setup}

    {\small \textcolor{gray}{2--3 hrs $\cdot$ Intermediate}}\\[12pt]

    \#\#\#\# Scenario  
The financial institution requires you to preprocess the raw data for analysis. This includes handling missing values, encoding categorical variables, and scaling numerical values.  

\#\#\#\# Student Assignment  
\begin{itemize}
    \item Impute missing values using appropriate techniques (e.g., median for numerical fields, mode for categorical fields).  
    \item Perform one-hot encoding for `Employment\_Status`.  
    \item Scale `Income`, `Debt`, and `Loan\_Amount` using MinMaxScaler.  
\end{itemize}

\#\#\#\# Deliverable  
A clean and preprocessed dataset ready for machine learning, saved as a `.csv` file.\par

    \vspace{16pt}

    \textcolor{cardBorder}{\hrule height 1pt}\vspace{8pt}

    \small \textcolor{gray}{When you are done with this task, mark it as complete to proceed to the next one.} \hfill \textbf{\textcolor{brandDark}{$\checkmark$ Completed}}

\end{taskbox}

\begin{taskbox}{Task 3: Build/Execute}

    {\small \textcolor{gray}{2--3 hrs $\cdot$ Intermediate}}\\[12pt]

    \#\#\#\# Scenario  
Your team needs an AI model to predict credit risk based on borrower profiles. You are tasked with building a machine learning model to classify borrowers as `Low Risk`, `Medium Risk`, or `High Risk`.  

\#\#\#\# Student Assignment  
\begin{itemize}
    \item Split the dataset into training and testing sets (80/20 split).  
    \item Train a decision tree classifier for credit risk assessment.  
    \item Evaluate model performance using accuracy, precision, recall, and F1 score metrics.  
\end{itemize}

\#\#\#\# Deliverable  
1. Trained decision tree model (`.pkl` file).  
2. Model evaluation report including metrics results.\par

    \vspace{16pt}

    \textcolor{cardBorder}{\hrule height 1pt}\vspace{8pt}

    \small \textcolor{gray}{When you are done with this task, mark it as complete to proceed to the next one.} \hfill \textbf{\textcolor{brandDark}{$\checkmark$ Completed}}

\end{taskbox}

\begin{taskbox}{Task 4: GenAI/Explanation}

    {\small \textcolor{gray}{2--3 hrs $\cdot$ Intermediate}}\\[12pt]

    \#\#\#\# Scenario  
The financial institution wants to ensure that the credit scoring model provides transparent and explainable decisions. You are required to use SHAP to interpret the model’s outputs.  

\#\#\#\# Student Assignment  
\begin{itemize}
    \item Use SHAP to explain which features contributed most to the model’s credit risk predictions.  
    \item Generate SHAP visualizations for individual predictions (e.g., force plots).  
    \item Identify any potential biases in the model’s decision-making process.  
\end{itemize}

\#\#\#\# Deliverable  
1. SHAP summary plot.  
2. Force plot for at least three individual predictions.  
3. Bias report highlighting potential concerns.\par

    \vspace{16pt}

    \textcolor{cardBorder}{\hrule height 1pt}\vspace{8pt}

    \small \textcolor{gray}{When you are done with this task, mark it as complete to proceed to the next one.} \hfill \textbf{\textcolor{brandDark}{$\checkmark$ Completed}}

\end{taskbox}

\begin{taskbox}{Task 5: Audit/Responsible AI}

    {\small \textcolor{gray}{2--3 hrs $\cdot$ Intermediate}}\\[12pt]

    \#\#\#\# Scenario  
To comply with responsible AI practices, the institution requests a fairness audit of the model. You need to evaluate whether the model disproportionately impacts specific groups (e.g., based on employment status or income levels).  

\#\#\#\# Student Assignment  
\begin{itemize}
    \item Identify if any demographic group (e.g., `Employment\_Status`) is unfairly penalized by the AI model.  
    \item Calculate metrics such as demographic parity and disparate impact ratio.  
    \item Suggest any necessary adjustments to ensure fairness.  
\end{itemize}

\#\#\#\# Deliverable  
1. Fairness audit report with calculated metrics.  
2. Recommendations for improving model fairness.\par

    \vspace{16pt}

    \textcolor{cardBorder}{\hrule height 1pt}\vspace{8pt}

    \small \textcolor{gray}{When you are done with this task, mark it as complete to proceed to the next one.} \hfill \textbf{\textcolor{brandDark}{$\checkmark$ Completed}}

\end{taskbox}

\begin{taskbox}{Task 6: Present Recommendation}

    {\small \textcolor{gray}{2--3 hrs $\cdot$ Intermediate}}\\[12pt]

    \#\#\#\# Scenario  
The financial advisory board is reviewing your model’s outputs. Prepare a comprehensive credit risk assessment report and presentation to justify your loan recommendations.  

\#\#\#\# Student Assignment  
\begin{itemize}
    \item Create a report summarizing the credit risk scoring methodology, results, and recommendations.  
    \item Prepare a presentation to outline your findings with visuals (charts, plots, SHAP explanations).  
    \item Simulate presenting your findings as if you were advising the board.  
\end{itemize}

\#\#\#\# Deliverable  
1. Credit risk assessment report (Word or PDF).  
2. Slide deck with visuals (PowerPoint or PDF).\par

    \vspace{16pt}

    \textcolor{cardBorder}{\hrule height 1pt}\vspace{8pt}

    \small \textcolor{gray}{When you are done with this task, mark it as complete to proceed to the next one.} \hfill \textbf{\textcolor{brandDark}{$\checkmark$ Completed}}

\end{taskbox}

\begin{completionbox}
    \textbf{Knowledge Assessment \& Verification:} Complete the online knowledge check, submit your code and presentation files for AI verification, and claim your \textbf{Skillzza Verified Certificate}.
\end{completionbox}

\newpage

% =============================================================================

% SIMULATION 18

% =============================================================================

\section*{Microfinance Analyst – AI Loan Eligibility Engine}

\addcontentsline{toc}{section}{Microfinance Analyst – AI Loan Eligibility Engine}

\begin{metabox}
    \begin{tabularx}{\textwidth}{@{}ll X@{}}
        \textbf{Industry:} & Technology & \textbf{Partner Enterprise:} \textbf{Skillzza Enterprise} \\
        \textbf{Duration:} & 12-18 hrs (Self-paced) & \textbf{Mode:} Individual, task-based simulation \\
        \textbf{Primary Skills:} & \multicolumn{2}{@{}l}{Python, AI, Data Analysis, Problem Solving}
    \end{tabularx}
\end{metabox}

\subsection*{Overview \& Problem Statement}

\#\#\# Scenario
Microfinance institutions play a critical role in providing financial access to underserved populations. However, traditional methods of assessing loan eligibility often exclude borrowers with limited credit histories. As a Microfinance Analyst specializing in AI-driven loan eligibility engines, you will help design a data-driven solution that leverages alternative signals (e.g., mobile usage, social behaviors, and spending patterns) to assess borrowers' risk profiles. 

\#\#\# Mission
Your mission is to develop, audit, and deploy a loan eligibility engine powered by AI models. By integrating advanced analytics and responsible AI principles, you will innovate the microfinance process to support inclusive lending decisions while maintaining financial stability for the institution.

\#\#\# Final Challenge
At the end of this simulation, you will present your AI-powered loan eligibility engine to a panel of stakeholders, showcasing insights derived from borrower data, your risk assessment model, and recommendations for ethical implementation.

---

\subsection*{Tasks \& Deliverables}

\begin{taskbox}{Task 1: Understand – Microfinance Domain and Dataset Exploration}

    {\small \textcolor{gray}{2--3 hrs $\cdot$ Intermediate}}\\[12pt]

    \#\#\#\# Scenario:
You are tasked with understanding the basics of microfinance lending and the dataset provided. The borrower dataset contains anonymized information, including demographics, mobile usage, spending patterns, and repayment history.

\#\#\#\# Student Assignment:
Review the provided dataset and documentation. Identify key features and their relevance to loan eligibility. Conduct exploratory data analysis (EDA) to uncover trends, correlations, and anomalies.

\#\#\#\# Deliverable:
\begin{itemize}
    \item A Jupyter Notebook with EDA visualizations.
    \item An insight summary of key trends and findings (300-500 words).\par
\end{itemize}

    \vspace{16pt}

    \textcolor{cardBorder}{\hrule height 1pt}\vspace{8pt}

    \small \textcolor{gray}{When you are done with this task, mark it as complete to proceed to the next one.} \hfill \textbf{\textcolor{brandDark}{$\checkmark$ Completed}}

\end{taskbox}

\begin{taskbox}{Task 2: Data/Setup – Preprocessing Borrower Data}

    {\small \textcolor{gray}{2--3 hrs $\cdot$ Intermediate}}\\[12pt]

    \#\#\#\# Scenario:
The borrower dataset contains missing values, outliers, and categorical data that need to be cleaned and prepared for modeling.

\#\#\#\# Student Assignment:
Preprocess the dataset by:
\begin{itemize}
    \item Handling missing values and outliers.
    \item Encoding categorical variables using one-hot encoding.
    \item Scaling numerical features using MinMaxScaler.
\end{itemize}

\#\#\#\# Deliverable:
\begin{itemize}
    \item A cleaned and preprocessed dataset (.csv or .xlsx format).
    \item Documentation of the preprocessing steps.\par
\end{itemize}

    \vspace{16pt}

    \textcolor{cardBorder}{\hrule height 1pt}\vspace{8pt}

    \small \textcolor{gray}{When you are done with this task, mark it as complete to proceed to the next one.} \hfill \textbf{\textcolor{brandDark}{$\checkmark$ Completed}}

\end{taskbox}

\begin{taskbox}{Task 3: Build/Execute – Develop the Loan Eligibility Model}

    {\small \textcolor{gray}{2--3 hrs $\cdot$ Intermediate}}\\[12pt]

    \#\#\#\# Scenario:
Using the preprocessed data, you will train a machine learning model to classify borrowers as eligible or ineligible for loans. Your focus will be on optimizing the model for fairness and accuracy.

\#\#\#\# Student Assignment:
\begin{itemize}
    \item Split the data into train and test sets.
    \item Train a classification model (e.g., logistic regression, random forest, or XGBoost).
    \item Evaluate performance using precision, recall, and F1-score.
\end{itemize}

\#\#\#\# Deliverable:
\begin{itemize}
    \item Python script for model training and evaluation.
    \item A report comparing the performance of different models.\par
\end{itemize}

    \vspace{16pt}

    \textcolor{cardBorder}{\hrule height 1pt}\vspace{8pt}

    \small \textcolor{gray}{When you are done with this task, mark it as complete to proceed to the next one.} \hfill \textbf{\textcolor{brandDark}{$\checkmark$ Completed}}

\end{taskbox}

\begin{taskbox}{Task 4: GenAI/Explanation – Explain Loan Decisions Using Generative AI}

    {\small \textcolor{gray}{2--3 hrs $\cdot$ Intermediate}}\\[12pt]

    \#\#\#\# Scenario:
Stakeholders in microfinance institutions need transparency into how the AI model makes decisions. You will use Generative AI tools to create clear, non-technical explanations for loan eligibility scores.

\#\#\#\# Student Assignment:
\begin{itemize}
    \item Use SHAP to generate explanations for individual predictions.
    \item Write a Generative AI prompt to create a draft of a client-facing explanation for a sample loan decision.
\end{itemize}

\#\#\#\# Deliverable:
\begin{itemize}
    \item SHAP visualization for select borrowers.
    \item Draft of client-facing explanation (200 words).\par
\end{itemize}

    \vspace{16pt}

    \textcolor{cardBorder}{\hrule height 1pt}\vspace{8pt}

    \small \textcolor{gray}{When you are done with this task, mark it as complete to proceed to the next one.} \hfill \textbf{\textcolor{brandDark}{$\checkmark$ Completed}}

\end{taskbox}

\begin{taskbox}{Task 5: Audit/Responsible AI – Bias Detection and Fairness Evaluation}

    {\small \textcolor{gray}{2--3 hrs $\cdot$ Intermediate}}\\[12pt]

    \#\#\#\# Scenario:
Addressing bias in AI models is critical for ethical lending. You will audit the loan eligibility model for fairness across demographic groups (e.g., gender, income level).

\#\#\#\# Student Assignment:
\begin{itemize}
    \item Use FairLearn or similar tools to evaluate bias in the model.
    \item Suggest actionable steps to mitigate any detected bias.
\end{itemize}

\#\#\#\# Deliverable:
\begin{itemize}
    \item Bias audit report with visualizations.
    \item Recommendations for improving fairness.\par
\end{itemize}

    \vspace{16pt}

    \textcolor{cardBorder}{\hrule height 1pt}\vspace{8pt}

    \small \textcolor{gray}{When you are done with this task, mark it as complete to proceed to the next one.} \hfill \textbf{\textcolor{brandDark}{$\checkmark$ Completed}}

\end{taskbox}

\begin{taskbox}{Task 6: Present Recommendation – Stakeholder Presentation}

    {\small \textcolor{gray}{2--3 hrs $\cdot$ Intermediate}}\\[12pt]

    \#\#\#\# Scenario:
Microfinance stakeholders are evaluating your AI engine. You need to present your findings, model performance, and ethical considerations.

\#\#\#\# Student Assignment:
\begin{itemize}
    \item Create a 5-slide presentation summarizing your work:
    \item Problem overview.
    \item Key insights from the data.
    \item Model performance and fairness evaluation.
    \item Recommendations for implementation.
    \item Deliver your presentation in a simulated stakeholder meeting.
\end{itemize}

\#\#\#\# Deliverable:
\begin{itemize}
    \item Presentation deck (.ppt or .pdf format).
    \item Script for presenting key points (300-500 words).\par
\end{itemize}

    \vspace{16pt}

    \textcolor{cardBorder}{\hrule height 1pt}\vspace{8pt}

    \small \textcolor{gray}{When you are done with this task, mark it as complete to proceed to the next one.} \hfill \textbf{\textcolor{brandDark}{$\checkmark$ Completed}}

\end{taskbox}

\begin{completionbox}
    \textbf{Knowledge Assessment \& Verification:} Complete the online knowledge check, submit your code and presentation files for AI verification, and claim your \textbf{Skillzza Verified Certificate}.
\end{completionbox}

\newpage

% =============================================================================

% SIMULATION 19

% =============================================================================

\section*{Financial Analyst – AI-Powered Financial Modelling}

\addcontentsline{toc}{section}{Financial Analyst – AI-Powered Financial Modelling}

\begin{metabox}
    \begin{tabularx}{\textwidth}{@{}ll X@{}}
        \textbf{Industry:} & Technology & \textbf{Partner Enterprise:} \textbf{Skillzza Enterprise} \\
        \textbf{Duration:} & 12-18 hrs (Self-paced) & \textbf{Mode:} Individual, task-based simulation \\
        \textbf{Primary Skills:} & \multicolumn{2}{@{}l}{Python, AI, Data Analysis, Problem Solving}
    \end{tabularx}
\end{metabox}

\subsection*{Overview \& Problem Statement}

\#\#\# \textbf{Scenario:}  
You are hired as a Financial Analyst at a fast-growing FinTech startup. The company wants to leverage AI for building dynamic financial models that can predict profitability under various scenarios. Your role is to analyze the company's financial statements, build AI-powered predictive models, and translate outputs into actionable insights for senior management.

\#\#\# \textbf{Mission:}  
Your ultimate mission is to design, test, and deploy an AI-powered financial model that can forecast revenue, expenses, and profitability under different market conditions. You'll present your findings and recommendations to the management team to help drive strategic decisions.

\#\#\# \textbf{Final Challenge:}  
Build a scenario-based financial model using AI, simulate revenue and expense forecasts under multiple conditions, and generate a management report with actionable insights.

---

\subsection*{Tasks \& Deliverables}

\begin{taskbox}{Task 1: Understand the Financial Fundamentals}

    {\small \textcolor{gray}{2--3 hrs $\cdot$ Intermediate}}\\[12pt]

      
Your manager hands you a spreadsheet containing the company’s last three years of financial statements. You are asked to analyze the key metrics and trends to set the foundation for predictive modeling.  


\vspace{8pt}
{\large \textbf{\textcolor{brandLight}{What you'll learn}}}\\[4pt]  
\begin{itemize}
    \item Review the Income Statement, Balance Sheet, and Cash Flow Statement.  
    \item Identify trends in revenue growth, operating expenses, and net profit margin.  
    \item Highlight any anomalies or irregularities in the data.  
\end{itemize}

\vspace{8pt}
{\large \textbf{\textcolor{brandLight}{What you'll do}}}\\[4pt]  
Submit a summary report highlighting 5 key financial trends and observations. Use graphs and charts for visualization.\par

    \vspace{16pt}

    \textcolor{cardBorder}{\hrule height 1pt}\vspace{8pt}

    \small \textcolor{gray}{When you are done with this task, mark it as complete to proceed to the next one.} \hfill \textbf{\textcolor{brandDark}{$\checkmark$ Completed}}

\end{taskbox}

\begin{taskbox}{Task 2: Prepare the Dataset for AI Modeling}

    {\small \textcolor{gray}{2--3 hrs $\cdot$ Intermediate}}\\[12pt]

      
You are tasked with cleaning and structuring the financial data for AI analysis. The raw dataset contains missing values, duplicate entries, and inconsistent formatting.  


\vspace{8pt}
{\large \textbf{\textcolor{brandLight}{What you'll learn}}}\\[4pt]  
\begin{itemize}
    \item Clean the dataset to remove errors and handle missing values.  
    \item Structure the dataset to ensure compatibility with AI tools.  
    \item Extract key features (e.g., revenue, expenses, EBITDA) required for predictive modeling.  
\end{itemize}

\vspace{8pt}
{\large \textbf{\textcolor{brandLight}{What you'll do}}}\\[4pt]  
Submit the cleaned, structured dataset along with a brief report describing the steps you took.\par

    \vspace{16pt}

    \textcolor{cardBorder}{\hrule height 1pt}\vspace{8pt}

    \small \textcolor{gray}{When you are done with this task, mark it as complete to proceed to the next one.} \hfill \textbf{\textcolor{brandDark}{$\checkmark$ Completed}}

\end{taskbox}

\begin{taskbox}{Task 3: Build and Execute AI-Powered Financial Model}

    {\small \textcolor{gray}{2--3 hrs $\cdot$ Intermediate}}\\[12pt]

      
Using the cleaned dataset, you need to train a predictive model to forecast revenue and expenses for the next fiscal year.  


\vspace{8pt}
{\large \textbf{\textcolor{brandLight}{What you'll learn}}}\\[4pt]  
\begin{itemize}
    \item Use regression algorithms from Scikit-learn to build the model.  
    \item Train the model with the historical financial data provided.  
    \item Evaluate the model’s accuracy using metrics like RMSE (Root Mean Squared Error).  
\end{itemize}

\vspace{8pt}
{\large \textbf{\textcolor{brandLight}{What you'll do}}}\\[4pt]  
Submit the Python code for the trained model and a summary of its performance metrics.\par

    \vspace{16pt}

    \textcolor{cardBorder}{\hrule height 1pt}\vspace{8pt}

    \small \textcolor{gray}{When you are done with this task, mark it as complete to proceed to the next one.} \hfill \textbf{\textcolor{brandDark}{$\checkmark$ Completed}}

\end{taskbox}

\begin{taskbox}{Task 4: Leverage Generative AI for Management Insights}

    {\small \textcolor{gray}{2--3 hrs $\cdot$ Intermediate}}\\[12pt]

      
Once the AI model generates forecasts, the management team requires a concise report summarizing the insights. You will use Generative AI tools to automate the creation of this report.  


\vspace{8pt}
{\large \textbf{\textcolor{brandLight}{What you'll learn}}}\\[4pt]  
\begin{itemize}
    \item Input the model’s outputs into a Generative AI tool like ChatGPT.  
    \item Generate a management-friendly summary that includes graphs, recommendations, and insights.  
    \item Ensure the report is free of biases and aligns with Responsible AI principles.  
\end{itemize}

\vspace{8pt}
{\large \textbf{\textcolor{brandLight}{What you'll do}}}\\[4pt]  
Submit the AI-generated management report in PDF format.\par

    \vspace{16pt}

    \textcolor{cardBorder}{\hrule height 1pt}\vspace{8pt}

    \small \textcolor{gray}{When you are done with this task, mark it as complete to proceed to the next one.} \hfill \textbf{\textcolor{brandDark}{$\checkmark$ Completed}}

\end{taskbox}

\begin{taskbox}{Task 5: Audit Model Performance and Responsible AI}

    {\small \textcolor{gray}{2--3 hrs $\cdot$ Intermediate}}\\[12pt]

      
The management team raises concerns about the ethical implications and reliability of AI-driven forecasts. You are asked to audit the model for Responsible AI principles.  


\vspace{8pt}
{\large \textbf{\textcolor{brandLight}{What you'll learn}}}\\[4pt]  
\begin{itemize}
    \item Evaluate the model’s performance under different scenarios for bias and accuracy.  
    \item Document any limitations or risks associated with the model.  
    \item Suggest improvements for ensuring Responsible AI compliance.  
\end{itemize}

\vspace{8pt}
{\large \textbf{\textcolor{brandLight}{What you'll do}}}\\[4pt]  
Submit an audit report summarizing biases, risks, and proposed improvements.\par

    \vspace{16pt}

    \textcolor{cardBorder}{\hrule height 1pt}\vspace{8pt}

    \small \textcolor{gray}{When you are done with this task, mark it as complete to proceed to the next one.} \hfill \textbf{\textcolor{brandDark}{$\checkmark$ Completed}}

\end{taskbox}

\begin{taskbox}{Task 6: Present Recommendations to Senior Management}

    {\small \textcolor{gray}{2--3 hrs $\cdot$ Intermediate}}\\[12pt]

      
You need to deliver a virtual presentation summarizing key financial insights, the predictive model’s outputs, and actionable recommendations for management.  


\vspace{8pt}
{\large \textbf{\textcolor{brandLight}{What you'll learn}}}\\[4pt]  
\begin{itemize}
    \item Create a slide deck summarizing financial trends, predictive model insights, and scenario analysis.  
    \item Record a 5-minute video explaining your recommendations to senior management.  
\end{itemize}

\vspace{8pt}
{\large \textbf{\textcolor{brandLight}{What you'll do}}}\\[4pt]  
Submit the slide deck and video presentation.\par

    \vspace{16pt}

    \textcolor{cardBorder}{\hrule height 1pt}\vspace{8pt}

    \small \textcolor{gray}{When you are done with this task, mark it as complete to proceed to the next one.} \hfill \textbf{\textcolor{brandDark}{$\checkmark$ Completed}}

\end{taskbox}

\begin{completionbox}
    \textbf{Knowledge Assessment \& Verification:} Complete the online knowledge check, submit your code and presentation files for AI verification, and claim your \textbf{Skillzza Verified Certificate}.
\end{completionbox}

\newpage

% =============================================================================

% SIMULATION 20

% =============================================================================

\section*{Advisory Analyst – Business Turnaround Case}

\addcontentsline{toc}{section}{Advisory Analyst – Business Turnaround Case}

\begin{metabox}
    \begin{tabularx}{\textwidth}{@{}ll X@{}}
        \textbf{Industry:} & Technology & \textbf{Partner Enterprise:} \textbf{Skillzza Enterprise} \\
        \textbf{Duration:} & 12-18 hrs (Self-paced) & \textbf{Mode:} Individual, task-based simulation \\
        \textbf{Primary Skills:} & \multicolumn{2}{@{}l}{Python, AI, Data Analysis, Problem Solving}
    \end{tabularx}
\end{metabox}

\subsection*{Overview \& Problem Statement}

\#\#\# \textbf{Scenario:}  
You are hired as an Advisory Analyst by a boutique consulting firm specializing in business turnarounds for struggling companies. Your client, \textbf{AlphaTech Solutions}, is a mid-sized technology company experiencing declining revenues, profitability, and employee morale. The CEO suspects operational inefficiencies, weak market strategy, and financial mismanagement are to blame. You are tasked with analyzing the situation and delivering actionable recommendations to turn the business around within six months.  

\#\#\# \textbf{Mission:}  
Your mission is to:  
1. Perform a deep dive into AlphaTech’s financial and operational data.  
2. Identify root causes of declining performance using structured analysis.  
3. Create a comprehensive turnaround plan focusing on operational efficiency, financial recovery, and strategic growth.  

\#\#\# \textbf{Final Challenge:}  
Present your findings and recommendations to AlphaTech’s executive board as a structured turnaround roadmap document and a 10-slide presentation. Your solution will need to balance quick wins with long-term strategies while demonstrating an understanding of financial modeling, operational scalability, and strategic advisory.  

---

\subsection*{Tasks \& Deliverables}

\begin{taskbox}{Task 1: Understand – Business Diagnostic}

    {\small \textcolor{gray}{2--3 hrs $\cdot$ Intermediate}}\\[12pt]

    \#\#\#\#   
The CEO of AlphaTech Solutions has provided you with an overview of the company’s challenges. You also have access to a dataset containing financial statements for the last three years, operational KPIs, and market reports.  

\#\#\#\# \\[10pt]
\vspace{8pt}
{\large \textbf{\textcolor{brandLight}{What you'll learn}}}\\[4pt]  
\begin{itemize}
    \item Review the company profile and the provided dataset.  
    \item Identify key areas of concern from the initial diagnostic (e.g., declining profitability, high operating costs).  
    \item Summarize challenges into three main categories: financial, operational, strategic.  
\end{itemize}

\#\#\#\# \vspace{8pt}
{\large \textbf{\textcolor{brandLight}{What you'll do}}}\\[4pt]  
A one-page summary identifying the key areas of concern and challenges faced by AlphaTech Solutions.\par

    \vspace{16pt}

    \textcolor{cardBorder}{\hrule height 1pt}\vspace{8pt}

    \small \textcolor{gray}{When you are done with this task, mark it as complete to proceed to the next one.} \hfill \textbf{\textcolor{brandDark}{$\checkmark$ Completed}}

\end{taskbox}

\begin{taskbox}{Task 2: Data/Setup – Financial and Operational Analysis}

    {\small \textcolor{gray}{2--3 hrs $\cdot$ Intermediate}}\\[12pt]

    \#\#\#\#   
You must analyze AlphaTech’s financial and operational data to uncover trends and anomalies. The dataset includes:  
\begin{itemize}
    \item Financial statements (Income Statement, Balance Sheet, Cash Flow).  
    \item Operational KPIs (Customer acquisition cost, churn rate, productivity metrics).  
\end{itemize}

\#\#\#\# \\[10pt]
\vspace{8pt}
{\large \textbf{\textcolor{brandLight}{What you'll learn}}}\\[4pt]  
\begin{itemize}
    \item Calculate financial ratios: profitability, liquidity, and efficiency.  
    \item Identify trends (e.g., declining gross margin, increasing operating expenses).  
    \item Benchmark operational KPIs against industry standards.  
\end{itemize}

\#\#\#\# \vspace{8pt}
{\large \textbf{\textcolor{brandLight}{What you'll do}}}\\[4pt]  
An Excel file or Google Sheet with calculated financial ratios, operational benchmarks, and trend insights.\par

    \vspace{16pt}

    \textcolor{cardBorder}{\hrule height 1pt}\vspace{8pt}

    \small \textcolor{gray}{When you are done with this task, mark it as complete to proceed to the next one.} \hfill \textbf{\textcolor{brandDark}{$\checkmark$ Completed}}

\end{taskbox}

\begin{taskbox}{Task 3: Build/Execute – Root Cause Analysis}

    {\small \textcolor{gray}{2--3 hrs $\cdot$ Intermediate}}\\[12pt]

    \#\#\#\#   
Using your findings, you must perform a root cause analysis to pinpoint the underlying reasons for AlphaTech’s declining performance.  

\#\#\#\# \\[10pt]
\vspace{8pt}
{\large \textbf{\textcolor{brandLight}{What you'll learn}}}\\[4pt]  
\begin{itemize}
    \item Apply frameworks like Ishikawa diagrams or the 5 Whys technique to break down root causes.  
    \item Categorize findings into controllable (e.g., internal operations) and uncontrollable (e.g., market trends) factors.  
\end{itemize}

\#\#\#\# \vspace{8pt}
{\large \textbf{\textcolor{brandLight}{What you'll do}}}\\[4pt]  
A structured document or slide deck outlining:  
1. The root causes.  
2. The methodology used for identification.\par

    \vspace{16pt}

    \textcolor{cardBorder}{\hrule height 1pt}\vspace{8pt}

    \small \textcolor{gray}{When you are done with this task, mark it as complete to proceed to the next one.} \hfill \textbf{\textcolor{brandDark}{$\checkmark$ Completed}}

\end{taskbox}

\begin{taskbox}{Task 4: GenAI/Explanation – Turnaround Strategy Planning}

    {\small \textcolor{gray}{2--3 hrs $\cdot$ Intermediate}}\\[12pt]

    \#\#\#\#   
AlphaTech’s CEO has requested a detailed turnaround strategy that includes quick wins and long-term initiatives. You need to draft this strategy.  

\#\#\#\# \\[10pt]
\vspace{8pt}
{\large \textbf{\textcolor{brandLight}{What you'll learn}}}\\[4pt]  
\begin{itemize}
    \item Use Generative AI tools (e.g., ChatGPT or Bard) to brainstorm innovative solutions.  
    \item Create a three-phase turnaround plan: stabilize finances, optimize operations, and drive growth.  
    \item Prioritize initiatives based on their ROI and implementation feasibility.  
\end{itemize}

\#\#\#\# \vspace{8pt}
{\large \textbf{\textcolor{brandLight}{What you'll do}}}\\[4pt]  
A three-phase turnaround strategy document with a detailed explanation of each phase, including:  
\begin{itemize}
    \item Key actions.  
    \item Required resources.  
    \item Expected outcomes.\par
\end{itemize}

    \vspace{16pt}

    \textcolor{cardBorder}{\hrule height 1pt}\vspace{8pt}

    \small \textcolor{gray}{When you are done with this task, mark it as complete to proceed to the next one.} \hfill \textbf{\textcolor{brandDark}{$\checkmark$ Completed}}

\end{taskbox}

\begin{taskbox}{Task 5: Audit/Responsible Advisory}

    {\small \textcolor{gray}{2--3 hrs $\cdot$ Intermediate}}\\[12pt]

    \#\#\#\#   
Your recommendations must align with ethical business practices and consider stakeholder impacts. The CEO has asked you to ensure your proposals are socially and environmentally responsible.  

\#\#\#\# \\[10pt]
\vspace{8pt}
{\large \textbf{\textcolor{brandLight}{What you'll learn}}}\\[4pt]  
\begin{itemize}
    \item Audit the turnaround plan for ethical considerations (e.g., employee layoffs, environmental sustainability).  
    \item Modify recommendations to balance profitability with stakeholder well-being.  
\end{itemize}

\#\#\#\# \vspace{8pt}
{\large \textbf{\textcolor{brandLight}{What you'll do}}}\\[4pt]  
An updated turnaround strategy document with a section on ethical considerations and adjustments.\par

    \vspace{16pt}

    \textcolor{cardBorder}{\hrule height 1pt}\vspace{8pt}

    \small \textcolor{gray}{When you are done with this task, mark it as complete to proceed to the next one.} \hfill \textbf{\textcolor{brandDark}{$\checkmark$ Completed}}

\end{taskbox}

\begin{taskbox}{Task 6: Present Recommendation – Final Challenge}

    {\small \textcolor{gray}{2--3 hrs $\cdot$ Intermediate}}\\[12pt]

    \#\#\#\#   
You are presenting your turnaround proposal to AlphaTech’s executive board. The presentation must be concise, visually engaging, and persuasive.  

\#\#\#\# \\[10pt]
\vspace{8pt}
{\large \textbf{\textcolor{brandLight}{What you'll learn}}}\\[4pt]  
\begin{itemize}
    \item Create a 10-slide presentation summarizing your analysis, findings, and recommendations.  
    \item Include data visualizations, ROI projections, and a clear implementation timeline.  
\end{itemize}

\#\#\#\# \vspace{8pt}
{\large \textbf{\textcolor{brandLight}{What you'll do}}}\\[4pt]  
A presentation deck in PowerPoint or Google Slides format.\par

    \vspace{16pt}

    \textcolor{cardBorder}{\hrule height 1pt}\vspace{8pt}

    \small \textcolor{gray}{When you are done with this task, mark it as complete to proceed to the next one.} \hfill \textbf{\textcolor{brandDark}{$\checkmark$ Completed}}

\end{taskbox}

\begin{completionbox}
    \textbf{Knowledge Assessment \& Verification:} Complete the online knowledge check, submit your code and presentation files for AI verification, and claim your \textbf{Skillzza Verified Certificate}.
\end{completionbox}

\newpage

% =============================================================================

% SIMULATION 21

% =============================================================================

\section*{Insurance AI Analyst – Claims Risk Prediction}

\addcontentsline{toc}{section}{Insurance AI Analyst – Claims Risk Prediction}

\begin{metabox}
    \begin{tabularx}{\textwidth}{@{}ll X@{}}
        \textbf{Industry:} & Technology & \textbf{Partner Enterprise:} \textbf{Skillzza Enterprise} \\
        \textbf{Duration:} & 12-18 hrs (Self-paced) & \textbf{Mode:} Individual, task-based simulation \\
        \textbf{Primary Skills:} & \multicolumn{2}{@{}l}{Python, AI, Data Analysis, Problem Solving}
    \end{tabularx}
\end{metabox}

\subsection*{Overview \& Problem Statement}

\#\#\# \textbf{Scenario}
The insurance industry is rapidly embracing AI to reduce fraud, predict risks, and improve claims management. Your role as an \textbf{Insurance AI Analyst} is to assess potential risks in claims, identify fraudulent patterns, and develop actionable insights to improve underwriting and claims processing. 

You have been hired by \textbf{SecureSure Insurance}, a mid-sized insurance company that provides auto, health, and property insurance. The company faces challenges with rising claim frauds and wants to leverage AI to predict risk levels associated with incoming claims. Your mission is to analyze historical claims data, build a risk prediction model, and recommend actionable interventions to minimize fraud while ensuring claims are processed efficiently.

\#\#\# \textbf{Mission}
\begin{itemize}
    \item Use historical claims data to identify key risk factors.
    \item Build an AI model to predict the risk score for new claims.
    \item Provide actionable insights to reduce fraudulent claims and improve operations.
\end{itemize}

\#\#\# \textbf{Final Challenge}
Design and present a comprehensive report that includes:
\begin{itemize}
    \item Analysis of risk patterns in the claims data.
    \item A trained and evaluated risk prediction model.
    \item A summary of interventions to reduce fraudulent claims and optimize claims processing.
\end{itemize}

---

\subsection*{Tasks \& Deliverables}

\begin{taskbox}{Task 1: Understand the Insurance Claims Scenario}

    {\small \textcolor{gray}{2--3 hrs $\cdot$ Intermediate}}\\[12pt]

    \#\#\#\# \textbf{Scenario} 
You’ve been onboarded as an Insurance AI Analyst. Your first task is to understand the problem statement and familiarize yourself with the claims data. SecureSure Insurance has provided a dataset containing historical claims information, some of which are labeled as fraudulent. 

\#\#\#\# \textbf{Student Assignment}
\begin{itemize}
    \item Review the dataset fields, including policy details, claim amounts, claimant characteristics, and fraud indicators.
    \item Identify any missing values, categorical fields, and numeric fields.
    \item Document the problem statement and the project’s goals.
\end{itemize}

\#\#\#\# \textbf{Deliverable}
\begin{itemize}
    \item A 1-2 page document summarizing:
    \item Key fields in the dataset.
    \item Challenges in claims risk prediction.
    \item Goals of your analysis.\par
\end{itemize}

    \vspace{16pt}

    \textcolor{cardBorder}{\hrule height 1pt}\vspace{8pt}

    \small \textcolor{gray}{When you are done with this task, mark it as complete to proceed to the next one.} \hfill \textbf{\textcolor{brandDark}{$\checkmark$ Completed}}

\end{taskbox}

\begin{taskbox}{Task 2: Prepare and Explore the Claims Data}

    {\small \textcolor{gray}{2--3 hrs $\cdot$ Intermediate}}\\[12pt]

    \#\#\#\# \textbf{Scenario}
The dataset has inconsistencies such as missing values, outliers, and mixed data types. Your task is to clean and preprocess the dataset and perform exploratory data analysis (EDA) to uncover trends and patterns.

\#\#\#\# \textbf{Student Assignment}
\begin{itemize}
    \item Handle missing values using imputation techniques.
    \item Encode categorical variables (e.g., One-Hot Encoding, Label Encoding).
    \item Visualize key distributions (e.g., claims amount, ages, policy types) and correlations.
    \item Identify any anomalies and key trends in fraudulent claims.
\end{itemize}

\#\#\#\# \textbf{Deliverable}
\begin{itemize}
    \item A Jupyter Notebook with:
    \item Data cleaning and preprocessing steps.
    \item Visualizations of key data distributions.
    \item A summary of insights from the EDA.\par
\end{itemize}

    \vspace{16pt}

    \textcolor{cardBorder}{\hrule height 1pt}\vspace{8pt}

    \small \textcolor{gray}{When you are done with this task, mark it as complete to proceed to the next one.} \hfill \textbf{\textcolor{brandDark}{$\checkmark$ Completed}}

\end{taskbox}

\begin{taskbox}{Task 3: Build and Train the Risk Prediction Model}

    {\small \textcolor{gray}{2--3 hrs $\cdot$ Intermediate}}\\[12pt]

    \#\#\#\# \textbf{Scenario}
SecureSure Insurance wants you to develop a machine learning model to predict the risk of new claims. Claims with higher risk scores will undergo additional scrutiny for fraud detection.

\#\#\#\# \textbf{Student Assignment}
\begin{itemize}
    \item Split the dataset into training and testing sets.
    \item Train a machine learning model (e.g., Logistic Regression, Random Forest, or XGBoost) to predict claim risk.
    \item Optimize the model using hyperparameter tuning.
    \item Evaluate the model using metrics like accuracy, precision, recall, F1-score, and AUC-ROC.
\end{itemize}

\#\#\#\# \textbf{Deliverable}
\begin{itemize}
    \item A trained and saved model file.
    \item A Jupyter Notebook containing the model training and evaluation process.\par
\end{itemize}

    \vspace{16pt}

    \textcolor{cardBorder}{\hrule height 1pt}\vspace{8pt}

    \small \textcolor{gray}{When you are done with this task, mark it as complete to proceed to the next one.} \hfill \textbf{\textcolor{brandDark}{$\checkmark$ Completed}}

\end{taskbox}

\begin{taskbox}{Task 4: Apply Explainable AI for Risk Interpretation}

    {\small \textcolor{gray}{2--3 hrs $\cdot$ Intermediate}}\\[12pt]

    \#\#\#\# \textbf{Scenario}
The underwriting team at SecureSure Insurance has requested an explanation of how the AI model determines risk scores, so they can justify decisions to stakeholders and regulators.

\#\#\#\# \textbf{Student Assignment}
\begin{itemize}
    \item Use the SHAP library to interpret the predictions of your risk model.
    \item Identify the most important features influencing the risk scores.
    \item Create a visualization (e.g., SHAP summary plot) to communicate top risk factors.
\end{itemize}

\#\#\#\# \textbf{Deliverable}
\begin{itemize}
    \item A Jupyter Notebook with SHAP analysis and visualizations.
    \item A brief report describing the top features and their impact on risk predictions.\par
\end{itemize}

    \vspace{16pt}

    \textcolor{cardBorder}{\hrule height 1pt}\vspace{8pt}

    \small \textcolor{gray}{When you are done with this task, mark it as complete to proceed to the next one.} \hfill \textbf{\textcolor{brandDark}{$\checkmark$ Completed}}

\end{taskbox}

\begin{taskbox}{Task 5: Audit Model Fairness and Ensure Responsible AI}

    {\small \textcolor{gray}{2--3 hrs $\cdot$ Intermediate}}\\[12pt]

    \#\#\#\# \textbf{Scenario}
To meet regulatory requirements and ensure fairness, SecureSure Insurance has asked you to audit the model for biases. For example, does the model unfairly assign higher risk scores based on gender or income levels?

\#\#\#\# \textbf{Student Assignment}
\begin{itemize}
    \item Analyze the model’s predictions for potential biases across sensitive attributes (e.g., gender, age group, income).
    \item Calculate fairness metrics such as Demographic Parity and Equal Opportunity Difference.
    \item Suggest improvements to mitigate any detected biases.
\end{itemize}

\#\#\#\# \textbf{Deliverable}
\begin{itemize}
    \item A Jupyter Notebook with:
    \item Fairness analysis and metrics.
    \item Suggestions for bias mitigation.\par
\end{itemize}

    \vspace{16pt}

    \textcolor{cardBorder}{\hrule height 1pt}\vspace{8pt}

    \small \textcolor{gray}{When you are done with this task, mark it as complete to proceed to the next one.} \hfill \textbf{\textcolor{brandDark}{$\checkmark$ Completed}}

\end{taskbox}

\begin{taskbox}{Task 6: Present Recommendations to Stakeholders}

    {\small \textcolor{gray}{2--3 hrs $\cdot$ Intermediate}}\\[12pt]

    \#\#\#\# \textbf{Scenario}
The leadership team at SecureSure Insurance will use your findings to make strategic decisions. Your job is to present your results, model performance, key risk factors, and recommendations for operational changes to reduce fraudulent claims.

\#\#\#\# \textbf{Student Assignment}
\begin{itemize}
    \item Create a presentation summarizing:
    \item Key insights from EDA.
    \item Model performance and explainability results.
    \item Identified risk patterns and interventions to reduce fraud.
    \item Include visualizations and key metrics to support your recommendations.
\end{itemize}

\#\#\#\# \textbf{Deliverable}
\begin{itemize}
    \item A slide deck (e.g., PowerPoint, Google Slides) with:
    \item Key findings and visualizations.
    \item Model results and explainability insights.
    \item Recommendations for improving claims management.\par
\end{itemize}

    \vspace{16pt}

    \textcolor{cardBorder}{\hrule height 1pt}\vspace{8pt}

    \small \textcolor{gray}{When you are done with this task, mark it as complete to proceed to the next one.} \hfill \textbf{\textcolor{brandDark}{$\checkmark$ Completed}}

\end{taskbox}

\begin{completionbox}
    \textbf{Knowledge Assessment \& Verification:} Complete the online knowledge check, submit your code and presentation files for AI verification, and claim your \textbf{Skillzza Verified Certificate}.
\end{completionbox}

\newpage

% =============================================================================

% SIMULATION 22

% =============================================================================

\section*{Insurance Fraud Analyst – Detect Suspicious Claims}

\addcontentsline{toc}{section}{Insurance Fraud Analyst – Detect Suspicious Claims}

\begin{metabox}
    \begin{tabularx}{\textwidth}{@{}ll X@{}}
        \textbf{Industry:} & Technology & \textbf{Partner Enterprise:} \textbf{Skillzza Enterprise} \\
        \textbf{Duration:} & 12-18 hrs (Self-paced) & \textbf{Mode:} Individual, task-based simulation \\
        \textbf{Primary Skills:} & \multicolumn{2}{@{}l}{Python, AI, Data Analysis, Problem Solving}
    \end{tabularx}
\end{metabox}

\subsection*{Overview \& Problem Statement}

\#\#\# \textbf{Scenario}  
You are hired as an Insurance Fraud Analyst at an InsurTech company specializing in streamlining claims processing and fraud detection. Insurance fraud costs the industry billions annually, and your role is critical in helping insurers identify suspicious claims through advanced analytics and AI techniques. You’ll step into the shoes of a fraud investigator, working on raw claims data to build, refine, and interpret AI anomaly detection models that flag potentially fraudulent claims.

\#\#\# \textbf{Mission}  
Your mission is to help the company detect fraudulent claims in its database using anomaly detection techniques and AI-driven insights. You will analyze patterns, identify red flags, and recommend investigative strategies to mitigate financial risks.

\#\#\# \textbf{Final Challenge}  
By the end of this simulation, you will present an evidence-based report outlining the suspicious claims detected, the methodology used, and actionable recommendations for the investigation team. Your insights will contribute to the company’s fraud prevention strategy.

---

\subsection*{Tasks \& Deliverables}

\begin{taskbox}{Task 1: Understand the Role and Data}

    {\small \textcolor{gray}{2--3 hrs $\cdot$ Intermediate}}\\[12pt]

    \#\#\#\# \textbf{Scenario}  
The company's fraud detection team receives thousands of claims daily. Your first task is to familiarize yourself with the claims dataset and understand the key fraud indicators.  

\#\#\#\# \textbf{Student Assignment}  
\begin{itemize}
    \item Review a sample insurance claims dataset containing fields such as `Claim\_ID`, `Policy\_Holder\_Age`, `Incident\_Date`, `Claim\_Amount`, `Claim\_Type`, and `Claim\_Status`.  
    \item Research common fraud schemes (e.g., duplicate claims, exaggerated damages).  
    \item Summarize the role of AI and analytics in fraud detection.  
\end{itemize}

\#\#\#\# \textbf{Deliverable}  
Submit a 300-word summary explaining the dataset, fraud schemes, and the role of analytics in detecting suspicious claims.\par

    \vspace{16pt}

    \textcolor{cardBorder}{\hrule height 1pt}\vspace{8pt}

    \small \textcolor{gray}{When you are done with this task, mark it as complete to proceed to the next one.} \hfill \textbf{\textcolor{brandDark}{$\checkmark$ Completed}}

\end{taskbox}

\begin{taskbox}{Task 2: Data Exploration and Setup}

    {\small \textcolor{gray}{2--3 hrs $\cdot$ Intermediate}}\\[12pt]

    \#\#\#\# \textbf{Scenario}  
You need to prepare the dataset for analysis by cleaning, transforming, and exploring its features. This step ensures the data is ready for anomaly detection.  

\#\#\#\# \textbf{Student Assignment}  
\begin{itemize}
    \item Clean the dataset by handling missing values and formatting inconsistencies.  
    \item Perform exploratory data analysis (EDA) to identify patterns and potential fraud indicators (e.g., unusually high claim amounts).  
    \item Visualize distributions and correlations using Matplotlib or Seaborn.  
\end{itemize}

\#\#\#\# \textbf{Deliverable}  
Submit a Jupyter Notebook containing clean data, EDA visualizations, and a short analysis of patterns or anomalies.\par

    \vspace{16pt}

    \textcolor{cardBorder}{\hrule height 1pt}\vspace{8pt}

    \small \textcolor{gray}{When you are done with this task, mark it as complete to proceed to the next one.} \hfill \textbf{\textcolor{brandDark}{$\checkmark$ Completed}}

\end{taskbox}

\begin{taskbox}{Task 3: Build an Anomaly Detection Model}

    {\small \textcolor{gray}{2--3 hrs $\cdot$ Intermediate}}\\[12pt]

    \#\#\#\# \textbf{Scenario}  
You’ll create an AI model to identify outliers in claims data that could indicate fraud. The goal is to leverage machine learning techniques to detect anomalies.  

\#\#\#\# \textbf{Student Assignment}  
\begin{itemize}
    \item Split the dataset into training and testing sets.  
    \item Train an anomaly detection model using Isolation Forest or an Autoencoder.  
    \item Evaluate model performance using metrics like precision, recall, and ROC-AUC.  
\end{itemize}

\#\#\#\# \textbf{Deliverable}  
Submit your trained model and a report explaining the methodology, performance metrics, and initial findings.\par

    \vspace{16pt}

    \textcolor{cardBorder}{\hrule height 1pt}\vspace{8pt}

    \small \textcolor{gray}{When you are done with this task, mark it as complete to proceed to the next one.} \hfill \textbf{\textcolor{brandDark}{$\checkmark$ Completed}}

\end{taskbox}

\begin{taskbox}{Task 4: Interpret Results Using Generative AI}

    {\small \textcolor{gray}{2--3 hrs $\cdot$ Intermediate}}\\[12pt]

    \#\#\#\# \textbf{Scenario}  
Your team needs an explanation of why certain claims were flagged as suspicious by the model. You’ll use Generative AI to write a detailed and human-readable interpretation of the fraud indicators.  

\#\#\#\# \textbf{Student Assignment}  
\begin{itemize}
    \item Use Generative AI platforms (e.g., ChatGPT) to explain why specific claims were flagged as anomalies.  
    \item Translate technical outputs into non-technical language for stakeholders.  
    \item Create example scenarios that show why certain claims might be fraudulent.  
\end{itemize}

\#\#\#\# \textbf{Deliverable}  
Submit a document containing Generative AI-generated explanations and real-world implications for flagged claims.\par

    \vspace{16pt}

    \textcolor{cardBorder}{\hrule height 1pt}\vspace{8pt}

    \small \textcolor{gray}{When you are done with this task, mark it as complete to proceed to the next one.} \hfill \textbf{\textcolor{brandDark}{$\checkmark$ Completed}}

\end{taskbox}

\begin{taskbox}{Task 5: Audit for Responsible AI Practices}

    {\small \textcolor{gray}{2--3 hrs $\cdot$ Intermediate}}\\[12pt]

    \#\#\#\# \textbf{Scenario}  
Your anomaly detection model must be fair and unbiased. Your task is to ensure responsible AI practices are adhered to by auditing the model for potential biases.  

\#\#\#\# \textbf{Student Assignment}  
\begin{itemize}
    \item Analyze the flagged claims for demographic bias (e.g., age, gender, location).  
    \item Adjust the model or dataset to minimize bias.  
    \item Document any ethical concerns and propose solutions.  
\end{itemize}

\#\#\#\# \textbf{Deliverable}  
Submit a report outlining bias audit results, adjustments made, and recommendations for maintaining ethical AI practices.\par

    \vspace{16pt}

    \textcolor{cardBorder}{\hrule height 1pt}\vspace{8pt}

    \small \textcolor{gray}{When you are done with this task, mark it as complete to proceed to the next one.} \hfill \textbf{\textcolor{brandDark}{$\checkmark$ Completed}}

\end{taskbox}

\begin{taskbox}{Task 6: Present Your Fraud Detection Report}

    {\small \textcolor{gray}{2--3 hrs $\cdot$ Intermediate}}\\[12pt]

    \#\#\#\# \textbf{Scenario}  
Now that your fraud detection model is complete, you need to present your findings to the investigation team, including actionable recommendations.  

\#\#\#\# \textbf{Student Assignment}  
\begin{itemize}
    \item Summarize your workflow from data exploration to model building and auditing.  
    \item Present top suspicious claims and explain the reasoning behind the anomalies.  
    \item Recommend next steps for investigating flagged claims.  
\end{itemize}

\#\#\#\# \textbf{Deliverable}  
Submit a slide deck (5-7 slides) including visuals, methodology, results, and recommendations.\par

    \vspace{16pt}

    \textcolor{cardBorder}{\hrule height 1pt}\vspace{8pt}

    \small \textcolor{gray}{When you are done with this task, mark it as complete to proceed to the next one.} \hfill \textbf{\textcolor{brandDark}{$\checkmark$ Completed}}

\end{taskbox}

\begin{completionbox}
    \textbf{Knowledge Assessment \& Verification:} Complete the online knowledge check, submit your code and presentation files for AI verification, and claim your \textbf{Skillzza Verified Certificate}.
\end{completionbox}

\newpage

% =============================================================================

% SIMULATION 23

% =============================================================================

\section*{InsurTech Product Analyst – Build AI Claims Assistant}

\addcontentsline{toc}{section}{InsurTech Product Analyst – Build AI Claims Assistant}

\begin{metabox}
    \begin{tabularx}{\textwidth}{@{}ll X@{}}
        \textbf{Industry:} & Technology & \textbf{Partner Enterprise:} \textbf{Skillzza Enterprise} \\
        \textbf{Duration:} & 12-18 hrs (Self-paced) & \textbf{Mode:} Individual, task-based simulation \\
        \textbf{Primary Skills:} & \multicolumn{2}{@{}l}{Python, AI, Data Analysis, Problem Solving}
    \end{tabularx}
\end{metabox}

\subsection*{Overview \& Problem Statement}

\#\#\# \textbf{Scenario}  
You’ve been hired as a Product Analyst at an innovative InsurTech company aiming to revolutionize the claims process. Customers are frustrated with the lengthy and inefficient claims journey, leading to dissatisfaction and churn. The company has tasked you with analyzing the claims journey, identifying pain points, and designing an AI-powered Claims Assistant prototype that streamlines the experience while improving customer satisfaction.

\#\#\# \textbf{Mission}  
Your mission is to deeply understand the claims process, identify bottlenecks and inefficiencies, and leverage AI tools to design a prototype solution. You will analyze claims data, train a chatbot model, define performance KPIs, and present your findings and recommendations to the company’s leadership team.

\#\#\# \textbf{Final Challenge}  
At the end of this virtual internship, you will present an AI Claims Assistant prototype tailored to customer needs. You’ll justify your solution with data-driven insights, demonstrate the model’s functionality, and showcase how your solution addresses customer pain points while meeting organizational KPIs.

---

\subsection*{Tasks \& Deliverables}

\begin{taskbox}{Task 1: Understand the Insurance Claims Process}

    {\small \textcolor{gray}{2--3 hrs $\cdot$ Intermediate}}\\[12pt]

    \#\#\#\# \textbf{Scenario}  
You’ve been provided with anonymized historical claims data and customer feedback reports. Your first task is to analyze the claims journey and identify critical pain points that reduce customer satisfaction.  

\#\#\#\# \textbf{Student Assignment}  
\begin{itemize}
    \item Study the provided dataset and review customer feedback.  
    \item Map the claims journey, highlighting steps such as filing, review, approval, and resolution.  
    \item Identify bottlenecks and inefficiencies in the process (e.g., long approval times, lack of status updates).  
\end{itemize}

\#\#\#\# \textbf{Deliverable}  
\begin{itemize}
    \item A visual claims journey map identifying pain points (e.g., delays, unclear communication).  
    \item A summary document detailing customer pain points and their impact on satisfaction metrics.\par
\end{itemize}

    \vspace{16pt}

    \textcolor{cardBorder}{\hrule height 1pt}\vspace{8pt}

    \small \textcolor{gray}{When you are done with this task, mark it as complete to proceed to the next one.} \hfill \textbf{\textcolor{brandDark}{$\checkmark$ Completed}}

\end{taskbox}

\begin{taskbox}{Task 2: Data Preprocessing and Setup}

    {\small \textcolor{gray}{2--3 hrs $\cdot$ Intermediate}}\\[12pt]

    \#\#\#\# \textbf{Scenario}  
Now that you understand the claims journey, you need to clean and preprocess the raw claims data. This data will be used to train your AI Claims Assistant.  

\#\#\#\# \textbf{Student Assignment}  
\begin{itemize}
    \item Use Python and Pandas to clean the dataset (handle missing values, standardize text fields, and remove duplicates).  
    \item Identify key fields such as claim type, customer sentiment, and claim status for training the AI model.  
\end{itemize}

\#\#\#\# \textbf{Deliverable}  
\begin{itemize}
    \item Cleaned and preprocessed dataset ready for analysis.  
    \item A Jupyter Notebook with documented code and insights from data preprocessing.\par
\end{itemize}

    \vspace{16pt}

    \textcolor{cardBorder}{\hrule height 1pt}\vspace{8pt}

    \small \textcolor{gray}{When you are done with this task, mark it as complete to proceed to the next one.} \hfill \textbf{\textcolor{brandDark}{$\checkmark$ Completed}}

\end{taskbox}

\begin{taskbox}{Task 3: Build Your AI Claims Assistant Prototype}

    {\small \textcolor{gray}{2--3 hrs $\cdot$ Intermediate}}\\[12pt]

    \#\#\#\# \textbf{Scenario}  
Using the OpenAI GPT API, you will create a conversational AI prototype that can answer customer queries related to their claim status, coverage, and process. This chatbot will be the foundation for your AI Claims Assistant.  

\#\#\#\# \textbf{Student Assignment}  
\begin{itemize}
    \item Train an AI chatbot using OpenAI GPT API to handle sample queries like:
    \item “What is the status of my claim?”  
    \item “How long will it take to process my claim?”  
    \item “What documents do I need to submit for my car insurance claim?”  
    \item Test the chatbot with sample inputs and iteratively improve responses.  
\end{itemize}

\#\#\#\# \textbf{Deliverable}  
\begin{itemize}
    \item Python script for the chatbot prototype.  
    \item A document showcasing sample queries and chatbot responses.\par
\end{itemize}

    \vspace{16pt}

    \textcolor{cardBorder}{\hrule height 1pt}\vspace{8pt}

    \small \textcolor{gray}{When you are done with this task, mark it as complete to proceed to the next one.} \hfill \textbf{\textcolor{brandDark}{$\checkmark$ Completed}}

\end{taskbox}

\begin{taskbox}{Task 4: Explain AI Functionality Using GenAI}

    {\small \textcolor{gray}{2--3 hrs $\cdot$ Intermediate}}\\[12pt]

    \#\#\#\# \textbf{Scenario}  
Your stakeholders have requested a detailed explanation of how the AI Claims Assistant works. They want to ensure that the solution is user-friendly and compliant with industry regulations.  

\#\#\#\# \textbf{Student Assignment}  
\begin{itemize}
    \item Use Generative AI tools like ChatGPT to generate a user-friendly explanation of the AI model’s functionality.  
    \item Include details on data inputs, training process, response generation, and safeguards against misinformation.  
\end{itemize}

\#\#\#\# \textbf{Deliverable}  
\begin{itemize}
    \item A comprehensive document explaining the AI Claims Assistant's design and functionality.  
    \item A visual flowchart illustrating the chatbot’s operation.\par
\end{itemize}

    \vspace{16pt}

    \textcolor{cardBorder}{\hrule height 1pt}\vspace{8pt}

    \small \textcolor{gray}{When you are done with this task, mark it as complete to proceed to the next one.} \hfill \textbf{\textcolor{brandDark}{$\checkmark$ Completed}}

\end{taskbox}

\begin{taskbox}{Task 5: Audit the AI Model for Responsible AI Practices}

    {\small \textcolor{gray}{2--3 hrs $\cdot$ Intermediate}}\\[12pt]

    \#\#\#\# \textbf{Scenario}  
Before presenting your prototype, you must ensure that it meets ethical AI standards. Specifically, you need to verify that the model avoids bias and adheres to privacy guidelines.  

\#\#\#\# \textbf{Student Assignment}  
\begin{itemize}
    \item Test the chatbot for potential biases (e.g., handling diverse customer queries).  
    \item Review compliance with data privacy regulations such as GDPR.  
    \item Propose solutions to address any identified ethical concerns.  
\end{itemize}

\#\#\#\# \textbf{Deliverable}  
\begin{itemize}
    \item An audit report on the AI model’s compliance with Responsible AI practices.  
    \item Recommendations for improving model fairness and transparency.\par
\end{itemize}

    \vspace{16pt}

    \textcolor{cardBorder}{\hrule height 1pt}\vspace{8pt}

    \small \textcolor{gray}{When you are done with this task, mark it as complete to proceed to the next one.} \hfill \textbf{\textcolor{brandDark}{$\checkmark$ Completed}}

\end{taskbox}

\begin{taskbox}{Task 6: Present Your Recommendation}

    {\small \textcolor{gray}{2--3 hrs $\cdot$ Intermediate}}\\[12pt]

    \#\#\#\# \textbf{Scenario}  
You are now ready to present your findings and prototype to the company’s leadership. Your presentation must include the claims journey analysis, AI prototype demonstration, and KPI improvements.  

\#\#\#\# \textbf{Student Assignment}  
\begin{itemize}
    \item Create a slide deck summarizing the claims journey pain points, AI solution, and expected impact on KPIs.  
    \item Provide a live demo or video demo of your AI Claims Assistant prototype.  
    \item Justify your solution using data insights and industry benchmarks.  
\end{itemize}

\#\#\#\# \textbf{Deliverable}  
\begin{itemize}
    \item Slide deck (5-7 slides).  
    \item Video demo or live demo of AI Claims Assistant prototype.  
    \item KPI comparison showing expected improvement in customer satisfaction and operational efficiency.\par
\end{itemize}

    \vspace{16pt}

    \textcolor{cardBorder}{\hrule height 1pt}\vspace{8pt}

    \small \textcolor{gray}{When you are done with this task, mark it as complete to proceed to the next one.} \hfill \textbf{\textcolor{brandDark}{$\checkmark$ Completed}}

\end{taskbox}

\begin{completionbox}
    \textbf{Knowledge Assessment \& Verification:} Complete the online knowledge check, submit your code and presentation files for AI verification, and claim your \textbf{Skillzza Verified Certificate}.
\end{completionbox}

\newpage

% =============================================================================

% SIMULATION 24

% =============================================================================

\section*{Cybersecurity Analyst – Investigate a Security Incident}

\addcontentsline{toc}{section}{Cybersecurity Analyst – Investigate a Security Incident}

\begin{metabox}
    \begin{tabularx}{\textwidth}{@{}ll X@{}}
        \textbf{Industry:} & Technology & \textbf{Partner Enterprise:} \textbf{Skillzza Enterprise} \\
        \textbf{Duration:} & 12-18 hrs (Self-paced) & \textbf{Mode:} Individual, task-based simulation \\
        \textbf{Primary Skills:} & \multicolumn{2}{@{}l}{Python, AI, Data Analysis, Problem Solving}
    \end{tabularx}
\end{metabox}

\subsection*{Overview \& Problem Statement}

\#\#\# \textbf{Scenario}  
You are a Cybersecurity Analyst at SecureTech Solutions, a mid-sized company that specializes in providing IT and security services to clients in various industries. One morning, your team receives an alert from the Security Information and Event Management (SIEM) system indicating unusual activity within the internal network. It appears that a potential breach has occurred, and you are assigned to investigate the situation.

\#\#\# \textbf{Mission}  
Your mission is to investigate the security incident, identify the root cause, and provide a detailed report along with recommendations to prevent similar incidents in the future. You will work with logs, analyze traffic patterns, and apply your cybersecurity knowledge to ensure the organization’s systems remain secure.

\#\#\# \textbf{Final Challenge}  
Your final challenge will be to present your findings in a well-structured incident report to your supervisor and suggest actionable recommendations to improve the organization’s security posture.

---

\subsection*{Tasks \& Deliverables}

\begin{taskbox}{Task 1: Understand the Incident}

    {\small \textcolor{gray}{2--3 hrs $\cdot$ Intermediate}}\\[12pt]

    \#\#\#\# \textbf{Scenario}  
Your manager has provided details of an alert from the SIEM system that flagged suspicious activity on the network. A user account attempted multiple failed logins in a short time, followed by successful access from an unknown IP address. This unusual behavior must be investigated.

\#\#\#\# \textbf{Student Assignment}  
1. Review the SIEM alert summary to understand the context of the incident.  
2. Identify the affected user account and note the timestamps of the alert.  
3. Outline potential causes of this activity (e.g., brute force attack, phishing).  

\#\#\#\# \textbf{Deliverable}  
A summary document (300-400 words) that includes:  
\begin{itemize}
    \item Description of the incident.  
    \item Possible causes of the suspicious behavior.  
    \item A list of questions or areas to investigate further.\par
\end{itemize}

    \vspace{16pt}

    \textcolor{cardBorder}{\hrule height 1pt}\vspace{8pt}

    \small \textcolor{gray}{When you are done with this task, mark it as complete to proceed to the next one.} \hfill \textbf{\textcolor{brandDark}{$\checkmark$ Completed}}

\end{taskbox}

\begin{taskbox}{Task 2: Analyze Logs and Set Up Investigation}

    {\small \textcolor{gray}{2--3 hrs $\cdot$ Intermediate}}\\[12pt]

    \#\#\#\# \textbf{Scenario}  
Your manager has provided access to system and network logs during the time of the incident. You need to analyze these logs to identify potential Indicators of Compromise (IoCs).

\#\#\#\# \textbf{Student Assignment}  
1. Import the provided logs into Splunk or Kibana.  
2. Analyze login attempts, IP addresses, and timestamps.  
3. Filter logs to identify anomalies such as unusual login times or unknown IPs.  

\#\#\#\# \textbf{Deliverable}  
\begin{itemize}
    \item A filtered log file highlighting anomalies (e.g., suspicious IPs, unusual timestamps).  
    \item A brief explanation (150-200 words) of the anomalies identified and their potential significance.\par
\end{itemize}

    \vspace{16pt}

    \textcolor{cardBorder}{\hrule height 1pt}\vspace{8pt}

    \small \textcolor{gray}{When you are done with this task, mark it as complete to proceed to the next one.} \hfill \textbf{\textcolor{brandDark}{$\checkmark$ Completed}}

\end{taskbox}

\begin{taskbox}{Task 3: Investigate the Root Cause}

    {\small \textcolor{gray}{2--3 hrs $\cdot$ Intermediate}}\\[12pt]

    \#\#\#\# \textbf{Scenario}  
After analyzing the logs, you notice that the suspicious IP address is associated with a known malware command-and-control server. Additionally, the user account was used to access sensitive files on the company server.

\#\#\#\# \textbf{Student Assignment}  
1. Use Wireshark to analyze the network packet capture (PCAP) file provided.  
2. Identify the types of traffic originating from the suspicious IP.  
3. Research the identified malware and describe its typical behavior.  

\#\#\#\# \textbf{Deliverable}  
\begin{itemize}
    \item A list of key indicators from the PCAP analysis (e.g., protocols, payloads, IPs).  
    \item A malware profile document (200–300 words) describing how it operates and its potential impact.\par
\end{itemize}

    \vspace{16pt}

    \textcolor{cardBorder}{\hrule height 1pt}\vspace{8pt}

    \small \textcolor{gray}{When you are done with this task, mark it as complete to proceed to the next one.} \hfill \textbf{\textcolor{brandDark}{$\checkmark$ Completed}}

\end{taskbox}

\begin{taskbox}{Task 4: Use AI to Identify Patterns}

    {\small \textcolor{gray}{2--3 hrs $\cdot$ Intermediate}}\\[12pt]

    \#\#\#\# \textbf{Scenario}  
You decide to leverage AI for pattern recognition to determine if similar incidents have occurred in the past. Your team has access to a small dataset of past security incidents logged over the last year.

\#\#\#\# \textbf{Student Assignment}  
1. Use a Python-based AI tool (e.g., TensorFlow or pandas library) to analyze the historical incident dataset provided.  
2. Identify patterns or trends (e.g., repeated attacks from the same IP range).  
3. Visualize the results using a graph or chart.  

\#\#\#\# \textbf{Deliverable}  
\begin{itemize}
    \item A Python script or notebook showing your analysis.  
    \item A visualization (e.g., bar chart, heatmap) summarizing the identified patterns.\par
\end{itemize}

    \vspace{16pt}

    \textcolor{cardBorder}{\hrule height 1pt}\vspace{8pt}

    \small \textcolor{gray}{When you are done with this task, mark it as complete to proceed to the next one.} \hfill \textbf{\textcolor{brandDark}{$\checkmark$ Completed}}

\end{taskbox}

\begin{taskbox}{Task 5: Audit and Responsible AI}

    {\small \textcolor{gray}{2--3 hrs $\cdot$ Intermediate}}\\[12pt]

    \#\#\#\# \textbf{Scenario}  
Your findings will be used to update the company's automated threat detection system. However, it's critical to ensure that the AI model does not produce biased or false-positive results.

\#\#\#\# \textbf{Student Assignment}  
1. Review the AI model's dataset and results for potential biases (e.g., flagging certain IP ranges disproportionately).  
2. Suggest improvements to ensure the detection system is fair and accurate.  
3. Document the ethical implications of using AI for cybersecurity.  

\#\#\#\# \textbf{Deliverable}  
\begin{itemize}
    \item A short report (200–300 words) on potential biases in the AI model and steps for mitigation.  
    \item Recommendations for improving the AI model’s accuracy and fairness.\par
\end{itemize}

    \vspace{16pt}

    \textcolor{cardBorder}{\hrule height 1pt}\vspace{8pt}

    \small \textcolor{gray}{When you are done with this task, mark it as complete to proceed to the next one.} \hfill \textbf{\textcolor{brandDark}{$\checkmark$ Completed}}

\end{taskbox}

\begin{taskbox}{Task 6: Present Your Recommendations}

    {\small \textcolor{gray}{2--3 hrs $\cdot$ Intermediate}}\\[12pt]

    \#\#\#\# \textbf{Scenario}  
You are required to create a final incident report summarizing your investigation and recommendations. This report will be presented to the company’s security team for review.

\#\#\#\# \textbf{Student Assignment}  
1. Compile your findings, including the root cause, IoCs, and analysis results.  
2. Provide actionable recommendations to prevent similar incidents in the future.  
3. Create a presentation slide deck summarizing your report.  

\#\#\#\# \textbf{Deliverable}  
\begin{itemize}
    \item A detailed incident report (600–800 words).  
    \item A slide deck (5–7 slides) summarizing your findings and recommendations.\par
\end{itemize}

    \vspace{16pt}

    \textcolor{cardBorder}{\hrule height 1pt}\vspace{8pt}

    \small \textcolor{gray}{When you are done with this task, mark it as complete to proceed to the next one.} \hfill \textbf{\textcolor{brandDark}{$\checkmark$ Completed}}

\end{taskbox}

\begin{completionbox}
    \textbf{Knowledge Assessment \& Verification:} Complete the online knowledge check, submit your code and presentation files for AI verification, and claim your \textbf{Skillzza Verified Certificate}.
\end{completionbox}

\newpage

% =============================================================================

% SIMULATION 25

% =============================================================================

\section*{Cyber Threat Intelligence Analyst – Investigate an Attack}

\addcontentsline{toc}{section}{Cyber Threat Intelligence Analyst – Investigate an Attack}

\begin{metabox}
    \begin{tabularx}{\textwidth}{@{}ll X@{}}
        \textbf{Industry:} & Technology & \textbf{Partner Enterprise:} \textbf{Skillzza Enterprise} \\
        \textbf{Duration:} & 12-18 hrs (Self-paced) & \textbf{Mode:} Individual, task-based simulation \\
        \textbf{Primary Skills:} & \multicolumn{2}{@{}l}{Python, AI, Data Analysis, Problem Solving}
    \end{tabularx}
\end{metabox}

\subsection*{Overview \& Problem Statement}

In an increasingly digital world, cyber threats are evolving at an unprecedented pace. As a \textbf{Cyber Threat Intelligence (CTI) Analyst}, your job is to analyze, detect, and mitigate these threats, ensuring organizations remain secure against potential attacks. 

In this virtual internship, you will step into the shoes of a CTI Analyst to investigate a simulated cyberattack. Using real-world tools and techniques, you will trace an attack from its initial indicators of compromise (IoCs) to uncover the attacker’s identity, motivations, and methods. Finally, you will recommend actionable mitigations to enhance organizational security.  

\#\#\# \textbf{Scenario}
A multinational e-commerce company, \textbf{ShopSphere}, has reported unusual activity on its network. Several employees have flagged suspicious emails, and the IT team has identified data exfiltration attempts from their payment processing servers. ShopSphere has asked you to investigate the incident, analyze threat intelligence, and recommend a mitigation strategy.

\#\#\# \textbf{Mission}
\begin{itemize}
    \item Identify and analyze indicators of compromise (IoCs) from the attack.
    \item Map the threats to the MITRE ATT\&CK framework to uncover their place in the kill chain.
    \item Investigate potential threat actors using open-source intelligence (OSINT) tools.
    \item Develop a mitigation strategy to prevent future attacks.
\end{itemize}

\#\#\# \textbf{Final Challenge}
Your investigation will culminate in a comprehensive Threat Intelligence Report. This report will include IoC analysis, threat actor profiling, attack chain mapping, and a set of actionable recommendations for improving ShopSphere’s cybersecurity defenses.

---

\subsection*{Tasks \& Deliverables}

\begin{taskbox}{Task 1: Understand the Incident}

    {\small \textcolor{gray}{2--3 hrs $\cdot$ Intermediate}}\\[12pt]

    \textbf{Scenario}: ShopSphere’s IT team has provided you with a brief on the suspicious activity, including phishing emails received by employees and packet capture (PCAP) files from their network.  
\textbf{Student Assignment}: 
\begin{itemize}
    \item Read the incident brief and familiarize yourself with the provided data, including email headers, PCAP files, and a list of anomalies flagged by the IT team.
    \item Identify potential IoCs such as suspicious email addresses, IPs, or file hashes.  
\textbf{Deliverable}: A list of at least 5 IoCs, including email addresses, IP addresses, or domain names, with a brief explanation of why each is suspicious.\par
\end{itemize}

    \vspace{16pt}

    \textcolor{cardBorder}{\hrule height 1pt}\vspace{8pt}

    \small \textcolor{gray}{When you are done with this task, mark it as complete to proceed to the next one.} \hfill \textbf{\textcolor{brandDark}{$\checkmark$ Completed}}

\end{taskbox}

\begin{taskbox}{Task 2: Prepare the Dataset for Analysis}

    {\small \textcolor{gray}{2--3 hrs $\cdot$ Intermediate}}\\[12pt]

    \textbf{Scenario}: After identifying IoCs, you need to process and organize the data for deeper analysis. This involves examining phishing emails and decrypting network traffic.  
\textbf{Student Assignment}:
\begin{itemize}
    \item Use Wireshark to analyze the provided PCAP files, filtering out irrelevant traffic and highlighting suspicious patterns.
    \item Extract URLs and file hashes from phishing emails.  
\textbf{Deliverable}: A cleaned dataset containing key IoCs, along with details (e.g., timestamps, source/destination IPs, and extracted URLs).\par
\end{itemize}

    \vspace{16pt}

    \textcolor{cardBorder}{\hrule height 1pt}\vspace{8pt}

    \small \textcolor{gray}{When you are done with this task, mark it as complete to proceed to the next one.} \hfill \textbf{\textcolor{brandDark}{$\checkmark$ Completed}}

\end{taskbox}

\begin{taskbox}{Task 3: Map Threats Using the MITRE ATT\&CK Framework}

    {\small \textcolor{gray}{2--3 hrs $\cdot$ Intermediate}}\\[12pt]

    \textbf{Scenario}: Now that you have IoCs, your next step is to analyze the tactics, techniques, and procedures (TTPs) used in the attack.  
\textbf{Student Assignment}:  
\begin{itemize}
    \item Use the MITRE ATT\&CK Navigator to map the IoCs to specific tactics and techniques.  
    \item Identify where these techniques fit into the cyber kill chain (e.g., initial access, lateral movement, exfiltration).  
\textbf{Deliverable}: A complete MITRE ATT\&CK mapping for the attack, including identified techniques and their descriptions.\par
\end{itemize}

    \vspace{16pt}

    \textcolor{cardBorder}{\hrule height 1pt}\vspace{8pt}

    \small \textcolor{gray}{When you are done with this task, mark it as complete to proceed to the next one.} \hfill \textbf{\textcolor{brandDark}{$\checkmark$ Completed}}

\end{taskbox}

\begin{taskbox}{Task 4: Investigate the Threat Actor Using OSINT}

    {\small \textcolor{gray}{2--3 hrs $\cdot$ Intermediate}}\\[12pt]

    \textbf{Scenario}: With insights into TTPs, you can now investigate potential threat actors behind the attack.  
\textbf{Student Assignment}:  
\begin{itemize}
    \item Use Shodan to investigate IP addresses from the IoC list.  
    \item Cross-reference IPs, hashes, and domains with MISP or VirusTotal to identify connections to known threat actors.  
\textbf{Deliverable}: A threat actor profile, including name (if identifiable), known motivations, and prior activities.\par
\end{itemize}

    \vspace{16pt}

    \textcolor{cardBorder}{\hrule height 1pt}\vspace{8pt}

    \small \textcolor{gray}{When you are done with this task, mark it as complete to proceed to the next one.} \hfill \textbf{\textcolor{brandDark}{$\checkmark$ Completed}}

\end{taskbox}

\begin{taskbox}{Task 5: Audit and Responsible Cybersecurity Practices}

    {\small \textcolor{gray}{2--3 hrs $\cdot$ Intermediate}}\\[12pt]

    \textbf{Scenario}: Before presenting your findings, ensure that your investigation adheres to ethical and responsible practices.  
\textbf{Student Assignment}:  
\begin{itemize}
    \item Identify potential biases or gaps in your analysis (e.g., false positives, incomplete data).  
    \item Verify that your proposed actions comply with cybersecurity laws and organizational policies.  
\textbf{Deliverable}: A one-page Responsible Cybersecurity Checklist with potential issues and their resolutions.\par
\end{itemize}

    \vspace{16pt}

    \textcolor{cardBorder}{\hrule height 1pt}\vspace{8pt}

    \small \textcolor{gray}{When you are done with this task, mark it as complete to proceed to the next one.} \hfill \textbf{\textcolor{brandDark}{$\checkmark$ Completed}}

\end{taskbox}

\begin{taskbox}{Task 6: Present Your Findings and Recommendations}

    {\small \textcolor{gray}{2--3 hrs $\cdot$ Intermediate}}\\[12pt]

    \textbf{Scenario}: As a final step, you will consolidate your findings into a report for ShopSphere’s executive team.  
\textbf{Student Assignment}:  
\begin{itemize}
    \item Summarize the attack chain, threat actor profile, and key findings.  
    \item Recommend at least three actionable mitigations to prevent similar attacks in the future.  
\textbf{Deliverable}: A polished Threat Intelligence Report (PDF or PPT) with visuals (e.g., MITRE ATT\&CK heatmap).\par
\end{itemize}

    \vspace{16pt}

    \textcolor{cardBorder}{\hrule height 1pt}\vspace{8pt}

    \small \textcolor{gray}{When you are done with this task, mark it as complete to proceed to the next one.} \hfill \textbf{\textcolor{brandDark}{$\checkmark$ Completed}}

\end{taskbox}

\begin{completionbox}
    \textbf{Knowledge Assessment \& Verification:} Complete the online knowledge check, submit your code and presentation files for AI verification, and claim your \textbf{Skillzza Verified Certificate}.
\end{completionbox}

\newpage

% =============================================================================

% SIMULATION 26

% =============================================================================

\section*{Cloud Security Analyst – Secure an AI Workload}

\addcontentsline{toc}{section}{Cloud Security Analyst – Secure an AI Workload}

\begin{metabox}
    \begin{tabularx}{\textwidth}{@{}ll X@{}}
        \textbf{Industry:} & Technology & \textbf{Partner Enterprise:} \textbf{Skillzza Enterprise} \\
        \textbf{Duration:} & 12-18 hrs (Self-paced) & \textbf{Mode:} Individual, task-based simulation \\
        \textbf{Primary Skills:} & \multicolumn{2}{@{}l}{Python, AI, Data Analysis, Problem Solving}
    \end{tabularx}
\end{metabox}

\subsection*{Overview \& Problem Statement}

\#\#\# \textbf{Scenario}
Organizations are increasingly deploying AI workloads in the cloud to leverage scalability, flexibility, and cost efficiency. As a Cloud Security Analyst, your role is to ensure these AI workloads are secure from potential threats, vulnerabilities, and compliance risks. You will work on securing a sensitive AI workload hosted on a public cloud platform, with a focus on architecture design, threat modeling, security controls, and providing actionable recommendations.

\#\#\# \textbf{Mission}
Your mission is to design and implement a secure cloud architecture for an AI workload that processes sensitive customer data. You will identify vulnerabilities, apply security controls, and create a comprehensive security recommendation report for stakeholders. The goal is to ensure confidentiality, integrity, and availability of the AI workload while adhering to industry-best practices and compliance requirements.

\#\#\# \textbf{Final Challenge}
By the end of this simulation, you will present a detailed security recommendation report to a simulated board of executives, explaining the security improvements implemented, trade-offs considered, and how the architecture aligns with compliance standards (e.g., GDPR, ISO 27001).

---

\subsection*{Tasks \& Deliverables}

\begin{taskbox}{Task 1: Understand the Cloud AI Workload}

    {\small \textcolor{gray}{2--3 hrs $\cdot$ Intermediate}}\\[12pt]

    \#\#\#\# \textbf{Scenario}
You are given a cloud-hosted AI model that processes customer financial data to predict credit scores. The workload consists of an inference pipeline deployed on a public cloud platform.

\#\#\#\# \textbf{Student Assignment}
\begin{itemize}
    \item Study the architecture of the AI workload.
    \item Identify sensitive data flows and critical components (e.g., API gateways, storage buckets).
    \item Pinpoint areas where security risks may arise (e.g., data exposure, unauthorized access).
\end{itemize}

\#\#\#\# \textbf{Deliverable}
Submit a report mapping the sensitive data flow within the workload along with a list of potential security risks.\par

    \vspace{16pt}

    \textcolor{cardBorder}{\hrule height 1pt}\vspace{8pt}

    \small \textcolor{gray}{When you are done with this task, mark it as complete to proceed to the next one.} \hfill \textbf{\textcolor{brandDark}{$\checkmark$ Completed}}

\end{taskbox}

\begin{taskbox}{Task 2: Set Up the Environment}

    {\small \textcolor{gray}{2--3 hrs $\cdot$ Intermediate}}\\[12pt]

    \#\#\#\# \textbf{Scenario}
You need to configure the cloud environment to host the AI workload securely. The workload will be deployed using containers and connected to a storage bucket for data input.

\#\#\#\# \textbf{Student Assignment}
\begin{itemize}
    \item Set up IAM roles and policies for secure access control.
    \item Configure VPCs (Virtual Private Clouds) and subnets for isolating the workload.
    \item Enable encryption for storage buckets and data in transit.
\end{itemize}

\#\#\#\# \textbf{Deliverable}
Submit screenshots of IAM configurations, VPC settings, and encryption setup.\par

    \vspace{16pt}

    \textcolor{cardBorder}{\hrule height 1pt}\vspace{8pt}

    \small \textcolor{gray}{When you are done with this task, mark it as complete to proceed to the next one.} \hfill \textbf{\textcolor{brandDark}{$\checkmark$ Completed}}

\end{taskbox}

\begin{taskbox}{Task 3: Build Threat Model}

    {\small \textcolor{gray}{2--3 hrs $\cdot$ Intermediate}}\\[12pt]

    \#\#\#\# \textbf{Scenario}
Use a threat modeling tool to identify vulnerabilities in the AI workload architecture.

\#\#\#\# \textbf{Student Assignment}
\begin{itemize}
    \item Create a data flow diagram (DFD) for the workload.
    \item Identify potential threats using STRIDE methodology (Spoofing, Tampering, Repudiation, Information Disclosure, Denial of Service, Elevation of Privilege).
    \item Suggest mitigations for each identified threat.
\end{itemize}

\#\#\#\# \textbf{Deliverable}
Submit the threat model diagram and a table listing threats with corresponding mitigations.\par

    \vspace{16pt}

    \textcolor{cardBorder}{\hrule height 1pt}\vspace{8pt}

    \small \textcolor{gray}{When you are done with this task, mark it as complete to proceed to the next one.} \hfill \textbf{\textcolor{brandDark}{$\checkmark$ Completed}}

\end{taskbox}

\begin{taskbox}{Task 4: Apply GenAI for Explanation}

    {\small \textcolor{gray}{2--3 hrs $\cdot$ Intermediate}}\\[12pt]

    \#\#\#\# \textbf{Scenario}
A junior stakeholder asks for an explanation of the security design in simple terms. Use a Generative AI tool to create an explanatory document summarizing the security measures.

\#\#\#\# \textbf{Student Assignment}
\begin{itemize}
    \item Input the security architecture into a generative AI tool (e.g., ChatGPT or Bard).
    \item Create a stakeholder-friendly document that explains the security measures in non-technical language.
    \item Include analogies or visual aids.
\end{itemize}

\#\#\#\# \textbf{Deliverable}
Submit a 1-page stakeholder document explaining the security design in simple terms.\par

    \vspace{16pt}

    \textcolor{cardBorder}{\hrule height 1pt}\vspace{8pt}

    \small \textcolor{gray}{When you are done with this task, mark it as complete to proceed to the next one.} \hfill \textbf{\textcolor{brandDark}{$\checkmark$ Completed}}

\end{taskbox}

\begin{taskbox}{Task 5: Perform Security Audit}

    {\small \textcolor{gray}{2--3 hrs $\cdot$ Intermediate}}\\[12pt]

    \#\#\#\# \textbf{Scenario}
You need to ensure the implemented security controls align with industry standards (e.g., GDPR, ISO 27001).

\#\#\#\# \textbf{Student Assignment}
\begin{itemize}
    \item Audit the environment using a compliance checker tool.
    \item Identify gaps in compliance and propose solutions.
    \item Test the effectiveness of encryption and IAM policies.
\end{itemize}

\#\#\#\# \textbf{Deliverable}
Submit a compliance audit report highlighting gaps and solutions.\par

    \vspace{16pt}

    \textcolor{cardBorder}{\hrule height 1pt}\vspace{8pt}

    \small \textcolor{gray}{When you are done with this task, mark it as complete to proceed to the next one.} \hfill \textbf{\textcolor{brandDark}{$\checkmark$ Completed}}

\end{taskbox}

\begin{taskbox}{Task 6: Present Security Recommendations}

    {\small \textcolor{gray}{2--3 hrs $\cdot$ Intermediate}}\\[12pt]

    \#\#\#\# \textbf{Scenario}
You are presenting your findings to the board of executives. They need to understand the security improvements, trade-offs, and compliance alignment.

\#\#\#\# \textbf{Student Assignment}
\begin{itemize}
    \item Create a presentation summarizing the architecture, threat model, applied controls, and compliance results.
    \item Highlight trade-offs (e.g., performance vs. security) and recommendations for future improvements.
\end{itemize}

\#\#\#\# \textbf{Deliverable}
Submit a PowerPoint presentation with speaker notes.\par

    \vspace{16pt}

    \textcolor{cardBorder}{\hrule height 1pt}\vspace{8pt}

    \small \textcolor{gray}{When you are done with this task, mark it as complete to proceed to the next one.} \hfill \textbf{\textcolor{brandDark}{$\checkmark$ Completed}}

\end{taskbox}

\begin{completionbox}
    \textbf{Knowledge Assessment \& Verification:} Complete the online knowledge check, submit your code and presentation files for AI verification, and claim your \textbf{Skillzza Verified Certificate}.
\end{completionbox}

\newpage

% =============================================================================

% SIMULATION 27

% =============================================================================

\section*{AI Developer – Build a RAG Application}

\addcontentsline{toc}{section}{AI Developer – Build a RAG Application}

\begin{metabox}
    \begin{tabularx}{\textwidth}{@{}ll X@{}}
        \textbf{Industry:} & Technology & \textbf{Partner Enterprise:} \textbf{Skillzza Enterprise} \\
        \textbf{Duration:} & 12-18 hrs (Self-paced) & \textbf{Mode:} Individual, task-based simulation \\
        \textbf{Primary Skills:} & \multicolumn{2}{@{}l}{Python, AI, Data Analysis, Problem Solving}
    \end{tabularx}
\end{metabox}

\subsection*{Overview \& Problem Statement}

\#\#\# Scenario  
You have been hired as an AI Developer at a cutting-edge technology company specializing in generative AI applications. Your team has been tasked with building a Retrieval-Augmented Generation (RAG) application for a client in the legal domain. The application will enable lawyers to query large legal document datasets and instantly retrieve relevant information, powered by generative AI capabilities.  

\#\#\# Mission  
Your mission is to design, implement, and deploy a working prototype of the RAG system. This involves chunking legal documents, generating embeddings, implementing a retrieval mechanism, integrating a generative AI model for answering questions, and evaluating the quality of generated answers.  

\#\#\# Final Challenge  
At the end of the internship, you will present your RAG application prototype to stakeholders. You will demonstrate its functionality, explain the underlying AI architecture, and showcase how effectively it retrieves and generates answers to complex legal queries.  

---

\subsection*{Tasks \& Deliverables}

\begin{taskbox}{Task 1: Understand}

    {\small \textcolor{gray}{2--3 hrs $\cdot$ Intermediate}}\\[12pt]

    \#\#\#\# Scenario  
Your client is a legal firm with thousands of case files, contracts, and regulatory documents. They need a system that allows lawyers to query this dataset and retrieve concise, accurate answers to legal questions. You are tasked to understand the architecture and requirements of a RAG system.  

\#\#\#\# Student Assignment  
\begin{itemize}
    \item Research and document how RAG systems work, including the retrieval and generation pipeline.  
    \item Identify key challenges in building a RAG system for domain-specific applications like legal documents.  
\end{itemize}

\#\#\#\# Deliverable  
Submit a 1-page technical document explaining the architecture of a RAG system and challenges specific to the legal domain.\par

    \vspace{16pt}

    \textcolor{cardBorder}{\hrule height 1pt}\vspace{8pt}

    \small \textcolor{gray}{When you are done with this task, mark it as complete to proceed to the next one.} \hfill \textbf{\textcolor{brandDark}{$\checkmark$ Completed}}

\end{taskbox}

\begin{taskbox}{Task 2: Data/Setup}

    {\small \textcolor{gray}{2--3 hrs $\cdot$ Intermediate}}\\[12pt]

    \#\#\#\# Scenario  
You have been provided with a dataset containing legal documents in PDF format. Your first step is to preprocess these documents for analysis.  

\#\#\#\# Student Assignment  
\begin{itemize}
    \item Extract text from the provided PDFs using PyPDF2 or a similar library.  
    \item Implement a chunking algorithm to split large documents into smaller, contextually coherent chunks (e.g., paragraphs or sections).  
\end{itemize}

\#\#\#\# Deliverable  
Submit a Jupyter Notebook containing the code for text extraction and document chunking, along with a summary of how many chunks were created per document.\par

    \vspace{16pt}

    \textcolor{cardBorder}{\hrule height 1pt}\vspace{8pt}

    \small \textcolor{gray}{When you are done with this task, mark it as complete to proceed to the next one.} \hfill \textbf{\textcolor{brandDark}{$\checkmark$ Completed}}

\end{taskbox}

\begin{taskbox}{Task 3: Build/Execute}

    {\small \textcolor{gray}{2--3 hrs $\cdot$ Intermediate}}\\[12pt]

    \#\#\#\# Scenario  
To enable retrieval, the chunks need to be converted into vector embeddings. You must generate embeddings using a pre-trained model and store them in a vector database.  

\#\#\#\# Student Assignment  
\begin{itemize}
    \item Use a pre-trained model like OpenAI's text-embedding-ada-002 or Hugging Face's Sentence Transformers to generate embeddings for the document chunks.  
    \item Store the embeddings in FAISS or Pinecone and implement a retrieval mechanism to fetch relevant chunks based on cosine similarity.  
\end{itemize}

\#\#\#\# Deliverable  
Submit the code for embedding generation and retrieval, along with a sample query and its retrieved document chunks.\par

    \vspace{16pt}

    \textcolor{cardBorder}{\hrule height 1pt}\vspace{8pt}

    \small \textcolor{gray}{When you are done with this task, mark it as complete to proceed to the next one.} \hfill \textbf{\textcolor{brandDark}{$\checkmark$ Completed}}

\end{taskbox}

\begin{taskbox}{Task 4: GenAI/Explanation}

    {\small \textcolor{gray}{2--3 hrs $\cdot$ Intermediate}}\\[12pt]

    \#\#\#\# Scenario  
Once retrieval is operational, integrate a generative AI model (e.g., GPT-4) to provide concise answers based on the retrieved chunks.  

\#\#\#\# Student Assignment  
\begin{itemize}
    \item Design a pipeline that passes the retrieved chunks as context to the generative AI model.  
    \item Fine-tune the prompts for the generative model to ensure that answers are accurate, domain-specific, and well-structured.  
\end{itemize}

\#\#\#\# Deliverable  
Submit the code for the generative pipeline along with sample queries and their generated answers.\par

    \vspace{16pt}

    \textcolor{cardBorder}{\hrule height 1pt}\vspace{8pt}

    \small \textcolor{gray}{When you are done with this task, mark it as complete to proceed to the next one.} \hfill \textbf{\textcolor{brandDark}{$\checkmark$ Completed}}

\end{taskbox}

\begin{taskbox}{Task 5: Audit/Responsible AI}

    {\small \textcolor{gray}{2--3 hrs $\cdot$ Intermediate}}\\[12pt]

    \#\#\#\# Scenario  
The client emphasizes ethical and responsible AI usage, particularly in the legal domain where accuracy is critical. You are tasked with auditing the system for biases and inaccuracies.  

\#\#\#\# Student Assignment  
\begin{itemize}
    \item Evaluate the accuracy of generated answers by comparing them to human-curated answers.  
    \item Document any instances of hallucination or biased outputs.  
    \item Propose strategies for mitigating these issues.  
\end{itemize}

\#\#\#\# Deliverable  
Submit an audit report highlighting the system’s performance, any identified issues, and proposed solutions.\par

    \vspace{16pt}

    \textcolor{cardBorder}{\hrule height 1pt}\vspace{8pt}

    \small \textcolor{gray}{When you are done with this task, mark it as complete to proceed to the next one.} \hfill \textbf{\textcolor{brandDark}{$\checkmark$ Completed}}

\end{taskbox}

\begin{taskbox}{Task 6: Present Recommendation}

    {\small \textcolor{gray}{2--3 hrs $\cdot$ Intermediate}}\\[12pt]

    \#\#\#\# Scenario  
You are now ready to present your RAG system prototype to the client. Explain how the system works, its benefits, and address any concerns raised during testing.  

\#\#\#\# Student Assignment  
\begin{itemize}
    \item Prepare a presentation that includes the system architecture, sample outputs, and performance metrics.  
    \item Include explanations of how you ensured responsible AI practices.  
\end{itemize}

\#\#\#\# Deliverable  
Submit a slide deck (PowerPoint or PDF) with at least 8 slides covering the above points.\par

    \vspace{16pt}

    \textcolor{cardBorder}{\hrule height 1pt}\vspace{8pt}

    \small \textcolor{gray}{When you are done with this task, mark it as complete to proceed to the next one.} \hfill \textbf{\textcolor{brandDark}{$\checkmark$ Completed}}

\end{taskbox}

\begin{completionbox}
    \textbf{Knowledge Assessment \& Verification:} Complete the online knowledge check, submit your code and presentation files for AI verification, and claim your \textbf{Skillzza Verified Certificate}.
\end{completionbox}

\newpage

% =============================================================================

% SIMULATION 28

% =============================================================================

\section*{LangChain Developer – Build an Enterprise AI Workflow}

\addcontentsline{toc}{section}{LangChain Developer – Build an Enterprise AI Workflow}

\begin{metabox}
    \begin{tabularx}{\textwidth}{@{}ll X@{}}
        \textbf{Industry:} & Technology & \textbf{Partner Enterprise:} \textbf{Skillzza Enterprise} \\
        \textbf{Duration:} & 12-18 hrs (Self-paced) & \textbf{Mode:} Individual, task-based simulation \\
        \textbf{Primary Skills:} & \multicolumn{2}{@{}l}{Python, AI, Data Analysis, Problem Solving}
    \end{tabularx}
\end{metabox}

\subsection*{Overview \& Problem Statement}

\#\#\# \textbf{Scenario}
You’ve been hired as a \textbf{LangChain Developer} by a leading AI consulting firm. Your task is to design and implement a cutting-edge \textbf{enterprise AI workflow} using \textbf{LangChain} to streamline customer support processes for a large e-commerce company. The company aims to use generative AI to provide intelligent responses to customer queries, fetch relevant product details, and escalate complex issues to human agents while maintaining context across interactions.

\#\#\# \textbf{Mission}
Your mission is to build a \textbf{multi-step conversational AI system} leveraging LangChain. This system will integrate \textbf{tools} like external APIs for fetching product information, utilize \textbf{memory} to store conversation context, and generate structured outputs for issue tracking and analytics. You will also evaluate the system’s performance using automated benchmarks and responsible AI principles.

\#\#\# \textbf{Final Challenge}
Design, implement, and present an \textbf{end-to-end LangChain prototype} that connects chains, tools, memory, and structured outputs. Your solution must demonstrate \textbf{high reliability}, \textbf{accuracy}, and \textbf{ethical considerations} in its responses.

---

\subsection*{Tasks \& Deliverables}

\begin{taskbox}{Task 1: Understand}

    {\small \textcolor{gray}{2--3 hrs $\cdot$ Intermediate}}\\[12pt]

    \#\#\#\# Scenario:
The e-commerce company needs an AI workflow to handle customer queries, fetch product details, and escalate issues. You must understand the requirements for building a \textbf{multi-step conversational chain}.

\#\#\#\# Student Assignment:
\begin{itemize}
    \item Research \textbf{LangChain components}: chains, tools, memory, and structured outputs.
    \item Study a sample dataset containing customer queries and product details (fields: `query\_id`, `customer\_query`, `product\_id`, `product\_name`, `price`, `availability`, `issue\_description`).
\end{itemize}

\#\#\#\# Deliverable:
\begin{itemize}
    \item Submit a document summarizing how LangChain components can be used to meet the requirements.
    \item Identify potential challenges (e.g., ambiguous queries, memory limitations).\par
\end{itemize}

    \vspace{16pt}

    \textcolor{cardBorder}{\hrule height 1pt}\vspace{8pt}

    \small \textcolor{gray}{When you are done with this task, mark it as complete to proceed to the next one.} \hfill \textbf{\textcolor{brandDark}{$\checkmark$ Completed}}

\end{taskbox}

\begin{taskbox}{Task 2: Data/Setup}

    {\small \textcolor{gray}{2--3 hrs $\cdot$ Intermediate}}\\[12pt]

    \#\#\#\# Scenario:
You will set up the environment and load the e-commerce dataset. The dataset includes fields like `query\_id`, `customer\_query`, `product\_id`, `product\_name`, `price`, and `availability`.

\#\#\#\# Student Assignment:
\begin{itemize}
    \item Install \textbf{LangChain} and related dependencies in your Python environment.
    \item Load the sample dataset into a Pandas DataFrame for preprocessing.
    \item Write a Python script to test API connectivity for product information retrieval (e.g., mock API for `GET product details`).
\end{itemize}

\#\#\#\# Deliverable:
\begin{itemize}
    \item Submit a Python script that loads the dataset and retrieves product details via API calls.
    \item Provide a screenshot of successful API responses.\par
\end{itemize}

    \vspace{16pt}

    \textcolor{cardBorder}{\hrule height 1pt}\vspace{8pt}

    \small \textcolor{gray}{When you are done with this task, mark it as complete to proceed to the next one.} \hfill \textbf{\textcolor{brandDark}{$\checkmark$ Completed}}

\end{taskbox}

\begin{taskbox}{Task 3: Build/Execute}

    {\small \textcolor{gray}{2--3 hrs $\cdot$ Intermediate}}\\[12pt]

    \#\#\#\# Scenario:
Using LangChain, you will implement a \textbf{multi-step chain} that processes customer queries, fetches product details, and escalates unresolved issues.

\#\#\#\# Student Assignment:
\begin{itemize}
    \item Build a \textbf{SequentialChain} in LangChain to handle:
  1. Parsing customer queries.
  2. Fetching product information using the API.
  3. Escalating complex issues.
    \item Integrate \textbf{memory} to store context across multiple queries from the same customer.
\end{itemize}

\#\#\#\# Deliverable:
\begin{itemize}
    \item Submit your Python code for the SequentialChain.
    \item Provide logs showing sample input queries and corresponding outputs.\par
\end{itemize}

    \vspace{16pt}

    \textcolor{cardBorder}{\hrule height 1pt}\vspace{8pt}

    \small \textcolor{gray}{When you are done with this task, mark it as complete to proceed to the next one.} \hfill \textbf{\textcolor{brandDark}{$\checkmark$ Completed}}

\end{taskbox}

\begin{taskbox}{Task 4: GenAI/Explanation}

    {\small \textcolor{gray}{2--3 hrs $\cdot$ Intermediate}}\\[12pt]

    \#\#\#\# Scenario:
Your AI workflow must generate \textbf{structured responses} (e.g., JSON format) to customer queries, including product details and escalation status.

\#\#\#\# Student Assignment:
\begin{itemize}
    \item Use \textbf{prompt engineering} to design templates for structured responses.
    \item Implement structured outputs in JSON format using LangChain.
    \item Write a short explanation of how prompt engineering improves response accuracy.
\end{itemize}

\#\#\#\# Deliverable:
\begin{itemize}
    \item Submit the updated Python script with structured output generation.
    \item Provide 3 sample queries and their JSON responses.\par
\end{itemize}

    \vspace{16pt}

    \textcolor{cardBorder}{\hrule height 1pt}\vspace{8pt}

    \small \textcolor{gray}{When you are done with this task, mark it as complete to proceed to the next one.} \hfill \textbf{\textcolor{brandDark}{$\checkmark$ Completed}}

\end{taskbox}

\begin{taskbox}{Task 5: Audit/Responsible AI}

    {\small \textcolor{gray}{2--3 hrs $\cdot$ Intermediate}}\\[12pt]

    \#\#\#\# Scenario:
The workflow must comply with responsible AI practices, ensuring accuracy, fairness, and ethical handling of customer data.

\#\#\#\# Student Assignment:
\begin{itemize}
    \item Implement an \textbf{Evaluation Chain} in LangChain to benchmark the workflow.
    \item Check for biases in product recommendations and escalation decisions.
    \item Write a report outlining responsible AI improvements.
\end{itemize}

\#\#\#\# Deliverable:
\begin{itemize}
    \item Submit the evaluation results and responsible AI report.
    \item Highlight areas for improvement (e.g., response accuracy or ethical compliance).\par
\end{itemize}

    \vspace{16pt}

    \textcolor{cardBorder}{\hrule height 1pt}\vspace{8pt}

    \small \textcolor{gray}{When you are done with this task, mark it as complete to proceed to the next one.} \hfill \textbf{\textcolor{brandDark}{$\checkmark$ Completed}}

\end{taskbox}

\begin{taskbox}{Task 6: Present Recommendation}

    {\small \textcolor{gray}{2--3 hrs $\cdot$ Intermediate}}\\[12pt]

    \#\#\#\# Scenario:
You will present your finished prototype to stakeholders. The presentation must include workflow architecture, sample customer interactions, and evaluation results.

\#\#\#\# Student Assignment:
\begin{itemize}
    \item Create a presentation showcasing:
  1. Workflow overview (chains, tools, memory, and outputs).
  2. Sample interactions and JSON outputs.
  3. Evaluation metrics and responsible AI practices.
    \item Record a video demo of the prototype in action.
\end{itemize}

\#\#\#\# Deliverable:
\begin{itemize}
    \item Submit the presentation slides (PDF) and video demo link.
    \item Provide a summary document with key recommendations.\par
\end{itemize}

    \vspace{16pt}

    \textcolor{cardBorder}{\hrule height 1pt}\vspace{8pt}

    \small \textcolor{gray}{When you are done with this task, mark it as complete to proceed to the next one.} \hfill \textbf{\textcolor{brandDark}{$\checkmark$ Completed}}

\end{taskbox}

\begin{completionbox}
    \textbf{Knowledge Assessment \& Verification:} Complete the online knowledge check, submit your code and presentation files for AI verification, and claim your \textbf{Skillzza Verified Certificate}.
\end{completionbox}

\newpage

% =============================================================================

% SIMULATION 29

% =============================================================================

\section*{LangGraph Developer – Build an Agentic Workflow}

\addcontentsline{toc}{section}{LangGraph Developer – Build an Agentic Workflow}

\begin{metabox}
    \begin{tabularx}{\textwidth}{@{}ll X@{}}
        \textbf{Industry:} & Technology & \textbf{Partner Enterprise:} \textbf{Skillzza Enterprise} \\
        \textbf{Duration:} & 12-18 hrs (Self-paced) & \textbf{Mode:} Individual, task-based simulation \\
        \textbf{Primary Skills:} & \multicolumn{2}{@{}l}{Python, AI, Data Analysis, Problem Solving}
    \end{tabularx}
\end{metabox}

\subsection*{Overview \& Problem Statement}

\#\#\# \textbf{Scenario}
The rise of Agentic AI (autonomous agents capable of performing multi-step, goal-driven tasks) has created a demand for skilled engineers who can design and build workflows that enable AI agents to interact with tools, humans, and external systems seamlessly. You’ve been hired as a \textbf{LangGraph Developer} at \textbf{FlowAI}, a cutting-edge organization specializing in the development of agent-based systems. 

Your mission is to design, implement, and test a modular LangGraph-based workflow that integrates state management, tools, and human approval checkpoints. This workflow will serve as the backbone for an AI-powered assistant capable of autonomously completing complex tasks such as scheduling, document drafting, and task prioritization.

\#\#\# \textbf{Mission}
As the LangGraph Developer, you’ll be tasked with creating a robust agentic workflow using LangChain’s graph-based architecture. You will define the state management structure, design agent interactions, connect external tools (e.g., APIs), and implement a human-in-the-loop approval mechanism. Finally, you will test and audit the system for responsible AI deployment.

\#\#\# \textbf{Final Challenge}
Your final challenge will be to deliver a fully functional LangGraph workflow for an AI assistant that can:
1. Process a user's task request (e.g., "Organize a meeting with John, draft an agenda, and send confirmation").
2. Break down the task into sub-tasks using agents.
3. Execute sub-tasks by interacting with external tools like Google Calendar and email APIs.
4. Request human approval for sensitive actions (e.g., sending an email).
5. Return the final output as a cohesive response.

---

\subsection*{Tasks \& Deliverables}

\begin{taskbox}{Task 1: Understand the Role of LangGraphs}

    {\small \textcolor{gray}{2--3 hrs $\cdot$ Intermediate}}\\[12pt]

    \#\#\#\# \textbf{Scenario}
You’re starting your journey as a LangGraph Developer at FlowAI. Your first task is to understand the foundational concepts of LangGraph, including nodes, directed edges, and how states are managed.

\#\#\#\# \textbf{Student Assignment}
\begin{itemize}
    \item Study the documentation for LangChain's LangGraph module and its role in AI workflows.
    \item Write a brief summary of:
  1. What a LangGraph is.
  2. The difference between a Tool Node, Agent Node, and Human Approval Node.
\end{itemize}

\#\#\#\# \textbf{Deliverable}
Submit a 300-word summary explaining LangGraphs and their components. Include examples of how each node type can be used in a workflow.\par

    \vspace{16pt}

    \textcolor{cardBorder}{\hrule height 1pt}\vspace{8pt}

    \small \textcolor{gray}{When you are done with this task, mark it as complete to proceed to the next one.} \hfill \textbf{\textcolor{brandDark}{$\checkmark$ Completed}}

\end{taskbox}

\begin{taskbox}{Task 2: Define and Set Up the Workflow}

    {\small \textcolor{gray}{2--3 hrs $\cdot$ Intermediate}}\\[12pt]

    \#\#\#\# \textbf{Scenario}
Your manager has requested a LangGraph workflow blueprint for an AI assistant that automates scheduling tasks, drafts emails, and requires human approval for sending sensitive emails.

\#\#\#\# \textbf{Student Assignment}
\begin{itemize}
    \item Create a LangGraph blueprint with:
  1. Nodes for task input, agents, tools, and human approval.
  2. State management to track sub-task progress.
\end{itemize}

\begin{itemize}
    \item Use pseudocode or a Python skeleton to outline your implementation steps.
\end{itemize}

\#\#\#\# \textbf{Deliverable}
Submit a LangGraph blueprint (diagram or detailed explanation) and a Python script skeleton for your workflow.\par

    \vspace{16pt}

    \textcolor{cardBorder}{\hrule height 1pt}\vspace{8pt}

    \small \textcolor{gray}{When you are done with this task, mark it as complete to proceed to the next one.} \hfill \textbf{\textcolor{brandDark}{$\checkmark$ Completed}}

\end{taskbox}

\begin{taskbox}{Task 3: Build and Execute the LangGraph}

    {\small \textcolor{gray}{2--3 hrs $\cdot$ Intermediate}}\\[12pt]

    \#\#\#\# \textbf{Scenario}
With your blueprint approved, it’s time to build the workflow in Python using LangChain. You’ll implement agents for task decomposition and tool interactions.

\#\#\#\# \textbf{Student Assignment}
\begin{itemize}
    \item Implement the LangGraph in Python.
    \item Create agents for:
  1. Task decomposition.
  2. Scheduling with Google Calendar API.
  3. Drafting emails using OpenAI GPT.
    \item Ensure all outputs are logged for debugging.
\end{itemize}

\#\#\#\# \textbf{Deliverable}
Submit your Python code for the fully functional LangGraph workflow, along with a README file detailing how to run the script.\par

    \vspace{16pt}

    \textcolor{cardBorder}{\hrule height 1pt}\vspace{8pt}

    \small \textcolor{gray}{When you are done with this task, mark it as complete to proceed to the next one.} \hfill \textbf{\textcolor{brandDark}{$\checkmark$ Completed}}

\end{taskbox}

\begin{taskbox}{Task 4: Integrate Generative AI for Contextual Explanations}

    {\small \textcolor{gray}{2--3 hrs $\cdot$ Intermediate}}\\[12pt]

    \#\#\#\# \textbf{Scenario}
The AI assistant must provide meaningful explanations for its actions to increase user trust and transparency.

\#\#\#\# \textbf{Student Assignment}
\begin{itemize}
    \item Use the OpenAI API to generate natural language explanations for:
  1. Why a task was decomposed into specific sub-tasks.
  2. Why a particular scheduling slot was chosen.
  3. Why human approval is required for certain actions.
\end{itemize}

\#\#\#\# \textbf{Deliverable}
Submit Python code showcasing the integration of generative AI for explanations. Provide example outputs for the 3 scenarios above.\par

    \vspace{16pt}

    \textcolor{cardBorder}{\hrule height 1pt}\vspace{8pt}

    \small \textcolor{gray}{When you are done with this task, mark it as complete to proceed to the next one.} \hfill \textbf{\textcolor{brandDark}{$\checkmark$ Completed}}

\end{taskbox}

\begin{taskbox}{Task 5: Audit for Responsible AI Practices}

    {\small \textcolor{gray}{2--3 hrs $\cdot$ Intermediate}}\\[12pt]

    \#\#\#\# \textbf{Scenario}
Before your LangGraph workflow goes live, you need to audit it for potential errors, biases, and ethical concerns. This includes ensuring that the AI assistant respects user preferences and privacy.

\#\#\#\# \textbf{Student Assignment}
\begin{itemize}
    \item Audit the LangGraph for:
  1. Potential bias in task decomposition and decision-making.
  2. Data privacy concerns in tool integrations.
  3. Failures to explain actions effectively.
    \item Propose solutions to address the identified issues.
\end{itemize}

\#\#\#\# \textbf{Deliverable}
Submit an audit report (500 words) detailing your findings and recommendations for improvement.\par

    \vspace{16pt}

    \textcolor{cardBorder}{\hrule height 1pt}\vspace{8pt}

    \small \textcolor{gray}{When you are done with this task, mark it as complete to proceed to the next one.} \hfill \textbf{\textcolor{brandDark}{$\checkmark$ Completed}}

\end{taskbox}

\begin{taskbox}{Task 6: Present Your Recommendation}

    {\small \textcolor{gray}{2--3 hrs $\cdot$ Intermediate}}\\[12pt]

    \#\#\#\# \textbf{Scenario}
It’s time to present your LangGraph workflow to FlowAI’s leadership team. They expect a clear explanation of how the system works, its benefits, and its reliability.

\#\#\#\# \textbf{Student Assignment}
\begin{itemize}
    \item Prepare a presentation including:
  1. An overview of your LangGraph workflow.
  2. Key features: task decomposition, tool integration, and human approval.
  3. Results from your audit and how issues were addressed.
  4. A demonstration of the workflow in action.
\end{itemize}

\#\#\#\# \textbf{Deliverable}
Submit a presentation deck (10-12 slides) with speaker notes and a recorded demo (3-5 minutes) of your LangGraph workflow in action.\par

    \vspace{16pt}

    \textcolor{cardBorder}{\hrule height 1pt}\vspace{8pt}

    \small \textcolor{gray}{When you are done with this task, mark it as complete to proceed to the next one.} \hfill \textbf{\textcolor{brandDark}{$\checkmark$ Completed}}

\end{taskbox}

\begin{completionbox}
    \textbf{Knowledge Assessment \& Verification:} Complete the online knowledge check, submit your code and presentation files for AI verification, and claim your \textbf{Skillzza Verified Certificate}.
\end{completionbox}

\newpage

% =============================================================================

% SIMULATION 30

% =============================================================================

\section*{AI Agent Developer – Customer Support Agent}

\addcontentsline{toc}{section}{AI Agent Developer – Customer Support Agent}

\begin{metabox}
    \begin{tabularx}{\textwidth}{@{}ll X@{}}
        \textbf{Industry:} & Technology & \textbf{Partner Enterprise:} \textbf{Skillzza Enterprise} \\
        \textbf{Duration:} & 12-18 hrs (Self-paced) & \textbf{Mode:} Individual, task-based simulation \\
        \textbf{Primary Skills:} & \multicolumn{2}{@{}l}{Python, AI, Data Analysis, Problem Solving}
    \end{tabularx}
\end{metabox}

\subsection*{Overview \& Problem Statement}

\#\#\# \textbf{Scenario:}
The future of customer service relies on intelligent automation. As an AI Agent Developer specializing in customer support, you are tasked with designing, training, and deploying an AI-powered virtual assistant that efficiently handles customer queries, resolves common issues, and escalates complex requests to human agents. This simulation places you in the role of a developer at a SaaS company that wants to automate FAQs and streamline escalation workflows while maintaining high customer satisfaction.

\#\#\# \textbf{Mission:}
Your mission is to design and deploy a conversational AI agent capable of handling customer FAQs, integrating with support tools, escalating unresolved tickets, and ensuring seamless communication while adhering to responsible AI principles. The ultimate goal is to enhance the company’s support experience and deliver measurable improvements in response time and resolution accuracy.

\#\#\# \textbf{Final Challenge:}
You will present a fully functional AI customer support agent prototype to the leadership team, showcasing how it addresses FAQs, handles escalations, and improves KPIs such as average response time, first-contact resolution rate, and customer satisfaction score.

---

\subsection*{Tasks \& Deliverables}

\begin{taskbox}{Task 1: Understand – Analyzing Customer Queries}

    {\small \textcolor{gray}{2--3 hrs $\cdot$ Intermediate}}\\[12pt]

    \#\#\#\# 
You’ve been provided with a dataset containing historical customer interaction logs. Your first step is to analyze the dataset, identify commonly asked questions, and categorize them into themes (e.g., billing, technical support, account management).

\#\#\#\# \\[10pt]
\vspace{8pt}
{\large \textbf{\textcolor{brandLight}{What you'll learn}}}\\[4pt]
\begin{itemize}
    \item Load the dataset and perform exploratory analysis to understand customer query patterns.
    \item Identify the top 10 FAQs based on frequency.
    \item Categorize FAQs into logical groups.
\end{itemize}

\#\#\#\# \vspace{8pt}
{\large \textbf{\textcolor{brandLight}{What you'll do}}}\\[4pt]
Submit a report containing:
1. Visualizations showing the distribution of query types.
2. A list of top 10 FAQs and their respective categories.\par

    \vspace{16pt}

    \textcolor{cardBorder}{\hrule height 1pt}\vspace{8pt}

    \small \textcolor{gray}{When you are done with this task, mark it as complete to proceed to the next one.} \hfill \textbf{\textcolor{brandDark}{$\checkmark$ Completed}}

\end{taskbox}

\begin{taskbox}{Task 2: Data/Setup – Preparing the FAQ Dataset}

    {\small \textcolor{gray}{2--3 hrs $\cdot$ Intermediate}}\\[12pt]

    \#\#\#\# 
You need to format the FAQ data for training a conversational AI model. This involves cleaning the text, removing duplicates, and ensuring FAQs are mapped to clear response templates.

\#\#\#\# \\[10pt]
\vspace{8pt}
{\large \textbf{\textcolor{brandLight}{What you'll learn}}}\\[4pt]
\begin{itemize}
    \item Preprocess the FAQ dataset (remove noise, duplicates, and irrelevant data).
    \item Structure the data into a Q\&A format (e.g., JSON or CSV format) suitable for training.
    \item Create sample response templates for each FAQ category.
\end{itemize}

\#\#\#\# \vspace{8pt}
{\large \textbf{\textcolor{brandLight}{What you'll do}}}\\[4pt]
Submit the cleaned and structured FAQ dataset along with response templates.\par

    \vspace{16pt}

    \textcolor{cardBorder}{\hrule height 1pt}\vspace{8pt}

    \small \textcolor{gray}{When you are done with this task, mark it as complete to proceed to the next one.} \hfill \textbf{\textcolor{brandDark}{$\checkmark$ Completed}}

\end{taskbox}

\begin{taskbox}{Task 3: Build/Execute – Developing the AI Agent}

    {\small \textcolor{gray}{2--3 hrs $\cdot$ Intermediate}}\\[12pt]

    \#\#\#\# 
It’s time to build the AI agent! Using OpenAI GPT (or an alternative NLP platform), train the model using the FAQ dataset and create dialogue flows to handle customer queries.

\#\#\#\# \\[10pt]
\vspace{8pt}
{\large \textbf{\textcolor{brandLight}{What you'll learn}}}\\[4pt]
\begin{itemize}
    \item Train the conversational AI model using the prepared dataset.
    \item Build logical dialogue flows for handling different FAQ categories.
    \item Test the model locally to ensure it responds accurately to queries.
\end{itemize}

\#\#\#\# \vspace{8pt}
{\large \textbf{\textcolor{brandLight}{What you'll do}}}\\[4pt]
Submit the Python code (or equivalent) for training and deploying the AI model, along with screenshots of test conversations.\par

    \vspace{16pt}

    \textcolor{cardBorder}{\hrule height 1pt}\vspace{8pt}

    \small \textcolor{gray}{When you are done with this task, mark it as complete to proceed to the next one.} \hfill \textbf{\textcolor{brandDark}{$\checkmark$ Completed}}

\end{taskbox}

\begin{taskbox}{Task 4: GenAI/Explanation – Handling Escalations}

    {\small \textcolor{gray}{2--3 hrs $\cdot$ Intermediate}}\\[12pt]

    \#\#\#\# 
Not all customer queries can be resolved by the AI agent. Create an escalation workflow where unresolved queries are handed over to human agents via API integration (e.g., Zendesk or Slack).

\#\#\#\# \\[10pt]
\vspace{8pt}
{\large \textbf{\textcolor{brandLight}{What you'll learn}}}\\[4pt]
\begin{itemize}
    \item Develop logic to identify unresolved queries (e.g., confidence threshold or “fallback” intents).
    \item Integrate APIs to escalate tickets to human agents.
    \item Test the escalation workflow with sample queries.
\end{itemize}

\#\#\#\# \vspace{8pt}
{\large \textbf{\textcolor{brandLight}{What you'll do}}}\\[4pt]
Submit the Python script for API integration and a flow diagram explaining the escalation process.\par

    \vspace{16pt}

    \textcolor{cardBorder}{\hrule height 1pt}\vspace{8pt}

    \small \textcolor{gray}{When you are done with this task, mark it as complete to proceed to the next one.} \hfill \textbf{\textcolor{brandDark}{$\checkmark$ Completed}}

\end{taskbox}

\begin{taskbox}{Task 5: Audit/Responsible AI – Ensuring Ethical AI Standards}

    {\small \textcolor{gray}{2--3 hrs $\cdot$ Intermediate}}\\[12pt]

    \#\#\#\# 
Your AI agent must comply with responsible AI principles. This includes ensuring unbiased responses, avoiding harmful outputs, and providing transparency in escalation workflows.

\#\#\#\# \\[10pt]
\vspace{8pt}
{\large \textbf{\textcolor{brandLight}{What you'll learn}}}\\[4pt]
\begin{itemize}
    \item Audit the AI model for bias in responses (e.g., gender, race, or cultural bias).
    \item Implement safeguards to prevent harmful outputs.
    \item Create a responsible AI checklist for your project.
\end{itemize}

\#\#\#\# \vspace{8pt}
{\large \textbf{\textcolor{brandLight}{What you'll do}}}\\[4pt]
Submit the audit report, safeguards implemented, and the responsible AI checklist.\par

    \vspace{16pt}

    \textcolor{cardBorder}{\hrule height 1pt}\vspace{8pt}

    \small \textcolor{gray}{When you are done with this task, mark it as complete to proceed to the next one.} \hfill \textbf{\textcolor{brandDark}{$\checkmark$ Completed}}

\end{taskbox}

\begin{taskbox}{Task 6: Present Recommendation – Improving KPIs}

    {\small \textcolor{gray}{2--3 hrs $\cdot$ Intermediate}}\\[12pt]

    \#\#\#\# 
Present your AI agent prototype to the leadership team, emphasizing how it improves customer support KPIs like average response time, CSAT, and FCR.

\#\#\#\# \\[10pt]
\vspace{8pt}
{\large \textbf{\textcolor{brandLight}{What you'll learn}}}\\[4pt]
\begin{itemize}
    \item Prepare a presentation with metrics showing the agent’s impact on KPIs.
    \item Create a live demo showcasing the agent handling queries and escalating tickets.
    \item Recommend further improvements to the AI agent.
\end{itemize}

\#\#\#\# \vspace{8pt}
{\large \textbf{\textcolor{brandLight}{What you'll do}}}\\[4pt]
Submit the presentation slides and demo video.\par

    \vspace{16pt}

    \textcolor{cardBorder}{\hrule height 1pt}\vspace{8pt}

    \small \textcolor{gray}{When you are done with this task, mark it as complete to proceed to the next one.} \hfill \textbf{\textcolor{brandDark}{$\checkmark$ Completed}}

\end{taskbox}

\begin{completionbox}
    \textbf{Knowledge Assessment \& Verification:} Complete the online knowledge check, submit your code and presentation files for AI verification, and claim your \textbf{Skillzza Verified Certificate}.
\end{completionbox}

\newpage

% =============================================================================

% SIMULATION 31

% =============================================================================

\section*{AI Agent Developer – Sales Qualification Agent}

\addcontentsline{toc}{section}{AI Agent Developer – Sales Qualification Agent}

\begin{metabox}
    \begin{tabularx}{\textwidth}{@{}ll X@{}}
        \textbf{Industry:} & Technology & \textbf{Partner Enterprise:} \textbf{Skillzza Enterprise} \\
        \textbf{Duration:} & 12-18 hrs (Self-paced) & \textbf{Mode:} Individual, task-based simulation \\
        \textbf{Primary Skills:} & \multicolumn{2}{@{}l}{Python, AI, Data Analysis, Problem Solving}
    \end{tabularx}
\end{metabox}

\subsection*{Overview \& Problem Statement}

\#\#\# \textbf{Scenario}  
You are hired as an AI Agent Developer at a fast-paced software company that specializes in building conversational AI solutions for sales teams. Your mission is to design and deploy an AI-powered Sales Qualification Agent that automates lead qualification, interacts with prospects, and updates the CRM system. The company is launching a new product and expects a surge in inbound leads. Your agent must efficiently score and qualify these leads, ensuring sales teams focus only on high-value opportunities.

\#\#\# \textbf{Mission}  
Your mission is to design, build, and test an advanced conversational AI agent capable of:  
\begin{itemize}
    \item Interpreting lead data.  
    \item Conducting initial qualification conversations.  
    \item Scoring leads based on predefined criteria.  
    \item Updating the CRM system with qualified leads.  
\end{itemize}

By the end of the simulation, you will present a fully functional Sales Qualification Agent prototype, along with a detailed report explaining its design logic, performance metrics, and ability to align with responsible AI principles.

\#\#\# \textbf{Final Challenge}  
Deploy your Sales Qualification Agent in a simulated environment where it processes real-world lead data. Ensure the agent generates accurate scores, provides clear qualification logic, and seamlessly integrates with a CRM system. Your report will be evaluated for clarity, technical accuracy, and ethical considerations.

---

\begin{completionbox}
    \textbf{Knowledge Assessment \& Verification:} Complete the online knowledge check, submit your code and presentation files for AI verification, and claim your \textbf{Skillzza Verified Certificate}.
\end{completionbox}

\newpage

% =============================================================================

% SIMULATION 32

% =============================================================================

\section*{AI Agent Developer – HR Recruitment Agent}

\addcontentsline{toc}{section}{AI Agent Developer – HR Recruitment Agent}

\begin{metabox}
    \begin{tabularx}{\textwidth}{@{}ll X@{}}
        \textbf{Industry:} & Technology & \textbf{Partner Enterprise:} \textbf{Skillzza Enterprise} \\
        \textbf{Duration:} & 12-18 hrs (Self-paced) & \textbf{Mode:} Individual, task-based simulation \\
        \textbf{Primary Skills:} & \multicolumn{2}{@{}l}{Python, AI, Data Analysis, Problem Solving}
    \end{tabularx}
\end{metabox}

\subsection*{Overview \& Problem Statement}

\#\#\# \textbf{Scenario}  
You are hired by "TalentOptimize Inc.," a global HR consulting firm that specializes in using AI technologies to streamline recruitment processes. Your role as an AI Agent Developer is to design, implement, and evaluate a recruitment agent powered by generative AI and recommendation algorithms. The agent should parse job descriptions (JDs), analyze candidate profiles, rank them for suitability, and incorporate bias control mechanisms to ensure fair and ethical hiring practices.  

\#\#\# \textbf{Mission}  
Your mission is to develop and deploy an AI-based recruitment agent capable of automating candidate screening and ranking while adhering to Responsible AI principles in HR. You will work with structured and unstructured datasets, leverage pre-trained generative AI models for text analysis, and implement bias detection and control mechanisms.  

\#\#\# \textbf{Final Challenge}  
By the end of this internship, you will present a fully functional AI HR Recruitment Agent prototype, complete with its screening and ranking logic, a bias audit report, and recommendations for ethical deployment in hiring workflows.  

---

\subsection*{Tasks \& Deliverables}

\begin{taskbox}{Task 1: Understand}

    {\small \textcolor{gray}{2--3 hrs $\cdot$ Intermediate}}\\[12pt]

     You are provided with a set of job descriptions (JDs) and candidate profiles. Your first task is to understand the structure and nuances of HR data.  


\vspace{8pt}
{\large \textbf{\textcolor{brandLight}{What you'll learn}}}\\[4pt]  
\begin{itemize}
    \item Analyze the provided dataset, which includes 20 job descriptions and 200 anonymized candidate profiles.  
    \item Identify key fields in the data (e.g., skills, experience, qualifications).  
    \item Summarize trends in JD requirements and candidate attributes.  
\end{itemize}

\vspace{8pt}
{\large \textbf{\textcolor{brandLight}{What you'll do}}}\\[4pt]  
A report detailing:  
\begin{itemize}
    \item Key trends in JD requirements (e.g., top skills, common qualifications).  
    \item Candidate profile diversity (e.g., experience levels, skillset coverage).\par
\end{itemize}

    \vspace{16pt}

    \textcolor{cardBorder}{\hrule height 1pt}\vspace{8pt}

    \small \textcolor{gray}{When you are done with this task, mark it as complete to proceed to the next one.} \hfill \textbf{\textcolor{brandDark}{$\checkmark$ Completed}}

\end{taskbox}

\begin{taskbox}{Task 2: Data/Setup}

    {\small \textcolor{gray}{2--3 hrs $\cdot$ Intermediate}}\\[12pt]

     Your team needs you to preprocess the HR dataset for AI model training.  


\vspace{8pt}
{\large \textbf{\textcolor{brandLight}{What you'll learn}}}\\[4pt]  
\begin{itemize}
    \item Clean and preprocess the JD and candidate profile datasets (e.g., remove duplicates, standardize fields).  
    \item Tokenize text fields (skills, experience, qualifications) using NLP techniques.  
    \item Split the dataset into training and testing subsets.  
\end{itemize}

\vspace{8pt}
{\large \textbf{\textcolor{brandLight}{What you'll do}}}\\[4pt]  
A cleaned and tokenized dataset saved as CSV files, ready for AI model input.\par

    \vspace{16pt}

    \textcolor{cardBorder}{\hrule height 1pt}\vspace{8pt}

    \small \textcolor{gray}{When you are done with this task, mark it as complete to proceed to the next one.} \hfill \textbf{\textcolor{brandDark}{$\checkmark$ Completed}}

\end{taskbox}

\begin{taskbox}{Task 3: Build/Execute}

    {\small \textcolor{gray}{2--3 hrs $\cdot$ Intermediate}}\\[12pt]

     You need to build the AI recruitment agent that matches JDs to candidate profiles and ranks them.  


\vspace{8pt}
{\large \textbf{\textcolor{brandLight}{What you'll learn}}}\\[4pt]  
\begin{itemize}
    \item Use OpenAI GPT for parsing job descriptions and summarizing key requirements.  
    \item Implement a candidate ranking algorithm based on a scoring system (e.g., matching skills, experience levels).  
    \item Generate top-10 ranked candidates for each JD.  
\end{itemize}

\vspace{8pt}
{\large \textbf{\textcolor{brandLight}{What you'll do}}}\\[4pt]  
Python scripts that:  
\begin{itemize}
    \item Parse and summarize JDs using generative AI.  
    \item Rank candidates and output top-10 matches per JD.\par
\end{itemize}

    \vspace{16pt}

    \textcolor{cardBorder}{\hrule height 1pt}\vspace{8pt}

    \small \textcolor{gray}{When you are done with this task, mark it as complete to proceed to the next one.} \hfill \textbf{\textcolor{brandDark}{$\checkmark$ Completed}}

\end{taskbox}

\begin{taskbox}{Task 4: GenAI/Explanation}

    {\small \textcolor{gray}{2--3 hrs $\cdot$ Intermediate}}\\[12pt]

     Your manager asks you to automate explanations for ranking decisions using generative AI.  


\vspace{8pt}
{\large \textbf{\textcolor{brandLight}{What you'll learn}}}\\[4pt]  
\begin{itemize}
    \item Use GPT to generate a justification for each candidate’s ranking based on matching criteria.  
    \item Ensure explanations are clear and align with the scoring logic.  
\end{itemize}

\vspace{8pt}
{\large \textbf{\textcolor{brandLight}{What you'll do}}}\\[4pt]  
A JSON file containing:  
\begin{itemize}
    \item Candidate rankings per JD.  
    \item GPT-generated explanations for each rank.\par
\end{itemize}

    \vspace{16pt}

    \textcolor{cardBorder}{\hrule height 1pt}\vspace{8pt}

    \small \textcolor{gray}{When you are done with this task, mark it as complete to proceed to the next one.} \hfill \textbf{\textcolor{brandDark}{$\checkmark$ Completed}}

\end{taskbox}

\begin{taskbox}{Task 5: Audit/Responsible AI}

    {\small \textcolor{gray}{2--3 hrs $\cdot$ Intermediate}}\\[12pt]

     You are tasked with auditing the AI recruitment agent for potential biases.  


\vspace{8pt}
{\large \textbf{\textcolor{brandLight}{What you'll learn}}}\\[4pt]  
\begin{itemize}
    \item Use IBM AI Fairness 360 to audit the ranking system for gender and ethnicity biases.  
    \item Identify any disparities in candidate rankings based on sensitive attributes.  
\end{itemize}

\vspace{8pt}
{\large \textbf{\textcolor{brandLight}{What you'll do}}}\\[4pt]  
A bias audit report that includes:  
\begin{itemize}
    \item Metrics evaluating fairness in rankings.  
    \item Recommendations for mitigating any biases found.\par
\end{itemize}

    \vspace{16pt}

    \textcolor{cardBorder}{\hrule height 1pt}\vspace{8pt}

    \small \textcolor{gray}{When you are done with this task, mark it as complete to proceed to the next one.} \hfill \textbf{\textcolor{brandDark}{$\checkmark$ Completed}}

\end{taskbox}

\begin{taskbox}{Task 6: Present Recommendation}

    {\small \textcolor{gray}{2--3 hrs $\cdot$ Intermediate}}\\[12pt]

     You need to present your recruitment agent prototype and bias audit findings to the HR leadership team.  


\vspace{8pt}
{\large \textbf{\textcolor{brandLight}{What you'll learn}}}\\[4pt]  
\begin{itemize}
    \item Create a dashboard using Streamlit to visualize JD summaries, candidate rankings, and bias audit results.  
    \item Prepare a 5-minute pitch explaining your prototype and its Responsible AI features.  
\end{itemize}

\vspace{8pt}
{\large \textbf{\textcolor{brandLight}{What you'll do}}}\\[4pt]  
\begin{itemize}
    \item A Streamlit dashboard URL showcasing your recruitment agent.  
    \item A slide deck summarizing your findings and recommendations.\par
\end{itemize}

    \vspace{16pt}

    \textcolor{cardBorder}{\hrule height 1pt}\vspace{8pt}

    \small \textcolor{gray}{When you are done with this task, mark it as complete to proceed to the next one.} \hfill \textbf{\textcolor{brandDark}{$\checkmark$ Completed}}

\end{taskbox}

\begin{completionbox}
    \textbf{Knowledge Assessment \& Verification:} Complete the online knowledge check, submit your code and presentation files for AI verification, and claim your \textbf{Skillzza Verified Certificate}.
\end{completionbox}

\newpage

% =============================================================================

% SIMULATION 33

% =============================================================================

\section*{Omnichannel CX Analyst – Map the Customer Journey}

\addcontentsline{toc}{section}{Omnichannel CX Analyst – Map the Customer Journey}

\begin{metabox}
    \begin{tabularx}{\textwidth}{@{}ll X@{}}
        \textbf{Industry:} & Technology & \textbf{Partner Enterprise:} \textbf{Skillzza Enterprise} \\
        \textbf{Duration:} & 12-18 hrs (Self-paced) & \textbf{Mode:} Individual, task-based simulation \\
        \textbf{Primary Skills:} & \multicolumn{2}{@{}l}{Python, AI, Data Analysis, Problem Solving}
    \end{tabularx}
\end{metabox}

\subsection*{Overview \& Problem Statement}

\#\#\# \textbf{Scenario}
Businesses today operate across multiple channels such as physical stores, websites, apps, social media, and customer service hotlines. To deliver seamless customer experiences, organizations need analysts who can track, evaluate, and optimize customer journeys across these touchpoints. As an \textbf{Omnichannel CX Analyst}, you'll help brands understand how their customers interact across channels and identify pain points that impact satisfaction and loyalty.

\#\#\# \textbf{Mission}
Your mission is to map the customer journey for a fictional retail brand, \textbf{Trendify}, which operates both online and offline. You will identify key touchpoints, analyze customer behavior across channels, uncover friction areas, and propose actionable recommendations to improve the experience.

\#\#\# \textbf{Final Challenge}
The internship culminates in presenting a detailed \textbf{Customer Journey Map}, channel analysis, and strategic recommendations to Trendify’s executive team. Your report will detail pain points, analytics findings, and actionable fixes to improve omnichannel customer experience.

---

\begin{completionbox}
    \textbf{Knowledge Assessment \& Verification:} Complete the online knowledge check, submit your code and presentation files for AI verification, and claim your \textbf{Skillzza Verified Certificate}.
\end{completionbox}

\newpage

% =============================================================================

% SIMULATION 34

% =============================================================================

\section*{AI Customer Service Manager – Design a GenAI Contact Centre}

\addcontentsline{toc}{section}{AI Customer Service Manager – Design a GenAI Contact Centre}

\begin{metabox}
    \begin{tabularx}{\textwidth}{@{}ll X@{}}
        \textbf{Industry:} & Technology & \textbf{Partner Enterprise:} \textbf{Skillzza Enterprise} \\
        \textbf{Duration:} & 12-18 hrs (Self-paced) & \textbf{Mode:} Individual, task-based simulation \\
        \textbf{Primary Skills:} & \multicolumn{2}{@{}l}{Python, AI, Data Analysis, Problem Solving}
    \end{tabularx}
\end{metabox}

\subsection*{Overview \& Problem Statement}

\#\#\# \textbf{Scenario}  
As businesses scale, customer service is often their first line of defense to retain customers, resolve complaints, and build brand loyalty. However, traditional contact centres face challenges like long response times, high agent turnover, and inconsistent customer experiences. Imagine being brought on board by a leading e-commerce company, "ShopSphere," to revolutionize their customer service operations using Generative AI (GenAI). Your mission is to design and implement an AI-powered contact centre that improves efficiency, enhances customer satisfaction, and integrates seamlessly with human agents.  

\#\#\# \textbf{Mission}  
Your goal is to identify inefficiencies in ShopSphere’s current contact centre workflow, implement AI-powered solutions to address gaps, and design an agent-assist system using GenAI to optimize operations. You will evaluate KPIs before and after integration and present your findings to senior leadership.  

\#\#\# \textbf{Final Challenge}  
Build a GenAI-powered contact centre framework for ShopSphere by leveraging AI tools for process automation, conversation modeling, and sentiment analysis. You will need to produce a detailed workflow diagram, a GenAI implementation plan, and a KPI improvement report for final submission.

---

\subsection*{Tasks \& Deliverables}

\begin{taskbox}{Task 1: Understand}

    {\small \textcolor{gray}{2--3 hrs $\cdot$ Intermediate}}\\[12pt]

    \#\#\#\# \textbf{Scenario}  
ShopSphere’s leadership team has provided the current contact centre workflow and customer complaints data. You need to analyze the workflow and identify pain points affecting service quality.  

\#\#\#\# \textbf{Student Assignment}  
\begin{itemize}
    \item Study the provided workflow diagram and identify bottlenecks.  
    \item Analyze customer complaints data to spot trends (e.g., delays, unresolved issues).  
    \item List areas where AI solutions can improve efficiency.  
\end{itemize}

\#\#\#\# \textbf{Deliverable}  
\begin{itemize}
    \item A report summarizing pain points and AI opportunities.  
    \item Annotated workflow diagram highlighting inefficiencies.\par
\end{itemize}

    \vspace{16pt}

    \textcolor{cardBorder}{\hrule height 1pt}\vspace{8pt}

    \small \textcolor{gray}{When you are done with this task, mark it as complete to proceed to the next one.} \hfill \textbf{\textcolor{brandDark}{$\checkmark$ Completed}}

\end{taskbox}

\begin{taskbox}{Task 2: Data/Setup}

    {\small \textcolor{gray}{2--3 hrs $\cdot$ Intermediate}}\\[12pt]

    \#\#\#\# \textbf{Scenario}  
You will prepare the dataset for AI implementation. This involves cleaning the data, performing exploratory data analysis (EDA), and identifying meaningful features like response times, resolution rates, and sentiment scores.  

\#\#\#\# \textbf{Student Assignment}  
\begin{itemize}
    \item Perform EDA on the customer complaints dataset.  
    \item Clean and preprocess data, ensuring it’s ready for AI modeling.  
    \item Generate sentiment scores for each complaint using a pre-trained sentiment analysis model.  
\end{itemize}

\#\#\#\# \textbf{Deliverable}  
\begin{itemize}
    \item Preprocessed dataset with added sentiment scores.  
    \item EDA report with visualizations showing key trends (e.g., recurring complaint types).\par
\end{itemize}

    \vspace{16pt}

    \textcolor{cardBorder}{\hrule height 1pt}\vspace{8pt}

    \small \textcolor{gray}{When you are done with this task, mark it as complete to proceed to the next one.} \hfill \textbf{\textcolor{brandDark}{$\checkmark$ Completed}}

\end{taskbox}

\begin{taskbox}{Task 3: Build/Execute}

    {\small \textcolor{gray}{2--3 hrs $\cdot$ Intermediate}}\\[12pt]

    \#\#\#\# \textbf{Scenario}  
Using OpenAI GPT APIs, you will create a chatbot that can handle customer inquiries, classify complaints, and provide resolution suggestions. Integrate the chatbot with a basic agent-assist system to help human agents with real-time recommendations.  

\#\#\#\# \textbf{Student Assignment}  
\begin{itemize}
    \item Create a Python script to integrate with OpenAI GPT API for generating responses.  
    \item Train the chatbot to classify complaints into categories (e.g., payment issues, delivery delays).  
    \item Build a basic agent-assist system that suggests responses based on customer sentiment and query type.  
\end{itemize}

\#\#\#\# \textbf{Deliverable}  
\begin{itemize}
    \item Python script for chatbot integration.  
    \item Agent-assist prototype with sample queries and responses.\par
\end{itemize}

    \vspace{16pt}

    \textcolor{cardBorder}{\hrule height 1pt}\vspace{8pt}

    \small \textcolor{gray}{When you are done with this task, mark it as complete to proceed to the next one.} \hfill \textbf{\textcolor{brandDark}{$\checkmark$ Completed}}

\end{taskbox}

\begin{taskbox}{Task 4: GenAI/Explanation}

    {\small \textcolor{gray}{2--3 hrs $\cdot$ Intermediate}}\\[12pt]

    \#\#\#\# \textbf{Scenario}  
Explain the rationale behind your GenAI implementation to ShopSphere’s leadership team. You need to justify how the chatbot and agent-assist system will improve customer service efficiency and satisfaction.  

\#\#\#\# \textbf{Student Assignment}  
\begin{itemize}
    \item Create a presentation explaining the GenAI model, its features, and functionality.  
    \item Use case studies/examples to demonstrate its capabilities.  
    \item Highlight the impact on KPIs like CSAT, AHT, and FCR.  
\end{itemize}

\#\#\#\# \textbf{Deliverable}  
\begin{itemize}
    \item Presentation deck with visuals (graphs, process flows, case studies).\par
\end{itemize}

    \vspace{16pt}

    \textcolor{cardBorder}{\hrule height 1pt}\vspace{8pt}

    \small \textcolor{gray}{When you are done with this task, mark it as complete to proceed to the next one.} \hfill \textbf{\textcolor{brandDark}{$\checkmark$ Completed}}

\end{taskbox}

\begin{taskbox}{Task 5: Audit/Responsible AI}

    {\small \textcolor{gray}{2--3 hrs $\cdot$ Intermediate}}\\[12pt]

    \#\#\#\# \textbf{Scenario}  
To ensure ethical AI implementation, ShopSphere requires you to audit the chatbot system for bias, fairness, and transparency. You will also propose fail-safe mechanisms for cases where human intervention is necessary.  

\#\#\#\# \textbf{Student Assignment}  
\begin{itemize}
    \item Analyze chatbot responses for potential bias or inaccuracies.  
    \item Develop guidelines for ethical AI use in customer service.  
    \item Suggest mechanisms for agent override when the AI fails or escalates issues.  
\end{itemize}

\#\#\#\# \textbf{Deliverable}  
\begin{itemize}
    \item Audit report on fairness and bias.  
    \item Ethical AI guidelines document.\par
\end{itemize}

    \vspace{16pt}

    \textcolor{cardBorder}{\hrule height 1pt}\vspace{8pt}

    \small \textcolor{gray}{When you are done with this task, mark it as complete to proceed to the next one.} \hfill \textbf{\textcolor{brandDark}{$\checkmark$ Completed}}

\end{taskbox}

\begin{taskbox}{Task 6: Present Recommendation}

    {\small \textcolor{gray}{2--3 hrs $\cdot$ Intermediate}}\\[12pt]

    \#\#\#\# \textbf{Scenario}  
Prepare a final presentation for ShopSphere’s executives, summarizing your findings, proposed GenAI solutions, and projected KPI improvements.  

\#\#\#\# \textbf{Student Assignment}  
\begin{itemize}
    \item Consolidate all previous tasks into a cohesive report.  
    \item Prepare a visual presentation highlighting the proposed GenAI contact centre framework.  
    \item Include KPI projections and a roadmap for implementation.  
\end{itemize}

\#\#\#\# \textbf{Deliverable}  
\begin{itemize}
    \item Final report (PDF).  
    \item Presentation deck (PPT/Google Slides).\par
\end{itemize}

    \vspace{16pt}

    \textcolor{cardBorder}{\hrule height 1pt}\vspace{8pt}

    \small \textcolor{gray}{When you are done with this task, mark it as complete to proceed to the next one.} \hfill \textbf{\textcolor{brandDark}{$\checkmark$ Completed}}

\end{taskbox}

\begin{completionbox}
    \textbf{Knowledge Assessment \& Verification:} Complete the online knowledge check, submit your code and presentation files for AI verification, and claim your \textbf{Skillzza Verified Certificate}.
\end{completionbox}

\newpage

% =============================================================================

% SIMULATION 35

% =============================================================================

\section*{AI Avatar Designer – Build a Virtual Receptionist}

\addcontentsline{toc}{section}{AI Avatar Designer – Build a Virtual Receptionist}

\begin{metabox}
    \begin{tabularx}{\textwidth}{@{}ll X@{}}
        \textbf{Industry:} & Technology & \textbf{Partner Enterprise:} \textbf{Skillzza Enterprise} \\
        \textbf{Duration:} & 12-18 hrs (Self-paced) & \textbf{Mode:} Individual, task-based simulation \\
        \textbf{Primary Skills:} & \multicolumn{2}{@{}l}{Python, AI, Data Analysis, Problem Solving}
    \end{tabularx}
\end{metabox}

\subsection*{Overview \& Problem Statement}

\textbf{Scenario:}  
Imagine you are part of an innovative AI startup tasked with designing a virtual receptionist for a global coworking space brand. Your goal is to create an AI-powered avatar that can welcome visitors, answer frequently asked questions, and assist with scheduling meeting rooms. The avatar must be engaging, professional, and capable of holding natural conversations with users.

\textbf{Mission:}  
In this job simulation, you will design, build, and test a conversational AI avatar. You'll start by defining a persona for the virtual receptionist, build its knowledge base, integrate conversational logic, create a script for its interactions, and finally test and improve its performance.

\textbf{Final Challenge:}  
You will deliver a fully functional AI-powered virtual receptionist that can handle at least 10 different visitor scenarios, including answering questions, addressing common complaints, and performing basic scheduling. Your final deliverable will include the avatar’s persona document, conversation flow design, scripts, and deployment test results.

---

\subsection*{Tasks \& Deliverables}

\begin{taskbox}{Task 1: Understand the Problem and Define the Avatar Persona}

    {\small \textcolor{gray}{2--3 hrs $\cdot$ Intermediate}}\\[12pt]

      
The coworking space brand has asked for a friendly and professional virtual receptionist that can handle interactions with users in English. The receptionist must reflect the brand’s values: professionalism, inclusivity, and innovation.


\vspace{8pt}
{\large \textbf{\textcolor{brandLight}{What you'll learn}}}\\[4pt]  
\begin{itemize}
    \item Study the company background and target audience.  
    \item Define a detailed persona for the virtual receptionist, including its appearance, tone, and key personality traits.  
    \item Draft a document that outlines the avatar’s objective and scope of interactions (e.g., greeting visitors, answering FAQs, scheduling appointments).  
\end{itemize}

\vspace{8pt}
{\large \textbf{\textcolor{brandLight}{What you'll do}}}\\[4pt]  
\begin{itemize}
    \item A detailed persona document with at least the following fields: Name, Appearance, Personality Traits, Tone of Voice, and Key Responsibilities.\par
\end{itemize}

    \vspace{16pt}

    \textcolor{cardBorder}{\hrule height 1pt}\vspace{8pt}

    \small \textcolor{gray}{When you are done with this task, mark it as complete to proceed to the next one.} \hfill \textbf{\textcolor{brandDark}{$\checkmark$ Completed}}

\end{taskbox}

\begin{taskbox}{Task 2: Build the Knowledge Base for the Avatar}

    {\small \textcolor{gray}{2--3 hrs $\cdot$ Intermediate}}\\[12pt]

      
The virtual receptionist will need to answer FAQs and provide assistance to users. You are tasked with building an initial knowledge base that will serve as the foundation for the avatar’s conversational abilities.


\vspace{8pt}
{\large \textbf{\textcolor{brandLight}{What you'll learn}}}\\[4pt]  
\begin{itemize}
    \item Research common FAQs for coworking spaces (e.g., pricing, amenities, booking rules).  
    \item Organize the questions and answers into a structured format using Google Sheets or Airtable.  
    \item Include at least 20 questions and answers in the knowledge base.  
\end{itemize}

\vspace{8pt}
{\large \textbf{\textcolor{brandLight}{What you'll do}}}\\[4pt]  
\begin{itemize}
    \item A knowledge base in Google Sheets or Airtable with fields: Question, Answer, and Tags (e.g., pricing, scheduling).\par
\end{itemize}

    \vspace{16pt}

    \textcolor{cardBorder}{\hrule height 1pt}\vspace{8pt}

    \small \textcolor{gray}{When you are done with this task, mark it as complete to proceed to the next one.} \hfill \textbf{\textcolor{brandDark}{$\checkmark$ Completed}}

\end{taskbox}

\begin{taskbox}{Task 3: Configure the Avatar’s Chatbot and Train the AI}

    {\small \textcolor{gray}{2--3 hrs $\cdot$ Intermediate}}\\[12pt]

      
It’s time to bring the virtual receptionist to life! You will set up a chatbot platform, import the knowledge base, and train the AI to handle user queries.


\vspace{8pt}
{\large \textbf{\textcolor{brandLight}{What you'll learn}}}\\[4pt]  
\begin{itemize}
    \item Choose a conversational AI platform (e.g., Dialogflow, Rasa).  
    \item Create intents and entities based on the knowledge base.  
    \item Train the AI using the imported data and refine its ability to understand user queries.  
\end{itemize}

\vspace{8pt}
{\large \textbf{\textcolor{brandLight}{What you'll do}}}\\[4pt]  
\begin{itemize}
    \item A working chatbot model that can handle the 20 FAQs with appropriate responses.  
    \item Screenshots or a short video demonstrating the chatbot’s functionality.\par
\end{itemize}

    \vspace{16pt}

    \textcolor{cardBorder}{\hrule height 1pt}\vspace{8pt}

    \small \textcolor{gray}{When you are done with this task, mark it as complete to proceed to the next one.} \hfill \textbf{\textcolor{brandDark}{$\checkmark$ Completed}}

\end{taskbox}

\begin{taskbox}{Task 4: Generate AI Avatar Scripts for Conversations}

    {\small \textcolor{gray}{2--3 hrs $\cdot$ Intermediate}}\\[12pt]

      
The virtual receptionist’s conversations must feel natural and engaging. You will design scripted responses and use a generative AI tool to generate variations of these scripts.


\vspace{8pt}
{\large \textbf{\textcolor{brandLight}{What you'll learn}}}\\[4pt]  
\begin{itemize}
    \item Write a list of 10 scripted dialogues for common scenarios (e.g., greeting visitors, answering FAQs, handling complaints).  
    \item Use a generative AI tool like OpenAI’s GPT to create variations of these scripts.  
    \item Ensure the scripts align with the persona’s tone and personality.  
\end{itemize}

\vspace{8pt}
{\large \textbf{\textcolor{brandLight}{What you'll do}}}\\[4pt]  
\begin{itemize}
    \item A document with the original scripts and AI-generated variations for the 10 scenarios.\par
\end{itemize}

    \vspace{16pt}

    \textcolor{cardBorder}{\hrule height 1pt}\vspace{8pt}

    \small \textcolor{gray}{When you are done with this task, mark it as complete to proceed to the next one.} \hfill \textbf{\textcolor{brandDark}{$\checkmark$ Completed}}

\end{taskbox}

\begin{taskbox}{Task 5: Test and Audit the Avatar for Responsible AI Practices}

    {\small \textcolor{gray}{2--3 hrs $\cdot$ Intermediate}}\\[12pt]

      
Before launching the virtual receptionist, you need to test its performance and ensure it aligns with ethical AI principles, such as avoiding bias and ensuring inclusivity.


\vspace{8pt}
{\large \textbf{\textcolor{brandLight}{What you'll learn}}}\\[4pt]  
\begin{itemize}
    \item Test the avatar with 10 different user inputs, including edge cases (e.g., ambiguous or offensive queries).  
    \item Document the avatar’s responses and identify areas for improvement.  
    \item Write a short report on how the avatar adheres to responsible AI principles.
\end{itemize}

\vspace{8pt}
{\large \textbf{\textcolor{brandLight}{What you'll do}}}\\[4pt]  
\begin{itemize}
    \item A testing report with user input, avatar output, and identified issues.  
    \item A brief summary of the ethical considerations addressed in the avatar’s design.\par
\end{itemize}

    \vspace{16pt}

    \textcolor{cardBorder}{\hrule height 1pt}\vspace{8pt}

    \small \textcolor{gray}{When you are done with this task, mark it as complete to proceed to the next one.} \hfill \textbf{\textcolor{brandDark}{$\checkmark$ Completed}}

\end{taskbox}

\begin{taskbox}{Task 6: Present Your Final Avatar and Recommendations}

    {\small \textcolor{gray}{2--3 hrs $\cdot$ Intermediate}}\\[12pt]

      
The client is excited to see the final virtual receptionist in action. You will present your work, explain your design choices, and outline recommendations for future improvements.


\vspace{8pt}
{\large \textbf{\textcolor{brandLight}{What you'll learn}}}\\[4pt]  
\begin{itemize}
    \item Create a presentation (slides or a video) showcasing the avatar’s persona, knowledge base, conversation flow, and testing results.  
    \item Include a demo of the avatar handling at least 5 different scenarios.  
    \item Suggest improvements or future features for the avatar.  
\end{itemize}

\vspace{8pt}
{\large \textbf{\textcolor{brandLight}{What you'll do}}}\\[4pt]  
\begin{itemize}
    \item A presentation (PDF or video) with the avatar’s demo and recommendations for future iterations.\par
\end{itemize}

    \vspace{16pt}

    \textcolor{cardBorder}{\hrule height 1pt}\vspace{8pt}

    \small \textcolor{gray}{When you are done with this task, mark it as complete to proceed to the next one.} \hfill \textbf{\textcolor{brandDark}{$\checkmark$ Completed}}

\end{taskbox}

\begin{completionbox}
    \textbf{Knowledge Assessment \& Verification:} Complete the online knowledge check, submit your code and presentation files for AI verification, and claim your \textbf{Skillzza Verified Certificate}.
\end{completionbox}

\newpage

% =============================================================================

% SIMULATION 36

% =============================================================================

\section*{AI Avatar Product Manager – Design an AI Digital Human}

\addcontentsline{toc}{section}{AI Avatar Product Manager – Design an AI Digital Human}

\begin{metabox}
    \begin{tabularx}{\textwidth}{@{}ll X@{}}
        \textbf{Industry:} & Technology & \textbf{Partner Enterprise:} \textbf{Skillzza Enterprise} \\
        \textbf{Duration:} & 12-18 hrs (Self-paced) & \textbf{Mode:} Individual, task-based simulation \\
        \textbf{Primary Skills:} & \multicolumn{2}{@{}l}{Python, AI, Data Analysis, Problem Solving}
    \end{tabularx}
\end{metabox}

\subsection*{Overview \& Problem Statement}

\#\#\# Scenario  
Imagine a future where digital human avatars handle customer support, provide personalized services, and even act as virtual companions. As an AI Avatar Product Manager, you’ll lead the development of a cutting-edge AI Digital Human. Your goal is to bridge the gap between user needs and technical capabilities by defining the use case, crafting the persona, designing conversational interactions, and ensuring the AI avatar meets key performance benchmarks.

\#\#\# Mission  
Your mission is to conceptualize, design, and evaluate an AI-powered digital human capable of holding natural, engaging, and contextually relevant conversations. You will work across the lifecycle of product development, from ideation to KPI measurement, ensuring the avatar meets both business and user expectations.

\#\#\# Final Challenge  
At the end of this simulation, you will present a product pitch for your AI Digital Human. Your pitch will include the use case, target user persona, core capabilities, conversation architecture, and the designed KPIs to measure success. You’ll also address ethical considerations and responsible AI practices in your design.

---

\subsection*{Tasks \& Deliverables}

\begin{taskbox}{Task 1: Understand – Define the Use Case}

    {\small \textcolor{gray}{2--3 hrs $\cdot$ Intermediate}}\\[12pt]

    \#\#\#\# Scenario:  
You have been hired by a healthcare startup to design an AI Digital Human that can act as a mental health support assistant for users. The CEO wants the avatar to provide empathetic support, answer common mental health questions, and encourage users to seek professional help if needed.  

\#\#\#\# Student Assignment:  
\begin{itemize}
    \item Research the mental health chatbot use case.  
    \item Define the target audience and the specific value the AI avatar will provide.  
    \item Identify the key challenges this type of AI avatar might face (e.g., ethical concerns, privacy, accuracy).  
\end{itemize}

\#\#\#\# Deliverable:  
A 300-word use case document outlining:  
\begin{itemize}
    \item Problem statement.  
    \item Target audience and their pain points.  
    \item Expected value and outcomes.\par
\end{itemize}

    \vspace{16pt}

    \textcolor{cardBorder}{\hrule height 1pt}\vspace{8pt}

    \small \textcolor{gray}{When you are done with this task, mark it as complete to proceed to the next one.} \hfill \textbf{\textcolor{brandDark}{$\checkmark$ Completed}}

\end{taskbox}

\begin{taskbox}{Task 2: Data/Setup – Create the Persona}

    {\small \textcolor{gray}{2--3 hrs $\cdot$ Intermediate}}\\[12pt]

    \#\#\#\# Scenario:  
Your team has identified the need for a trustworthy and empathetic digital persona. It’s now your job to create a detailed persona that aligns with the avatar’s role as a mental health support assistant.  

\#\#\#\# Student Assignment:  
\begin{itemize}
    \item Create a persona profile that includes:  
    \item Name, age, and visual appearance.  
    \item Communication style (e.g., tone, vocabulary).  
    \item Emotional intelligence capabilities.  
    \item Use design software (e.g., Figma or Canva) to create a visual representation of the avatar.  
\end{itemize}

\#\#\#\# Deliverable:  
1. A one-page persona profile document.  
2. A visual design of the AI avatar.\par

    \vspace{16pt}

    \textcolor{cardBorder}{\hrule height 1pt}\vspace{8pt}

    \small \textcolor{gray}{When you are done with this task, mark it as complete to proceed to the next one.} \hfill \textbf{\textcolor{brandDark}{$\checkmark$ Completed}}

\end{taskbox}

\begin{taskbox}{Task 3: Build/Execute – Design the Conversation Architecture}

    {\small \textcolor{gray}{2--3 hrs $\cdot$ Intermediate}}\\[12pt]

    \#\#\#\# Scenario:  
Now that the persona is ready, it’s time to design the avatar’s conversational capabilities. The goal is to ensure the AI provides empathetic, accurate, and on-brand responses while maintaining user engagement.  

\#\#\#\# Student Assignment:  
\begin{itemize}
    \item Outline a conversation flow for the following scenarios:  
  1. User reports feeling anxious.  
  2. User asks for tips to improve sleep.  
  3. User requests information about professional counseling.  
    \item Define intents, entities, and sample utterances using a conversational AI platform like DialogFlow.  
\end{itemize}

\#\#\#\# Deliverable:  
1. A conversation flow diagram (PDF or image).  
2. DialogFlow project export file containing defined intents and entities.\par

    \vspace{16pt}

    \textcolor{cardBorder}{\hrule height 1pt}\vspace{8pt}

    \small \textcolor{gray}{When you are done with this task, mark it as complete to proceed to the next one.} \hfill \textbf{\textcolor{brandDark}{$\checkmark$ Completed}}

\end{taskbox}

\begin{taskbox}{Task 4: GenAI/Explanation – Implement Generative AI for Dynamic Responses}

    {\small \textcolor{gray}{2--3 hrs $\cdot$ Intermediate}}\\[12pt]

    \#\#\#\# Scenario:  
To make the AI avatar more dynamic, you decide to integrate a generative AI model like GPT-4 for open-ended conversations. Your task is to ensure it generates responses that are empathetic, accurate, and contextually relevant.  

\#\#\#\# Student Assignment:  
\begin{itemize}
    \item Connect GPT-4 API to the conversational AI platform.  
    \item Fine-tune the response generation by writing a set of prompts.  
    \item Test sample conversations to ensure the avatar’s responses align with the defined persona.  
\end{itemize}

\#\#\#\# Deliverable:  
1. A set of 5 GPT-4 prompts optimized for empathetic responses.  
2. A recorded or text-based dialogue demonstrating the generative AI in action.\par

    \vspace{16pt}

    \textcolor{cardBorder}{\hrule height 1pt}\vspace{8pt}

    \small \textcolor{gray}{When you are done with this task, mark it as complete to proceed to the next one.} \hfill \textbf{\textcolor{brandDark}{$\checkmark$ Completed}}

\end{taskbox}

\begin{taskbox}{Task 5: Audit/Responsible AI – Evaluate for Bias and Responsiveness}

    {\small \textcolor{gray}{2--3 hrs $\cdot$ Intermediate}}\\[12pt]

    \#\#\#\# Scenario:  
Before the avatar goes live, you need to ensure it adheres to responsible AI principles. This involves testing for bias, inclusivity, and response accuracy.  

\#\#\#\# Student Assignment:  
\begin{itemize}
    \item Test the avatar across diverse user scenarios (e.g., different age groups, genders, and cultural backgrounds).  
    \item Identify potential biases in the responses.  
    \item Suggest improvements to make the avatar more inclusive.  
\end{itemize}

\#\#\#\# Deliverable:  
1. A bias evaluation report (500 words).  
2. Spreadsheet of test results (e.g., responsiveness, accuracy, inclusivity ratings).\par

    \vspace{16pt}

    \textcolor{cardBorder}{\hrule height 1pt}\vspace{8pt}

    \small \textcolor{gray}{When you are done with this task, mark it as complete to proceed to the next one.} \hfill \textbf{\textcolor{brandDark}{$\checkmark$ Completed}}

\end{taskbox}

\begin{taskbox}{Task 6: Present Recommendation – Final Product Pitch}

    {\small \textcolor{gray}{2--3 hrs $\cdot$ Intermediate}}\\[12pt]

    \#\#\#\# Scenario:  
You’re ready to present your AI Digital Human to the healthcare startup’s leadership team. Your pitch needs to cover the avatar’s use case, persona, conversation architecture, generative AI integration, and KPIs.  

\#\#\#\# Student Assignment:  
\begin{itemize}
    \item Create a 5-slide presentation covering:  
    \item Use case and target audience.  
    \item Persona design.  
    \item Conversation architecture and generative AI integration.  
    \item KPIs and performance metrics.  
    \item Ethical considerations and audit results.  
    \item Record a 5-minute video presenting your pitch.  
\end{itemize}

\#\#\#\# Deliverable:  
1. A 5-slide presentation (PDF or PPT format).  
2. A 5-minute video pitch recording.\par

    \vspace{16pt}

    \textcolor{cardBorder}{\hrule height 1pt}\vspace{8pt}

    \small \textcolor{gray}{When you are done with this task, mark it as complete to proceed to the next one.} \hfill \textbf{\textcolor{brandDark}{$\checkmark$ Completed}}

\end{taskbox}

\begin{completionbox}
    \textbf{Knowledge Assessment \& Verification:} Complete the online knowledge check, submit your code and presentation files for AI verification, and claim your \textbf{Skillzza Verified Certificate}.
\end{completionbox}

\newpage

% =============================================================================

% SIMULATION 37

% =============================================================================

\section*{Social Media Intelligence Analyst – Decode a Brand Conversation}

\addcontentsline{toc}{section}{Social Media Intelligence Analyst – Decode a Brand Conversation}

\begin{metabox}
    \begin{tabularx}{\textwidth}{@{}ll X@{}}
        \textbf{Industry:} & Technology & \textbf{Partner Enterprise:} \textbf{Skillzza Enterprise} \\
        \textbf{Duration:} & 12-18 hrs (Self-paced) & \textbf{Mode:} Individual, task-based simulation \\
        \textbf{Primary Skills:} & \multicolumn{2}{@{}l}{Python, AI, Data Analysis, Problem Solving}
    \end{tabularx}
\end{metabox}

\subsection*{Overview \& Problem Statement}

\#\#\# \textbf{Scenario}
You’ve been hired as a Social Media Intelligence Analyst for a leading digital marketing agency. Your client is a fast-growing brand in the lifestyle and fitness industry that wants to understand how customers perceive its products online. The brand also wants to compare its social media performance with key competitors and identify actionable insights to improve its online reputation.

Social media platforms like Twitter, Instagram, and forums are flooded with conversations about the brand, its competitors, and related topics. Your mission is to decode these conversations using social media analytics techniques and provide insights that drive business decisions.

---

\#\#\# \textbf{Mission}
Your role is to analyze the brand’s social media presence, interpret sentiment, identify recurring themes, and benchmark its performance against competitors. Your work will directly influence the brand’s digital marketing strategy.

---

\#\#\# \textbf{Final Challenge}
By the end of this internship, you will present a comprehensive social media intelligence report covering:
1. Customer sentiment analysis for the brand.
2. Thematic analysis of conversations.
3. Competitive benchmarking insights.
4. Actionable recommendations for improving the brand’s social media strategy.

---

\begin{completionbox}
    \textbf{Knowledge Assessment \& Verification:} Complete the online knowledge check, submit your code and presentation files for AI verification, and claim your \textbf{Skillzza Verified Certificate}.
\end{completionbox}

\newpage

% =============================================================================

% SIMULATION 38

% =============================================================================

\section*{UI/UX Designer – Design an AI Product Experience}

\addcontentsline{toc}{section}{UI/UX Designer – Design an AI Product Experience}

\begin{metabox}
    \begin{tabularx}{\textwidth}{@{}ll X@{}}
        \textbf{Industry:} & Technology & \textbf{Partner Enterprise:} \textbf{Skillzza Enterprise} \\
        \textbf{Duration:} & 12-18 hrs (Self-paced) & \textbf{Mode:} Individual, task-based simulation \\
        \textbf{Primary Skills:} & \multicolumn{2}{@{}l}{Python, AI, Data Analysis, Problem Solving}
    \end{tabularx}
\end{metabox}

\subsection*{Overview \& Problem Statement}

\textbf{Scenario:}  
Imagine you’ve been hired as a UI/UX Designer for an innovative AI startup developing a cutting-edge product: an AI-powered mental wellness assistant. Your goal is to ensure the product provides an intuitive, user-friendly, and visually appealing experience while addressing end-user needs effectively.

\textbf{Mission:}  
Your mission is to guide the design process for this AI product by conducting user research, developing personas, designing user flows, creating wireframes, and testing usability. By the end of this simulation, you will deliver a prototype-ready design that balances functionality, accessibility, and aesthetics.

\textbf{Final Challenge:}  
Present your design recommendations, supported by research findings, design decisions, and usability test outcomes, to the product team. Convince stakeholders that your design ensures a delightful user experience while meeting business goals.

---

\subsection*{Tasks \& Deliverables}

\begin{taskbox}{Task 1: Understand – Conduct User Research}

    {\small \textcolor{gray}{2--3 hrs $\cdot$ Intermediate}}\\[12pt]

      
You are tasked with understanding the target audience for the AI mental wellness assistant. You need to uncover user pain points, preferences, and expectations for an AI-powered product that supports mental health.


\vspace{8pt}
{\large \textbf{\textcolor{brandLight}{What you'll learn}}}\\[4pt]  
\begin{itemize}
    \item Create a Google Form survey with 10 targeted questions to gather insights on user needs (e.g., "What challenges do you face in maintaining mental wellness-).  
    \item Conduct 3 virtual interviews with potential users to uncover qualitative insights.
\end{itemize}

\vspace{8pt}
{\large \textbf{\textcolor{brandLight}{What you'll do}}}\\[4pt]  
Submit a short report summarizing the survey results and key findings from interviews. Include charts or graphs representing quantitative data.\par

    \vspace{16pt}

    \textcolor{cardBorder}{\hrule height 1pt}\vspace{8pt}

    \small \textcolor{gray}{When you are done with this task, mark it as complete to proceed to the next one.} \hfill \textbf{\textcolor{brandDark}{$\checkmark$ Completed}}

\end{taskbox}

\begin{taskbox}{Task 2: Data/Setup – Develop Personas}

    {\small \textcolor{gray}{2--3 hrs $\cdot$ Intermediate}}\\[12pt]

      
Based on your user research, you need to create detailed personas that represent the core segments of your audience. These personas will guide your design decisions throughout the project.


\vspace{8pt}
{\large \textbf{\textcolor{brandLight}{What you'll learn}}}\\[4pt]  
\begin{itemize}
    \item Use Miro or a similar tool to create 3 personas. Include key details: demographics, goals, frustrations, and behaviors.  
    \item Highlight how each persona interacts with AI products and their expectations.
\end{itemize}

\vspace{8pt}
{\large \textbf{\textcolor{brandLight}{What you'll do}}}\\[4pt]  
Submit a PDF or image export of your personas from Miro, ensuring each persona is well-defined with visuals.\par

    \vspace{16pt}

    \textcolor{cardBorder}{\hrule height 1pt}\vspace{8pt}

    \small \textcolor{gray}{When you are done with this task, mark it as complete to proceed to the next one.} \hfill \textbf{\textcolor{brandDark}{$\checkmark$ Completed}}

\end{taskbox}

\begin{taskbox}{Task 3: Build/Execute – Design the User Flow}

    {\small \textcolor{gray}{2--3 hrs $\cdot$ Intermediate}}\\[12pt]

      
You need to map out the user journey for interacting with the AI mental wellness assistant, from onboarding to using its features (e.g., mood tracking, chatbot conversations).


\vspace{8pt}
{\large \textbf{\textcolor{brandLight}{What you'll learn}}}\\[4pt]  
\begin{itemize}
    \item Use Miro to create a user flow diagram that highlights key paths like onboarding, accessing features, and receiving insights.  
    \item Define decision points and actions at each step (e.g., "User selects mood tracker → AI provides analysis").
\end{itemize}

\vspace{8pt}
{\large \textbf{\textcolor{brandLight}{What you'll do}}}\\[4pt]  
Submit a user flow diagram exported from Miro. Include annotations explaining each step and decision point.\par

    \vspace{16pt}

    \textcolor{cardBorder}{\hrule height 1pt}\vspace{8pt}

    \small \textcolor{gray}{When you are done with this task, mark it as complete to proceed to the next one.} \hfill \textbf{\textcolor{brandDark}{$\checkmark$ Completed}}

\end{taskbox}

\begin{taskbox}{Task 4: GenAI/Explanation – Create Wireframes}

    {\small \textcolor{gray}{2--3 hrs $\cdot$ Intermediate}}\\[12pt]

      
It’s time to convert your user flow into wireframes for core screens of the AI product. Focus on intuitive layouts and ensuring users can easily interact with AI features.


\vspace{8pt}
{\large \textbf{\textcolor{brandLight}{What you'll learn}}}\\[4pt]  
\begin{itemize}
    \item Design wireframes for 3 primary screens: Onboarding, Mood Tracking Dashboard, and Chatbot Interface.  
    \item Ensure your designs follow accessibility guidelines (e.g., WCAG contrast standards).
\end{itemize}

\vspace{8pt}
{\large \textbf{\textcolor{brandLight}{What you'll do}}}\\[4pt]  
Submit high-quality wireframes created in Figma or Adobe XD. Include annotations describing design choices.\par

    \vspace{16pt}

    \textcolor{cardBorder}{\hrule height 1pt}\vspace{8pt}

    \small \textcolor{gray}{When you are done with this task, mark it as complete to proceed to the next one.} \hfill \textbf{\textcolor{brandDark}{$\checkmark$ Completed}}

\end{taskbox}

\begin{taskbox}{Task 5: Audit/Responsible AI – Conduct Usability Testing}

    {\small \textcolor{gray}{2--3 hrs $\cdot$ Intermediate}}\\[12pt]

      
You need to ensure your wireframes deliver a seamless experience. Conduct usability tests to gather feedback and identify areas for improvement.


\vspace{8pt}
{\large \textbf{\textcolor{brandLight}{What you'll learn}}}\\[4pt]  
\begin{itemize}
    \item Share your wireframes with 3 test users and ask them to perform specific tasks (e.g., "Navigate to the Mood Tracker").  
    \item Gather feedback on usability issues and suggestions for improvement.
\end{itemize}

\vspace{8pt}
{\large \textbf{\textcolor{brandLight}{What you'll do}}}\\[4pt]  
Submit a usability test report summarizing participant feedback, including a prioritized list of design improvements.\par

    \vspace{16pt}

    \textcolor{cardBorder}{\hrule height 1pt}\vspace{8pt}

    \small \textcolor{gray}{When you are done with this task, mark it as complete to proceed to the next one.} \hfill \textbf{\textcolor{brandDark}{$\checkmark$ Completed}}

\end{taskbox}

\begin{taskbox}{Task 6: Present Recommendation – Final Design Proposal}

    {\small \textcolor{gray}{2--3 hrs $\cdot$ Intermediate}}\\[12pt]

      
You’re presenting your design recommendations to the product team. Use data from research, personas, and usability tests to justify your decisions.


\vspace{8pt}
{\large \textbf{\textcolor{brandLight}{What you'll learn}}}\\[4pt]  
\begin{itemize}
    \item Create a presentation deck showcasing your design process: research findings, personas, user flow, wireframes, and test results.  
    \item Include actionable recommendations for the product team.
\end{itemize}

\vspace{8pt}
{\large \textbf{\textcolor{brandLight}{What you'll do}}}\\[4pt]  
Submit your presentation as a PDF or PowerPoint file. Ensure it’s visually engaging and includes clear, data-driven insights.\par

    \vspace{16pt}

    \textcolor{cardBorder}{\hrule height 1pt}\vspace{8pt}

    \small \textcolor{gray}{When you are done with this task, mark it as complete to proceed to the next one.} \hfill \textbf{\textcolor{brandDark}{$\checkmark$ Completed}}

\end{taskbox}

\begin{completionbox}
    \textbf{Knowledge Assessment \& Verification:} Complete the online knowledge check, submit your code and presentation files for AI verification, and claim your \textbf{Skillzza Verified Certificate}.
\end{completionbox}

\newpage

% =============================================================================

% SIMULATION 39

% =============================================================================

\section*{AWS AI Architect – Design a GenAI Enterprise Solution}

\addcontentsline{toc}{section}{AWS AI Architect – Design a GenAI Enterprise Solution}

\begin{metabox}
    \begin{tabularx}{\textwidth}{@{}ll X@{}}
        \textbf{Industry:} & Technology & \textbf{Partner Enterprise:} \textbf{Skillzza Enterprise} \\
        \textbf{Duration:} & 12-18 hrs (Self-paced) & \textbf{Mode:} Individual, task-based simulation \\
        \textbf{Primary Skills:} & \multicolumn{2}{@{}l}{Python, AI, Data Analysis, Problem Solving}
    \end{tabularx}
\end{metabox}

\subsection*{Overview \& Problem Statement}

\#\#\# \textbf{Scenario}  
Imagine you are an AWS AI Architect working for a global consulting firm. Your client, a leading ecommerce company, wants to deploy a Generative AI (GenAI) solution to enhance customer experience by providing intelligent product recommendations and real-time customer support. They require a solution that is scalable, secure, cost-effective, and compliant with industry regulations.  

Your mission is to design this enterprise GenAI solution using AWS services, ensuring seamless integration, responsible AI practices, and optimal performance.  

\#\#\# \textbf{Mission}  
You will navigate through real-world challenges such as defining business requirements, designing a cloud-based architecture, integrating GenAI tools, addressing security concerns, and estimating operational costs.  

\#\#\# \textbf{Final Challenge}  
At the end of this internship, you will present a comprehensive AWS-based GenAI enterprise solution to the client’s stakeholders. Your solution must include architecture diagrams, security measures, cost analysis, and a justification of your design decisions.  

---

\subsection*{Tasks \& Deliverables}

\begin{taskbox}{Task 1: Understand Business Requirements}

    {\small \textcolor{gray}{2--3 hrs $\cdot$ Intermediate}}\\[12pt]

    \textbf{Scenario}: The client has provided a business summary document outlining their objectives for implementing GenAI. Your first task is to analyze their requirements and translate them into technical specifications.  

\textbf{Student Assignment}:  
\begin{itemize}
    \item Review the business summary.  
    \item Identify key requirements like scalability, compliance, and AI capabilities.  
    \item Map out the functional and non-functional requirements.  
\end{itemize}

\textbf{Deliverable}: A one-page technical requirement document summarizing the client’s goals and expectations.\par

    \vspace{16pt}

    \textcolor{cardBorder}{\hrule height 1pt}\vspace{8pt}

    \small \textcolor{gray}{When you are done with this task, mark it as complete to proceed to the next one.} \hfill \textbf{\textcolor{brandDark}{$\checkmark$ Completed}}

\end{taskbox}

\begin{taskbox}{Task 2: Data Pipeline and Environment Setup}

    {\small \textcolor{gray}{2--3 hrs $\cdot$ Intermediate}}\\[12pt]

    \textbf{Scenario}: Now that the requirements are clear, you need to prepare the AWS environment and set up the data pipeline for the solution.  

\textbf{Student Assignment}:  
\begin{itemize}
    \item Configure the necessary AWS services: Amazon S3, DynamoDB, and IAM roles.  
    \item Design the data pipeline to feed inputs into the GenAI model (e.g., real-time customer queries).  
    \item Ensure scalability by implementing auto-scaling policies.  
\end{itemize}

\textbf{Deliverable}: A diagram of the data pipeline architecture and screenshots of the AWS environment setup.\par

    \vspace{16pt}

    \textcolor{cardBorder}{\hrule height 1pt}\vspace{8pt}

    \small \textcolor{gray}{When you are done with this task, mark it as complete to proceed to the next one.} \hfill \textbf{\textcolor{brandDark}{$\checkmark$ Completed}}

\end{taskbox}

\begin{taskbox}{Task 3: Build and Deploy GenAI Model}

    {\small \textcolor{gray}{2--3 hrs $\cdot$ Intermediate}}\\[12pt]

    \textbf{Scenario}: You need to build a GenAI model using AWS services to process natural language queries and return intelligent responses.  

\textbf{Student Assignment}:  
\begin{itemize}
    \item Use Amazon Bedrock to deploy foundation models for GenAI.  
    \item Integrate SageMaker for custom model training if needed.  
    \item Implement Retrieval-Augmented Generation (RAG) to enhance the model’s accuracy.  
\end{itemize}

\textbf{Deliverable}: A deployed GenAI solution with sample inputs and outputs showcasing its functionality.\par

    \vspace{16pt}

    \textcolor{cardBorder}{\hrule height 1pt}\vspace{8pt}

    \small \textcolor{gray}{When you are done with this task, mark it as complete to proceed to the next one.} \hfill \textbf{\textcolor{brandDark}{$\checkmark$ Completed}}

\end{taskbox}

\begin{taskbox}{Task 4: Implement Responsible AI and Security Measures}

    {\small \textcolor{gray}{2--3 hrs $\cdot$ Intermediate}}\\[12pt]

    \textbf{Scenario}: To ensure compliance and ethical AI practices, you must integrate responsible AI principles and security protocols into the design.  

\textbf{Student Assignment}:  
\begin{itemize}
    \item Configure AWS IAM policies for role-based access control (RBAC).  
    \item Implement encryption using AWS KMS for data security.  
    \item Add filters to prevent bias or inappropriate GenAI outputs.  
\end{itemize}

\textbf{Deliverable}: Documentation describing the security and responsible AI measures implemented, along with AWS configuration screenshots.\par

    \vspace{16pt}

    \textcolor{cardBorder}{\hrule height 1pt}\vspace{8pt}

    \small \textcolor{gray}{When you are done with this task, mark it as complete to proceed to the next one.} \hfill \textbf{\textcolor{brandDark}{$\checkmark$ Completed}}

\end{taskbox}

\begin{taskbox}{Task 5: Audit and Optimize Solution}

    {\small \textcolor{gray}{2--3 hrs $\cdot$ Intermediate}}\\[12pt]

    \textbf{Scenario}: The client needs assurance that the solution is cost-effective and optimized for performance.  

\textbf{Student Assignment}:  
\begin{itemize}
    \item Use Amazon CloudWatch to monitor resource usage and identify bottlenecks.  
    \item Analyze operational costs using AWS Cost Explorer and suggest optimizations.  
    \item Identify risks and propose mitigation strategies.  
\end{itemize}

\textbf{Deliverable}: A performance audit report and cost optimization plan.\par

    \vspace{16pt}

    \textcolor{cardBorder}{\hrule height 1pt}\vspace{8pt}

    \small \textcolor{gray}{When you are done with this task, mark it as complete to proceed to the next one.} \hfill \textbf{\textcolor{brandDark}{$\checkmark$ Completed}}

\end{taskbox}

\begin{taskbox}{Task 6: Present Enterprise Solution}

    {\small \textcolor{gray}{2--3 hrs $\cdot$ Intermediate}}\\[12pt]

    \textbf{Scenario}: Your final task is to present the GenAI enterprise solution to the client’s stakeholders, explaining design decisions, cost estimates, and security measures.  

\textbf{Student Assignment}:  
\begin{itemize}
    \item Create a PowerPoint presentation summarizing the solution architecture, costs, and responsible AI practices.  
    \item Justify your design choices using data and client requirements.  
\end{itemize}

\textbf{Deliverable}: A professional presentation deck (min. 10 slides) and a recorded video walkthrough (optional).\par

    \vspace{16pt}

    \textcolor{cardBorder}{\hrule height 1pt}\vspace{8pt}

    \small \textcolor{gray}{When you are done with this task, mark it as complete to proceed to the next one.} \hfill \textbf{\textcolor{brandDark}{$\checkmark$ Completed}}

\end{taskbox}

\begin{completionbox}
    \textbf{Knowledge Assessment \& Verification:} Complete the online knowledge check, submit your code and presentation files for AI verification, and claim your \textbf{Skillzza Verified Certificate}.
\end{completionbox}

\newpage

% =============================================================================

% SIMULATION 40

% =============================================================================

\section*{Data Lake Architect – Build an Enterprise Data Foundation}

\addcontentsline{toc}{section}{Data Lake Architect – Build an Enterprise Data Foundation}

\begin{metabox}
    \begin{tabularx}{\textwidth}{@{}ll X@{}}
        \textbf{Industry:} & Technology & \textbf{Partner Enterprise:} \textbf{Skillzza Enterprise} \\
        \textbf{Duration:} & 12-18 hrs (Self-paced) & \textbf{Mode:} Individual, task-based simulation \\
        \textbf{Primary Skills:} & \multicolumn{2}{@{}l}{Python, AI, Data Analysis, Problem Solving}
    \end{tabularx}
\end{metabox}

\subsection*{Overview \& Problem Statement}

\#\#\# \textbf{Scenario:}
In the modern enterprise landscape, data is the lifeblood of decision-making. As a Data Lake Architect, you are tasked with designing a scalable and efficient data lake solution to serve as the backbone for enterprise-wide analytics. Your mission? To integrate disparate data sources, establish robust ingestion pipelines, and ensure governance while optimizing the architecture for analytics workloads.

\#\#\# \textbf{Mission:}
You are hired as a Data Lake Architect at a multinational retail company, "RetailSphere Inc." The company is struggling to unify its data environment. They have transactional data from point-of-sale systems, customer profiles from CRM tools, and inventory data from supply chain systems—all siloed in disparate formats. Your goal is to design an enterprise data lake architecture that consolidates these data streams, supports analytics, and ensures regulatory compliance.

\#\#\# \textbf{Final Challenge:}
Present a fully operational enterprise data lake blueprint, complete with ingestion pipelines, governance strategy, and analytics integration. Your solution will be judged on scalability, performance, compliance, and clarity of recommendations.

---

\subsection*{Tasks \& Deliverables}

\begin{taskbox}{Task 1: Understand the Data Landscape}

    {\small \textcolor{gray}{2--3 hrs $\cdot$ Intermediate}}\\[12pt]

    \#\#\#\#   
RetailSphere Inc. collects data from multiple sources:  
1. \textbf{Sales Transactions:} Point-of-sale data in CSV format.  
2. \textbf{Customer Profiles:} CRM data in JSON format.  
3. \textbf{Inventory:} Supply chain data in MySQL databases.  

You need to analyze the data sources and identify challenges in integration.

\#\#\#\# \\[10pt]
\vspace{8pt}
{\large \textbf{\textcolor{brandLight}{What you'll learn}}}\\[4pt]  
\begin{itemize}
    \item Review provided datasets (sample CSV, JSON, and SQL dump).  
    \item Identify schema mismatches, duplication, and data quality issues.  
    \item Map out the relationships between datasets.  
\end{itemize}

\#\#\#\# \vspace{8pt}
{\large \textbf{\textcolor{brandLight}{What you'll do}}}\\[4pt]  
A report summarizing:  
\begin{itemize}
    \item Data source types and formats.  
    \item Challenges in unifying the data.  
    \item A preliminary schema for integration.\par
\end{itemize}

    \vspace{16pt}

    \textcolor{cardBorder}{\hrule height 1pt}\vspace{8pt}

    \small \textcolor{gray}{When you are done with this task, mark it as complete to proceed to the next one.} \hfill \textbf{\textcolor{brandDark}{$\checkmark$ Completed}}

\end{taskbox}

\begin{taskbox}{Task 2: Setup the Data Lake Environment}

    {\small \textcolor{gray}{2--3 hrs $\cdot$ Intermediate}}\\[12pt]

    \#\#\#\#   
RetailSphere Inc. wants to use AWS as the cloud platform for their data lake. You need to configure the foundational infrastructure.

\#\#\#\# \\[10pt]
\vspace{8pt}
{\large \textbf{\textcolor{brandLight}{What you'll learn}}}\\[4pt]  
\begin{itemize}
    \item Provision an S3 bucket for the data lake.  
    \item Set up AWS Glue for ETL pipelines.  
    \item Create IAM roles and policies for secure access.  
\end{itemize}

\#\#\#\# \vspace{8pt}
{\large \textbf{\textcolor{brandLight}{What you'll do}}}\\[4pt]  
\begin{itemize}
    \item Screenshots of the S3 bucket and Glue configurations.  
    \item A document outlining IAM policy settings.\par
\end{itemize}

    \vspace{16pt}

    \textcolor{cardBorder}{\hrule height 1pt}\vspace{8pt}

    \small \textcolor{gray}{When you are done with this task, mark it as complete to proceed to the next one.} \hfill \textbf{\textcolor{brandDark}{$\checkmark$ Completed}}

\end{taskbox}

\begin{taskbox}{Task 3: Build and Execute Data Ingestion Pipelines}

    {\small \textcolor{gray}{2--3 hrs $\cdot$ Intermediate}}\\[12pt]

    \#\#\#\#   
Your task is to ingest data from all sources into the data lake while ensuring scalability and fault tolerance.

\#\#\#\# \\[10pt]
\vspace{8pt}
{\large \textbf{\textcolor{brandLight}{What you'll learn}}}\\[4pt]  
\begin{itemize}
    \item Use AWS Glue to create ingestion pipelines for CSV, JSON, and SQL data.  
    \item Perform schema unification during ingestion.  
    \item Implement error handling for incomplete or corrupt data.  
\end{itemize}

\#\#\#\# \vspace{8pt}
{\large \textbf{\textcolor{brandLight}{What you'll do}}}\\[4pt]  
\begin{itemize}
    \item Python/AWS Glue scripts for ingestion.  
    \item Logs showing successful ingestion.  
    \item Unified schema stored in the data lake.\par
\end{itemize}

    \vspace{16pt}

    \textcolor{cardBorder}{\hrule height 1pt}\vspace{8pt}

    \small \textcolor{gray}{When you are done with this task, mark it as complete to proceed to the next one.} \hfill \textbf{\textcolor{brandDark}{$\checkmark$ Completed}}

\end{taskbox}

\begin{taskbox}{Task 4: AI-Powered Metadata Tagging}

    {\small \textcolor{gray}{2--3 hrs $\cdot$ Intermediate}}\\[12pt]

    \#\#\#\#   
RetailSphere Inc. needs a searchable data catalog to support analytics. Use AI tools to generate metadata tags and automate schema evolution.

\#\#\#\# \\[10pt]
\vspace{8pt}
{\large \textbf{\textcolor{brandLight}{What you'll learn}}}\\[4pt]  
\begin{itemize}
    \item Use AWS Glue Data Catalog or Apache Atlas to tag datasets with metadata.  
    \item Implement AI-driven schema evolution to handle new data.  
    \item Provide a queryable view of metadata for analysts.  
\end{itemize}

\#\#\#\# \vspace{8pt}
{\large \textbf{\textcolor{brandLight}{What you'll do}}}\\[4pt]  
\begin{itemize}
    \item Metadata tagging scripts.  
    \item Screenshots of the catalog interface.  
    \item A report on schema evolution automation.\par
\end{itemize}

    \vspace{16pt}

    \textcolor{cardBorder}{\hrule height 1pt}\vspace{8pt}

    \small \textcolor{gray}{When you are done with this task, mark it as complete to proceed to the next one.} \hfill \textbf{\textcolor{brandDark}{$\checkmark$ Completed}}

\end{taskbox}

\begin{taskbox}{Task 5: Audit and Responsible Data Governance}

    {\small \textcolor{gray}{2--3 hrs $\cdot$ Intermediate}}\\[12pt]

    \#\#\#\#   
RetailSphere Inc. needs to comply with GDPR and CCPA regulations. Implement governance policies for data handling.

\#\#\#\# \\[10pt]
\vspace{8pt}
{\large \textbf{\textcolor{brandLight}{What you'll learn}}}\\[4pt]  
\begin{itemize}
    \item Define governance policies for data retention and access control.  
    \item Use Immuta or AWS Lake Formation to enforce compliance.  
    \item Audit data access logs to ensure compliance.  
\end{itemize}

\#\#\#\# \vspace{8pt}
{\large \textbf{\textcolor{brandLight}{What you'll do}}}\\[4pt]  
\begin{itemize}
    \item Governance policy document.  
    \item Screenshots of compliance tool configurations.  
    \item Audit logs showing adherence to policies.\par
\end{itemize}

    \vspace{16pt}

    \textcolor{cardBorder}{\hrule height 1pt}\vspace{8pt}

    \small \textcolor{gray}{When you are done with this task, mark it as complete to proceed to the next one.} \hfill \textbf{\textcolor{brandDark}{$\checkmark$ Completed}}

\end{taskbox}

\begin{taskbox}{Task 6: Present Recommendations to Stakeholders}

    {\small \textcolor{gray}{2--3 hrs $\cdot$ Intermediate}}\\[12pt]

    \#\#\#\#   
RetailSphere Inc. executives need a presentation summarizing your data lake architecture and its business impact.

\#\#\#\# \\[10pt]
\vspace{8pt}
{\large \textbf{\textcolor{brandLight}{What you'll learn}}}\\[4pt]  
\begin{itemize}
    \item Create a presentation highlighting architecture design, governance strategy, and analytics potential.  
    \item Include ROI calculations (e.g., cost savings from improved data accessibility).  
    \item Prepare to answer questions on scalability and compliance.  
\end{itemize}

\#\#\#\# \vspace{8pt}
{\large \textbf{\textcolor{brandLight}{What you'll do}}}\\[4pt]  
\begin{itemize}
    \item PowerPoint presentation (10 slides max).  
    \item ROI calculation spreadsheet.  
    \item Recorded video or written script for key points.\par
\end{itemize}

    \vspace{16pt}

    \textcolor{cardBorder}{\hrule height 1pt}\vspace{8pt}

    \small \textcolor{gray}{When you are done with this task, mark it as complete to proceed to the next one.} \hfill \textbf{\textcolor{brandDark}{$\checkmark$ Completed}}

\end{taskbox}

\begin{completionbox}
    \textbf{Knowledge Assessment \& Verification:} Complete the online knowledge check, submit your code and presentation files for AI verification, and claim your \textbf{Skillzza Verified Certificate}.
\end{completionbox}

\newpage

% =============================================================================

% SIMULATION 41

% =============================================================================

\section*{AI Data Engineer – Build a RAG Data Pipeline}

\addcontentsline{toc}{section}{AI Data Engineer – Build a RAG Data Pipeline}

\begin{metabox}
    \begin{tabularx}{\textwidth}{@{}ll X@{}}
        \textbf{Industry:} & Technology & \textbf{Partner Enterprise:} \textbf{Skillzza Enterprise} \\
        \textbf{Duration:} & 12-18 hrs (Self-paced) & \textbf{Mode:} Individual, task-based simulation \\
        \textbf{Primary Skills:} & \multicolumn{2}{@{}l}{Python, AI, Data Analysis, Problem Solving}
    \end{tabularx}
\end{metabox}

\subsection*{Overview \& Problem Statement}

\#\#\# \textbf{Scenario}  
Imagine you're an AI Data Engineer at a fast-growing tech startup that specializes in building intelligent search and retrieval systems. Your team has been tasked with creating a \textbf{Retrieval-Augmented Generation (RAG) pipeline} for a business client who needs a robust solution for document search and question answering. The client has a large repository of internal documents (e.g., policy manuals, reports, and technical documents) and wants to leverage generative AI to provide accurate and contextual answers to user queries.

\#\#\# \textbf{Mission}  
Your mission is to design and implement a fully functional \textbf{RAG Data Pipeline}, which involves preprocessing raw documents, generating embeddings, storing them in a \textbf{Vector Database}, retrieving relevant documents for a user query, and integrating this with a generative AI model for answer generation.

\#\#\# \textbf{Final Challenge}  
At the end of the internship, you will present a fully operational RAG pipeline to simulate a real-world project delivery. You'll showcase your pipeline's ability to answer user queries accurately while ensuring scalability, efficiency, and adherence to ethical AI principles.

---

\subsection*{Tasks \& Deliverables}

\begin{taskbox}{Task 1: Understand the Problem and Define Pipeline Architecture}

    {\small \textcolor{gray}{2--3 hrs $\cdot$ Intermediate}}\\[12pt]

     Your client has provided a dataset of internal documents (PDFs and TXT files) containing critical company information. You need to design a RAG pipeline that processes these documents and integrates with a generative AI model to provide query-based answers.  


\vspace{8pt}
{\large \textbf{\textcolor{brandLight}{What you'll learn}}}\\[4pt]  
\begin{itemize}
    \item Research the RAG architecture.  
    \item Define the components of your pipeline (e.g., preprocessing, embedding generation, vector DB, retrieval, GenAI integration).  
    \item Document the flow of data through your pipeline.  
\end{itemize}

\vspace{8pt}
{\large \textbf{\textcolor{brandLight}{What you'll do}}}\\[4pt]  
\begin{itemize}
    \item Submit a flowchart or diagram of your pipeline architecture with a brief explanation of each component.\par
\end{itemize}

    \vspace{16pt}

    \textcolor{cardBorder}{\hrule height 1pt}\vspace{8pt}

    \small \textcolor{gray}{When you are done with this task, mark it as complete to proceed to the next one.} \hfill \textbf{\textcolor{brandDark}{$\checkmark$ Completed}}

\end{taskbox}

\begin{taskbox}{Task 2: Data Preprocessing}

    {\small \textcolor{gray}{2--3 hrs $\cdot$ Intermediate}}\\[12pt]

     The provided dataset contains raw documents with inconsistent formatting, typos, and irrelevant information. Before you can generate embeddings, these documents need to be cleaned and structured.  


\vspace{8pt}
{\large \textbf{\textcolor{brandLight}{What you'll learn}}}\\[4pt]  
\begin{itemize}
    \item Write a Python script to preprocess the documents:  
    \item Remove special characters and stopwords.  
    \item Split long documents into smaller, meaningful chunks (e.g., 200-300 words).  
    \item Save the processed chunks into a structured format (e.g., JSON or CSV).  
\end{itemize}

\vspace{8pt}
{\large \textbf{\textcolor{brandLight}{What you'll do}}}\\[4pt]  
\begin{itemize}
    \item Upload the cleaned and structured dataset in JSON or CSV format.  
    \item Submit your Python code with comments explaining each step.\par
\end{itemize}

    \vspace{16pt}

    \textcolor{cardBorder}{\hrule height 1pt}\vspace{8pt}

    \small \textcolor{gray}{When you are done with this task, mark it as complete to proceed to the next one.} \hfill \textbf{\textcolor{brandDark}{$\checkmark$ Completed}}

\end{taskbox}

\begin{taskbox}{Task 3: Generate Embeddings and Store in a Vector Database}

    {\small \textcolor{gray}{2--3 hrs $\cdot$ Intermediate}}\\[12pt]

     Now that the documents are preprocessed, the next step is to generate vector embeddings and store them in a Vector Database. These embeddings will allow you to efficiently retrieve relevant chunks during query time.  


\vspace{8pt}
{\large \textbf{\textcolor{brandLight}{What you'll learn}}}\\[4pt]  
\begin{itemize}
    \item Use a pre-trained embedding model (e.g., Sentence Transformers or OpenAI Embedding API) to generate embeddings for your document chunks.  
    \item Set up a Vector Database (e.g., Pinecone, Weaviate, or FAISS) and store the embeddings for efficient search and retrieval.  
\end{itemize}

\vspace{8pt}
{\large \textbf{\textcolor{brandLight}{What you'll do}}}\\[4pt]  
\begin{itemize}
    \item Upload the Python script used to generate embeddings and interact with the Vector Database.  
    \item Include a sample output showcasing stored embeddings in the Vector Database.\par
\end{itemize}

    \vspace{16pt}

    \textcolor{cardBorder}{\hrule height 1pt}\vspace{8pt}

    \small \textcolor{gray}{When you are done with this task, mark it as complete to proceed to the next one.} \hfill \textbf{\textcolor{brandDark}{$\checkmark$ Completed}}

\end{taskbox}

\begin{taskbox}{Task 4: Integrate Generative AI for Query-Based Retrieval}

    {\small \textcolor{gray}{2--3 hrs $\cdot$ Intermediate}}\\[12pt]

     After storing embeddings in the Vector Database, the client now wants to test the pipeline's ability to retrieve relevant document chunks and generate contextual answers to user queries.  


\vspace{8pt}
{\large \textbf{\textcolor{brandLight}{What you'll learn}}}\\[4pt]  
\begin{itemize}
    \item Write a Python script to:  
    \item Accept a user query as input.  
    \item Use the Vector Database to retrieve the top 3 relevant document chunks for the query.  
    \item Pass these chunks to a Generative AI model (e.g., OpenAI GPT-4) to generate a coherent answer.  
\end{itemize}

\vspace{8pt}
{\large \textbf{\textcolor{brandLight}{What you'll do}}}\\[4pt]  
\begin{itemize}
    \item Submit your Python script and a sample output showcasing a query and the AI-generated answer.  
    \item Provide a brief explanation of how retrieval and generation were integrated.\par
\end{itemize}

    \vspace{16pt}

    \textcolor{cardBorder}{\hrule height 1pt}\vspace{8pt}

    \small \textcolor{gray}{When you are done with this task, mark it as complete to proceed to the next one.} \hfill \textbf{\textcolor{brandDark}{$\checkmark$ Completed}}

\end{taskbox}

\begin{taskbox}{Task 5: Audit for Accuracy and Responsible AI Practices}

    {\small \textcolor{gray}{2--3 hrs $\cdot$ Intermediate}}\\[12pt]

     The client has concerns about the accuracy and fairness of the AI-generated answers. You need to audit the RAG pipeline to ensure it adheres to Responsible AI principles and minimizes bias or misinformation.  


\vspace{8pt}
{\large \textbf{\textcolor{brandLight}{What you'll learn}}}\\[4pt]  
\begin{itemize}
    \item Identify potential risks in the pipeline (e.g., misinformation, biased embeddings, hallucinations).  
    \item Perform an accuracy check for a set of test queries by comparing AI-generated answers to ground truth data.  
    \item Suggest strategies to mitigate risks (e.g., adding a fact-checking layer or confidence thresholds).  
\end{itemize}

\vspace{8pt}
{\large \textbf{\textcolor{brandLight}{What you'll do}}}\\[4pt]  
\begin{itemize}
    \item Submit a report highlighting identified risks, audit results, and proposed mitigation strategies.\par
\end{itemize}

    \vspace{16pt}

    \textcolor{cardBorder}{\hrule height 1pt}\vspace{8pt}

    \small \textcolor{gray}{When you are done with this task, mark it as complete to proceed to the next one.} \hfill \textbf{\textcolor{brandDark}{$\checkmark$ Completed}}

\end{taskbox}

\begin{taskbox}{Task 6: Present Your RAG Pipeline}

    {\small \textcolor{gray}{2--3 hrs $\cdot$ Intermediate}}\\[12pt]

     It’s time to showcase your RAG pipeline to the client. Demonstrate its functionality, explain the architecture, and highlight how it meets the client’s requirements.  


\vspace{8pt}
{\large \textbf{\textcolor{brandLight}{What you'll learn}}}\\[4pt]  
\begin{itemize}
    \item Prepare a presentation that includes:  
    \item An overview of the pipeline architecture.  
    \item A walkthrough of the preprocessing, embedding generation, and retrieval processes.  
    \item Live demo of the pipeline answering a sample query.  
    \item Responsible AI measures implemented.  
\end{itemize}

\vspace{8pt}
{\large \textbf{\textcolor{brandLight}{What you'll do}}}\\[4pt]  
\begin{itemize}
    \item Submit a recording of your presentation (5-7 minutes).  
    \item Include your presentation slides and any supporting documentation.\par
\end{itemize}

    \vspace{16pt}

    \textcolor{cardBorder}{\hrule height 1pt}\vspace{8pt}

    \small \textcolor{gray}{When you are done with this task, mark it as complete to proceed to the next one.} \hfill \textbf{\textcolor{brandDark}{$\checkmark$ Completed}}

\end{taskbox}

\begin{completionbox}
    \textbf{Knowledge Assessment \& Verification:} Complete the online knowledge check, submit your code and presentation files for AI verification, and claim your \textbf{Skillzza Verified Certificate}.
\end{completionbox}

\newpage

% =============================================================================

% SIMULATION 42

% =============================================================================

\section*{AI IT Service Manager – Design an Agentic IT Helpdesk}

\addcontentsline{toc}{section}{AI IT Service Manager – Design an Agentic IT Helpdesk}

\begin{metabox}
    \begin{tabularx}{\textwidth}{@{}ll X@{}}
        \textbf{Industry:} & Technology & \textbf{Partner Enterprise:} \textbf{Skillzza Enterprise} \\
        \textbf{Duration:} & 12-18 hrs (Self-paced) & \textbf{Mode:} Individual, task-based simulation \\
        \textbf{Primary Skills:} & \multicolumn{2}{@{}l}{Python, AI, Data Analysis, Problem Solving}
    \end{tabularx}
\end{metabox}

\subsection*{Overview \& Problem Statement}

\#\#\# \textbf{Scenario:}  
You are an AI IT Service Manager tasked with transforming a traditional IT helpdesk into a cutting-edge, agentic AI-powered helpdesk. The goal is to enhance ticket resolution efficiency, improve the user experience, and implement governance policies to ensure responsible AI practices. Your stakeholders include IT administrators, end-users, and company leadership, all relying on your expertise to deliver a seamless automated helpdesk system.  

\#\#\# \textbf{Mission:}  
Your mission is to design an AI-driven IT helpdesk system that integrates ticket workflows, intelligent agents, a knowledge base, and automated processes. Additionally, you will implement governance measures to ensure ethical and responsible AI usage. This immersive simulation will test your ability to optimize IT service management using AI tools while balancing automation with human oversight.  

\#\#\# \textbf{Final Challenge:}  
At the end of this internship, you will present a detailed IT helpdesk solution that includes a defined ticket workflow, agentic AI configurations, an enriched knowledge base, automation scripts, governance policies, and data-driven recommendations.  

---

\subsection*{Tasks \& Deliverables}

\begin{taskbox}{Task 1: Understand the IT Helpdesk System}

    {\small \textcolor{gray}{2--3 hrs $\cdot$ Intermediate}}\\[12pt]

    \#\#\#\#   
Your company processes hundreds of IT support tickets daily. Currently, resolution times are slow due to repetitive queries and inefficient workflows. Leadership has asked you to analyze the existing system and identify areas for AI intervention.  

\#\#\#\# \\[10pt]
\vspace{8pt}
{\large \textbf{\textcolor{brandLight}{What you'll learn}}}\\[4pt]  
\begin{itemize}
    \item Study the provided dataset of historical IT tickets (fields: ticket ID, category, issue description, resolution time, agent notes).  
    \item Categorize tickets into “simple,” “complex,” and “critical” based on resolution time and issue type.  
    \item Identify areas where agentic AI could reduce resolution time.  
\end{itemize}

\#\#\#\# \vspace{8pt}
{\large \textbf{\textcolor{brandLight}{What you'll do}}}\\[4pt]  
Submit a report categorizing tickets and highlighting inefficiencies in the current IT helpdesk system. Include 3 suggested areas for AI intervention.\par

    \vspace{16pt}

    \textcolor{cardBorder}{\hrule height 1pt}\vspace{8pt}

    \small \textcolor{gray}{When you are done with this task, mark it as complete to proceed to the next one.} \hfill \textbf{\textcolor{brandDark}{$\checkmark$ Completed}}

\end{taskbox}

\begin{taskbox}{Task 2: Build the Ticket Workflow \& Knowledge Base}

    {\small \textcolor{gray}{2--3 hrs $\cdot$ Intermediate}}\\[12pt]

    \#\#\#\#   
You need to design a workflow for handling IT tickets more efficiently. Additionally, a centralized knowledge base must be created to assist AI agents in resolving simple queries autonomously.  

\#\#\#\# \\[10pt]
\vspace{8pt}
{\large \textbf{\textcolor{brandLight}{What you'll learn}}}\\[4pt]  
\begin{itemize}
    \item Create a ticket workflow diagram with stages (e.g., ticket creation, classification, agent assignment, resolution).  
    \item Build a knowledge base using sample queries and solutions provided (fields: query, solution, tags).  
    \item Integrate the knowledge base with Zendesk.  
\end{itemize}

\#\#\#\# \vspace{8pt}
{\large \textbf{\textcolor{brandLight}{What you'll do}}}\\[4pt]  
Submit the ticket workflow diagram and a functional knowledge base in Zendesk. Include screenshots of the integration process.\par

    \vspace{16pt}

    \textcolor{cardBorder}{\hrule height 1pt}\vspace{8pt}

    \small \textcolor{gray}{When you are done with this task, mark it as complete to proceed to the next one.} \hfill \textbf{\textcolor{brandDark}{$\checkmark$ Completed}}

\end{taskbox}

\begin{taskbox}{Task 3: Configure Agentic AI for IT Queries}

    {\small \textcolor{gray}{2--3 hrs $\cdot$ Intermediate}}\\[12pt]

    \#\#\#\#   
To automate IT support, you will configure conversational AI agents capable of handling simple and repetitive queries. These AI agents must escalate complex or critical tickets to human agents.  

\#\#\#\# \\[10pt]
\vspace{8pt}
{\large \textbf{\textcolor{brandLight}{What you'll learn}}}\\[4pt]  
\begin{itemize}
    \item Use Dialogflow CX to create conversational flows for sample IT queries (e.g., password reset, software installation).  
    \item Configure escalation triggers for complex issues.  
    \item Test the AI agents with sample queries.  
\end{itemize}

\#\#\#\# \vspace{8pt}
{\large \textbf{\textcolor{brandLight}{What you'll do}}}\\[4pt]  
Submit a video recording of the conversational flows and escalation triggers in action. Provide the Dialogflow CX configuration files.\par

    \vspace{16pt}

    \textcolor{cardBorder}{\hrule height 1pt}\vspace{8pt}

    \small \textcolor{gray}{When you are done with this task, mark it as complete to proceed to the next one.} \hfill \textbf{\textcolor{brandDark}{$\checkmark$ Completed}}

\end{taskbox}

\begin{taskbox}{Task 4: Automate IT Processes Using AI}

    {\small \textcolor{gray}{2--3 hrs $\cdot$ Intermediate}}\\[12pt]

    \#\#\#\#   
Certain IT processes, such as ticket triage and resolution updates, are repetitive and time-consuming. Automating these tasks will improve efficiency and reduce human error.  

\#\#\#\# \\[10pt]
\vspace{8pt}
{\large \textbf{\textcolor{brandLight}{What you'll learn}}}\\[4pt]  
\begin{itemize}
    \item Use Power Automate to script workflows for ticket triage and resolution updates.  
    \item Automate notifications for ticket status changes.  
    \item Test the automation with sample tickets.  
\end{itemize}

\#\#\#\# \vspace{8pt}
{\large \textbf{\textcolor{brandLight}{What you'll do}}}\\[4pt]  
Submit the Power Automate scripts and a demonstration video showing ticket triage and resolution updates.\par

    \vspace{16pt}

    \textcolor{cardBorder}{\hrule height 1pt}\vspace{8pt}

    \small \textcolor{gray}{When you are done with this task, mark it as complete to proceed to the next one.} \hfill \textbf{\textcolor{brandDark}{$\checkmark$ Completed}}

\end{taskbox}

\begin{taskbox}{Task 5: Audit the AI System for Responsible Usage}

    {\small \textcolor{gray}{2--3 hrs $\cdot$ Intermediate}}\\[12pt]

    \#\#\#\#   
Leadership has emphasized the importance of responsible AI. You are tasked with auditing the AI-powered IT helpdesk system to ensure it adheres to governance policies.  

\#\#\#\# \\[10pt]
\vspace{8pt}
{\large \textbf{\textcolor{brandLight}{What you'll learn}}}\\[4pt]  
\begin{itemize}
    \item Use the AI Governance Toolkit to identify biases in ticket classification and agent responses.  
    \item Develop a policy document outlining ethical AI usage in IT service management.  
    \item Recommend changes to improve fairness and transparency.  
\end{itemize}

\#\#\#\# \vspace{8pt}
{\large \textbf{\textcolor{brandLight}{What you'll do}}}\\[4pt]  
Submit the audit report and the policy document. Include screenshots or logs from the AI Governance Toolkit.\par

    \vspace{16pt}

    \textcolor{cardBorder}{\hrule height 1pt}\vspace{8pt}

    \small \textcolor{gray}{When you are done with this task, mark it as complete to proceed to the next one.} \hfill \textbf{\textcolor{brandDark}{$\checkmark$ Completed}}

\end{taskbox}

\begin{taskbox}{Task 6: Present Your AI IT Helpdesk Solution}

    {\small \textcolor{gray}{2--3 hrs $\cdot$ Intermediate}}\\[12pt]

    \#\#\#\#   
You need to present your IT helpdesk solution to company leadership. Your presentation should demonstrate the workflow, agentic AI capabilities, automation scripts, and governance policies.  

\#\#\#\# \\[10pt]
\vspace{8pt}
{\large \textbf{\textcolor{brandLight}{What you'll learn}}}\\[4pt]  
\begin{itemize}
    \item Create a PowerPoint presentation summarizing your solution.  
    \item Include visuals of the ticket workflow, AI agent configurations, and automation scripts.  
    \item Explain how governance policies ensure responsible AI usage.  
\end{itemize}

\#\#\#\# \vspace{8pt}
{\large \textbf{\textcolor{brandLight}{What you'll do}}}\\[4pt]  
Submit the PowerPoint presentation and a 5-minute recorded video of your pitch.\par

    \vspace{16pt}

    \textcolor{cardBorder}{\hrule height 1pt}\vspace{8pt}

    \small \textcolor{gray}{When you are done with this task, mark it as complete to proceed to the next one.} \hfill \textbf{\textcolor{brandDark}{$\checkmark$ Completed}}

\end{taskbox}

\begin{completionbox}
    \textbf{Knowledge Assessment \& Verification:} Complete the online knowledge check, submit your code and presentation files for AI verification, and claim your \textbf{Skillzza Verified Certificate}.
\end{completionbox}

\newpage

% =============================================================================

% SIMULATION 43

% =============================================================================

\section*{AI Project Manager – Rescue a Delayed AI Project}

\addcontentsline{toc}{section}{AI Project Manager – Rescue a Delayed AI Project}

\begin{metabox}
    \begin{tabularx}{\textwidth}{@{}ll X@{}}
        \textbf{Industry:} & Technology & \textbf{Partner Enterprise:} \textbf{Skillzza Enterprise} \\
        \textbf{Duration:} & 12-18 hrs (Self-paced) & \textbf{Mode:} Individual, task-based simulation \\
        \textbf{Primary Skills:} & \multicolumn{2}{@{}l}{Python, AI, Data Analysis, Problem Solving}
    \end{tabularx}
\end{metabox}

\subsection*{Overview \& Problem Statement}

\#\#\# \textbf{Scenario}
You are an AI Project Manager at a mid-sized artificial intelligence development firm. One of your critical projects—a machine learning model designed to predict product demand for a major retail client—has fallen behind schedule by three months due to scope creep, resource shortages, and technical bottlenecks. The client is losing patience and is considering terminating the contract unless immediate action is taken.

Your task? Rescue the project by identifying the root causes of the delays, formulating a recovery plan, and regaining client trust through transparent communication and effective leadership.

---

\#\#\# \textbf{Mission}
Your mission is to:
1. Assess the project data to analyze the reasons for the delay.
2. Develop actionable strategies to overcome bottlenecks and resource constraints.
3. Create a comprehensive recovery plan and communicate it effectively to stakeholders.

This simulation will immerse you in the high-stakes world of project management for AI development, teaching you how to manage complex timelines, technical challenges, and client relations simultaneously.

---

\#\#\# \textbf{Final Challenge}
Deliver a polished recovery plan in the form of a project roadmap and a stakeholder communication memo. Demonstrate your ability to manage risks, resources, and expectations in a high-pressure project rescue scenario.

---

\subsection*{Tasks \& Deliverables}

\begin{taskbox}{Task 1: Understand – Analyze the Current State of the AI Project}

    {\small \textcolor{gray}{2--3 hrs $\cdot$ Intermediate}}\\[12pt]

\begin{itemize}
    \item You’ve been handed the current project status reports, including Gantt charts, team performance metrics, and client feedback logs. Your first task is to assess the health of the project to identify key areas of concern.

\end{itemize}
\vspace{8pt}
{\large \textbf{\textcolor{brandLight}{What you'll learn}}}\\[4pt] 
\begin{itemize}
    \item Analyze the Gantt chart to identify overdue milestones.
    \item Review the client feedback to understand their concerns.
    \item Examine team performance metrics and resource allocation.

\end{itemize}
{\large \textbf{\textcolor{brandLight}{What you'll do}}}\\[4pt] 
\begin{itemize}
    \item A one-page summary identifying at least three root causes of the project delay, supported by data from the reports.\par
\end{itemize}

    \vspace{16pt}

    \textcolor{cardBorder}{\hrule height 1pt}\vspace{8pt}

    \small \textcolor{gray}{When you are done with this task, mark it as complete to proceed to the next one.} \hfill \textbf{\textcolor{brandDark}{$\checkmark$ Completed}}

\end{taskbox}

\begin{taskbox}{Task 2: Data/Setup – Prioritize Recovery Objectives}

    {\small \textcolor{gray}{2--3 hrs $\cdot$ Intermediate}}\\[12pt]

\begin{itemize}
    \item After identifying the root causes, you need to set clear recovery objectives. The client wants the project delivered within 6 weeks, but your team has limited resources.

\end{itemize}
\vspace{8pt}
{\large \textbf{\textcolor{brandLight}{What you'll learn}}}\\[4pt] 
\begin{itemize}
    \item Use provided project data (e.g., resource availability, task dependencies) to prioritize tasks.
    \item Define SMART (Specific, Measurable, Achievable, Relevant, Time-bound) recovery objectives.

\end{itemize}
{\large \textbf{\textcolor{brandLight}{What you'll do}}}\\[4pt] 
\begin{itemize}
    \item A prioritized task list and a set of 3–5 SMART objectives for the recovery plan.\par
\end{itemize}

    \vspace{16pt}

    \textcolor{cardBorder}{\hrule height 1pt}\vspace{8pt}

    \small \textcolor{gray}{When you are done with this task, mark it as complete to proceed to the next one.} \hfill \textbf{\textcolor{brandDark}{$\checkmark$ Completed}}

\end{taskbox}

\begin{taskbox}{Task 3: Build/Execute – Develop a Recovery Plan}

    {\small \textcolor{gray}{2--3 hrs $\cdot$ Intermediate}}\\[12pt]

\begin{itemize}
    \item Now it’s time to create a detailed recovery plan. You’ll need to reallocate resources, adjust timelines, and propose additional measures to meet the client’s deadline.

\end{itemize}
\vspace{8pt}
{\large \textbf{\textcolor{brandLight}{What you'll learn}}}\\[4pt] 
\begin{itemize}
    \item Use JIRA to create a sprint plan for the next 6 weeks.
    \item Develop a revised Gantt chart with updated timelines.
    \item Propose solutions to address the identified bottlenecks (e.g., outsourcing, overtime, automation).

\end{itemize}
{\large \textbf{\textcolor{brandLight}{What you'll do}}}\\[4pt] 
\begin{itemize}
    \item A detailed recovery plan document that includes the revised Gantt chart, sprint plan, and proposed solutions.\par
\end{itemize}

    \vspace{16pt}

    \textcolor{cardBorder}{\hrule height 1pt}\vspace{8pt}

    \small \textcolor{gray}{When you are done with this task, mark it as complete to proceed to the next one.} \hfill \textbf{\textcolor{brandDark}{$\checkmark$ Completed}}

\end{taskbox}

\begin{taskbox}{Task 4: GenAI/Explanation – Use Generative AI to Draft Stakeholder Communication}

    {\small \textcolor{gray}{2--3 hrs $\cdot$ Intermediate}}\\[12pt]

\begin{itemize}
    \item You must communicate the recovery plan to the client and stakeholders. Miscommunication could lead to further misunderstandings or mistrust.

\end{itemize}
\vspace{8pt}
{\large \textbf{\textcolor{brandLight}{What you'll learn}}}\\[4pt] 
\begin{itemize}
    \item Use a Generative AI tool (like ChatGPT) to draft a stakeholder memo explaining the recovery plan.
    \item Ensure the memo is clear, concise, and aligned with the client’s priorities.

\end{itemize}
{\large \textbf{\textcolor{brandLight}{What you'll do}}}\\[4pt] 
\begin{itemize}
    \item A polished, client-ready stakeholder memo generated using GenAI, explaining the recovery steps, timelines, and expected outcomes.\par
\end{itemize}

    \vspace{16pt}

    \textcolor{cardBorder}{\hrule height 1pt}\vspace{8pt}

    \small \textcolor{gray}{When you are done with this task, mark it as complete to proceed to the next one.} \hfill \textbf{\textcolor{brandDark}{$\checkmark$ Completed}}

\end{taskbox}

\begin{taskbox}{Task 5: Audit/Responsible AI – Ensure Ethical and Responsible AI Practices}

    {\small \textcolor{gray}{2--3 hrs $\cdot$ Intermediate}}\\[12pt]

\begin{itemize}
    \item The project involves sensitive customer data, and the client is concerned about data privacy and ethical AI practices.

\end{itemize}
\vspace{8pt}
{\large \textbf{\textcolor{brandLight}{What you'll learn}}}\\[4pt] 
\begin{itemize}
    \item Review the project scope and ensure compliance with data privacy laws (e.g., GDPR, CCPA).
    \item Identify any ethical or bias-related risks in the AI model.
    \item Propose mitigation steps for the identified risks.

\end{itemize}
{\large \textbf{\textcolor{brandLight}{What you'll do}}}\\[4pt] 
\begin{itemize}
    \item A checklist of compliance and ethical considerations, along with a brief report on how these will be addressed in the recovery plan.\par
\end{itemize}

    \vspace{16pt}

    \textcolor{cardBorder}{\hrule height 1pt}\vspace{8pt}

    \small \textcolor{gray}{When you are done with this task, mark it as complete to proceed to the next one.} \hfill \textbf{\textcolor{brandDark}{$\checkmark$ Completed}}

\end{taskbox}

\begin{taskbox}{Task 6: Present Recommendation – Deliver the Recovery Plan}

    {\small \textcolor{gray}{2--3 hrs $\cdot$ Intermediate}}\\[12pt]

\begin{itemize}
    \item The client has scheduled a meeting to review your recovery plan. Your goal is to gain their approval and rebuild trust.

\end{itemize}
\vspace{8pt}
{\large \textbf{\textcolor{brandLight}{What you'll learn}}}\\[4pt] 
\begin{itemize}
    \item Create a PowerPoint presentation summarizing the recovery plan.
    \item Include key insights from the root cause analysis, revised timelines, and risk mitigation strategies.
    \item Practice delivering the presentation using a clear and professional tone.

\end{itemize}
{\large \textbf{\textcolor{brandLight}{What you'll do}}}\\[4pt] 
\begin{itemize}
    \item A 6–8 slide PowerPoint presentation, including a script or speaker notes.\par
\end{itemize}

    \vspace{16pt}

    \textcolor{cardBorder}{\hrule height 1pt}\vspace{8pt}

    \small \textcolor{gray}{When you are done with this task, mark it as complete to proceed to the next one.} \hfill \textbf{\textcolor{brandDark}{$\checkmark$ Completed}}

\end{taskbox}

\begin{completionbox}
    \textbf{Knowledge Assessment \& Verification:} Complete the online knowledge check, submit your code and presentation files for AI verification, and claim your \textbf{Skillzza Verified Certificate}.
\end{completionbox}

\newpage

% =============================================================================

% SIMULATION 44

% =============================================================================

\section*{Sustainability Analyst – Build a Corporate Carbon Dashboard}

\addcontentsline{toc}{section}{Sustainability Analyst – Build a Corporate Carbon Dashboard}

\begin{metabox}
    \begin{tabularx}{\textwidth}{@{}ll X@{}}
        \textbf{Industry:} & Technology & \textbf{Partner Enterprise:} \textbf{Skillzza Enterprise} \\
        \textbf{Duration:} & 12-18 hrs (Self-paced) & \textbf{Mode:} Individual, task-based simulation \\
        \textbf{Primary Skills:} & \multicolumn{2}{@{}l}{Python, AI, Data Analysis, Problem Solving}
    \end{tabularx}
\end{metabox}

\subsection*{Overview \& Problem Statement}

\#\#\# \textbf{Scenario}  
You have been hired as a Sustainability Analyst for a multinational corporation aiming to reduce its carbon footprint and improve environmental impact transparency. The company operates across various industries, generating Scope 1, Scope 2, and Scope 3 emissions. Your mission is to design a corporate carbon dashboard that visualizes emissions data and provides actionable insights for a reduction plan.

\#\#\# \textbf{Mission}  
Your mission is to analyze corporate emissions data across three scopes, identify key contributors, and create a data-driven dashboard that empowers stakeholders to make informed decisions. You will also propose a sustainability roadmap based on your findings to help the company align its operations with global climate goals, such as those outlined in the Paris Agreement.

\#\#\# \textbf{Final Challenge}  
Your final task is to present a polished carbon dashboard to key stakeholders, explaining the data insights, reduction strategies, and how the dashboard can drive measurable impact. You will need to defend your recommendations using analytics and responsible AI principles.

---

\subsection*{Tasks \& Deliverables}

\begin{taskbox}{Task 1: Understand the Carbon Emissions Landscape}

    {\small \textcolor{gray}{2--3 hrs $\cdot$ Intermediate}}\\[12pt]

    \#\#\#\# \textbf{Scenario}  
Your manager hands you a briefing document outlining the company’s operations and emissions footprint. You need to identify the key differentiators between Scope 1, Scope 2, and Scope 3 emissions and their impact on corporate sustainability goals.

\#\#\#\# \textbf{Student Assignment}
\begin{itemize}
    \item Read the provided document on carbon emissions and the Greenhouse Gas Protocol.  
    \item Summarize the differences between Scope 1, Scope 2, and Scope 3 emissions.  
    \item Identify three challenges companies face when tracking Scope 3 emissions.
\end{itemize}

\#\#\#\# \textbf{Deliverable}
Submit a one-page summary defining Scope 1, Scope 2, and Scope 3 emissions, including examples for each scope and key challenges.\par

    \vspace{16pt}

    \textcolor{cardBorder}{\hrule height 1pt}\vspace{8pt}

    \small \textcolor{gray}{When you are done with this task, mark it as complete to proceed to the next one.} \hfill \textbf{\textcolor{brandDark}{$\checkmark$ Completed}}

\end{taskbox}

\begin{taskbox}{Task 2: Prepare and Clean Emissions Data}

    {\small \textcolor{gray}{2--3 hrs $\cdot$ Intermediate}}\\[12pt]

    \#\#\#\# \textbf{Scenario}  
You have received raw emissions data from the company’s various departments, including energy consumption, transportation, and supplier data. The data is inconsistent and requires cleaning before analysis.

\#\#\#\# \textbf{Student Assignment}
\begin{itemize}
    \item Import the raw emissions dataset into Excel or Google Sheets.  
    \item Perform the following data cleaning tasks:  
    \item Remove duplicates.  
    \item Standardize units of measurement (e.g., convert kWh to MWh).  
    \item Handle missing values using interpolation or replacement strategies.  
    \item Categorize the cleaned data into Scope 1, Scope 2, and Scope 3 emissions.
\end{itemize}

\#\#\#\# \textbf{Deliverable}
Submit the cleaned and categorized dataset in .xlsx or Google Sheets format.\par

    \vspace{16pt}

    \textcolor{cardBorder}{\hrule height 1pt}\vspace{8pt}

    \small \textcolor{gray}{When you are done with this task, mark it as complete to proceed to the next one.} \hfill \textbf{\textcolor{brandDark}{$\checkmark$ Completed}}

\end{taskbox}

\begin{taskbox}{Task 3: Build the Carbon Dashboard}

    {\small \textcolor{gray}{2--3 hrs $\cdot$ Intermediate}}\\[12pt]

    \#\#\#\# \textbf{Scenario}  
After cleaning the data, your next task is to build a visual dashboard that highlights the company’s carbon emissions across Scope 1, Scope 2, and Scope 3. The dashboard should include actionable insights and trends.

\#\#\#\# \textbf{Student Assignment}
\begin{itemize}
    \item Use Tableau or Power BI to create the dashboard.  
    \item Include the following visualizations:  
    \item Total emissions by scope (bar chart).  
    \item Monthly emissions trend (line graph).  
    \item Percentage contribution of Scope 3 emissions (pie chart).  
    \item Ensure that the dashboard is clear, interactive, and stakeholder-friendly.
\end{itemize}

\#\#\#\# \textbf{Deliverable}
Submit a link or export of your dashboard in Tableau/Power BI format, with screenshots explaining each visualization.\par

    \vspace{16pt}

    \textcolor{cardBorder}{\hrule height 1pt}\vspace{8pt}

    \small \textcolor{gray}{When you are done with this task, mark it as complete to proceed to the next one.} \hfill \textbf{\textcolor{brandDark}{$\checkmark$ Completed}}

\end{taskbox}

\begin{taskbox}{Task 4: Generate Insights Using GenAI}

    {\small \textcolor{gray}{2--3 hrs $\cdot$ Intermediate}}\\[12pt]

    \#\#\#\# \textbf{Scenario}  
The company wants to leverage AI for deeper insights into emissions data. You are tasked with using a Generative AI tool to generate a narrative analysis of the data and suggest areas for improvement.

\#\#\#\# \textbf{Student Assignment}
\begin{itemize}
    \item Use a Generative AI tool (e.g., ChatGPT or Bard) to analyze the cleaned dataset.  
    \item Provide insights on:  
    \item Which scope contributes the most to emissions and why.  
    \item Predictive trends for emissions in the next 5 years.  
    \item Three reduction strategies based on the data.  
\end{itemize}

\#\#\#\# \textbf{Deliverable}  
Submit the AI-generated narrative analysis and highlight how the GenAI tool helped derive actionable insights.\par

    \vspace{16pt}

    \textcolor{cardBorder}{\hrule height 1pt}\vspace{8pt}

    \small \textcolor{gray}{When you are done with this task, mark it as complete to proceed to the next one.} \hfill \textbf{\textcolor{brandDark}{$\checkmark$ Completed}}

\end{taskbox}

\begin{taskbox}{Task 5: Audit the Dashboard for Responsible AI}

    {\small \textcolor{gray}{2--3 hrs $\cdot$ Intermediate}}\\[12pt]

    \#\#\#\# \textbf{Scenario}  
Your manager is concerned about the ethical implications of using AI for sustainability reporting. You need to audit your dashboard and explain how responsible AI principles have been incorporated.

\#\#\#\# \textbf{Student Assignment}
\begin{itemize}
    \item Perform an audit of your dashboard and AI narrative analysis.  
    \item Evaluate potential biases or inaccuracies in the data and AI-generated insights.  
    \item Suggest improvements to enhance transparency and trust.
\end{itemize}

\#\#\#\# \textbf{Deliverable}  
Submit a 500-word audit report highlighting the responsible AI measures taken and areas for improvement.\par

    \vspace{16pt}

    \textcolor{cardBorder}{\hrule height 1pt}\vspace{8pt}

    \small \textcolor{gray}{When you are done with this task, mark it as complete to proceed to the next one.} \hfill \textbf{\textcolor{brandDark}{$\checkmark$ Completed}}

\end{taskbox}

\begin{taskbox}{Task 6: Present Your Recommendations}

    {\small \textcolor{gray}{2--3 hrs $\cdot$ Intermediate}}\\[12pt]

    \#\#\#\# \textbf{Scenario}  
You are invited to present your carbon dashboard and reduction strategy at the company’s sustainability meeting. Prepare a presentation that effectively communicates your findings and recommendations.

\#\#\#\# \textbf{Student Assignment}
\begin{itemize}
    \item Create a slide deck with the following sections:  
    \item Company’s emissions overview.  
    \item Dashboard insights and visualizations.  
    \item Proposed reduction strategies with cost implications and potential ROI.  
\end{itemize}

\#\#\#\# \textbf{Deliverable}  
Submit your presentation deck in .pptx or Google Slides format.\par

    \vspace{16pt}

    \textcolor{cardBorder}{\hrule height 1pt}\vspace{8pt}

    \small \textcolor{gray}{When you are done with this task, mark it as complete to proceed to the next one.} \hfill \textbf{\textcolor{brandDark}{$\checkmark$ Completed}}

\end{taskbox}

\begin{completionbox}
    \textbf{Knowledge Assessment \& Verification:} Complete the online knowledge check, submit your code and presentation files for AI verification, and claim your \textbf{Skillzza Verified Certificate}.
\end{completionbox}

\newpage

% =============================================================================

% SIMULATION 45

% =============================================================================

\section*{Real Estate AI Analyst – Predict Property Demand}

\addcontentsline{toc}{section}{Real Estate AI Analyst – Predict Property Demand}

\begin{metabox}
    \begin{tabularx}{\textwidth}{@{}ll X@{}}
        \textbf{Industry:} & Technology & \textbf{Partner Enterprise:} \textbf{Skillzza Enterprise} \\
        \textbf{Duration:} & 12-18 hrs (Self-paced) & \textbf{Mode:} Individual, task-based simulation \\
        \textbf{Primary Skills:} & \multicolumn{2}{@{}l}{Python, AI, Data Analysis, Problem Solving}
    \end{tabularx}
\end{metabox}

\subsection*{Overview \& Problem Statement}

\#\#\# \textbf{Scenario:}
The real estate industry is undergoing a massive transformation, driven by AI and data analytics. As a Real Estate AI Analyst in a PropTech startup, your role is to leverage cutting-edge AI models to predict property demand across different locations. Your insights will help the company identify high-demand areas, optimize investments, and provide data-driven recommendations to clients, including property developers and investors.

\#\#\# \textbf{Mission:}
Your mission is to build and deploy an AI-driven Property Demand Prediction Model using market data, evaluate location-specific demand metrics, and present your insights to key stakeholders for strategic decision-making. You’ll need to ensure that your recommendations are ethically sound, explainable, and aligned with sustainable urban planning principles.

\#\#\# \textbf{Final Challenge:}
Deploy a predictive demand model for a new metropolitan area, analyze patterns across neighborhoods, and create a visually compelling report to recommend the top 3 high-demand areas for development. Justify your selections with data insights, responsible AI considerations, and potential business implications.

---

\subsection*{Tasks \& Deliverables}

\begin{taskbox}{Task 1: Understand the Real Estate Market}

    {\small \textcolor{gray}{2--3 hrs $\cdot$ Intermediate}}\\[12pt]

    \#\#\#\# 
Your manager has provided a dataset containing real estate market data, including historical property prices, population growth, household income, crime rates, and proximity to schools and public transport. Your first task is to explore the dataset and identify key factors that influence property demand.

\#\#\#\# \\[10pt]
\vspace{8pt}
{\large \textbf{\textcolor{brandLight}{What you'll learn}}}\\[4pt]
\begin{itemize}
    \item Analyze the dataset and identify relevant trends and patterns.
    \item Calculate correlations between features like property prices, population density, and proximity to amenities.
    \item Present a summary of your findings.
\end{itemize}

\#\#\#\# \vspace{8pt}
{\large \textbf{\textcolor{brandLight}{What you'll do}}}\\[4pt]
\begin{itemize}
    \item A Jupyter Notebook with exploratory data analysis (EDA) including visualizations and summary statistics.\par
\end{itemize}

    \vspace{16pt}

    \textcolor{cardBorder}{\hrule height 1pt}\vspace{8pt}

    \small \textcolor{gray}{When you are done with this task, mark it as complete to proceed to the next one.} \hfill \textbf{\textcolor{brandDark}{$\checkmark$ Completed}}

\end{taskbox}

\begin{taskbox}{Task 2: Data Cleaning and Feature Engineering}

    {\small \textcolor{gray}{2--3 hrs $\cdot$ Intermediate}}\\[12pt]

    \#\#\#\# 
The dataset contains missing values and inconsistent formats. You’ve also noticed some features that could be engineered to better capture demand signals, such as "Distance to Nearest School" or "Crime Rate Index."

\#\#\#\# \\[10pt]
\vspace{8pt}
{\large \textbf{\textcolor{brandLight}{What you'll learn}}}\\[4pt]
\begin{itemize}
    \item Handle missing data (e.g., imputation or removal).
    \item Engineer new features that could improve model prediction accuracy.
    \item Normalize and scale features for better model performance.
\end{itemize}

\#\#\#\# \vspace{8pt}
{\large \textbf{\textcolor{brandLight}{What you'll do}}}\\[4pt]
\begin{itemize}
    \item A cleaned and feature-engineered dataset in CSV format.
    \item A Python script detailing all transformations.\par
\end{itemize}

    \vspace{16pt}

    \textcolor{cardBorder}{\hrule height 1pt}\vspace{8pt}

    \small \textcolor{gray}{When you are done with this task, mark it as complete to proceed to the next one.} \hfill \textbf{\textcolor{brandDark}{$\checkmark$ Completed}}

\end{taskbox}

\begin{taskbox}{Task 3: Build and Train Your Property Demand Model}

    {\small \textcolor{gray}{2--3 hrs $\cdot$ Intermediate}}\\[12pt]

    \#\#\#\# 
Using the cleaned dataset, you’ll now build a machine learning model to predict property demand in different areas. Your manager has suggested starting with regression models but encourages experimentation with other algorithms like Decision Trees or Gradient Boosting.

\#\#\#\# \\[10pt]
\vspace{8pt}
{\large \textbf{\textcolor{brandLight}{What you'll learn}}}\\[4pt]
\begin{itemize}
    \item Train a regression model to predict demand (target variable: Demand Score).
    \item Evaluate the model using appropriate metrics (e.g., RMSE, MAE).
    \item Experiment with at least one advanced algorithm and compare results.
\end{itemize}

\#\#\#\# \vspace{8pt}
{\large \textbf{\textcolor{brandLight}{What you'll do}}}\\[4pt]
\begin{itemize}
    \item A trained model file (e.g., `.pkl`) and evaluation report.
    \item A Jupyter Notebook showing the model-building process.\par
\end{itemize}

    \vspace{16pt}

    \textcolor{cardBorder}{\hrule height 1pt}\vspace{8pt}

    \small \textcolor{gray}{When you are done with this task, mark it as complete to proceed to the next one.} \hfill \textbf{\textcolor{brandDark}{$\checkmark$ Completed}}

\end{taskbox}

\begin{taskbox}{Task 4: Geospatial Analysis and Demand Visualization}

    {\small \textcolor{gray}{2--3 hrs $\cdot$ Intermediate}}\\[12pt]

    \#\#\#\# 
Your trained model is ready, but stakeholders need to see the results in an actionable format. Use geospatial mapping techniques to visualize property demand across different neighborhoods.

\#\#\#\# \\[10pt]
\vspace{8pt}
{\large \textbf{\textcolor{brandLight}{What you'll learn}}}\\[4pt]
\begin{itemize}
    \item Apply your model to predict demand scores for all neighborhoods in the dataset.
    \item Create a heatmap or choropleth map to visually represent demand distribution.
    \item Identify high-demand and low-demand areas with actionable insights.
\end{itemize}

\#\#\#\# \vspace{8pt}
{\large \textbf{\textcolor{brandLight}{What you'll do}}}\\[4pt]
\begin{itemize}
    \item A geospatial visualization (e.g., Folium map or heatmap) highlighting demand patterns.
    \item A short report summarizing insights from the visualization.\par
\end{itemize}

    \vspace{16pt}

    \textcolor{cardBorder}{\hrule height 1pt}\vspace{8pt}

    \small \textcolor{gray}{When you are done with this task, mark it as complete to proceed to the next one.} \hfill \textbf{\textcolor{brandDark}{$\checkmark$ Completed}}

\end{taskbox}

\begin{taskbox}{Task 5: Audit for Responsible AI}

    {\small \textcolor{gray}{2--3 hrs $\cdot$ Intermediate}}\\[12pt]

    \#\#\#\# 
A stakeholder raised concerns about potential biases in your model. For instance, could the model disadvantage certain neighborhoods (e.g., low-income areas) due to biased data?

\#\#\#\# \\[10pt]
\vspace{8pt}
{\large \textbf{\textcolor{brandLight}{What you'll learn}}}\\[4pt]
\begin{itemize}
    \item Analyze the model for biases using fairness metrics.
    \item Document potential risks and propose solutions to mitigate bias (e.g., reweighting data, excluding problematic features).
    \item Ensure model predictions are explainable to non-technical stakeholders.
\end{itemize}

\#\#\#\# \vspace{8pt}
{\large \textbf{\textcolor{brandLight}{What you'll do}}}\\[4pt]
\begin{itemize}
    \item A bias audit report including metrics and mitigation strategies.
    \item A presentation slide summarizing Responsible AI considerations.\par
\end{itemize}

    \vspace{16pt}

    \textcolor{cardBorder}{\hrule height 1pt}\vspace{8pt}

    \small \textcolor{gray}{When you are done with this task, mark it as complete to proceed to the next one.} \hfill \textbf{\textcolor{brandDark}{$\checkmark$ Completed}}

\end{taskbox}

\begin{taskbox}{Task 6: Present Your Recommendations}

    {\small \textcolor{gray}{2--3 hrs $\cdot$ Intermediate}}\\[12pt]

    \#\#\#\# 
The final challenge is to present your findings to the executive team. They are interested in the top 3 neighborhoods for investment, along with data-driven justifications and an explanation of the model’s predictions.

\#\#\#\# \\[10pt]
\vspace{8pt}
{\large \textbf{\textcolor{brandLight}{What you'll learn}}}\\[4pt]
\begin{itemize}
    \item Create a presentation summarizing your findings, including:
    \item Key features influencing demand.
    \item Geospatial visualization of high-demand areas.
    \item Top 3 neighborhoods for investment, with reasoning.
    \item Responsible AI considerations.
    \item Record a 5-minute video presenting your recommendations.
\end{itemize}

\#\#\#\# \vspace{8pt}
{\large \textbf{\textcolor{brandLight}{What you'll do}}}\\[4pt]
\begin{itemize}
    \item A PowerPoint or PDF presentation.
    \item A recorded video (5 minutes max) presenting your findings and recommendations.\par
\end{itemize}

    \vspace{16pt}

    \textcolor{cardBorder}{\hrule height 1pt}\vspace{8pt}

    \small \textcolor{gray}{When you are done with this task, mark it as complete to proceed to the next one.} \hfill \textbf{\textcolor{brandDark}{$\checkmark$ Completed}}

\end{taskbox}

\begin{completionbox}
    \textbf{Knowledge Assessment \& Verification:} Complete the online knowledge check, submit your code and presentation files for AI verification, and claim your \textbf{Skillzza Verified Certificate}.
\end{completionbox}

\newpage

% =============================================================================

% SIMULATION 46

% =============================================================================

\section*{LegalTech Analyst – Contract Risk Detection}

\addcontentsline{toc}{section}{LegalTech Analyst – Contract Risk Detection}

\begin{metabox}
    \begin{tabularx}{\textwidth}{@{}ll X@{}}
        \textbf{Industry:} & Technology & \textbf{Partner Enterprise:} \textbf{Skillzza Enterprise} \\
        \textbf{Duration:} & 12-18 hrs (Self-paced) & \textbf{Mode:} Individual, task-based simulation \\
        \textbf{Primary Skills:} & \multicolumn{2}{@{}l}{Python, AI, Data Analysis, Problem Solving}
    \end{tabularx}
\end{metabox}

\subsection*{Overview \& Problem Statement}

\#\#\# \textbf{Scenario}  
In the fast-evolving LegalTech industry, companies are leveraging AI to automate contract analysis, improve legal workflows, and detect risks hidden in contract clauses. You’ve just been hired as a LegalTech Analyst at a leading AI-powered legal services firm, \textbf{ClauseGuard Technologies}. Your mission is to analyze legal contracts, identify risky clauses, and provide actionable recommendations to mitigate potential legal and business risks.

\#\#\# \textbf{Mission}  
Your task is to assist the legal team by building a scalable AI pipeline to:  
\begin{itemize}
    \item Extract relevant clauses from contracts.  
    \item Detect and flag risky clauses using pre-defined risk categories.  
    \item Summarize risks and provide actionable recommendations for resolution.  
\end{itemize}

\#\#\# \textbf{Final Challenge}  
For your final assignment, you will analyze a real-world contract, use AI tools to detect risks, explain the rationale behind the flagged risks, and deliver a concise summary report with actionable recommendations to senior stakeholders.

---

\subsection*{Tasks \& Deliverables}

\begin{taskbox}{Task 1: Understand}

    {\small \textcolor{gray}{2--3 hrs $\cdot$ Intermediate}}\\[12pt]

    \#\#\#\# \textbf{Scenario}  
You are given a legal contract and a risk taxonomy document that defines common risk types (e.g., termination, indemnity, and liability risks). Your first step is to thoroughly understand risk categories and their implications.  

\#\#\#\# \textbf{Student Assignment}  
\begin{itemize}
    \item Read the provided \textbf{Risk Taxonomy Document} and identify key risk types.  
    \item Review the sample contract and highlight clauses that may correspond to risk categories.  
\end{itemize}

\#\#\#\# \textbf{Deliverable}  
\begin{itemize}
    \item A 1-page summary of the risk categories and examples of clauses associated with each category.\par
\end{itemize}

    \vspace{16pt}

    \textcolor{cardBorder}{\hrule height 1pt}\vspace{8pt}

    \small \textcolor{gray}{When you are done with this task, mark it as complete to proceed to the next one.} \hfill \textbf{\textcolor{brandDark}{$\checkmark$ Completed}}

\end{taskbox}

\begin{taskbox}{Task 2: Data/Setup}

    {\small \textcolor{gray}{2--3 hrs $\cdot$ Intermediate}}\\[12pt]

    \#\#\#\# \textbf{Scenario}  
To begin automated analysis, you need to preprocess the contract text and prepare it for clause extraction.  

\#\#\#\# \textbf{Student Assignment}  
1. Load the sample contract into a Python script.  
2. Use \textbf{spaCy} to split the contract into clauses and tokenize the text.  
3. Annotate a subset of clauses manually to create a training dataset.  

\#\#\#\# \textbf{Deliverable}  
\begin{itemize}
    \item A CSV file with tokenized clauses and annotated risk categories for at least 20 clauses.\par
\end{itemize}

    \vspace{16pt}

    \textcolor{cardBorder}{\hrule height 1pt}\vspace{8pt}

    \small \textcolor{gray}{When you are done with this task, mark it as complete to proceed to the next one.} \hfill \textbf{\textcolor{brandDark}{$\checkmark$ Completed}}

\end{taskbox}

\begin{taskbox}{Task 3: Build/Execute}

    {\small \textcolor{gray}{2--3 hrs $\cdot$ Intermediate}}\\[12pt]

    \#\#\#\# \textbf{Scenario}  
Now, it’s time to build an AI model to classify clauses into risk categories.  

\#\#\#\# \textbf{Student Assignment}  
1. Fine-tune a \textbf{Hugging Face Transformers} model on the annotated clause dataset.  
2. Implement a pipeline to classify clauses into risk categories.  
3. Evaluate the model's performance (precision, recall, F1-score).  

\#\#\#\# \textbf{Deliverable}  
\begin{itemize}
    \item A Jupyter Notebook with model training code, evaluation metrics, and sample predictions.\par
\end{itemize}

    \vspace{16pt}

    \textcolor{cardBorder}{\hrule height 1pt}\vspace{8pt}

    \small \textcolor{gray}{When you are done with this task, mark it as complete to proceed to the next one.} \hfill \textbf{\textcolor{brandDark}{$\checkmark$ Completed}}

\end{taskbox}

\begin{taskbox}{Task 4: GenAI/Explanation}

    {\small \textcolor{gray}{2--3 hrs $\cdot$ Intermediate}}\\[12pt]

    \#\#\#\# \textbf{Scenario}  
Your AI model has flagged several risky clauses. Your task is to explain the flagged risks in simple terms to non-technical legal stakeholders.  

\#\#\#\# \textbf{Student Assignment}  
\begin{itemize}
    \item Use \textbf{LangChain} or a similar framework to generate explanations for flagged clauses.  
    \item Provide a human-readable summary of why each clause was flagged.  
\end{itemize}

\#\#\#\# \textbf{Deliverable}  
\begin{itemize}
    \item A document with 5 flagged clauses, their risk categories, and AI-generated explanations.\par
\end{itemize}

    \vspace{16pt}

    \textcolor{cardBorder}{\hrule height 1pt}\vspace{8pt}

    \small \textcolor{gray}{When you are done with this task, mark it as complete to proceed to the next one.} \hfill \textbf{\textcolor{brandDark}{$\checkmark$ Completed}}

\end{taskbox}

\begin{taskbox}{Task 5: Audit/Responsible AI}

    {\small \textcolor{gray}{2--3 hrs $\cdot$ Intermediate}}\\[12pt]

    \#\#\#\# \textbf{Scenario}  
Your legal team is concerned about potential bias or errors in the AI system. You need to conduct an audit to ensure the AI is reliable and fair.  

\#\#\#\# \textbf{Student Assignment}  
1. Analyze false positives and false negatives in the model's predictions.  
2. Propose improvements for reducing model bias (e.g., expanding the training dataset).  
3. Document how your model complies with Responsible AI guidelines.  

\#\#\#\# \textbf{Deliverable}  
\begin{itemize}
    \item An audit report with error analysis and recommendations for improvement.\par
\end{itemize}

    \vspace{16pt}

    \textcolor{cardBorder}{\hrule height 1pt}\vspace{8pt}

    \small \textcolor{gray}{When you are done with this task, mark it as complete to proceed to the next one.} \hfill \textbf{\textcolor{brandDark}{$\checkmark$ Completed}}

\end{taskbox}

\begin{taskbox}{Task 6: Present Recommendation}

    {\small \textcolor{gray}{2--3 hrs $\cdot$ Intermediate}}\\[12pt]

    \#\#\#\# \textbf{Scenario}  
You will present your findings and recommendations to senior legal stakeholders in a concise, actionable format.  

\#\#\#\# \textbf{Student Assignment}  
\begin{itemize}
    \item Create a risk summary report for the analyzed contract.  
    \item Include key risks, their implications, and actionable recommendations.  
    \item Prepare a 5-slide presentation summarizing the results.  
\end{itemize}

\#\#\#\# \textbf{Deliverable}  
\begin{itemize}
    \item A risk summary report (PDF/Word) and a slide deck (PDF/PPT).\par
\end{itemize}

    \vspace{16pt}

    \textcolor{cardBorder}{\hrule height 1pt}\vspace{8pt}

    \small \textcolor{gray}{When you are done with this task, mark it as complete to proceed to the next one.} \hfill \textbf{\textcolor{brandDark}{$\checkmark$ Completed}}

\end{taskbox}

\begin{completionbox}
    \textbf{Knowledge Assessment \& Verification:} Complete the online knowledge check, submit your code and presentation files for AI verification, and claim your \textbf{Skillzza Verified Certificate}.
\end{completionbox}

\newpage

% =============================================================================

% SIMULATION 47

% =============================================================================

\section*{Citizen Services AI Agent – Resolve a Government Complaint}

\addcontentsline{toc}{section}{Citizen Services AI Agent – Resolve a Government Complaint}

\begin{metabox}
    \begin{tabularx}{\textwidth}{@{}ll X@{}}
        \textbf{Industry:} & Technology & \textbf{Partner Enterprise:} \textbf{Skillzza Enterprise} \\
        \textbf{Duration:} & 12-18 hrs (Self-paced) & \textbf{Mode:} Individual, task-based simulation \\
        \textbf{Primary Skills:} & \multicolumn{2}{@{}l}{Python, AI, Data Analysis, Problem Solving}
    \end{tabularx}
\end{metabox}

\subsection*{Overview \& Problem Statement}

\#\#\# \textbf{Scenario}  
Imagine working as an AI Specialist for a government innovation lab tasked with improving citizen services. Your mission is to enhance the efficiency of handling citizen complaints by designing an AI-based solution that automates complaint resolution, reduces response times, and ensures that unresolved complaints are escalated appropriately.  

You will be working with a dataset of citizen complaints categorized by department, issue type, and resolution status. Your job is to build a Citizen Services AI Agent capable of classifying complaints, retrieving relevant knowledge articles for resolution, and escalating unresolved complaints to the appropriate department.

\#\#\# \textbf{Mission}  
Your task is to design and execute a pipeline for the Citizen Services AI Agent that:  
1. \textbf{Classifies complaints} based on department and urgency.  
2. \textbf{Retrieves knowledge base articles} to offer automated resolutions.  
3. Determines when a complaint needs \textbf{human escalation} and routes it to the correct department.  

\#\#\# \textbf{Final Challenge}  
By the end of this simulation, you will present a functional AI pipeline alongside a report detailing classification accuracy, knowledge retrieval effectiveness, and escalation logic. Your solution will demonstrate how this system can improve citizen satisfaction and reduce service delays.

---

\subsection*{Tasks \& Deliverables}

\begin{taskbox}{Task 1: Understand the Problem}

    {\small \textcolor{gray}{2--3 hrs $\cdot$ Intermediate}}\\[12pt]

    \#\#\#\# \textbf{Scenario}  
A government services portal receives 10,000+ citizen complaints every month. These complaints vary from potholes and waste management to utility bill disputes. Your first step is to understand the dataset and problem statement.  

\#\#\#\# \textbf{Student Assignment}  
1. Explore a provided dataset containing complaints with fields:  
\begin{itemize}
    \item `Complaint\_ID`: Unique identifier.  
    \item `Citizen\_Name`: Name of complainant.  
    \item `Complaint\_Text`: The full description of the complaint.  
    \item `Department`: Department responsible for resolution (e.g., Waste Management, Public Works).  
    \item `Urgency`: Low, Medium, High.  
    \item `Status`: Resolved, Pending, Escalated.  
2. Analyze the dataset for patterns, such as the most common complaint types and departments with the highest complaint volumes.  
\end{itemize}

\#\#\#\# \textbf{Deliverable}  
Submit a short report (300 words) with:  
\begin{itemize}
    \item Key patterns in the dataset.  
    \item Suggestions for how AI can reduce backlog and improve response times.\par
\end{itemize}

    \vspace{16pt}

    \textcolor{cardBorder}{\hrule height 1pt}\vspace{8pt}

    \small \textcolor{gray}{When you are done with this task, mark it as complete to proceed to the next one.} \hfill \textbf{\textcolor{brandDark}{$\checkmark$ Completed}}

\end{taskbox}

\begin{taskbox}{Task 2: Prepare the Data}

    {\small \textcolor{gray}{2--3 hrs $\cdot$ Intermediate}}\\[12pt]

    \#\#\#\# \textbf{Scenario}  
You must clean and preprocess the complaint data to prepare it for text classification.  

\#\#\#\# \textbf{Student Assignment}  
1. Perform text preprocessing on the `Complaint\_Text` field, including:  
\begin{itemize}
    \item Tokenization.  
    \item Stopword removal.  
    \item Lemmatization.  
2. Split the dataset into training, validation, and test sets.  
3. Encode the labels (`Department`) for classification using one-hot encoding or label encoding.  
\end{itemize}

\#\#\#\# \textbf{Deliverable}  
Submit the cleaned dataset and a summary of your preprocessing steps.\par

    \vspace{16pt}

    \textcolor{cardBorder}{\hrule height 1pt}\vspace{8pt}

    \small \textcolor{gray}{When you are done with this task, mark it as complete to proceed to the next one.} \hfill \textbf{\textcolor{brandDark}{$\checkmark$ Completed}}

\end{taskbox}

\begin{taskbox}{Task 3: Build the Classification Model}

    {\small \textcolor{gray}{2--3 hrs $\cdot$ Intermediate}}\\[12pt]

    \#\#\#\# \textbf{Scenario}  
Now, you will build a machine learning model to classify complaints into the correct department.  

\#\#\#\# \textbf{Student Assignment}  
1. Train a text classification model using Hugging Face Transformers (e.g., DistilBERT).  
2. Evaluate the model using metrics like accuracy, precision, recall, and F1-score. Aim for an F1-score > 0.8.  
3. Save the trained model for later use in the pipeline.  

\#\#\#\# \textbf{Deliverable}  
Submit the trained model file and a report on model performance (include a confusion matrix).\par

    \vspace{16pt}

    \textcolor{cardBorder}{\hrule height 1pt}\vspace{8pt}

    \small \textcolor{gray}{When you are done with this task, mark it as complete to proceed to the next one.} \hfill \textbf{\textcolor{brandDark}{$\checkmark$ Completed}}

\end{taskbox}

\begin{taskbox}{Task 4: Integrate Knowledge Retrieval}

    {\small \textcolor{gray}{2--3 hrs $\cdot$ Intermediate}}\\[12pt]

    \#\#\#\# \textbf{Scenario}  
Once a complaint is classified, the AI Agent should retrieve relevant knowledge articles to suggest resolutions.  

\#\#\#\# \textbf{Student Assignment}  
1. Load a provided knowledge base with fields:  
\begin{itemize}
    \item `Article\_ID`: Unique identifier.  
    \item `Department`: Associated department.  
    \item `Keywords`: Tags for retrieval.  
    \item `Article\_Text`: Full text of the article.  
2. Use embeddings (e.g., SentenceTransformers) to generate vector representations for the articles and complaints.  
3. Implement a similarity search algorithm (e.g., cosine similarity) to retrieve the top 3 articles for each complaint.  
\end{itemize}

\#\#\#\# \textbf{Deliverable}  
Submit your knowledge retrieval script and an output file showing the top 3 articles for 10 test complaints.\par

    \vspace{16pt}

    \textcolor{cardBorder}{\hrule height 1pt}\vspace{8pt}

    \small \textcolor{gray}{When you are done with this task, mark it as complete to proceed to the next one.} \hfill \textbf{\textcolor{brandDark}{$\checkmark$ Completed}}

\end{taskbox}

\begin{taskbox}{Task 5: Test Responsible AI Practices}

    {\small \textcolor{gray}{2--3 hrs $\cdot$ Intermediate}}\\[12pt]

    \#\#\#\# \textbf{Scenario}  
AI systems for citizen services must be fair and unbiased. You need to audit your classification model for potential biases.  

\#\#\#\# \textbf{Student Assignment}  
1. Analyze model performance across different demographics (e.g., names indicating gender or region).  
2. Check for overrepresentation or neglect of certain complaint categories.  
3. Propose methods to address any bias found (e.g., rebalancing training data).  

\#\#\#\# \textbf{Deliverable}  
Submit an audit report (300 words) outlining your findings and proposed solutions.\par

    \vspace{16pt}

    \textcolor{cardBorder}{\hrule height 1pt}\vspace{8pt}

    \small \textcolor{gray}{When you are done with this task, mark it as complete to proceed to the next one.} \hfill \textbf{\textcolor{brandDark}{$\checkmark$ Completed}}

\end{taskbox}

\begin{taskbox}{Task 6: Present Final Recommendation}

    {\small \textcolor{gray}{2--3 hrs $\cdot$ Intermediate}}\\[12pt]

    \#\#\#\# \textbf{Scenario}  
You are pitching your solution to government stakeholders. Summarize your AI pipeline and demonstrate its impact on complaint resolution efficiency.  

\#\#\#\# \textbf{Student Assignment}  
1. Create a presentation (5 slides) covering:  
\begin{itemize}
    \item Problem statement and dataset analysis.  
    \item Classification model results.  
    \item Knowledge retrieval accuracy.  
    \item Responsible AI findings.  
    \item Overall impact and next steps.  
2. Record a 3-minute video walkthrough of your presentation.  
\end{itemize}

\#\#\#\# \textbf{Deliverable}  
Submit your slides and video walkthrough.\par

    \vspace{16pt}

    \textcolor{cardBorder}{\hrule height 1pt}\vspace{8pt}

    \small \textcolor{gray}{When you are done with this task, mark it as complete to proceed to the next one.} \hfill \textbf{\textcolor{brandDark}{$\checkmark$ Completed}}

\end{taskbox}

\begin{completionbox}
    \textbf{Knowledge Assessment \& Verification:} Complete the online knowledge check, submit your code and presentation files for AI verification, and claim your \textbf{Skillzza Verified Certificate}.
\end{completionbox}

\newpage

% =============================================================================

% SIMULATION 48

% =============================================================================

\section*{Drone Technology Analyst – AI-Based Infrastructure Inspection}

\addcontentsline{toc}{section}{Drone Technology Analyst – AI-Based Infrastructure Inspection}

\begin{metabox}
    \begin{tabularx}{\textwidth}{@{}ll X@{}}
        \textbf{Industry:} & Technology & \textbf{Partner Enterprise:} \textbf{Skillzza Enterprise} \\
        \textbf{Duration:} & 12-18 hrs (Self-paced) & \textbf{Mode:} Individual, task-based simulation \\
        \textbf{Primary Skills:} & \multicolumn{2}{@{}l}{Python, AI, Data Analysis, Problem Solving}
    \end{tabularx}
\end{metabox}

\subsection*{Overview \& Problem Statement}

\#\#\# \textbf{Scenario}  
In the rapidly evolving DroneTech industry, drones equipped with AI capabilities are transforming infrastructure inspections by delivering precise, efficient, and cost-effective results. As a Drone Technology Analyst, you'll play a critical role in ensuring infrastructure safety by leveraging drone-collected data and AI-driven models to detect structural risks in bridges, towers, and pipelines.  

\#\#\# \textbf{Mission}  
Your mission is to simulate a professional role in a DroneTech company tasked with analyzing drone-captured imagery and sensor data. You will develop AI models to detect potential infrastructure risks, audit model performance for responsible AI practices, and present actionable insights to stakeholders.  

\#\#\# \textbf{Final Challenge}  
At the end of this internship, you will deliver a comprehensive risk report based on drone imagery captured from a bridge inspection. The report will include AI-detected anomalies, risk categorization, and actionable recommendations for remediation.

---

\subsection*{Tasks \& Deliverables}

\begin{taskbox}{Task 1: Understand}

    {\small \textcolor{gray}{2--3 hrs $\cdot$ Intermediate}}\\[12pt]

    \textbf{Scenario}: Your company has been hired to inspect a bridge for structural vulnerabilities using drone technology. You receive a briefing about the inspection goals and the expected outcomes.  
\textbf{Student Assignment}: Review the inspection brief, learn about the structural risks commonly associated with bridges, and understand the data format captured by drones (images and sensor logs).  
\textbf{Deliverable}: Submit a summary document outlining the inspection objectives and risk indicators.\par

    \vspace{16pt}

    \textcolor{cardBorder}{\hrule height 1pt}\vspace{8pt}

    \small \textcolor{gray}{When you are done with this task, mark it as complete to proceed to the next one.} \hfill \textbf{\textcolor{brandDark}{$\checkmark$ Completed}}

\end{taskbox}

\begin{taskbox}{Task 2: Data/Setup}

    {\small \textcolor{gray}{2--3 hrs $\cdot$ Intermediate}}\\[12pt]

    \textbf{Scenario}: You have received raw drone data comprising high-resolution images and sensor readings. Before AI modeling, you need to preprocess and annotate the data.  
\textbf{Student Assignment}:  
\begin{itemize}
    \item Clean the image dataset to remove duplicates or irrelevant data.  
    \item Annotate images to highlight cracks, rust, or other anomalies using tools like LabelImg.  
    \item Split the dataset into training and testing subsets.  
\textbf{Deliverable}: Submit the preprocessed dataset and a summary of your annotation process.\par
\end{itemize}

    \vspace{16pt}

    \textcolor{cardBorder}{\hrule height 1pt}\vspace{8pt}

    \small \textcolor{gray}{When you are done with this task, mark it as complete to proceed to the next one.} \hfill \textbf{\textcolor{brandDark}{$\checkmark$ Completed}}

\end{taskbox}

\begin{taskbox}{Task 3: Build/Execute}

    {\small \textcolor{gray}{2--3 hrs $\cdot$ Intermediate}}\\[12pt]

    \textbf{Scenario}: Your team is tasked with training an AI model to detect anomalies in the drone imagery. You will use CNNs for image classification and anomaly detection.  
\textbf{Student Assignment}:  
\begin{itemize}
    \item Build a convolutional neural network (CNN) using TensorFlow or PyTorch.  
    \item Train the model using the annotated dataset.  
    \item Evaluate the model's accuracy and generate predictions on the test set.  
\textbf{Deliverable}: Submit the Python code, model evaluation metrics (accuracy, precision, recall), and predictions.\par
\end{itemize}

    \vspace{16pt}

    \textcolor{cardBorder}{\hrule height 1pt}\vspace{8pt}

    \small \textcolor{gray}{When you are done with this task, mark it as complete to proceed to the next one.} \hfill \textbf{\textcolor{brandDark}{$\checkmark$ Completed}}

\end{taskbox}

\begin{taskbox}{Task 4: GenAI/Explanation}

    {\small \textcolor{gray}{2--3 hrs $\cdot$ Intermediate}}\\[12pt]

    \textbf{Scenario}: Stakeholders need a clear explanation of the AI results and actionable recommendations based on detected risks. Use Generative AI tools to draft a professional report.  
\textbf{Student Assignment}:  
\begin{itemize}
    \item Input AI prediction results into a GPT-based tool to generate a stakeholder-friendly report.  
    \item Ensure the report includes detected anomalies, severity levels, and remediation steps.  
\textbf{Deliverable}: Submit the AI-generated report in PDF format.\par
\end{itemize}

    \vspace{16pt}

    \textcolor{cardBorder}{\hrule height 1pt}\vspace{8pt}

    \small \textcolor{gray}{When you are done with this task, mark it as complete to proceed to the next one.} \hfill \textbf{\textcolor{brandDark}{$\checkmark$ Completed}}

\end{taskbox}

\begin{taskbox}{Task 5: Audit/Responsible AI}

    {\small \textcolor{gray}{2--3 hrs $\cdot$ Intermediate}}\\[12pt]

    \textbf{Scenario}: Before releasing the AI model for operational use, you need to ensure it adheres to responsible AI practices. This includes auditing for bias and interpretability.  
\textbf{Student Assignment}:  
\begin{itemize}
    \item Use SHAP or Lime to interpret AI model predictions.  
    \item Verify that the model is free from bias (e.g., ensuring consistent detection across different image qualities).  
    \item Document the audit findings.  
\textbf{Deliverable}: Submit the audit report highlighting interpretability insights and bias checks.\par
\end{itemize}

    \vspace{16pt}

    \textcolor{cardBorder}{\hrule height 1pt}\vspace{8pt}

    \small \textcolor{gray}{When you are done with this task, mark it as complete to proceed to the next one.} \hfill \textbf{\textcolor{brandDark}{$\checkmark$ Completed}}

\end{taskbox}

\begin{taskbox}{Task 6: Present Recommendation}

    {\small \textcolor{gray}{2--3 hrs $\cdot$ Intermediate}}\\[12pt]

    \textbf{Scenario}: Your final task is to present actionable recommendations to the engineering team based on your analysis. Stakeholders expect concise and impactful insights.  
\textbf{Student Assignment}:  
\begin{itemize}
    \item Create a presentation summarizing the bridge inspection results, AI model outputs, and risk categorization.  
    \item Include visualizations and a clear call-to-action for remediation priorities.  
\textbf{Deliverable}: Submit the presentation in PowerPoint or PDF format.\par
\end{itemize}

    \vspace{16pt}

    \textcolor{cardBorder}{\hrule height 1pt}\vspace{8pt}

    \small \textcolor{gray}{When you are done with this task, mark it as complete to proceed to the next one.} \hfill \textbf{\textcolor{brandDark}{$\checkmark$ Completed}}

\end{taskbox}

\begin{completionbox}
    \textbf{Knowledge Assessment \& Verification:} Complete the online knowledge check, submit your code and presentation files for AI verification, and claim your \textbf{Skillzza Verified Certificate}.
\end{completionbox}

\newpage

% =============================================================================

% SIMULATION 49

% =============================================================================

\section*{Space Technology Analyst – AI for Satellite Intelligence}

\addcontentsline{toc}{section}{Space Technology Analyst – AI for Satellite Intelligence}

\begin{metabox}
    \begin{tabularx}{\textwidth}{@{}ll X@{}}
        \textbf{Industry:} & Technology & \textbf{Partner Enterprise:} \textbf{Skillzza Enterprise} \\
        \textbf{Duration:} & 12-18 hrs (Self-paced) & \textbf{Mode:} Individual, task-based simulation \\
        \textbf{Primary Skills:} & \multicolumn{2}{@{}l}{Python, AI, Data Analysis, Problem Solving}
    \end{tabularx}
\end{metabox}

\subsection*{Overview \& Problem Statement}

\#\#\# \textbf{Scenario}  
You are recruited as a Space Technology Analyst for a leading SpaceTech organization specializing in satellite intelligence. Your mission is to optimize the use of satellite imagery to detect anomalies, such as environmental changes, infrastructure damage, or illegal activities. You will leverage AI and machine learning algorithms to analyze geospatial data, detect patterns, identify anomalies, and generate actionable intelligence reports for stakeholders like government agencies, environmental organizations, and private enterprises.

\#\#\# \textbf{Mission}  
Your objective is to develop a satellite image analysis pipeline that processes raw geospatial data, identifies anomalies using AI models, and translates findings into actionable intelligence. You will simulate real-world scenarios where timely intelligence is critical for decision-making, such as disaster response, infrastructure monitoring, or environmental conservation.

\#\#\# \textbf{Final Challenge}  
The final challenge involves detecting anomalies in a satellite-generated dataset, validating your findings with a responsible AI framework, and presenting a high-stakes intelligence report to stakeholders. This is your chance to prove your expertise in SpaceTech analytics and AI-driven decision-making.

---

\subsection*{Tasks \& Deliverables}

\begin{taskbox}{Task 1: Understand}

    {\small \textcolor{gray}{2--3 hrs $\cdot$ Intermediate}}\\[12pt]

    \#\#\#\# Scenario  
Your team has received raw satellite imagery for a disaster-prone region. Stakeholders are counting on you to provide actionable insights to prevent future damage. Before diving into the data, you must understand the scope and types of satellite images provided.  

\#\#\#\# Student Assignment  
\begin{itemize}
    \item Research the types of satellite images (optical vs. radar).  
    \item Identify the key metadata fields in the dataset (e.g., resolution, timestamp, geolocation).  
\end{itemize}

\#\#\#\# Deliverable  
A short report summarizing the dataset characteristics, the types of anomalies to detect, and the potential challenges in processing satellite imagery.\par

    \vspace{16pt}

    \textcolor{cardBorder}{\hrule height 1pt}\vspace{8pt}

    \small \textcolor{gray}{When you are done with this task, mark it as complete to proceed to the next one.} \hfill \textbf{\textcolor{brandDark}{$\checkmark$ Completed}}

\end{taskbox}

\begin{taskbox}{Task 2: Data/Setup}

    {\small \textcolor{gray}{2--3 hrs $\cdot$ Intermediate}}\\[12pt]

    \#\#\#\# Scenario  
The raw satellite images are available in TIFF format, but they need to be preprocessed before analysis. Your task is to clean, preprocess, and prepare the data for AI-based anomaly detection.  

\#\#\#\# Student Assignment  
\begin{itemize}
    \item Use OpenCV or similar libraries to load and preprocess images (e.g., resize, grayscale conversion).  
    \item Extract metadata fields (e.g., coordinates, timestamps) for analysis.  
    \item Visualize a subset of images using Matplotlib for exploratory analysis.  
\end{itemize}

\#\#\#\# Deliverable  
A set of preprocessed images and a Jupyter notebook showing the preprocessing pipeline.\par

    \vspace{16pt}

    \textcolor{cardBorder}{\hrule height 1pt}\vspace{8pt}

    \small \textcolor{gray}{When you are done with this task, mark it as complete to proceed to the next one.} \hfill \textbf{\textcolor{brandDark}{$\checkmark$ Completed}}

\end{taskbox}

\begin{taskbox}{Task 3: Build/Execute}

    {\small \textcolor{gray}{2--3 hrs $\cdot$ Intermediate}}\\[12pt]

    \#\#\#\# Scenario  
You are tasked with building an anomaly detection model using AI. The goal is to identify regions in the satellite images where anomalies such as deforestation or infrastructure damage might have occurred.  

\#\#\#\# Student Assignment  
\begin{itemize}
    \item Train an anomaly detection model using unsupervised learning (e.g., autoencoders).  
    \item Test the model on preprocessed satellite images.  
    \item Fine-tune the model to minimize false positives and false negatives.  
\end{itemize}

\#\#\#\# Deliverable  
A trained AI model and a notebook documenting the training process, evaluation metrics, and sample outputs.\par

    \vspace{16pt}

    \textcolor{cardBorder}{\hrule height 1pt}\vspace{8pt}

    \small \textcolor{gray}{When you are done with this task, mark it as complete to proceed to the next one.} \hfill \textbf{\textcolor{brandDark}{$\checkmark$ Completed}}

\end{taskbox}

\begin{taskbox}{Task 4: GenAI/Explanation}

    {\small \textcolor{gray}{2--3 hrs $\cdot$ Intermediate}}\\[12pt]

    \#\#\#\# Scenario  
Stakeholders need to understand the anomalies detected without technical jargon. Using Generative AI tools, you must create a simplified explanation for non-technical audiences.  

\#\#\#\# Student Assignment  
\begin{itemize}
    \item Use Generative AI (e.g., ChatGPT or GPT-based tools) to draft a summary explaining the detected anomalies.  
    \item Include visuals (e.g., annotated satellite images) to support your explanation.  
\end{itemize}

\#\#\#\# Deliverable  
A stakeholder-friendly document summarizing the anomalies and their implications, accompanied by annotated visuals.\par

    \vspace{16pt}

    \textcolor{cardBorder}{\hrule height 1pt}\vspace{8pt}

    \small \textcolor{gray}{When you are done with this task, mark it as complete to proceed to the next one.} \hfill \textbf{\textcolor{brandDark}{$\checkmark$ Completed}}

\end{taskbox}

\begin{taskbox}{Task 5: Audit/Responsible AI}

    {\small \textcolor{gray}{2--3 hrs $\cdot$ Intermediate}}\\[12pt]

    \#\#\#\# Scenario  
Your anomaly detection model must be reviewed for bias and ethical implications. Ensure that your AI pipeline aligns with responsible AI practices.  

\#\#\#\# Student Assignment  
\begin{itemize}
    \item Evaluate the model for any biases (e.g., geographic or data selection bias).  
    \item Apply fairness metrics and transparency checks to your pipeline.  
    \item Suggest improvements to make the AI model more ethical.  
\end{itemize}

\#\#\#\# Deliverable  
A responsible AI audit report outlining the steps taken to ensure fairness and ethics in the analysis.\par

    \vspace{16pt}

    \textcolor{cardBorder}{\hrule height 1pt}\vspace{8pt}

    \small \textcolor{gray}{When you are done with this task, mark it as complete to proceed to the next one.} \hfill \textbf{\textcolor{brandDark}{$\checkmark$ Completed}}

\end{taskbox}

\begin{taskbox}{Task 6: Present Recommendation}

    {\small \textcolor{gray}{2--3 hrs $\cdot$ Intermediate}}\\[12pt]

    \#\#\#\# Scenario  
You are presenting your findings to a panel of stakeholders, including government officials and environmental agencies. Your goal is to deliver actionable recommendations based on the anomalies detected.  

\#\#\#\# Student Assignment  
\begin{itemize}
    \item Create a presentation summarizing the analysis, key findings, and recommendations.  
    \item Include insights on how stakeholders can act on the intelligence provided.  
\end{itemize}

\#\#\#\# Deliverable  
A presentation deck (PPT or PDF) with data visualizations, recommendations, and a clear action plan.\par

    \vspace{16pt}

    \textcolor{cardBorder}{\hrule height 1pt}\vspace{8pt}

    \small \textcolor{gray}{When you are done with this task, mark it as complete to proceed to the next one.} \hfill \textbf{\textcolor{brandDark}{$\checkmark$ Completed}}

\end{taskbox}

\begin{completionbox}
    \textbf{Knowledge Assessment \& Verification:} Complete the online knowledge check, submit your code and presentation files for AI verification, and claim your \textbf{Skillzza Verified Certificate}.
\end{completionbox}

\newpage

% =============================================================================

% SIMULATION 50

% =============================================================================

\section*{AI Startup Analyst – Evaluate an Early-Stage Startup}

\addcontentsline{toc}{section}{AI Startup Analyst – Evaluate an Early-Stage Startup}

\begin{metabox}
    \begin{tabularx}{\textwidth}{@{}ll X@{}}
        \textbf{Industry:} & Technology & \textbf{Partner Enterprise:} \textbf{Skillzza Enterprise} \\
        \textbf{Duration:} & 12-18 hrs (Self-paced) & \textbf{Mode:} Individual, task-based simulation \\
        \textbf{Primary Skills:} & \multicolumn{2}{@{}l}{Python, AI, Data Analysis, Problem Solving}
    \end{tabularx}
\end{metabox}

\subsection*{Overview \& Problem Statement}

\#\#\# \textbf{Scenario}  
You are hired as an AI Startup Analyst by a venture capital firm that specializes in investing in early-stage AI startups. Your firm is considering investing in a startup that claims to revolutionize predictive healthcare using AI. Your role is to evaluate the startup holistically—its pitch deck, market opportunity, product feasibility, technology stack, unit economics, and overall investment potential.  

\#\#\# \textbf{Mission}  
Your mission is to conduct a rigorous analysis of the startup's pitch deck and supporting materials to determine if it is a viable investment. You’ll assess the market demand, validate the product’s potential, evaluate the technological stack for scalability, and analyze its unit economics to write an investment memo recommending whether to fund the startup or not.  

\#\#\# \textbf{Final Challenge}  
Write a professional, data-driven investment memo that evaluates the startup across critical parameters. Your memo will be reviewed by a simulated Venture Capital Investment Committee for feedback.  

---

\subsection*{Tasks \& Deliverables}

\begin{taskbox}{Task 1: Understand – Analyze the AI Startup Pitch Deck}

    {\small \textcolor{gray}{2--3 hrs $\cdot$ Intermediate}}\\[12pt]

    \#\#\#\# \textbf{Scenario}  
The startup has shared its 12-slide pitch deck, which includes sections like Problem Statement, Solution, Market Opportunity, Financials, and Competitive Advantage. You need to analyze the deck’s structure and identify gaps.  

\#\#\#\# \textbf{Student Assignment}  
\begin{itemize}
    \item Review the pitch deck and identify key strengths and weaknesses.  
    \item Evaluate the clarity of the problem statement and solution.  
    \item Assess whether the competitive advantage and market opportunity are compelling.  
\end{itemize}

\#\#\#\# \textbf{Deliverable}  
Submit a short report highlighting the strengths, weaknesses, and missing elements in the pitch deck.\par

    \vspace{16pt}

    \textcolor{cardBorder}{\hrule height 1pt}\vspace{8pt}

    \small \textcolor{gray}{When you are done with this task, mark it as complete to proceed to the next one.} \hfill \textbf{\textcolor{brandDark}{$\checkmark$ Completed}}

\end{taskbox}

\begin{taskbox}{Task 2: Data/Setup – Market and Competitive Analysis}

    {\small \textcolor{gray}{2--3 hrs $\cdot$ Intermediate}}\\[12pt]

    \#\#\#\# \textbf{Scenario}  
The startup claims to operate in a market worth \$5 billion annually with only two direct competitors. You’re tasked to validate this claim using external data sources and benchmark competitors.  

\#\#\#\# \textbf{Student Assignment}  
\begin{itemize}
    \item Use tools like Crunchbase and Statista to research market size and trends.  
    \item Identify direct and indirect competitors. Compare their offerings and positioning.  
    \item Calculate TAM, SAM (Serviceable Addressable Market), and SOM (Serviceable Obtainable Market).  
\end{itemize}

\#\#\#\# \textbf{Deliverable}  
Submit a spreadsheet and a 1-page summary of market validation and competitive analysis.\par

    \vspace{16pt}

    \textcolor{cardBorder}{\hrule height 1pt}\vspace{8pt}

    \small \textcolor{gray}{When you are done with this task, mark it as complete to proceed to the next one.} \hfill \textbf{\textcolor{brandDark}{$\checkmark$ Completed}}

\end{taskbox}

\begin{taskbox}{Task 3: Build/Execute – Unit Economics Assessment}

    {\small \textcolor{gray}{2--3 hrs $\cdot$ Intermediate}}\\[12pt]

    \#\#\#\# \textbf{Scenario}  
The startup’s financials suggest a CAC of \$150 and an LTV of \$1,200, with a customer retention rate of 80\%. You need to verify these claims and assess their unit economics.  

\#\#\#\# \textbf{Student Assignment}  
\begin{itemize}
    \item Validate the startup’s CAC and LTV figures using provided financial data.  
    \item Build a simple financial model to assess scalability and break-even point.  
    \item Suggest ways to optimize CAC and improve profitability.  
\end{itemize}

\#\#\#\# \textbf{Deliverable}  
Submit the financial model in Excel or Google Sheets and a brief recommendation report.\par

    \vspace{16pt}

    \textcolor{cardBorder}{\hrule height 1pt}\vspace{8pt}

    \small \textcolor{gray}{When you are done with this task, mark it as complete to proceed to the next one.} \hfill \textbf{\textcolor{brandDark}{$\checkmark$ Completed}}

\end{taskbox}

\begin{taskbox}{Task 4: GenAI/Explanation – Evaluate AI Technology Stack}

    {\small \textcolor{gray}{2--3 hrs $\cdot$ Intermediate}}\\[12pt]

    \#\#\#\# \textbf{Scenario}  
The startup claims to use a proprietary AI algorithm for predictive healthcare analytics. You need to evaluate the feasibility of this claim and check for scalability and ethical risks.  

\#\#\#\# \textbf{Student Assignment}  
\begin{itemize}
    \item Review documentation on the startup’s AI technology stack (e.g., Python-based ML models, cloud hosting platforms).  
    \item Use a generative AI tool to draft a summary of potential challenges in scalability and ethical concerns (e.g., bias, data privacy).  
\end{itemize}

\#\#\#\# \textbf{Deliverable}  
Submit a 2-page evaluation of the AI technology stack, including identified risks and recommendations.\par

    \vspace{16pt}

    \textcolor{cardBorder}{\hrule height 1pt}\vspace{8pt}

    \small \textcolor{gray}{When you are done with this task, mark it as complete to proceed to the next one.} \hfill \textbf{\textcolor{brandDark}{$\checkmark$ Completed}}

\end{taskbox}

\begin{taskbox}{Task 5: Audit/Responsible AI – Perform an AI Ethics Audit}

    {\small \textcolor{gray}{2--3 hrs $\cdot$ Intermediate}}\\[12pt]

    \#\#\#\# \textbf{Scenario}  
The startup has no formal Responsible AI framework. You are tasked with identifying ethical risks and proposing mitigation strategies for its predictive healthcare solution.  

\#\#\#\# \textbf{Student Assignment}  
\begin{itemize}
    \item Audit the startup’s AI solution for potential biases (e.g., gender, ethnicity, socioeconomic).  
    \item Suggest actionable Responsible AI practices, including data governance and transparency measures.  
\end{itemize}

\#\#\#\# \textbf{Deliverable}  
Submit a Responsible AI Audit checklist and a 1-page summary of proposed mitigation strategies.\par

    \vspace{16pt}

    \textcolor{cardBorder}{\hrule height 1pt}\vspace{8pt}

    \small \textcolor{gray}{When you are done with this task, mark it as complete to proceed to the next one.} \hfill \textbf{\textcolor{brandDark}{$\checkmark$ Completed}}

\end{taskbox}

\begin{taskbox}{Task 6: Present Recommendation – Write an Investment Memo}

    {\small \textcolor{gray}{2--3 hrs $\cdot$ Intermediate}}\\[12pt]

    \#\#\#\# \textbf{Scenario}  
The Venture Capital Investment Committee requires a final investment memo summarizing your findings across market, product, technology, and financials.  

\#\#\#\# \textbf{Student Assignment}  
\begin{itemize}
    \item Consolidate findings from previous tasks into a formal investment memo.  
    \item Use visuals (charts, graphs) to strengthen your arguments.  
    \item Make a clear recommendation: Invest or Pass.  
\end{itemize}

\#\#\#\# \textbf{Deliverable}  
Submit a 5-page professional investment memo in PDF format, with supporting charts and financial models.\par

    \vspace{16pt}

    \textcolor{cardBorder}{\hrule height 1pt}\vspace{8pt}

    \small \textcolor{gray}{When you are done with this task, mark it as complete to proceed to the next one.} \hfill \textbf{\textcolor{brandDark}{$\checkmark$ Completed}}

\end{taskbox}

\begin{completionbox}
    \textbf{Knowledge Assessment \& Verification:} Complete the online knowledge check, submit your code and presentation files for AI verification, and claim your \textbf{Skillzza Verified Certificate}.
\end{completionbox}

\end{document}
